\documentclass[12pt,a4paper,openany]{book}
\usepackage[utf8]{inputenc}
\usepackage[T1]{fontenc}
\usepackage{lmodern}
\usepackage[margin=2.5cm]{geometry}
\usepackage{amsmath,amssymb,amsthm,mathrsfs,mathtools}
\usepackage{tikz}
\usetikzlibrary{arrows.meta,calc,positioning,patterns,decorations.markings,decorations.pathmorphing,cd}
\usepackage{pgfplots}
\pgfplotsset{compat=1.18}
\usepackage[most]{tcolorbox}
\tcbuselibrary{theorems,skins,breakable}
\usepackage[colorlinks=true,linkcolor=blue!60!black,citecolor=green!50!black]{hyperref}
\usepackage{makeidx}
\makeindex
\usepackage{enumitem}
\usepackage{fancyhdr}
\usepackage{caption}
\usepackage{booktabs}
\usepackage{microtype}
\usepackage{xcolor}

\pagestyle{fancy}
\fancyhf{}
\fancyhead[LE]{\slshape\leftmark}
\fancyhead[RO]{\slshape\rightmark}
\fancyfoot[C]{\thepage}

% ── Theorem environments (amsthm) ──
\theoremstyle{definition}
\newtheorem{definition}{Definition}[chapter]
\newtheorem{example}[definition]{Example}
\newtheorem{exercise}{Exercise}[chapter]

\theoremstyle{plain}
\newtheorem{theorem}[definition]{Theorem}
\newtheorem{proposition}[definition]{Proposition}
\newtheorem{lemma}[definition]{Lemma}
\newtheorem{corollary}[definition]{Corollary}

\theoremstyle{remark}
\newtheorem{remark}[definition]{Remark}

% ── tcolorbox wrapping (visual only) ──
\tcolorboxenvironment{definition}{enhanced,breakable,colback=blue!5,colframe=blue!70!black,fonttitle=\bfseries,before skip=10pt,after skip=10pt,left=8pt,right=8pt}
\tcolorboxenvironment{theorem}{enhanced,breakable,colback=red!5,colframe=red!70!black,fonttitle=\bfseries,before skip=10pt,after skip=10pt,left=8pt,right=8pt}
\tcolorboxenvironment{proposition}{enhanced,breakable,colback=orange!5,colframe=orange!70!black,fonttitle=\bfseries,before skip=10pt,after skip=10pt,left=8pt,right=8pt}
\tcolorboxenvironment{lemma}{enhanced,breakable,colback=yellow!5,colframe=yellow!50!black,fonttitle=\bfseries,before skip=10pt,after skip=10pt,left=8pt,right=8pt}
\tcolorboxenvironment{corollary}{enhanced,breakable,colback=green!5,colframe=green!50!black,fonttitle=\bfseries,before skip=10pt,after skip=10pt,left=8pt,right=8pt}
\tcolorboxenvironment{remark}{enhanced,breakable,colback=gray!5,colframe=gray!50!black,before skip=8pt,after skip=8pt,left=6pt,right=6pt}
\tcolorboxenvironment{example}{enhanced,breakable,colback=purple!5,colframe=purple!50!black,before skip=8pt,after skip=8pt,left=6pt,right=6pt}

% ── Custom commands ──
\newcommand{\C}{\mathbb{C}}
\newcommand{\R}{\mathbb{R}}
\newcommand{\Z}{\mathbb{Z}}
\newcommand{\N}{\mathbb{N}}
\newcommand{\Q}{\mathbb{Q}}
\newcommand{\D}{\mathbb{D}}
\newcommand{\T}{\mathbb{T}}
\DeclareMathOperator{\re}{Re}
\DeclareMathOperator{\im}{Im}
\DeclareMathOperator{\Arg}{Arg}
\DeclareMathOperator{\Log}{Log}
\DeclareMathOperator{\Res}{Res}
\DeclareMathOperator{\ind}{ind}
\newcommand{\dd}{\mathrm{d}}
\newcommand{\abs}[1]{\left|#1\right|}
\newcommand{\norm}[1]{\left\|#1\right\|}
\newcommand{\conj}[1]{\overline{#1}}
\DeclareMathOperator{\ord}{ord}
\DeclareMathOperator{\diam}{diam}

\title{\Huge\bfseries Complex Analysis \\ \Large Course Notes}
\author{L3 Mathematics}
\date{Academic Year 2025--2026}

\begin{document}

\frontmatter

% ══════════════════════════════════════════════
%  TITLE PAGE
% ══════════════════════════════════════════════
\begin{titlepage}
\begin{tikzpicture}[remember picture, overlay]
  \fill[left color=blue!15!white, right color=blue!5!white]
    (current page.south west) rectangle (current page.north east);
  \fill[color=orange!70!yellow]
    ([yshift=-2cm]current page.north west) rectangle ([yshift=-2.3cm]current page.north east);
  \fill[color=blue!80!black, rounded corners=8pt]
    ([xshift=3cm, yshift=-4cm]current page.north west) rectangle
    ([xshift=-3cm, yshift=-10cm]current page.north east);
  \node[anchor=center, text=white, font=\Huge\bfseries]
    at ([yshift=-6cm]current page.north) {Complex Analysis};
  \node[anchor=center, text=orange!80!yellow, font=\Large]
    at ([yshift=-7.5cm]current page.north) {Lecture Notes};
  \node[anchor=center, text=white!80, font=\large]
    at ([yshift=-8.8cm]current page.north) {Licence L3 --- 2025--2026};
  \node[anchor=center, text=blue!80!black, font=\Large\itshape]
    at ([yshift=-12cm]current page.north) {Ya\'e Ulrich Gaba};
  \draw[orange!70!yellow, line width=2pt]
    ([xshift=5cm, yshift=-13.5cm]current page.north west) --
    ([xshift=-5cm, yshift=-13.5cm]current page.north east);
  \node[anchor=center, text width=12cm, align=center, text=blue!60!black, font=\itshape]
    at ([yshift=-16cm]current page.north) {``The shortest path between two truths in the real domain passes through the complex domain.''\\[4pt]--- Jacques Hadamard};
  \node[anchor=south, text=gray, font=\small]
    at ([yshift=1.5cm]current page.south) {\today};
\end{tikzpicture}
\end{titlepage}

% ══════════════════════════════════════════════
%  PREFACE
% ══════════════════════════════════════════════
\chapter*{Preface}
\addcontentsline{toc}{chapter}{Preface}

This course is an introduction to the theory of functions of a single complex variable, aimed at third-year undergraduate students in mathematics. The subject is one of the great triumphs of nineteenth-century analysis, and its influence permeates nearly every branch of modern mathematics, from number theory and algebraic geometry to theoretical physics and engineering.

\medskip

\textbf{Why is complex differentiability so much richer than real differentiability?}

\medskip

In real analysis, a function $f:\R\to\R$ can be differentiable once without being differentiable twice; it can be infinitely differentiable without being analytic (i.e., without equalling its Taylor series). The landscape of real-differentiable functions is vast and somewhat wild.

In stark contrast, if a function $f:\C\to\C$ is complex-differentiable (holomorphic) on an open set, then it is automatically:
\begin{itemize}[leftmargin=2cm]
  \item \textbf{infinitely differentiable} --- all higher-order derivatives exist;
  \item \textbf{analytic} --- it equals its Taylor series in a neighbourhood of every point;
  \item \textbf{determined by its values on tiny sets} --- the identity principle tells us that if two holomorphic functions agree on a set with an accumulation point, they agree everywhere on their connected domain;
  \item \textbf{conformal} --- wherever $f'(z)\neq 0$, the mapping preserves angles between curves.
\end{itemize}

The source of this rigidity is the \emph{Cauchy--Riemann equations}: the requirement that the limit defining $f'(z_0)$ be the same regardless of the \emph{direction} from which $z$ approaches $z_0$ in the plane. In $\R$, there are only two directions (left and right); in $\C\simeq\R^2$, there are infinitely many. This single constraint binds the real and imaginary parts of $f$ into a system of partial differential equations that forces extraordinary regularity.

The crowning achievement is \textbf{Cauchy's integral formula}: the value of a holomorphic function at a point is completely determined by its values on any surrounding closed curve. From this single formula flow nearly all the major theorems of the subject: Liouville's theorem, the maximum modulus principle, the residue theorem, and the open mapping theorem.

\medskip

We have organised the material into the following chapters:
\begin{enumerate}
  \item \textbf{Complex Numbers} --- algebraic and geometric review, topology of the complex plane.
  \item \textbf{Holomorphic Functions} --- complex derivative, Cauchy--Riemann equations, conformality, harmonic functions.
  \item \textbf{Complex Integration} --- line integrals, Cauchy's theorem and integral formula.
  \item \textbf{Power Series and Analytic Functions} --- Taylor and Laurent series, isolated singularities.
  \item \textbf{The Residue Theorem and Applications} --- computation of real integrals, argument principle, Rouch\'e's theorem.
\end{enumerate}

Each chapter contains complete proofs, worked examples, and exercises graded by difficulty: one star~($\bigstar$) for routine practice, two stars~($\bigstar\!\bigstar$) for problems requiring synthesis, and three stars~($\bigstar\!\bigstar\!\bigstar$) for challenging or theoretical problems.

\medskip

\textbf{Prerequisites.} We assume familiarity with: linear algebra over $\R$; real analysis in one and several variables (limits, continuity, differentiation, Riemann integration); and basic point-set topology (open and closed sets, compactness, connectedness).

\vfill
\hfill\textit{The Author, 2025}

% ══════════════════════════════════════════════
%  NOTATION
% ══════════════════════════════════════════════
\chapter*{Notation}
\addcontentsline{toc}{chapter}{Notation}

\begin{tabular}{@{}ll@{}}
\toprule
\textbf{Symbol} & \textbf{Meaning} \\
\midrule
$\N,\;\Z,\;\Q,\;\R,\;\C$ & Natural numbers, integers, rationals, reals, complex numbers \\
$\re(z),\;\im(z)$ & Real and imaginary parts of $z\in\C$ \\
$\abs{z}$ & Modulus of $z$: $\abs{z}=\sqrt{\re(z)^2+\im(z)^2}$ \\
$\conj{z}$ & Complex conjugate of $z$: $\conj{x+iy}=x-iy$ \\
$\Arg(z)$ & Principal argument of $z$, taken in $(-\pi,\pi]$ \\
$\arg(z)$ & Multi-valued argument of $z$ \\
$\Log(z)$ & Principal branch of the logarithm: $\Log(z)=\ln\abs{z}+i\Arg(z)$ \\
$D(a,r)$ & Open disk of centre $a$ and radius $r$: $\{z\in\C:\abs{z-a}<r\}$ \\
$\conj{D}(a,r)$ & Closed disk: $\{z\in\C:\abs{z-a}\leq r\}$ \\
$C(a,r)$ & Circle of centre $a$ and radius $r$: $\{z\in\C:\abs{z-a}=r\}$ \\
$\D$ & Unit disk $D(0,1)$ \\
$\T$ & Unit circle $C(0,1)$ \\
$\C^*$ & $\C\setminus\{0\}$ \\
$\hat{\C}$ & Extended complex plane (Riemann sphere): $\C\cup\{\infty\}$ \\
$\mathscr{O}(\Omega)$ & Set of holomorphic functions on the open set $\Omega\subseteq\C$ \\
$f'(z_0)$ & Complex derivative of $f$ at $z_0$ \\
$\partial f/\partial z,\;\partial f/\partial\conj{z}$ & Wirtinger derivatives \\
$\dd z,\;\dd\conj{z}$ & Complex differentials \\
$\oint_\gamma f(z)\,\dd z$ & Complex line integral of $f$ along the curve $\gamma$ \\
$\ind(\gamma,a)$ & Winding number (index) of $\gamma$ about $a$ \\
$\Res(f,a)$ & Residue of $f$ at $a$ \\
$\Delta$ & Laplacian: $\Delta=\partial^2/\partial x^2+\partial^2/\partial y^2$ \\
\bottomrule
\end{tabular}

\tableofcontents

\mainmatter

% ══════════════════════════════════════════════════════════════════════
%  CHAPTER 1 — COMPLEX NUMBERS: REVIEW AND PLANE GEOMETRY
% ══════════════════════════════════════════════════════════════════════
\chapter{Complex Numbers --- Review and Plane Geometry}
\label{ch:complex-numbers}

\section{Algebraic Construction of \texorpdfstring{$\C$}{C}}
\label{sec:construction-C}

\begin{definition}[The field of complex numbers]
\label{def:complex-field}
The \emph{field of complex numbers} is the set $\C = \R^2 = \{(x,y) : x,y\in\R\}$ equipped with the operations
\[
  (x_1,y_1)+(x_2,y_2) = (x_1+x_2,\, y_1+y_2),
\]
\[
  (x_1,y_1)\cdot(x_2,y_2) = (x_1 x_2 - y_1 y_2,\; x_1 y_2 + y_1 x_2).
\]
We write $z = x + iy$ where $i = (0,1)$ satisfies $i^2 = -1$. The real number $x = \re(z)$ is the \emph{real part} and $y = \im(z)$ is the \emph{imaginary part} of $z$.
\end{definition}

\begin{proposition}[$\C$ is a field]
\label{prop:C-field}
$(\C,+,\cdot)$ is a commutative field. The additive identity is $0=(0,0)$, and the multiplicative identity is $1=(1,0)$. For $z=x+iy\neq 0$, the multiplicative inverse is
\[
  z^{-1} = \frac{x}{x^2+y^2} - i\,\frac{y}{x^2+y^2} = \frac{\conj{z}}{\abs{z}^2}.
\]
\end{proposition}

\begin{proof}
Commutativity and associativity of addition are inherited from $\R^2$. For multiplication, commutativity follows from the symmetry of the formula. Associativity is verified by direct computation: for $z_j = x_j+iy_j$,
\begin{align*}
(z_1 z_2)z_3 &= \bigl((x_1x_2-y_1y_2)+i(x_1y_2+y_1x_2)\bigr)(x_3+iy_3)\\
&= (x_1x_2x_3 - y_1y_2x_3 - x_1y_2y_3 - y_1x_2y_3) \\
&\quad + i(x_1x_2y_3 - y_1y_2y_3 + x_1y_2x_3 + y_1x_2x_3),
\end{align*}
which is symmetric under re-bracketing as $z_1(z_2 z_3)$. Distributivity is similarly verified. For the inverse, if $z = x+iy\neq 0$ then $\abs{z}^2 = x^2+y^2 > 0$ and
\[
  z \cdot \frac{\conj{z}}{\abs{z}^2} = \frac{z\conj{z}}{\abs{z}^2} = \frac{\abs{z}^2}{\abs{z}^2} = 1. \qedhere
\]
\end{proof}

\begin{remark}[The field $\C$ is not ordered]
There is no total order $\leq$ on $\C$ that is compatible with the field operations. Indeed, if such an order existed, then either $i > 0$ or $i < 0$. In the first case, $i^2 = -1 > 0$, a contradiction. In the second case, $(-i)>0$, so $(-i)^2 = -1 > 0$, the same contradiction.
\end{remark}

\begin{definition}[Conjugate]
\label{def:conjugate}
For $z = x+iy\in\C$, the \emph{complex conjugate} is $\conj{z} = x - iy$. The conjugation map $z\mapsto\conj{z}$ is the unique non-trivial field automorphism of $\C$ that fixes $\R$ pointwise.
\end{definition}

\begin{proposition}[Properties of conjugation]
\label{prop:conjugation-props}
For all $z,w\in\C$:
\begin{enumerate}[label=\textup{(\roman*)}]
  \item $\conj{z+w} = \conj{z}+\conj{w}$ and $\conj{zw} = \conj{z}\,\conj{w}$.
  \item $\conj{\conj{z}} = z$.
  \item $z\in\R$ if and only if $\conj{z}=z$; and $z$ is purely imaginary if and only if $\conj{z}=-z$.
  \item $\re(z) = \dfrac{z+\conj{z}}{2}$ and $\im(z) = \dfrac{z-\conj{z}}{2i}$.
\end{enumerate}
\end{proposition}

\begin{proof}
All statements follow by writing $z=x+iy$, $w=u+iv$ and computing directly. For instance, $\conj{zw} = \conj{(xu-yv)+i(xv+yu)} = (xu-yv)-i(xv+yu) = (x-iy)(u-iv) = \conj{z}\,\conj{w}$.
\end{proof}

\section{Modulus, Argument, and Exponential Form}
\label{sec:modulus-argument}

\begin{definition}[Modulus]
\label{def:modulus}
The \emph{modulus} of $z = x+iy$ is
\[
  \abs{z} = \sqrt{x^2 + y^2} = \sqrt{z\conj{z}}.
\]
Geometrically, $\abs{z}$ is the Euclidean distance from the origin to the point $(x,y)$ in $\R^2$.
\end{definition}

\begin{proposition}[Properties of the modulus]
\label{prop:modulus-props}
For all $z,w\in\C$:
\begin{enumerate}[label=\textup{(\roman*)}]
  \item $\abs{z}\geq 0$, with equality if and only if $z=0$.
  \item $\abs{zw} = \abs{z}\abs{w}$.
  \item $\abs{z+w}\leq\abs{z}+\abs{w}$ \textup{(triangle inequality)}.
  \item $\bigl|\abs{z}-\abs{w}\bigr| \leq \abs{z-w}$ \textup{(reverse triangle inequality)}.
  \item $\abs{\re(z)}\leq\abs{z}$ and $\abs{\im(z)}\leq\abs{z}$.
\end{enumerate}
\end{proposition}

\begin{proof}
(i) and (v) are immediate.

(ii) $\abs{zw}^2 = (zw)\conj{(zw)} = (zw)(\conj{z}\,\conj{w}) = (z\conj{z})(w\conj{w}) = \abs{z}^2\abs{w}^2$.

(iii) We compute:
\begin{align*}
\abs{z+w}^2 &= (z+w)\conj{(z+w)} = \abs{z}^2 + z\conj{w} + \conj{z}w + \abs{w}^2 \\
             &= \abs{z}^2 + 2\re(z\conj{w}) + \abs{w}^2 \\
             &\leq \abs{z}^2 + 2\abs{z\conj{w}} + \abs{w}^2 = \abs{z}^2 + 2\abs{z}\abs{w} + \abs{w}^2 \\
             &= (\abs{z}+\abs{w})^2.
\end{align*}

(iv) From (iii), $\abs{z} = \abs{(z-w)+w}\leq\abs{z-w}+\abs{w}$, so $\abs{z}-\abs{w}\leq\abs{z-w}$. Exchanging the roles of $z$ and $w$ gives $\abs{w}-\abs{z}\leq\abs{z-w}$, whence the result.
\end{proof}

\begin{definition}[Argument]
\label{def:argument}
For $z\in\C^*$, an \emph{argument} of $z$ is any real number $\theta$ such that
\[
  z = \abs{z}(\cos\theta + i\sin\theta).
\]
The set of all arguments of $z$ is $\arg(z) = \{\theta_0 + 2k\pi : k\in\Z\}$ for any particular argument $\theta_0$. The \emph{principal argument} $\Arg(z)\in(-\pi,\pi]$ is the unique argument in this interval:
\[
  \Arg(z) = \begin{cases}
    \arctan(y/x) & \text{if } x > 0,\\
    \arctan(y/x)+\pi & \text{if } x < 0,\; y \geq 0,\\
    \arctan(y/x)-\pi & \text{if } x < 0,\; y < 0,\\
    \pi/2 & \text{if } x = 0,\; y > 0,\\
    -\pi/2 & \text{if } x = 0,\; y < 0.
  \end{cases}
\]
\end{definition}

\begin{theorem}[Euler's formula]
\label{thm:euler}
For every $\theta\in\R$,
\[
  e^{i\theta} = \cos\theta + i\sin\theta.
\]
This can be taken as a definition (from the power series of $e^z$) or proved by showing that $f(\theta) = e^{-i\theta}(\cos\theta+i\sin\theta)$ satisfies $f'(\theta)=0$ and $f(0)=1$.
\end{theorem}

\begin{proof}
Consider the function $f:\R\to\C$ defined by $f(\theta) = e^{-i\theta}(\cos\theta+i\sin\theta)$. Differentiating:
\begin{align*}
f'(\theta) &= -ie^{-i\theta}(\cos\theta+i\sin\theta) + e^{-i\theta}(-\sin\theta+i\cos\theta) \\
            &= e^{-i\theta}\bigl[-i\cos\theta+\sin\theta-\sin\theta+i\cos\theta\bigr] = 0.
\end{align*}
Since $f$ is constant and $f(0) = 1\cdot(1+0) = 1$, we have $f(\theta)=1$ for all $\theta$, giving $e^{i\theta} = \cos\theta+i\sin\theta$.
\end{proof}

\begin{remark}
Using Euler's formula, every $z\in\C^*$ can be written in \emph{exponential form}:
\[
  z = r\,e^{i\theta}, \qquad r = \abs{z},\quad \theta\in\arg(z).
\]
Multiplication in exponential form is especially transparent: if $z_1 = r_1 e^{i\theta_1}$ and $z_2 = r_2 e^{i\theta_2}$, then
\[
  z_1 z_2 = r_1 r_2\, e^{i(\theta_1+\theta_2)}.
\]
Thus multiplication \emph{multiplies moduli and adds arguments}.
\end{remark}

\begin{example}[Computing powers]
\label{ex:powers}
Let $z = 1+i$. Then $\abs{z} = \sqrt{2}$ and $\Arg(z) = \pi/4$, so $z = \sqrt{2}\,e^{i\pi/4}$. Therefore
\[
  z^{10} = (\sqrt{2})^{10}\, e^{i\cdot 10\pi/4} = 2^5\, e^{i\cdot 5\pi/2} = 32\, e^{i\pi/2} = 32i.
\]
\end{example}

\begin{example}[De Moivre's formula]
For $n\in\Z$ and $\theta\in\R$:
\[
  (\cos\theta+i\sin\theta)^n = \cos(n\theta)+i\sin(n\theta).
\]
This follows immediately from $(e^{i\theta})^n = e^{in\theta}$. Taking $n=3$ and expanding the left side gives the triple-angle formulas:
\begin{align*}
  \cos 3\theta &= \cos^3\theta - 3\cos\theta\sin^2\theta = 4\cos^3\theta - 3\cos\theta,\\
  \sin 3\theta &= 3\cos^2\theta\sin\theta - \sin^3\theta = 3\sin\theta - 4\sin^3\theta.
\end{align*}
\end{example}

\section{Roots of Unity}
\label{sec:roots-unity}

\begin{theorem}[$n$-th roots of a complex number]
\label{thm:nth-roots}
Let $w = \rho\,e^{i\phi}\in\C^*$ and $n\geq 1$ be an integer. The equation $z^n = w$ has exactly $n$ distinct solutions in $\C$, given by
\[
  z_k = \rho^{1/n}\, \exp\!\left(i\,\frac{\phi + 2k\pi}{n}\right), \qquad k = 0, 1, \ldots, n-1.
\]
\end{theorem}

\begin{proof}
Write $z = r\,e^{i\theta}$. Then $z^n = w$ becomes $r^n e^{in\theta} = \rho\, e^{i\phi}$. Equating moduli: $r^n = \rho$, so $r = \rho^{1/n}$ (the unique positive real $n$-th root). Equating arguments: $n\theta = \phi + 2k\pi$ for some $k\in\Z$, giving $\theta = (\phi+2k\pi)/n$. For $k = 0,1,\ldots,n-1$, these give $n$ distinct values in $[0,2\pi)$; for other values of $k$, the angle differs by a multiple of $2\pi$ and so the same complex number is obtained.
\end{proof}

\begin{definition}[Roots of unity]
\label{def:roots-unity}
The \emph{$n$-th roots of unity} are the $n$ solutions of $z^n = 1$:
\[
  \omega_k = e^{2ik\pi/n}, \qquad k = 0, 1, \ldots, n-1.
\]
Setting $\omega = \omega_1 = e^{2i\pi/n}$, the roots of unity are $1, \omega, \omega^2, \ldots, \omega^{n-1}$. They form a cyclic group of order $n$ under multiplication, and they are the vertices of a regular $n$-gon inscribed in the unit circle.
\end{definition}

\begin{proposition}[Sum of roots of unity]
\label{prop:sum-roots}
For every integer $n\geq 2$,
\[
  \sum_{k=0}^{n-1}\omega^k = 0 \qquad\text{and}\qquad \prod_{k=0}^{n-1}\omega^k = (-1)^{n+1}.
\]
\end{proposition}

\begin{proof}
For the sum, since $\omega\neq 1$ we use the geometric series formula: $\sum_{k=0}^{n-1}\omega^k = (1-\omega^n)/(1-\omega) = 0$. For the product, $\prod_{k=0}^{n-1}\omega^k = \omega^{0+1+\cdots+(n-1)} = \omega^{n(n-1)/2} = e^{i\pi(n-1)} = (-1)^{n-1} = (-1)^{n+1}$.
\end{proof}

\begin{center}
\begin{tikzpicture}[scale=2.2]
  % Unit circle
  \draw[gray!40] (0,0) circle (1);
  \draw[gray!40,->] (-1.3,0)--(1.3,0) node[right]{$\re$};
  \draw[gray!40,->] (0,-1.3)--(0,1.3) node[above]{$\im$};

  % 6th roots of unity
  \foreach \k in {0,...,5} {
    \pgfmathsetmacro{\angle}{60*\k}
    \coordinate (R\k) at (\angle:1);
    \fill[blue!70!black] (R\k) circle (1.2pt);
  }
  % Polygon
  \draw[blue!70!black,thick] (R0)--(R1)--(R2)--(R3)--(R4)--(R5)--cycle;
  % Labels
  \node[right=2pt] at (R0) {$1$};
  \node[above right=1pt] at (R1) {$\omega$};
  \node[above left=1pt] at (R2) {$\omega^2$};
  \node[left=2pt] at (R3) {$\omega^3=-1$};
  \node[below left=1pt] at (R4) {$\omega^4$};
  \node[below right=1pt] at (R5) {$\omega^5$};

  \node[below] at (0,-1.5) {The $6$th roots of unity ($\omega = e^{i\pi/3}$)};
\end{tikzpicture}
\hspace{1.5cm}
\begin{tikzpicture}[scale=2.2]
  % Unit circle
  \draw[gray!40] (0,0) circle (1);
  \draw[gray!40,->] (-1.3,0)--(1.3,0) node[right]{$\re$};
  \draw[gray!40,->] (0,-1.3)--(0,1.3) node[above]{$\im$};

  % 5th roots of unity
  \foreach \k in {0,...,4} {
    \pgfmathsetmacro{\angle}{72*\k}
    \coordinate (S\k) at (\angle:1);
    \fill[red!70!black] (S\k) circle (1.2pt);
  }
  \draw[red!70!black,thick] (S0)--(S1)--(S2)--(S3)--(S4)--cycle;
  \node[right=2pt] at (S0) {$1$};
  \node[above right=1pt] at (S1) {$\zeta$};
  \node[above left=1pt] at (S2) {$\zeta^2$};
  \node[below left=1pt] at (S3) {$\zeta^3$};
  \node[below right=1pt] at (S4) {$\zeta^4$};

  \node[below] at (0,-1.5) {The $5$th roots of unity ($\zeta = e^{2i\pi/5}$)};
\end{tikzpicture}
\end{center}

\begin{example}[Cube roots of $-8i$]
\label{ex:cube-roots}
We wish to solve $z^3 = -8i$. First, write $-8i = 8\,e^{-i\pi/2}$ (since $\abs{-8i}=8$ and $\Arg(-8i)=-\pi/2$). Then
\[
  z_k = 8^{1/3}\exp\!\left(i\,\frac{-\pi/2+2k\pi}{3}\right) = 2\,\exp\!\left(i\,\frac{(4k-1)\pi}{6}\right), \quad k=0,1,2.
\]
Explicitly:
\begin{align*}
  z_0 &= 2e^{-i\pi/6} = 2\!\left(\frac{\sqrt{3}}{2} - \frac{i}{2}\right) = \sqrt{3}-i, \\
  z_1 &= 2e^{i\pi/2} = 2i, \\
  z_2 &= 2e^{i7\pi/6} = 2\!\left(-\frac{\sqrt{3}}{2} - \frac{i}{2}\right) = -\sqrt{3}-i.
\end{align*}
\end{example}

\section{Inversion and M\"obius Transformations}
\label{sec:mobius}

\begin{definition}[Complex inversion]
\label{def:inversion}
The \emph{inversion} map is $z\mapsto 1/z$, defined on $\C^*$. In exponential form, if $z=re^{i\theta}$ then $1/z = (1/r)\,e^{-i\theta}$: inversion inverts the modulus and negates the argument.
\end{definition}

\begin{proposition}[Inversion maps circles/lines to circles/lines]
\label{prop:inversion-circles}
The inversion $z\mapsto 1/z$ maps:
\begin{enumerate}[label=\textup{(\roman*)}]
  \item a circle not passing through the origin to a circle not passing through the origin;
  \item a circle passing through the origin to a line not passing through the origin;
  \item a line not passing through the origin to a circle passing through the origin;
  \item a line passing through the origin to itself \textup{(as a set)}.
\end{enumerate}
\end{proposition}

\begin{proof}
A general circle or line in $\C$ can be written
\[
  \alpha(x^2+y^2) + \beta x + \gamma y + \delta = 0, \qquad \alpha,\beta,\gamma,\delta\in\R,
\]
which in complex notation (using $x=(z+\conj{z})/2$, $y=(z-\conj{z})/(2i)$, $x^2+y^2=\abs{z}^2$) becomes
\[
  \alpha\abs{z}^2 + \re((\beta-i\gamma)z) + \delta = 0.
\]
Under the substitution $z\mapsto w=1/z$, i.e., $z=1/w$, and using $\abs{1/w}^2 = 1/\abs{w}^2$:
\[
  \alpha/\abs{w}^2 + \re((\beta-i\gamma)/w) + \delta = 0.
\]
Multiplying through by $\abs{w}^2$:
\[
  \alpha + \re((\beta-i\gamma)\conj{w}) + \delta\abs{w}^2 = 0.
\]
This is again a circle/line equation with parameters $(\delta,\beta,\gamma,\alpha)$. If $\alpha\neq 0$ and $\delta\neq 0$, both equations describe circles not through the origin. If $\alpha=0$ (the original is a line) then the image has $\delta$ in the role of the leading coefficient, giving a circle through the origin, etc.
\end{proof}

\begin{definition}[M\"obius transformation]
\label{def:mobius}
A \emph{M\"obius transformation} (or \emph{fractional linear transformation}) is a map $f:\hat{\C}\to\hat{\C}$ of the form
\[
  f(z) = \frac{az+b}{cz+d}, \qquad a,b,c,d\in\C,\quad ad-bc\neq 0.
\]
When $c\neq 0$, we set $f(-d/c)=\infty$ and $f(\infty)=a/c$. When $c=0$, we set $f(\infty)=\infty$. The condition $ad-bc\neq 0$ ensures that $f$ is not constant.
\end{definition}

\begin{proposition}[Basic properties of M\"obius transformations]
\label{prop:mobius-properties}
\leavevmode
\begin{enumerate}[label=\textup{(\roman*)}]
  \item Every M\"obius transformation is a bijection $\hat{\C}\to\hat{\C}$, with inverse
    \[
      f^{-1}(w) = \frac{dw-b}{-cw+a}.
    \]
  \item The composition of two M\"obius transformations is a M\"obius transformation, with the associated matrix being the product of the individual matrices:
    \[
      f \leftrightarrow \begin{pmatrix} a & b \\ c & d \end{pmatrix}, \qquad g\circ f \leftrightarrow \begin{pmatrix} a' & b' \\ c' & d' \end{pmatrix}\begin{pmatrix} a & b \\ c & d \end{pmatrix}.
    \]
  \item Every M\"obius transformation is a composition of translations $z\mapsto z+b$, dilations/rotations $z\mapsto az$ $(a\neq 0)$, and the inversion $z\mapsto 1/z$.
  \item M\"obius transformations map circles and lines in $\hat{\C}$ to circles and lines.
\end{enumerate}
\end{proposition}

\begin{proof}
(i) Direct verification: $f(f^{-1}(w))=w$.

(ii) If $f(z)=(az+b)/(cz+d)$ and $g(w)=(a'w+b')/(c'w+d')$, then
\[
  g(f(z)) = \frac{a'\frac{az+b}{cz+d}+b'}{c'\frac{az+b}{cz+d}+d'} = \frac{(a'a+b'c)z+(a'b+b'd)}{(c'a+d'c)z+(c'b+d'd)},
\]
which corresponds to the matrix product.

(iii) If $c=0$, then $f(z)=(a/d)z+b/d$ is an affine map (composition of dilation and translation). If $c\neq 0$, write
\[
  f(z) = \frac{a}{c} - \frac{ad-bc}{c^2}\cdot\frac{1}{z+d/c},
\]
which is a translation, then inversion, then dilation, then translation.

(iv) This follows from (iii), since translations and rotations obviously preserve circles/lines, and inversion does so by Proposition~\ref{prop:inversion-circles}.
\end{proof}

\begin{example}[Cayley transform]
\label{ex:cayley}
The M\"obius transformation
\[
  w = \frac{z-i}{z+i}
\]
maps the upper half-plane $\{z\in\C:\im(z)>0\}$ biholomorphically onto the unit disk $\D$. Indeed, for $z=x+iy$ with $y>0$:
\[
  \abs{w}^2 = \frac{\abs{z-i}^2}{\abs{z+i}^2} = \frac{x^2+(y-1)^2}{x^2+(y+1)^2} = \frac{x^2+y^2-2y+1}{x^2+y^2+2y+1} < 1,
\]
since $-2y<2y$ when $y>0$. The real axis maps to $\abs{w}=1$, i.e., the unit circle. The inverse (Cayley transform) is $z = i(1+w)/(1-w)$.
\end{example}

\begin{center}
\begin{tikzpicture}[scale=1.4]
  % Left: upper half-plane
  \fill[blue!8] (-2.5,0) rectangle (2.5,2.5);
  \draw[thick,->] (-2.7,0)--(2.7,0) node[right]{$\re$};
  \draw[thick,->] (0,-0.3)--(0,2.7) node[above]{$\im$};
  \draw[blue!60!black,thick] (-2.5,0)--(2.5,0);
  \node at (0,1.2) {$\im(z)>0$};
  \node[below] at (0,-0.6) {Upper half-plane};
  \fill[red] (0,1) circle (1.5pt) node[right=3pt]{$i$};

  % Arrow
  \draw[-{Stealth},thick] (3.2,1.2)--(4.5,1.2) node[midway,above]{$\dfrac{z-i}{z+i}$};

  % Right: unit disk
  \begin{scope}[xshift=7.2cm]
    \fill[green!8] (0,0) circle (1.5);
    \draw[thick,->] (-2,0)--(2,0) node[right]{$\re$};
    \draw[thick,->] (0,-2)--(0,2) node[above]{$\im$};
    \draw[green!50!black,thick] (0,0) circle (1.5);
    \node at (0,0) {$\D$};
    \node[below] at (0,-2.1) {Unit disk};
    \fill[red] (0,0) circle (1.5pt) node[above right=2pt]{$0$};
  \end{scope}
\end{tikzpicture}
\end{center}

\section{Topology of \texorpdfstring{$\C$}{C}}
\label{sec:topology-C}

Since $\C$ is identified with $\R^2$, it inherits the Euclidean metric $d(z,w) = \abs{z-w}$. We recall the essential topological notions.

\begin{definition}[Open and closed sets]
\label{def:open-closed}
Let $A\subseteq\C$.
\begin{itemize}
  \item $A$ is \emph{open} if for every $z\in A$, there exists $r>0$ such that $D(z,r)\subseteq A$.
  \item $A$ is \emph{closed} if $\C\setminus A$ is open, equivalently if $A$ contains all its limit points.
  \item The \emph{interior} of $A$ is $\mathring{A} = \bigcup\{U\subseteq A : U \text{ open}\}$, the largest open set contained in $A$.
  \item The \emph{closure} of $A$ is $\overline{A} = \bigcap\{F\supseteq A : F \text{ closed}\}$, the smallest closed set containing $A$.
  \item The \emph{boundary} of $A$ is $\partial A = \overline{A}\setminus\mathring{A}$.
\end{itemize}
\end{definition}

\begin{example}[Open disk and closed disk]
The open disk $D(a,r)=\{z:\abs{z-a}<r\}$ is open; the closed disk $\conj{D}(a,r)=\{z:\abs{z-a}\leq r\}$ is closed; the circle $C(a,r)=\{z:\abs{z-a}=r\}$ is the boundary of both.
\end{example}

\begin{definition}[Connected and path-connected sets]
\label{def:connected}
\leavevmode
\begin{itemize}
  \item $A\subseteq\C$ is \emph{connected} if it cannot be written as a disjoint union of two non-empty open subsets of $A$.
  \item $A$ is \emph{path-connected} if for every $z_0,z_1\in A$, there exists a continuous path $\gamma:[0,1]\to A$ with $\gamma(0)=z_0$ and $\gamma(1)=z_1$.
\end{itemize}
A non-empty connected open set in $\C$ is called a \emph{domain}.
\end{definition}

\begin{proposition}
\label{prop:path-connected}
In $\C$ (more generally in $\R^n$), a non-empty open set is connected if and only if it is path-connected.
\end{proposition}

\begin{proof}
Path-connected clearly implies connected (standard argument). For the converse, let $\Omega\subseteq\C$ be open and connected. Fix $z_0\in\Omega$ and let $U = \{z\in\Omega : z \text{ can be joined to } z_0 \text{ by a path in } \Omega\}$. Then $U$ is non-empty ($z_0\in U$). $U$ is open: if $z\in U$ then since $\Omega$ is open, there is a disk $D(z,r)\subseteq\Omega$; any point in this disk can be joined to $z$ by a line segment, hence to $z_0$ by concatenation. Similarly, $\Omega\setminus U$ is open: if $w\in\Omega\setminus U$ then a disk $D(w,r)\subseteq\Omega$ consists entirely of points not joinable to $z_0$ (otherwise $w$ would be joinable). Since $\Omega$ is connected and $U\neq\emptyset$, we must have $U = \Omega$.
\end{proof}

\begin{definition}[Simply connected domain]
\label{def:simply-connected}
A domain $\Omega\subseteq\C$ is \emph{simply connected} if every closed curve in $\Omega$ can be continuously deformed (within $\Omega$) to a point. Intuitively, $\Omega$ has ``no holes.'' Formally, for every continuous closed curve $\gamma:[0,1]\to\Omega$ with $\gamma(0)=\gamma(1)$, there exists a continuous map $H:[0,1]\times[0,1]\to\Omega$ with $H(t,0)=\gamma(t)$, $H(t,1)=\gamma(0)$ for all $t$, and $H(0,s)=H(1,s)$ for all $s$.
\end{definition}

\begin{example}
\leavevmode
\begin{enumerate}[label=\textup{(\roman*)}]
  \item $\C$, any open disk $D(a,r)$, any open half-plane, any convex open set: all simply connected.
  \item $\C^* = \C\setminus\{0\}$: connected but \emph{not} simply connected (the unit circle cannot be contracted to a point within $\C^*$).
  \item An annulus $\{z: r<\abs{z}<R\}$: connected but not simply connected.
\end{enumerate}
\end{example}

\begin{definition}[Compact sets]
\label{def:compact}
A set $K\subseteq\C$ is \emph{compact} if every open cover of $K$ has a finite subcover. By the Heine--Borel theorem, $K\subseteq\C$ is compact if and only if $K$ is closed and bounded.
\end{definition}

\subsection{The Extended Complex Plane and the Riemann Sphere}
\label{subsec:riemann-sphere}

\begin{definition}[Riemann sphere]
\label{def:riemann-sphere}
The \emph{extended complex plane} (or \emph{Riemann sphere}) is $\hat{\C} = \C\cup\{\infty\}$, where $\infty$ is a single additional point called the \emph{point at infinity}. We define a topology on $\hat{\C}$ by declaring the open sets to be:
\begin{itemize}
  \item all open subsets of $\C$, and
  \item sets of the form $\{z\in\C:\abs{z}>R\}\cup\{\infty\}$ for $R>0$ (neighbourhoods of $\infty$).
\end{itemize}
The space $\hat{\C}$ is homeomorphic to the $2$-sphere $S^2\subset\R^3$ via stereographic projection.
\end{definition}

\begin{proposition}[Stereographic projection]
\label{prop:stereographic}
Consider the unit sphere $S^2 = \{(X,Y,Z)\in\R^3 : X^2+Y^2+Z^2=1\}$ and identify $\C$ with the equatorial plane $\{Z=0\}$. The stereographic projection from the north pole $N=(0,0,1)$ is the map $\pi:S^2\setminus\{N\}\to\C$ given by
\[
  \pi(X,Y,Z) = \frac{X+iY}{1-Z}.
\]
Its inverse is
\[
  \pi^{-1}(z) = \left(\frac{2\re(z)}{\abs{z}^2+1},\; \frac{2\im(z)}{\abs{z}^2+1},\; \frac{\abs{z}^2-1}{\abs{z}^2+1}\right).
\]
Setting $\pi(N) = \infty$ extends $\pi$ to a homeomorphism $S^2\to\hat{\C}$.
\end{proposition}

\begin{proof}
A point $(X,Y,Z)\in S^2\setminus\{N\}$ and $N=(0,0,1)$ determine a line in $\R^3$:
\[
  \ell(t) = (1-t)N + t(X,Y,Z) = (tX,\; tY,\; 1-t+tZ).
\]
This intersects the plane $Z=0$ when $1-t+tZ=0$, i.e., $t=1/(1-Z)$. The intersection point is $(X/(1-Z),\; Y/(1-Z),\; 0)$, corresponding to $z=\frac{X+iY}{1-Z}$.

For the inverse, set $z=x+iy$ and solve: from $x=X/(1-Z)$, $y=Y/(1-Z)$ and $X^2+Y^2+Z^2=1$, we get $x^2+y^2 = (1-Z^2)/(1-Z)^2 = (1+Z)/(1-Z)$. Then $\abs{z}^2(1-Z) = 1+Z$, giving $Z = (\abs{z}^2-1)/(\abs{z}^2+1)$ and $1-Z = 2/(\abs{z}^2+1)$, whence $X = 2x/(\abs{z}^2+1)$ and $Y = 2y/(\abs{z}^2+1)$.
\end{proof}

\begin{center}
\begin{tikzpicture}[scale=2.0]
  % Sphere outline
  \draw[thick] (0,0) circle (1);
  % Equator (ellipse for 3d effect)
  \draw[gray,dashed] (0,0) ellipse (1 and 0.35);
  % North pole
  \fill[red] (0,1) circle (1.5pt) node[above right]{$N=(0,0,1)$};
  % A point on the sphere
  \coordinate (P) at (-0.5,0.3);
  \fill[blue!70!black] (P) circle (1.5pt) node[left=3pt]{$(X,Y,Z)$};
  % Projection line from N through P to equatorial plane
  \pgfmathsetmacro{\tval}{1/0.7}
  \pgfmathsetmacro{\projx}{-0.5*\tval}
  \draw[dashed,thick,red!60!black] (0,1)--(P)--(\projx,-0.35);
  \fill[green!50!black] (\projx,-0.35) circle (1.5pt) node[below left]{$z=\pi(X,Y,Z)$};
  % Equatorial plane line
  \draw[green!50!black,thick] (-1.5,-0.35)--(1.5,-0.35) node[right]{\small$\C$ (equatorial plane)};
\end{tikzpicture}
\end{center}

\begin{remark}[Chordal metric]
\label{rem:chordal}
The \emph{chordal metric} on $\hat{\C}$ is defined by
\[
  \chi(z,w) = \frac{2\abs{z-w}}{\sqrt{1+\abs{z}^2}\sqrt{1+\abs{w}^2}}, \qquad \chi(z,\infty) = \frac{2}{\sqrt{1+\abs{z}^2}}.
\]
This is the Euclidean distance between the stereographic images on $S^2$. The chordal metric makes $\hat{\C}$ into a compact metric space.
\end{remark}

\section{Exercises for Chapter 1}
\label{sec:exercises-ch1}

\begin{exercise}[$\bigstar$]
\label{exo:ch1-1}
Let $z=3+4i$ and $w=1-2i$. Compute $z+w$, $zw$, $z/w$, $\abs{z}$, $\abs{w}$, $\conj{z}w$, and $\Arg(z)$.
\end{exercise}

\begin{exercise}[$\bigstar$]
\label{exo:ch1-2}
Show that for all $z,w\in\C$,
\[
  \abs{z+w}^2 + \abs{z-w}^2 = 2(\abs{z}^2+\abs{w}^2) \qquad\text{(parallelogram law)}.
\]
Give a geometric interpretation.
\end{exercise}

\begin{exercise}[$\bigstar$]
\label{exo:ch1-3}
Find all $z\in\C$ such that $z^4 = -16$. Express the solutions in both exponential and Cartesian form, and plot them.
\end{exercise}

\begin{exercise}[$\bigstar$]
\label{exo:ch1-4}
Let $\omega = e^{2i\pi/7}$. Compute $\displaystyle\sum_{k=0}^{6}\omega^{3k}$.
\end{exercise}

\begin{exercise}[$\bigstar\!\bigstar$]
\label{exo:ch1-5}
Prove that for any $n\geq 1$ and $\theta\in\R$ with $\theta\neq 2k\pi/n$ for any $k\in\Z$:
\[
  \sum_{k=0}^{n-1}\cos(k\theta) = \frac{\sin(n\theta/2)}{\sin(\theta/2)}\cos\!\left(\frac{(n-1)\theta}{2}\right).
\]
\textit{Hint:} Compute $\re\!\left(\sum_{k=0}^{n-1}e^{ik\theta}\right)$.
\end{exercise}

\begin{exercise}[$\bigstar\!\bigstar$]
\label{exo:ch1-6}
Show that the M\"obius transformation $f(z) = (z-1)/(z+1)$ maps the right half-plane $\{\re(z)>0\}$ onto the unit disk $\D$. What is the image of the imaginary axis?
\end{exercise}

\begin{exercise}[$\bigstar\!\bigstar$]
\label{exo:ch1-7}
Let $a\in\D$. Show that the \emph{Blaschke factor}
\[
  \varphi_a(z) = \frac{z-a}{1-\conj{a}z}
\]
maps $\D$ bijectively onto $\D$ and $\T$ onto $\T$. Compute $\varphi_a^{-1}$.
\end{exercise}

\begin{exercise}[$\bigstar\!\bigstar$]
\label{exo:ch1-8}
Let $f(z) = (az+b)/(cz+d)$ with $ad-bc\neq 0$. Show that $f$ has exactly $1$ or $2$ fixed points in $\hat{\C}$ (i.e., solutions of $f(z)=z$). When does it have exactly one?
\end{exercise}

\begin{exercise}[$\bigstar\!\bigstar\!\bigstar$]
\label{exo:ch1-9}
Show that the group of M\"obius transformations preserving the unit disk $\D$ is
\[
  \mathrm{Aut}(\D) = \left\{z\mapsto e^{i\theta}\frac{z-a}{1-\conj{a}z} : \theta\in\R,\; a\in\D\right\}.
\]
\textit{Hint:} If $f$ is such a M\"obius transformation, let $a=f^{-1}(0)$ and consider $\varphi_a\circ f^{-1}$.
\end{exercise}

\begin{exercise}[$\bigstar\!\bigstar\!\bigstar$]
\label{exo:ch1-10}
\textbf{(Cross-ratio.)} For four distinct points $z_1,z_2,z_3,z_4\in\hat{\C}$, define the \emph{cross-ratio}
\[
  (z_1,z_2;z_3,z_4) = \frac{(z_1-z_3)(z_2-z_4)}{(z_1-z_4)(z_2-z_3)}.
\]
\begin{enumerate}[label=\textup{(\alph*)}]
  \item Show that the cross-ratio is invariant under M\"obius transformations: $(f(z_1),f(z_2);f(z_3),f(z_4))=(z_1,z_2;z_3,z_4)$.
  \item Show that four distinct points $z_1,z_2,z_3,z_4\in\hat{\C}$ lie on a circle or line if and only if $(z_1,z_2;z_3,z_4)\in\R$.
  \item Given three distinct points $z_2,z_3,z_4\in\hat{\C}$, show that $z\mapsto (z,z_2;z_3,z_4)$ is the unique M\"obius transformation sending $z_2\mapsto 1$, $z_3\mapsto 0$, $z_4\mapsto\infty$.
\end{enumerate}
\end{exercise}

% ══════════════════════════════════════════════════════════════════════
%  CHAPTER 2 — HOLOMORPHIC FUNCTIONS: CAUCHY–RIEMANN CONDITIONS
% ══════════════════════════════════════════════════════════════════════
\chapter{Holomorphic Functions --- Cauchy--Riemann Conditions}
\label{ch:holomorphic}

In this chapter we introduce the central concept of the course: \emph{holomorphic} (complex-differentiable) functions. We shall see that a single requirement --- that the complex derivative exist --- imposes far stronger constraints than real differentiability. The key link is the system of \emph{Cauchy--Riemann equations}, which connect the partial derivatives of the real and imaginary parts.

\section{The Complex Derivative}
\label{sec:complex-derivative}

\begin{definition}[Holomorphic function]
\label{def:holomorphic}
Let $\Omega\subseteq\C$ be open and $f:\Omega\to\C$. We say that $f$ is \emph{complex-differentiable} (or \emph{holomorphic}) at a point $z_0\in\Omega$ if the limit
\[
  f'(z_0) = \lim_{h\to 0}\frac{f(z_0+h)-f(z_0)}{h}
\]
exists in $\C$. Here $h\in\C^*$ and $h\to 0$ means $\abs{h}\to 0$ with $h$ approaching $0$ from \emph{any direction} in the plane. The number $f'(z_0)$ is called the \emph{complex derivative} of $f$ at $z_0$.

We say that $f$ is \emph{holomorphic on $\Omega$} if it is holomorphic at every point of $\Omega$. We write $f\in\mathscr{O}(\Omega)$.
\end{definition}

\begin{remark}[The crucial difference from $\R$]
\label{rem:crucial-difference}
In the real case, the variable $h$ in $f'(x_0) = \lim_{h\to 0}(f(x_0+h)-f(x_0))/h$ can only approach $0$ from the left or from the right --- two directions. In the complex case, $h$ can approach $0$ from \emph{any direction} in $\R^2$: along the real axis, the imaginary axis, spirals, or any other path. The requirement that the same limit be obtained for \emph{all} directions is extremely restrictive and is the source of the remarkable rigidity of holomorphic functions.
\end{remark}

\begin{proposition}[Holomorphic implies continuous]
\label{prop:holo-continuous}
If $f$ is holomorphic at $z_0$, then $f$ is continuous at $z_0$.
\end{proposition}

\begin{proof}
If $f'(z_0)$ exists, then for $h\neq 0$:
\[
  f(z_0+h) - f(z_0) = \frac{f(z_0+h)-f(z_0)}{h}\cdot h \to f'(z_0)\cdot 0 = 0 \quad\text{as } h\to 0. \qedhere
\]
\end{proof}

\begin{proposition}[Algebraic rules for the complex derivative]
\label{prop:derivative-rules}
Let $f,g\in\mathscr{O}(\Omega)$ and $\alpha\in\C$. Then:
\begin{enumerate}[label=\textup{(\roman*)}]
  \item $(\alpha f+g)'(z) = \alpha f'(z)+g'(z)$.
  \item $(fg)'(z) = f'(z)g(z)+f(z)g'(z)$ \textup{(product rule)}.
  \item If $g(z)\neq 0$, then $(f/g)'(z) = \bigl(f'(z)g(z)-f(z)g'(z)\bigr)/g(z)^2$ \textup{(quotient rule)}.
  \item If $f:\Omega\to\Omega'$ and $g:\Omega'\to\C$ are holomorphic, then $(g\circ f)'(z) = g'(f(z))\cdot f'(z)$ \textup{(chain rule)}.
\end{enumerate}
\end{proposition}

\begin{proof}
The proofs are identical to the real case. We give the chain rule as an example. Let $w_0 = f(z_0)$ and define
\[
  \varphi(w) = \begin{cases} \dfrac{g(w)-g(w_0)}{w-w_0} & \text{if } w\neq w_0, \\ g'(w_0) & \text{if } w = w_0. \end{cases}
\]
Then $\varphi$ is continuous at $w_0$ (since $g$ is holomorphic at $w_0$) and $g(w)-g(w_0) = \varphi(w)(w-w_0)$ for all $w$. Setting $w = f(z)$:
\[
  \frac{g(f(z))-g(f(z_0))}{z-z_0} = \varphi(f(z))\cdot\frac{f(z)-f(z_0)}{z-z_0}.
\]
As $z\to z_0$, $f(z)\to f(z_0)=w_0$ (by continuity of $f$), so $\varphi(f(z))\to\varphi(w_0)=g'(w_0)$, and the second factor tends to $f'(z_0)$.
\end{proof}

\begin{example}[Polynomials and rational functions]
\label{ex:polynomials}
\leavevmode
\begin{enumerate}[label=\textup{(\roman*)}]
  \item \textbf{Constant functions:} $f(z)=c$ has $f'(z)=0$ everywhere.
  \item \textbf{The identity:} $f(z)=z$ has $f'(z)=1$.
  \item \textbf{Powers:} $f(z)=z^n$ ($n\in\N$) has $f'(z)=nz^{n-1}$ (by induction using the product rule).
  \item \textbf{Polynomials:} $p(z) = \sum_{k=0}^n a_k z^k$ is holomorphic on $\C$ with $p'(z) = \sum_{k=1}^n k a_k z^{k-1}$.
  \item \textbf{Rational functions:} $f(z)=p(z)/q(z)$ is holomorphic on $\C\setminus\{z:q(z)=0\}$.
\end{enumerate}
\end{example}

\begin{example}[The exponential function]
\label{ex:exponential}
Define $e^z = e^x(\cos y + i\sin y)$ for $z = x+iy$. Then
\[
  \frac{e^{z+h}-e^z}{h} = e^z\cdot\frac{e^h-1}{h}.
\]
We claim $\lim_{h\to 0}(e^h-1)/h = 1$. Indeed, writing $h = s+it$ with $\abs{h}$ small:
\[
  e^h - 1 = e^s\cos t - 1 + ie^s\sin t = (e^s-1)\cos t + (\cos t-1) + ie^s\sin t.
\]
Using $e^s = 1+s+O(s^2)$, $\cos t = 1 - t^2/2 + O(t^4)$, $\sin t = t+O(t^3)$:
\[
  e^h-1 = s + it + O(\abs{h}^2) = h + O(\abs{h}^2),
\]
so $(e^h-1)/h \to 1$. Therefore $(e^z)' = e^z$ for all $z\in\C$.
\end{example}

\begin{example}[Conjugation is not holomorphic]
\label{ex:conjugation-not-holo}
The function $f(z) = \conj{z}$ is \emph{not} holomorphic anywhere. Indeed,
\[
  \frac{f(z+h)-f(z)}{h} = \frac{\conj{h}}{h}.
\]
If $h = t\in\R^*$, this equals $t/t = 1$. If $h = it$ with $t\in\R^*$, this equals $-it/(it) = -1$. Since the limit along the real axis gives $1$ and the limit along the imaginary axis gives $-1$, the limit does not exist.
\end{example}

\begin{example}[The modulus squared]
\label{ex:modulus-squared}
The function $f(z) = \abs{z}^2 = z\conj{z}$ is holomorphic \emph{only at the origin}. Indeed,
\[
  \frac{f(z+h)-f(z)}{h} = \frac{(z+h)\conj{(z+h)}-z\conj{z}}{h} = \conj{z} + z\frac{\conj{h}}{h} + \conj{h}.
\]
For this to have a limit as $h\to 0$, the term $z\cdot\conj{h}/h$ must have a limit. By the same argument as above, $\conj{h}/h$ has no limit unless $z=0$. When $z=0$, the expression becomes $\conj{h}$, which tends to $0$. So $f'(0)=0$ and $f$ is not holomorphic at any $z\neq 0$.
\end{example}

\section{The Cauchy--Riemann Equations}
\label{sec:cauchy-riemann}

The following theorem provides an analytic criterion for holomorphy in terms of the real and imaginary parts of $f$.

Write $f = u + iv$ where $u,v:\Omega\to\R$ and identify $z = x+iy$.

\begin{theorem}[Cauchy--Riemann equations --- necessary condition]
\label{thm:CR-necessary}
If $f = u+iv$ is holomorphic at $z_0 = x_0+iy_0$, then the partial derivatives $u_x$, $u_y$, $v_x$, $v_y$ all exist at $(x_0,y_0)$ and satisfy the \emph{Cauchy--Riemann equations}:
\[
  \boxed{\frac{\partial u}{\partial x} = \frac{\partial v}{\partial y}, \qquad \frac{\partial u}{\partial y} = -\frac{\partial v}{\partial x}.}
\]
Moreover, the complex derivative is given by
\[
  f'(z_0) = u_x(x_0,y_0) + iv_x(x_0,y_0) = v_y(x_0,y_0) - iu_y(x_0,y_0).
\]
\end{theorem}

\begin{proof}
Since $f'(z_0)$ exists, the limit $\lim_{h\to 0}(f(z_0+h)-f(z_0))/h$ is the same for all directions of approach.

\textbf{Approach along the real axis:} Set $h = t\in\R$, $t\to 0$:
\begin{align*}
  f'(z_0) &= \lim_{t\to 0}\frac{f(x_0+t+iy_0)-f(x_0+iy_0)}{t} \\
           &= \lim_{t\to 0}\frac{u(x_0+t,y_0)-u(x_0,y_0)}{t} + i\lim_{t\to 0}\frac{v(x_0+t,y_0)-v(x_0,y_0)}{t} \\
           &= u_x(x_0,y_0) + iv_x(x_0,y_0).
\end{align*}

\textbf{Approach along the imaginary axis:} Set $h = it$, $t\in\R$, $t\to 0$:
\begin{align*}
  f'(z_0) &= \lim_{t\to 0}\frac{f(x_0+i(y_0+t))-f(x_0+iy_0)}{it} \\
           &= \frac{1}{i}\left[u_y(x_0,y_0)+iv_y(x_0,y_0)\right] \\
           &= -iu_y(x_0,y_0) + v_y(x_0,y_0).
\end{align*}

Equating the two expressions for $f'(z_0)$:
\[
  u_x + iv_x = v_y - iu_y.
\]
Equating real and imaginary parts: $u_x = v_y$ and $v_x = -u_y$.
\end{proof}

\begin{remark}[The Cauchy--Riemann equations are not sufficient without regularity]
\label{rem:CR-not-sufficient}
The Cauchy--Riemann equations at a point $(x_0,y_0)$ are necessary but not sufficient for holomorphy. Consider
\[
  f(z) = \begin{cases} \exp(-1/z^4) & \text{if } z\neq 0, \\ 0 & \text{if } z=0. \end{cases}
\]
One can verify that $u_x = v_y = u_y = -v_x = 0$ at the origin, yet $f$ is not even continuous at $0$ (the limit depends on the direction of approach). The additional requirement is that the partial derivatives be \emph{continuous}.
\end{remark}

\begin{theorem}[Cauchy--Riemann equations --- sufficient condition]
\label{thm:CR-sufficient}
Let $\Omega\subseteq\C$ be open and $f = u+iv:\Omega\to\C$. Suppose that the partial derivatives $u_x$, $u_y$, $v_x$, $v_y$ exist in $\Omega$ and are \emph{continuous} at $z_0 = x_0+iy_0$, and that the Cauchy--Riemann equations hold at $(x_0,y_0)$:
\[
  u_x(x_0,y_0) = v_y(x_0,y_0), \qquad u_y(x_0,y_0) = -v_x(x_0,y_0).
\]
Then $f$ is holomorphic at $z_0$, with $f'(z_0) = u_x(x_0,y_0)+iv_x(x_0,y_0)$.
\end{theorem}

\begin{proof}
Since $u$ has continuous partial derivatives at $(x_0,y_0)$, it is (real) differentiable there:
\[
  u(x_0+s,y_0+t) = u(x_0,y_0) + u_x s + u_y t + \varepsilon_1(s,t)\sqrt{s^2+t^2},
\]
where $\varepsilon_1(s,t)\to 0$ as $(s,t)\to(0,0)$. Similarly for $v$:
\[
  v(x_0+s,y_0+t) = v(x_0,y_0) + v_x s + v_y t + \varepsilon_2(s,t)\sqrt{s^2+t^2}.
\]
Setting $h = s+it$ so that $\abs{h} = \sqrt{s^2+t^2}$:
\begin{align*}
  f(z_0+h)-f(z_0) &= (u_x s + u_y t + \varepsilon_1\abs{h}) + i(v_x s + v_y t + \varepsilon_2\abs{h}).
\end{align*}
Using the Cauchy--Riemann equations ($v_y = u_x$ and $v_x = -u_y$):
\begin{align*}
  f(z_0+h)-f(z_0) &= u_x s - v_x t + i(v_x s + u_x t) + (\varepsilon_1+i\varepsilon_2)\abs{h} \\
                    &= (u_x+iv_x)(s+it) + (\varepsilon_1+i\varepsilon_2)\abs{h} \\
                    &= (u_x+iv_x)\,h + (\varepsilon_1+i\varepsilon_2)\abs{h}.
\end{align*}
Dividing by $h$:
\[
  \frac{f(z_0+h)-f(z_0)}{h} = u_x+iv_x + (\varepsilon_1+i\varepsilon_2)\frac{\abs{h}}{h}.
\]
Since $\abs{\abs{h}/h} = 1$ and $\varepsilon_1,\varepsilon_2\to 0$, the last term vanishes as $h\to 0$. Therefore $f'(z_0) = u_x(x_0,y_0)+iv_x(x_0,y_0)$.
\end{proof}

\begin{corollary}
\label{cor:C1-holo}
If $f=u+iv:\Omega\to\C$ is of class $C^1$ (i.e., $u_x,u_y,v_x,v_y$ exist and are continuous on $\Omega$), then $f$ is holomorphic on $\Omega$ if and only if the Cauchy--Riemann equations hold at every point of $\Omega$.
\end{corollary}

\begin{example}[Verification for $e^z$]
\label{ex:CR-exp}
For $f(z) = e^z = e^x\cos y + ie^x\sin y$, we have $u(x,y)=e^x\cos y$ and $v(x,y)=e^x\sin y$. Then:
\[
  u_x = e^x\cos y = v_y, \qquad u_y = -e^x\sin y = -v_x.
\]
The Cauchy--Riemann equations hold everywhere, with continuous partial derivatives. Therefore $f$ is holomorphic on $\C$, with $f'(z) = u_x+iv_x = e^x\cos y + ie^x\sin y = e^z$.
\end{example}

\begin{example}[Verification for $\conj{z}$ --- failure of Cauchy--Riemann]
\label{ex:CR-conjugate}
For $f(z)=\conj{z} = x-iy$, we have $u=x$, $v=-y$. Then $u_x=1$ but $v_y=-1$, so $u_x\neq v_y$. The Cauchy--Riemann equations fail everywhere, confirming that $\conj{z}$ is nowhere holomorphic.
\end{example}

\begin{definition}[Wirtinger derivatives]
\label{def:wirtinger}
The \emph{Wirtinger operators} are defined by
\[
  \frac{\partial}{\partial z} = \frac{1}{2}\left(\frac{\partial}{\partial x}-i\frac{\partial}{\partial y}\right), \qquad
  \frac{\partial}{\partial\conj{z}} = \frac{1}{2}\left(\frac{\partial}{\partial x}+i\frac{\partial}{\partial y}\right).
\]
\end{definition}

\begin{proposition}[Holomorphy via Wirtinger derivatives]
\label{prop:wirtinger-holo}
A $C^1$ function $f:\Omega\to\C$ is holomorphic if and only if
\[
  \frac{\partial f}{\partial\conj{z}} = 0 \quad\text{on }\Omega.
\]
In this case, $f'(z) = \dfrac{\partial f}{\partial z}$.
\end{proposition}

\begin{proof}
Write $f = u+iv$. Then
\begin{align*}
  \frac{\partial f}{\partial\conj{z}} &= \frac{1}{2}\left(\frac{\partial}{\partial x}+i\frac{\partial}{\partial y}\right)(u+iv) \\
  &= \frac{1}{2}\bigl[(u_x-v_y)+i(v_x+u_y)\bigr].
\end{align*}
This vanishes if and only if $u_x = v_y$ and $u_y = -v_x$, which are the Cauchy--Riemann equations.
\end{proof}

\begin{remark}
The condition $\partial f/\partial\conj{z} = 0$ has an elegant interpretation: a holomorphic function is one that ``depends on $z$ but not on $\conj{z}$.'' This perspective is fundamental in several complex variables and in complex differential geometry.
\end{remark}

\section{Geometric Interpretation: Conformality}
\label{sec:conformality}

One of the most striking properties of holomorphic functions is that they preserve angles. This makes them fundamental in applications ranging from fluid dynamics to cartography.

\begin{definition}[Angle between curves]
\label{def:angle-curves}
Let $\gamma_1,\gamma_2:[0,1]\to\C$ be two smooth curves with $\gamma_1(0)=\gamma_2(0)=z_0$ and $\gamma_1'(0)\neq 0$, $\gamma_2'(0)\neq 0$. The \emph{angle} between $\gamma_1$ and $\gamma_2$ at $z_0$ is
\[
  \angle(\gamma_1,\gamma_2) = \arg\!\left(\frac{\gamma_2'(0)}{\gamma_1'(0)}\right) \pmod{2\pi}.
\]
This is the signed angle from the tangent of $\gamma_1$ to the tangent of $\gamma_2$ at $z_0$.
\end{definition}

\begin{theorem}[Holomorphic maps are conformal]
\label{thm:conformal}
Let $f:\Omega\to\C$ be holomorphic with $f'(z_0)\neq 0$. Then $f$ is \emph{conformal} at $z_0$: it preserves angles between curves and preserves orientation. More precisely, if $\gamma_1,\gamma_2$ are smooth curves through $z_0$, then the image curves $f\circ\gamma_1$ and $f\circ\gamma_2$ meet at $f(z_0)$ with the same angle:
\[
  \angle(f\circ\gamma_1,\, f\circ\gamma_2) = \angle(\gamma_1,\gamma_2).
\]
\end{theorem}

\begin{proof}
By the chain rule, $(f\circ\gamma_j)'(0) = f'(z_0)\cdot\gamma_j'(0)$ for $j=1,2$. Therefore
\[
  \frac{(f\circ\gamma_2)'(0)}{(f\circ\gamma_1)'(0)} = \frac{f'(z_0)\gamma_2'(0)}{f'(z_0)\gamma_1'(0)} = \frac{\gamma_2'(0)}{\gamma_1'(0)}.
\]
Taking arguments: $\angle(f\circ\gamma_1,f\circ\gamma_2) = \arg(\gamma_2'(0)/\gamma_1'(0)) = \angle(\gamma_1,\gamma_2)$.
\end{proof}

\begin{remark}
The condition $f'(z_0)\neq 0$ is essential. At a zero of $f'$ of order $k$ (i.e., $f(z)=f(z_0)+c(z-z_0)^{k+1}+\cdots$ with $c\neq 0$), angles are multiplied by $k+1$. For instance, $f(z)=z^2$ doubles angles at the origin.
\end{remark}

\begin{proposition}[Conformal $\Leftrightarrow$ holomorphic or anti-holomorphic]
\label{prop:conformal-iff}
Let $f:\Omega\to\C$ be a $C^1$ map with non-vanishing Jacobian. Then $f$ is angle-preserving (conformal) at every point if and only if $f$ is holomorphic with $f'\neq 0$ everywhere. If instead $f$ reverses angles (anti-conformal), then $\conj{f}$ is holomorphic.
\end{proposition}

\begin{proof}
The real Jacobian matrix of $f = u+iv$ is
\[
  J_f = \begin{pmatrix} u_x & u_y \\ v_x & v_y \end{pmatrix}.
\]
Conformality means $J_f$ is a positive scalar multiple of a rotation matrix. A $2\times 2$ matrix $\begin{pmatrix}a & b \\ c & d\end{pmatrix}$ is a positive scalar times a rotation if and only if $a=d$ and $c=-b$ (with $a^2+b^2>0$). This gives $u_x=v_y$ and $u_y=-v_x$, which are the Cauchy--Riemann equations. Anti-conformality similarly gives $u_x=-v_y$ and $u_y=v_x$, the Cauchy--Riemann equations for $\conj{f}$.
\end{proof}

\begin{example}[Conformal image of a grid under $z\mapsto z^2$]
\label{ex:conformal-z2}
The map $f(z)=z^2$ is conformal away from $z=0$ (where $f'(0)=0$). Let us see what happens to the standard grid of horizontal and vertical lines.

A horizontal line $z=t+ib$ ($t\in\R$) maps to $f(z)=(t+ib)^2 = t^2-b^2+2ibt$, i.e., the parametric curve $x=t^2-b^2$, $y=2bt$. Eliminating $t$: $t=y/(2b)$, so $x = y^2/(4b^2)-b^2$, a parabola opening to the right.

A vertical line $z=a+it$ maps to $(a+it)^2 = a^2-t^2+2iat$, giving $x=a^2-t^2$, $y=2at$, so $x=a^2-y^2/(4a^2)$, a parabola opening to the left.

These two families of parabolas intersect at right angles, illustrating the conformal property.
\end{example}

\begin{center}
\begin{tikzpicture}[scale=0.85]
  % Left: original grid
  \begin{scope}
    \draw[thick,->] (-2.5,0)--(2.5,0) node[right]{$x$};
    \draw[thick,->] (0,-2.5)--(0,2.5) node[above]{$y$};
    % Horizontal lines
    \foreach \b in {-2,-1.5,...,2} {
      \draw[blue!60,thin] (-2,\b)--(2,\b);
    }
    % Vertical lines
    \foreach \a in {-2,-1.5,...,2} {
      \draw[red!60,thin] (\a,-2)--(\a,2);
    }
    \node[below] at (0,-2.8) {$z$-plane};
  \end{scope}

  % Arrow
  \draw[-{Stealth},very thick] (3.2,0)--(4.8,0) node[midway,above]{$z\mapsto z^2$};

  % Right: image under z^2
  \begin{scope}[xshift=9cm]
    \draw[thick,->] (-4,0)--(4.5,0) node[right]{$u$};
    \draw[thick,->] (0,-4)--(0,4.5) node[above]{$v$};
    % Image of horizontal lines: parabolas x = t^2-b^2, y = 2bt
    \foreach \b in {0.3,0.6,0.9,1.2,1.5} {
      \draw[blue!60,thick,smooth,samples=60,domain=-1.8:1.8]
        plot ({\x*\x-\b*\b},{2*\b*\x});
      \draw[blue!60,thick,smooth,samples=60,domain=-1.8:1.8]
        plot ({\x*\x-\b*\b},{-2*\b*\x});
    }
    % Image of vertical lines: parabolas x = a^2-t^2, y = 2at
    \foreach \a in {0.3,0.6,0.9,1.2,1.5} {
      \draw[red!60,thick,smooth,samples=60,domain=-1.8:1.8]
        plot ({\a*\a-\x*\x},{2*\a*\x});
      \draw[red!60,thick,smooth,samples=60,domain=-1.8:1.8]
        plot ({\a*\a-\x*\x},{-2*\a*\x});
    }
    \node[below] at (0,-4.3) {$w$-plane ($w=z^2$)};
  \end{scope}
\end{tikzpicture}
\end{center}

\begin{example}[Conformal image under $z\mapsto e^z$]
\label{ex:conformal-exp}
The map $f(z) = e^z$ sends the horizontal strip $\{x+iy : x\in\R,\; 0<y<\pi\}$ conformally onto the upper half-plane. Indeed, $e^{x+iy} = e^x e^{iy}$ has modulus $e^x\in(0,\infty)$ and argument $y\in(0,\pi)$, which sweeps out all of $\{\im(w)>0\}$.

Horizontal lines $y=c$ map to rays from the origin at angle $c$, and vertical lines $x=a$ map to circles of radius $e^a$. These two families are orthogonal in the $w$-plane.
\end{example}

\begin{center}
\begin{tikzpicture}[scale=0.9]
  % Left: strip
  \begin{scope}
    \fill[blue!6] (-3,0) rectangle (3,3.14159);
    \draw[thick,->] (-3.3,0)--(3.3,0) node[right]{$x$};
    \draw[thick,->] (0,-0.5)--(0,3.8) node[above]{$y$};
    \draw[blue!60!black,thick] (-3,0)--(3,0);
    \draw[blue!60!black,thick] (-3,3.14159)--(3,3.14159);
    \node at (3.6,3.14159) {$y=\pi$};
    % Horizontal lines
    \foreach \c in {0.5,1.0,...,2.5} {
      \draw[blue!50,thin] (-3,\c)--(3,\c);
    }
    % Vertical lines
    \foreach \a in {-2.5,-2,...,2.5} {
      \draw[red!50,thin] (\a,0)--(\a,3.14159);
    }
    \node[below] at (0,-0.8) {$z$-plane (strip $0<y<\pi$)};
  \end{scope}

  \draw[-{Stealth},very thick] (4,1.5)--(5.5,1.5) node[midway,above]{$e^z$};

  % Right: upper half-plane
  \begin{scope}[xshift=9.5cm]
    \fill[green!6] (-3.5,0) rectangle (3.5,3.8);
    \draw[thick,->] (-3.7,0)--(3.7,0) node[right]{$u$};
    \draw[thick,->] (0,-0.5)--(0,3.8) node[above]{$v$};
    % Rays from origin (images of horizontal lines y=c)
    \foreach \c in {0.4,0.8,...,2.8} {
      \draw[blue!50,thick] (0,0)--({3.2*cos(\c r)},{3.2*sin(\c r)});
    }
    % Circles centred at origin (images of vertical lines x=a)
    \foreach \r in {0.5,1.0,1.5,2.0,2.5,3.0} {
      \draw[red!50,thick] (\r,0) arc (0:180:\r);
    }
    \node[below] at (0,-0.8) {$w$-plane (upper half-plane)};
  \end{scope}
\end{tikzpicture}
\end{center}

\section{Harmonic Functions and Their Connection to Holomorphy}
\label{sec:harmonic}

\begin{definition}[Harmonic function]
\label{def:harmonic}
A $C^2$ function $\varphi:\Omega\to\R$ (where $\Omega\subseteq\R^2$ is open) is \emph{harmonic} if it satisfies the \emph{Laplace equation}:
\[
  \Delta\varphi = \frac{\partial^2\varphi}{\partial x^2} + \frac{\partial^2\varphi}{\partial y^2} = 0.
\]
\end{definition}

\begin{theorem}[Real and imaginary parts of holomorphic functions are harmonic]
\label{thm:holo-harmonic}
If $f = u+iv$ is holomorphic on a domain $\Omega$, then $u$ and $v$ are harmonic on $\Omega$.
\end{theorem}

\begin{proof}
We shall use the fact (proved later via Cauchy's integral formula) that a holomorphic function is infinitely differentiable. Granted this, the Cauchy--Riemann equations $u_x = v_y$ and $u_y = -v_x$ can be differentiated:
\[
  u_{xx} = v_{yx}, \qquad u_{yy} = -v_{xy}.
\]
Since $v$ is $C^2$, Schwarz's theorem gives $v_{xy} = v_{yx}$, hence
\[
  \Delta u = u_{xx}+u_{yy} = v_{yx}-v_{xy} = 0.
\]
Similarly, differentiating the Cauchy--Riemann equations the other way:
\[
  v_{xx} = -u_{yx}, \qquad v_{yy} = u_{xy}.
\]
So $\Delta v = v_{xx}+v_{yy} = -u_{yx}+u_{xy} = 0$.
\end{proof}

\begin{definition}[Harmonic conjugate]
\label{def:harmonic-conjugate}
Let $u:\Omega\to\R$ be harmonic on a domain $\Omega$. A harmonic function $v:\Omega\to\R$ is called a \emph{harmonic conjugate} of $u$ if $f = u+iv$ is holomorphic on $\Omega$, i.e., if $u$ and $v$ satisfy the Cauchy--Riemann equations.
\end{definition}

\begin{theorem}[Existence of harmonic conjugate on simply connected domains]
\label{thm:harmonic-conjugate}
Let $\Omega\subseteq\C$ be a simply connected domain and $u:\Omega\to\R$ a harmonic function. Then $u$ has a harmonic conjugate $v$ on $\Omega$, unique up to an additive constant. Consequently, $f = u+iv$ is holomorphic on $\Omega$.
\end{theorem}

\begin{proof}
We need to find $v$ satisfying $v_x = -u_y$ and $v_y = u_x$. Fix a point $(x_0,y_0)\in\Omega$ and define
\[
  v(x,y) = \int_\gamma \left(-u_y\,\dd x + u_x\,\dd y\right),
\]
where $\gamma$ is any piecewise-smooth path in $\Omega$ from $(x_0,y_0)$ to $(x,y)$. For this to be well-defined (independent of the path), we need the $1$-form $\omega = -u_y\,\dd x + u_x\,\dd y$ to be closed, which means
\[
  \frac{\partial}{\partial y}(-u_y) = \frac{\partial}{\partial x}(u_x) \quad\Longleftrightarrow\quad -u_{yy} = u_{xx} \quad\Longleftrightarrow\quad \Delta u = 0.
\]
This holds since $u$ is harmonic. On a simply connected domain, a closed $1$-form is exact (Poincar\'e lemma), so the integral is path-independent and defines a well-defined function $v$.

By the fundamental theorem of calculus for line integrals, $v_x = -u_y$ and $v_y = u_x$, which are the Cauchy--Riemann equations. Uniqueness up to a constant follows because if $v_1,v_2$ are two harmonic conjugates, then $(v_1-v_2)_x = (v_1-v_2)_y = 0$, so $v_1-v_2$ is constant on the connected set $\Omega$.
\end{proof}

\begin{remark}
Simple connectivity is essential. On $\C^* = \C\setminus\{0\}$, the function $u(x,y)=\frac{1}{2}\ln(x^2+y^2)$ is harmonic, but its conjugate $v = \arctan(y/x)$ is multi-valued (the argument function). There is no single-valued harmonic conjugate on all of $\C^*$.
\end{remark}

\begin{example}[Finding a harmonic conjugate]
\label{ex:harmonic-conjugate}
Let $u(x,y) = x^2-y^2$ (which is harmonic since $u_{xx}+u_{yy}=2+(-2)=0$). We seek $v$ with $v_x = -u_y = 2y$ and $v_y = u_x = 2x$. From the first equation: $v(x,y) = 2xy + g(y)$. Substituting into the second: $2x + g'(y) = 2x$, so $g'(y)=0$, i.e., $g(y)=C$. Therefore $v(x,y)=2xy+C$, and
\[
  f(z) = (x^2-y^2)+i(2xy+C) = z^2+iC.
\]
\end{example}

\begin{example}[Harmonic conjugate of $\ln\abs{z}$]
\label{ex:harmonic-ln}
On $\Omega = \C\setminus(-\infty,0]$ (the slit plane), $u(x,y) = \frac{1}{2}\ln(x^2+y^2) = \ln\abs{z}$ is harmonic. The harmonic conjugate is $v = \Arg(z)$ (the principal argument), giving $f = u+iv = \Log(z)$, the principal branch of the logarithm.
\end{example}

\begin{proposition}[Laplacian in terms of Wirtinger derivatives]
\label{prop:laplacian-wirtinger}
For a $C^2$ function $\varphi:\Omega\to\R$,
\[
  \Delta\varphi = 4\,\frac{\partial^2\varphi}{\partial z\,\partial\conj{z}}.
\]
\end{proposition}

\begin{proof}
Using $\partial/\partial z = \frac{1}{2}(\partial_x - i\partial_y)$ and $\partial/\partial\conj{z} = \frac{1}{2}(\partial_x+i\partial_y)$:
\[
  4\frac{\partial^2\varphi}{\partial z\,\partial\conj{z}} = 4\cdot\frac{1}{2}(\partial_x-i\partial_y)\cdot\frac{1}{2}(\partial_x+i\partial_y)\varphi = (\partial_x-i\partial_y)(\partial_x+i\partial_y)\varphi = (\partial_x^2+\partial_y^2)\varphi = \Delta\varphi.
\]
(The cross terms cancel: $i\partial_x\partial_y - i\partial_y\partial_x = 0$ by Schwarz.)
\end{proof}

\begin{proposition}[Mean value property for harmonic functions]
\label{prop:mean-value-harmonic}
If $\varphi$ is harmonic on a domain $\Omega$ containing $\conj{D}(a,r)$, then
\[
  \varphi(a) = \frac{1}{2\pi}\int_0^{2\pi}\varphi(a+re^{i\theta})\,\dd\theta.
\]
That is, the value at the centre equals the average on any circle centred there.
\end{proposition}

\begin{proof}
This will be proved as a consequence of the Cauchy integral formula in Chapter~3. For now we state it to illustrate the power of harmonicity.
\end{proof}

\begin{proposition}[Maximum principle for harmonic functions]
\label{prop:max-principle-harmonic}
Let $\varphi$ be harmonic on a domain $\Omega$. If $\varphi$ attains a maximum or minimum value at an interior point of $\Omega$, then $\varphi$ is constant.
\end{proposition}

\begin{proof}
Suppose $\varphi(a) = \max_{z\in\Omega}\varphi(z)$ for some $a\in\Omega$. By the mean value property, for any small $r>0$:
\[
  \varphi(a) = \frac{1}{2\pi}\int_0^{2\pi}\varphi(a+re^{i\theta})\,\dd\theta.
\]
Since $\varphi(a+re^{i\theta})\leq\varphi(a)$ for all $\theta$ and the average equals $\varphi(a)$, we must have $\varphi(a+re^{i\theta})=\varphi(a)$ for all $\theta$ (a non-negative continuous function with zero integral is identically zero). Thus $\varphi$ is constant on $D(a,r)$. The set $\{z\in\Omega:\varphi(z)=\varphi(a)\}$ is both open (by the argument above) and closed (by continuity), hence equal to $\Omega$ by connectedness.
\end{proof}

\section{Entire Functions and First Glimpses Beyond}
\label{sec:entire}

\begin{definition}[Entire function]
\label{def:entire}
A function $f:\C\to\C$ that is holomorphic on all of $\C$ is called an \emph{entire function}. Examples include polynomials, $e^z$, $\sin z$, $\cos z$.
\end{definition}

\begin{proposition}[Trigonometric functions as entire functions]
\label{prop:trig-entire}
The functions
\[
  \sin z = \frac{e^{iz}-e^{-iz}}{2i}, \qquad \cos z = \frac{e^{iz}+e^{-iz}}{2}
\]
are entire, with $(\sin z)' = \cos z$ and $(\cos z)' = -\sin z$. Note that unlike their real counterparts, $\sin z$ and $\cos z$ are \emph{unbounded} on $\C$: for instance, $\sin(iy) = i\sinh(y)\to\infty$ as $y\to\infty$.
\end{proposition}

\begin{example}[The hyperbolic functions]
\label{ex:hyperbolic}
The hyperbolic functions are defined by
\[
  \sinh z = \frac{e^z-e^{-z}}{2}, \qquad \cosh z = \frac{e^z+e^{-z}}{2}.
\]
They are entire and satisfy $\sin(iz)=i\sinh z$ and $\cos(iz)=\cosh z$.
\end{example}

\begin{remark}[Preview: Liouville's theorem]
\label{rem:liouville-preview}
We shall prove in Chapter~3 that every \emph{bounded} entire function is constant (Liouville's theorem). This is a dramatic consequence of the rigidity of holomorphic functions, with no analogue in real analysis. As a corollary, one obtains the \emph{fundamental theorem of algebra}: every non-constant polynomial has a root in $\C$.
\end{remark}

\section{Exercises for Chapter 2}
\label{sec:exercises-ch2}

\begin{exercise}[$\bigstar$]
\label{exo:ch2-1}
Determine which of the following functions are holomorphic, and compute the derivative where it exists:
\begin{enumerate}[label=\textup{(\alph*)}]
  \item $f(z) = z^3-2z+i$
  \item $f(z) = \abs{z}^2$
  \item $f(z) = \re(z)$
  \item $f(z) = e^{z^2}$
  \item $f(z) = z\conj{z}+z^2$
\end{enumerate}
\end{exercise}

\begin{exercise}[$\bigstar$]
\label{exo:ch2-2}
Verify the Cauchy--Riemann equations for $f(z) = \sin z$ by writing $f = u+iv$ explicitly.
\end{exercise}

\begin{exercise}[$\bigstar$]
\label{exo:ch2-3}
Find the harmonic conjugate of $u(x,y) = e^x\cos y$ and express the resulting holomorphic function in terms of $z$.
\end{exercise}

\begin{exercise}[$\bigstar$]
\label{exo:ch2-4}
Show that $u(x,y) = x^3-3xy^2$ is harmonic and find a holomorphic function $f$ with $\re(f)=u$.
\end{exercise}

\begin{exercise}[$\bigstar\!\bigstar$]
\label{exo:ch2-5}
Let $f:\Omega\to\C$ be holomorphic and suppose $\re(f)$ is constant on $\Omega$ (a connected open set). Show that $f$ is constant.
\end{exercise}

\begin{exercise}[$\bigstar\!\bigstar$]
\label{exo:ch2-6}
Let $f:\Omega\to\C$ be holomorphic. Show that if $\abs{f}$ is constant on $\Omega$ (connected), then $f$ is constant.

\textit{Hint:} If $\abs{f}^2 = u^2+v^2 = c$ is constant, differentiate and use the Cauchy--Riemann equations.
\end{exercise}

\begin{exercise}[$\bigstar\!\bigstar$]
\label{exo:ch2-7}
Let $f = u+iv$ be holomorphic on a domain $\Omega$. Show that the level curves $\{u=\text{const}\}$ and $\{v=\text{const}\}$ are orthogonal at every point where $f'(z)\neq 0$.

\textit{Hint:} The gradient of $u$ is $\nabla u = (u_x,u_y)$ and that of $v$ is $\nabla v = (v_x,v_y)$. Compute $\nabla u\cdot\nabla v$ using the Cauchy--Riemann equations.
\end{exercise}

\begin{exercise}[$\bigstar\!\bigstar$]
\label{exo:ch2-8}
Show that $f(z) = \sqrt{\abs{xy}}$ (where $z=x+iy$) satisfies the Cauchy--Riemann equations at the origin but is \emph{not} holomorphic there.

\textit{Hint:} Compute the partial derivatives at the origin and check whether $f$ is differentiable (in the real sense) at the origin.
\end{exercise}

\begin{exercise}[$\bigstar\!\bigstar$]
\label{exo:ch2-9}
Describe the image of the following sets under $f(z)=z^2$:
\begin{enumerate}[label=\textup{(\alph*)}]
  \item The first quadrant $\{x>0, y>0\}$.
  \item The half-strip $\{0<x<1, y>0\}$.
  \item The sector $\{0<\arg z<\pi/4\}$.
\end{enumerate}
In each case, verify conformality by describing what happens to the boundary curves.
\end{exercise}

\begin{exercise}[$\bigstar\!\bigstar$]
\label{exo:ch2-10}
Let $f(z) = z + 1/z$ (defined on $\C^*$).
\begin{enumerate}[label=\textup{(\alph*)}]
  \item Show that $f$ maps the unit circle $\abs{z}=1$ onto the interval $[-2,2]$.
  \item Show that $f$ maps $\C\setminus\conj{\D}$ (the exterior of the closed unit disk) conformally onto $\C\setminus[-2,2]$.
  \item This is the \emph{Joukowski map}, fundamental in aerodynamics. Describe the images of circles $\abs{z}=r$ for $r>1$.
\end{enumerate}
\end{exercise}

\begin{exercise}[$\bigstar\!\bigstar\!\bigstar$]
\label{exo:ch2-11}
\textbf{(Harmonic functions and the Dirichlet problem.)} Let $u$ be a continuous function on $\conj{\D}$ that is harmonic on $\D$. Suppose $u=0$ on $\T = \partial\D$. Prove that $u\equiv 0$ on $\D$.

\textit{Hint:} Apply the maximum principle to both $u$ and $-u$.
\end{exercise}

\begin{exercise}[$\bigstar\!\bigstar\!\bigstar$]
\label{exo:ch2-12}
\textbf{(Looman--Menchov theorem preview.)} Show that if $f=u+iv$ is continuous on a domain $\Omega$, the partial derivatives $u_x,u_y,v_x,v_y$ exist everywhere in $\Omega$, and the Cauchy--Riemann equations hold at every point, then $f$ is holomorphic. (This is significantly harder than Theorem~\ref{thm:CR-sufficient} because we do not assume continuity of the partial derivatives. You may use without proof the fact that a continuous function satisfying the Cauchy--Riemann equations is holomorphic if it has only countably many points where it is not differentiable in the real sense.)
\end{exercise}

\begin{exercise}[$\bigstar\!\bigstar\!\bigstar$]
\label{exo:ch2-13}
Let $\Omega$ be a domain in $\C$ and $f:\Omega\to\C$ a holomorphic function. Show that if $f$ maps $\Omega$ into a straight line, then $f$ is constant.

\textit{Hint:} After composing with an affine transformation, reduce to the case $\im(f) = 0$.
\end{exercise}

\begin{exercise}[$\bigstar\!\bigstar\!\bigstar$]
\label{exo:ch2-14}
Let $\Omega$ be a domain and $f:\Omega\to\C$ holomorphic with $f'\neq 0$ on $\Omega$. Show that $f$ is an open mapping (images of open sets are open) by using the inverse function theorem for $\R^2$ and the fact that the Jacobian determinant of $f$ is $\abs{f'(z)}^2 > 0$.
\end{exercise}

%%%%%%%%%%%%%%%%%%%%%%%%%%%%%%%%%%%%%%%%%%%%%%%%%%%%%%%%%%%%
%%
%%  CHAPTER 3: Elementary Functions
%%
%%%%%%%%%%%%%%%%%%%%%%%%%%%%%%%%%%%%%%%%%%%%%%%%%%%%%%%%%%%%
\chapter{Elementary Functions}
\label{chap:elementary-functions}

\section*{Introduction}
\addcontentsline{toc}{section}{Introduction}

In real analysis, the elementary functions --- exponential, logarithm,
trigonometric and hyperbolic functions --- are old friends.  In the complex
setting they acquire a richer and sometimes surprising personality.  The
exponential function $e^z$ becomes periodic in the imaginary direction; the
logarithm splits into infinitely many \emph{branches}; and the trigonometric
functions lose their boundedness.  This chapter is devoted to a careful study
of these functions, their mapping properties, and the subtleties that arise
from multi-valuedness.

%% ════════════════════════════════════════════════════════════
\section{The Complex Exponential Function}
\label{sec:complex-exp}
%% ════════════════════════════════════════════════════════════

\subsection{Definition and first properties}

\begin{definition}[Complex exponential]\label{def:complex-exp}
For $z = x + iy \in \C$ (with $x,y\in\R$), the \textbf{complex exponential
function}\index{exponential function} is defined by
\[
  \exp(z) \;=\; e^z \;=\; e^x(\cos y + i\sin y).
\]
Equivalently, using the power series that converges for every $z\in\C$,
\[
  \exp(z) \;=\; \sum_{n=0}^{\infty} \frac{z^n}{n!}.
\]
\end{definition}

\begin{proposition}[Basic properties of $\exp$]\label{prop:exp-properties}
The complex exponential satisfies the following properties.
\begin{enumerate}[label=\textup{(\roman*)}]
  \item \textbf{Homomorphism.} For all $z,w\in\C$,
    $e^{z+w} = e^z\cdot e^w$.
  \item \textbf{Non-vanishing.} $e^z \ne 0$ for every $z\in\C$, and
    $(e^z)^{-1} = e^{-z}$.
  \item \textbf{Modulus and argument.}
    $\abs{e^z} = e^{\re z}$ and $\arg(e^z) = \im z + 2\pi\Z$.
  \item \textbf{Periodicity.} $e^{z+2\pi i} = e^z$ for all $z\in\C$;
    moreover $2\pi i$ is a \emph{fundamental period}: if $e^z = e^w$ then
    $z - w \in 2\pi i\Z$.
  \item \textbf{Holomorphicity.} $\exp$ is entire and
    $\dfrac{\dd}{\dd z}\,e^z = e^z$.
\end{enumerate}
\end{proposition}

\begin{proof}
\textbf{(i)} Using the power series, the Cauchy product gives
\[
  e^z \cdot e^w
  = \Bigl(\sum_{n\ge 0}\frac{z^n}{n!}\Bigr)
    \Bigl(\sum_{m\ge 0}\frac{w^m}{m!}\Bigr)
  = \sum_{k\ge 0}\frac{1}{k!}\sum_{j=0}^{k}\binom{k}{j}z^j w^{k-j}
  = \sum_{k\ge 0}\frac{(z+w)^k}{k!}
  = e^{z+w}.
\]
\textbf{(ii)} From (i), $e^z \cdot e^{-z} = e^0 = 1$, so $e^z\ne 0$.

\textbf{(iii)} Write $z = x+iy$.  Then
$\abs{e^z} = \abs{e^x(\cos y + i\sin y)} = e^x\sqrt{\cos^2 y+\sin^2 y} = e^x = e^{\re z}$.
The argument statement follows from $\cos y + i\sin y = e^{iy}$.

\textbf{(iv)} $e^{z+2\pi i} = e^z e^{2\pi i} = e^z\cdot 1 = e^z$.
If $e^z = e^w$, then $e^{z-w}=1$, so $e^{x'-x''}(\cos(y'-y'')+i\sin(y'-y''))=1$
where $z-w = (x'-x'')+i(y'-y'')$.  Taking moduli, $e^{x'-x''}=1$ so $x'=x''$.
Then $\cos(y'-y'')=1$ and $\sin(y'-y'')=0$, giving $y'-y''\in 2\pi\Z$.

\textbf{(v)} The power series $\sum z^n/n!$ has infinite radius of convergence,
so $\exp$ is entire.  Term-by-term differentiation gives
$\exp'(z) = \sum_{n\ge 1} z^{n-1}/(n-1)! = \exp(z)$.
\end{proof}

\begin{remark}[Surjectivity]\label{rem:exp-surjective}
The map $\exp\colon\C\to\C$ is \emph{not} surjective onto~$\C$ (the value $0$
is never attained), but it \emph{is} surjective onto $\C^* = \C\setminus\{0\}$.
Indeed, given $w = \rho e^{i\theta}\ne 0$, we can choose $z = \ln\rho + i\theta$
and verify that $e^z = w$.
\end{remark}

\subsection{Mapping properties of $\exp$}

The exponential maps horizontal lines to rays from the origin and vertical
lines to circles centered at the origin.  More precisely:

\begin{proposition}[Image of lines and strips]\label{prop:exp-mapping}
\leavevmode
\begin{enumerate}[label=\textup{(\roman*)}]
  \item The horizontal line $\{\,x+iy_0 : x\in\R\,\}$ is mapped onto the
    ray $\{\,re^{iy_0} : r>0\,\}$.
  \item The vertical line $\{\,x_0+iy : y\in\R\,\}$ is mapped onto the
    circle $\{\,e^{x_0}e^{iy} : y\in\R\,\} = \{w\in\C : \abs{w}=e^{x_0}\}$.
  \item The horizontal strip $S_\alpha = \{\,x+iy : x\in\R,\;\alpha < y < \alpha+2\pi\,\}$
    is mapped \emph{bijectively} onto $\C^*$.
\end{enumerate}
\end{proposition}

\begin{proof}
Parts (i) and (ii) are immediate from $e^{x+iy} = e^x e^{iy}$.  For (iii),
surjectivity follows from Remark~\ref{rem:exp-surjective}: given $w\in\C^*$,
write $w = \abs{w}e^{i\theta}$ with $\theta$ chosen uniquely in
$(\alpha,\alpha+2\pi)$, and set $z = \ln\abs{w}+i\theta\in S_\alpha$.
Injectivity follows from Proposition~\ref{prop:exp-properties}(iv): if
$e^{z_1}=e^{z_2}$ with $z_1,z_2\in S_\alpha$, then $z_1-z_2\in 2\pi i\Z$
and $\abs{\im(z_1-z_2)}<2\pi$, so $z_1=z_2$.
\end{proof}

\begin{center}
\begin{tikzpicture}[scale=0.85, >=Stealth]
  %% LEFT: z-plane
  \begin{scope}[shift={(-5,0)}]
    \fill[blue!8] (-3,0) rectangle (3,3.14159);
    \draw[thick, blue!60!black] (-3,0) -- (3,0);
    \draw[thick, blue!60!black] (-3,3.14159) -- (3,3.14159);
    \draw[dashed, blue!40!black] (-3,1.5708) -- (3,1.5708);
    \draw[->] (-3.3,0) -- (3.5,0) node[right]{$\re z$};
    \draw[->] (0,-0.5) -- (0,3.8) node[above]{$\im z$};
    \node[left, font=\small] at (-3,0) {$0$};
    \node[left, font=\small] at (-3,1.5708) {$\frac{\pi}{2}$};
    \node[left, font=\small] at (-3,3.14159) {$\pi$};
    \node[below, font=\small\bfseries] at (0,-1) {$z$-plane};
    \node[blue!60!black, font=\small] at (0,2.3) {$S_0$ (upper half)};
    %% Horizontal lines
    \draw[red!70!black, thick] (-2.5,0.6) -- (2.5,0.6);
    \draw[green!50!black, thick] (-2.5,2.2) -- (2.5,2.2);
    \node[red!70!black, font=\tiny, right] at (2.5,0.6) {$y=y_1$};
    \node[green!50!black, font=\tiny, right] at (2.5,2.2) {$y=y_2$};
  \end{scope}
  %% Arrow
  \draw[->, line width=1.5pt] (-1.2,1.5) -- (0.8,1.5)
    node[midway, above, font=\small]{$\exp$};
  %% RIGHT: w-plane
  \begin{scope}[shift={(4,0)}]
    \fill[blue!8] (0,0) -- (3,0) arc (0:180:3) -- cycle;
    \draw[->] (-3.5,0) -- (3.5,0) node[right]{$\re w$};
    \draw[->] (0,-0.5) -- (0,3.8) node[above]{$\im w$};
    %% Rays corresponding to horizontal lines
    \draw[red!70!black, thick] (0,0) -- ({3*cos(34.377)},{3*sin(34.377)});
    \draw[green!50!black, thick] (0,0) -- ({3*cos(126.1)},{3*sin(126.1)});
    \node[red!70!black, font=\tiny] at ({2.8*cos(34.377)+0.4},{2.8*sin(34.377)+0.2}) {$\arg=y_1$};
    \node[green!50!black, font=\tiny] at ({2.2*cos(126.1)-0.7},{2.2*sin(126.1)+0.2}) {$\arg=y_2$};
    %% Unit circle
    \draw[dashed, gray] (0,0) circle (1);
    \node[gray, font=\tiny] at (1.3,-0.3) {$\abs{w}=1$};
    \node[below, font=\small\bfseries] at (0,-1) {$w$-plane};
    \node[blue!60!black, font=\small] at (0.8,2.5) {upper half-plane};
  \end{scope}
\end{tikzpicture}
\captionof{figure}{The exponential maps the strip
$\{0<\im z<\pi\}$ bijectively onto the upper half of $\C^*$.
Horizontal lines become rays; vertical lines (not shown) become
circles.}
\label{fig:exp-strip-to-sector}
\end{center}

\begin{example}[Image of a rectangle]\label{ex:exp-rectangle}
Consider the rectangle $R = \{x+iy : 0\le x\le 1,\;0\le y\le\pi/2\}$.
Under $\exp$:
\begin{itemize}
  \item The bottom edge ($y=0$, $0\le x\le 1$) maps to the segment
    $[1,e]$ on the positive real axis.
  \item The top edge ($y=\pi/2$, $0\le x\le 1$) maps to the segment
    $[i,ei]$ on the positive imaginary axis.
  \item The left edge ($x=0$, $0\le y\le\pi/2$) maps to the arc
    $\{e^{iy}: 0\le y\le\pi/2\}$, a quarter of the unit circle.
  \item The right edge ($x=1$, $0\le y\le\pi/2$) maps to the arc
    $\{ee^{iy}: 0\le y\le\pi/2\}$, a quarter of the circle of radius~$e$.
\end{itemize}
The image is the sector
$\{w\in\C : 1\le\abs{w}\le e,\;0\le\arg w\le\pi/2\}$.
\end{example}

%% ════════════════════════════════════════════════════════════
\section{The Complex Logarithm}
\label{sec:complex-log}
%% ════════════════════════════════════════════════════════════

\subsection{Motivation and multi-valuedness}

We seek a function $\log$ that ``inverts'' the exponential: $e^{\log w} = w$
for $w\in\C^*$.  If $w = \rho e^{i\theta}$ with $\rho>0$ and
$\theta\in\R$, then $z = \ln\rho + i\theta$ satisfies $e^z = w$.  But so does
$z + 2k\pi i$ for every $k\in\Z$.  The complex logarithm is therefore
intrinsically \emph{multi-valued}.

\begin{definition}[Multi-valued logarithm]\label{def:log-multivalued}
For $w\in\C^*$, the \textbf{multi-valued logarithm}\index{logarithm!multi-valued}
is the set
\[
  \log w \;=\; \bigl\{\,\ln\abs{w} + i\arg w + 2k\pi i : k\in\Z\,\bigr\}
         \;=\; \ln\abs{w} + i\Arg w + 2\pi i\Z,
\]
where $\Arg w\in(-\pi,\pi]$ denotes the principal value of the argument.
Each element of this set is called a \textbf{value} (or \textbf{determination})
of $\log w$.
\end{definition}

\subsection{The principal branch}

\begin{definition}[Principal logarithm]\label{def:principal-log}
The \textbf{principal logarithm}\index{logarithm!principal branch}
(or \textbf{principal branch of the logarithm}) is the function
\[
  \Log\colon \C\setminus(-\infty,0] \;\longrightarrow\; \C, \qquad
  \Log w \;=\; \ln\abs{w} + i\Arg w,
\]
where $\Arg w$ is the unique argument of~$w$ in $(-\pi,\pi)$.
The set $(-\infty,0] = \{x\in\R : x\le 0\}$ is called the
\textbf{branch cut}\index{branch cut}.
\end{definition}

\begin{remark}\label{rem:log-branch-cut}
The branch cut $(-\infty,0]$ is a somewhat arbitrary choice.  One could
instead remove any ray from the origin to infinity and define a
corresponding branch.  The important point is that one must remove
\emph{some} curve from~$0$ to~$\infty$ in order to obtain a single-valued,
continuous logarithm.  The reason is topological: the argument function
cannot be defined continuously on all of $\C^*$, because $\C^*$ has a
``hole'' at the origin.
\end{remark}

\begin{theorem}[Properties of $\Log$]\label{thm:log-properties}
\leavevmode
\begin{enumerate}[label=\textup{(\roman*)}]
  \item $\Log$ is holomorphic on $\C\setminus(-\infty,0]$ and
    \[
      \frac{\dd}{\dd z}\,\Log z \;=\; \frac{1}{z}.
    \]
  \item $e^{\Log z} = z$ for all $z\in\C\setminus(-\infty,0]$.
  \item $\Log(e^z) = z$ if and only if $\im z\in(-\pi,\pi)$.
  \item The identity $\Log(z_1 z_2) = \Log z_1 + \Log z_2$ holds
    \emph{modulo $2\pi i$}: there exists $k\in\{-1,0,1\}$ such that
    $\Log(z_1 z_2) = \Log z_1 + \Log z_2 + 2k\pi i$.
\end{enumerate}
\end{theorem}

\begin{proof}
\textbf{(i)} Write $z = re^{i\theta}$ with $r>0$ and $\theta\in(-\pi,\pi)$.
In terms of $u = \re z$ and $v = \im z$ we have $r = \sqrt{u^2+v^2}$ and
$\theta = \arctan(v/u)$ (with appropriate adjustments for $u\le 0$, but we
have excluded the negative real axis).  Then
\[
  \Log z = \tfrac{1}{2}\ln(u^2+v^2) + i\arctan(v/u).
\]
Both components are smooth, and a direct computation verifies the Cauchy--Riemann
equations.  For the derivative, let $g = \Log$ and note
$e^{g(z)} = z$, so differentiating: $g'(z)\cdot e^{g(z)} = 1$, giving
$g'(z) = 1/e^{g(z)} = 1/z$.

\textbf{(ii)} Let $z = re^{i\theta}$ with $\theta\in(-\pi,\pi)$.
Then $e^{\Log z} = e^{\ln r + i\theta} = re^{i\theta} = z$.

\textbf{(iii)} $\Log(e^z) = \ln\abs{e^z}+i\Arg(e^z) = \re z + i\Arg(e^z)$.
This equals $z = \re z + i\,\im z$ if and only if $\Arg(e^z) = \im z$, which
holds precisely when $\im z \in (-\pi,\pi)$.

\textbf{(iv)} Both sides of $\Log(z_1 z_2) = \Log z_1 + \Log z_2$ are values
of $\log(z_1 z_2)$, so they differ by $2k\pi i$ for some $k\in\Z$.
Since arguments lie in $(-\pi,\pi)$, the difference of the two sides has
imaginary part in $(-3\pi,3\pi)$, forcing $k\in\{-1,0,1\}$.
\end{proof}

\subsection{General branches of the logarithm}

\begin{definition}[Branch of the logarithm]\label{def:log-branch}
Let $\Omega\subseteq\C^*$ be a domain (connected open set).  A
\textbf{branch of the logarithm}\index{logarithm!branch} on~$\Omega$ is a
continuous function $L\colon\Omega\to\C$ satisfying $e^{L(z)}=z$ for all
$z\in\Omega$.  Any such branch is necessarily holomorphic with
$L'(z)=1/z$, and any two branches on~$\Omega$ differ by a constant
$2k\pi i$ with $k\in\Z$.
\end{definition}

\begin{proposition}[Existence of branches]\label{prop:log-branch-existence}
A branch of the logarithm exists on a domain $\Omega\subseteq\C^*$ if and
only if the function $1/z$ has a primitive on~$\Omega$.  In particular, a branch
exists whenever $\Omega$ is \emph{simply connected} (i.e., every closed curve
in~$\Omega$ can be continuously deformed to a point within~$\Omega$).
\end{proposition}

\begin{proof}
If $L$ is a branch of $\log$ on $\Omega$, then $L'(z)=1/z$, so $L$ is a
primitive of $1/z$.  Conversely, if $F$ is a primitive of $1/z$ on $\Omega$,
fix $z_0\in\Omega$ and choose $c\in\C$ with $e^{F(z_0)+c}=z_0$ (this is
possible since $z_0\ne 0$).  Define $L(z) = F(z)+c$.  Then
$\frac{\dd}{\dd z}\bigl[z\,e^{-L(z)}\bigr] = e^{-L(z)} - z\cdot\frac{1}{z}
\cdot e^{-L(z)} = 0$, so $z\,e^{-L(z)}$ is constant on~$\Omega$, equal to
$z_0\,e^{-L(z_0)}=1$.  Hence $e^{L(z)}=z$.

The existence of primitives on simply connected domains will follow from
Cauchy's theorem (Chapter~\ref{chap:complex-integration}).
\end{proof}

\subsection{Intuition of Riemann surfaces}

\begin{remark}[The logarithmic Riemann surface]\label{rem:riemann-surface-log}
Rather than choosing a branch, one can ``unwind'' all the values of $\log z$
into a single geometric object: the \textbf{Riemann surface} of the logarithm.
Imagine infinitely many copies of $\C^*$, labelled by $k\in\Z$, each slit
along the negative real axis.  The upper lip of the $k$-th sheet is glued to the
lower lip of the $(k+1)$-th sheet.  The result is a \emph{helicoid}-like surface
on which the logarithm becomes a well-defined, single-valued holomorphic function.
\end{remark}

\begin{center}
\begin{tikzpicture}[scale=0.8, >=Stealth]
  %% Schematic of Riemann surface for log: stacked sheets
  \foreach \k/\lab/\col in {0/$k=0$/blue!50!black, 1/$k=1$/red!60!black, 2/$k=2$/green!50!black} {
    \begin{scope}[shift={(0,\k*1.8)}, opacity=0.85]
      \fill[\col!10] (-3,-0.3) -- (3,-0.3) -- (3.5,0.5) -- (-2.5,0.5) -- cycle;
      \draw[\col, thick] (-3,-0.3) -- (3,-0.3) -- (3.5,0.5) -- (-2.5,0.5) -- cycle;
      \draw[\col, thick, dashed] (-0.5,0.1) -- (3.2,0.1);
      \node[\col, font=\small] at (-2.0,0.1) {\lab};
    \end{scope}
  }
  %% Arrows showing connection
  \draw[->, thick, purple] (3.7,0.5) to[out=30,in=-30] (3.7,1.5);
  \draw[->, thick, purple] (3.7,2.3) to[out=30,in=-30] (3.7,3.3);
  \node[right, font=\small, purple] at (4.0,1.0) {glue};
  \node[right, font=\small, purple] at (4.0,2.8) {glue};
  \node[below, font=\small] at (0.5,-0.6) {Branch cuts along $(-\infty,0]$};
\end{tikzpicture}
\captionof{figure}{Schematic of the Riemann surface of $\log z$: infinitely
many sheets, each carrying one branch $\ln\abs{z}+i(\Arg z + 2k\pi)$,
glued along the branch cut.}
\label{fig:riemann-surface-log}
\end{center}

%% ════════════════════════════════════════════════════════════
\section{Complex Powers}
\label{sec:complex-powers}
%% ════════════════════════════════════════════════════════════

\begin{definition}[Complex power]\label{def:complex-power}
Let $\alpha\in\C$.  For $z\in\C\setminus(-\infty,0]$, the
\textbf{principal value}\index{complex power} of $z^\alpha$ is
\[
  z^\alpha \;=\; \exp\!\bigl(\alpha\,\Log z\bigr)
           \;=\; e^{\alpha\,\Log z}.
\]
More generally, using any branch $L$ of the logarithm on a domain
$\Omega\subseteq\C^*$, one defines $z^\alpha = e^{\alpha\, L(z)}$.
\end{definition}

\begin{proposition}[Properties of complex powers]\label{prop:powers-properties}
Let $\alpha,\beta\in\C$ and $z\in\C\setminus(-\infty,0]$.  Using the principal
branch:
\begin{enumerate}[label=\textup{(\roman*)}]
  \item $z^\alpha$ is holomorphic on $\C\setminus(-\infty,0]$ with
    $\dfrac{\dd}{\dd z}\,z^\alpha = \alpha\, z^{\alpha-1}$.
  \item $z^\alpha \cdot z^\beta = z^{\alpha+\beta}$ (principal values).
  \item $(z^\alpha)^\beta = z^{\alpha\beta}$ holds for principal values
    only up to a factor $e^{2k\pi i\beta}$ with $k\in\Z$.
  \item If $\alpha = n\in\Z$, then $z^n$ coincides with the usual
    integer power, and no branch cut is needed.
  \item If $\alpha = 1/n$ with $n\in\Z_{\ge 1}$, then $z^{1/n}$ gives
    the principal $n$-th root.
\end{enumerate}
\end{proposition}

\begin{proof}
\textbf{(i)} Since $\Log z$ and $\exp$ are holomorphic, so is their composition.
By the chain rule:
\[
  \frac{\dd}{\dd z}\,e^{\alpha\Log z}
  = e^{\alpha\Log z}\cdot\frac{\alpha}{z}
  = \alpha\,\frac{z^\alpha}{z}
  = \alpha\,z^{\alpha-1}.
\]
\textbf{(ii)}
$z^\alpha\cdot z^\beta = e^{\alpha\Log z}\cdot e^{\beta\Log z}
= e^{(\alpha+\beta)\Log z} = z^{\alpha+\beta}$.

\textbf{(iii)} $(z^\alpha)^\beta = e^{\beta\Log(z^\alpha)}$, and
$\Log(z^\alpha)$ may differ from $\alpha\Log z$ by $2k\pi i$.

\textbf{(iv)} When $\alpha=n\in\Z$, $e^{n\Log z} = e^{n\ln\abs{z}+in\Arg z}
= \abs{z}^n e^{in\Arg z} = z^n$ by Euler's formula.  Since all branches of $\log$
give $e^{n\cdot(L(z))}=z^n$ (the $2k\pi i n$ disappears under $\exp$), no branch
cut is needed.

\textbf{(v)} Follows from the definition with $\alpha = 1/n$.
\end{proof}

\begin{example}[Powers of $i$]\label{ex:powers-of-i}
Since $\Log i = \ln 1 + i\,\frac{\pi}{2} = \frac{i\pi}{2}$, we have:
\begin{itemize}
  \item $i^i = e^{i\Log i} = e^{i\cdot i\pi/2} = e^{-\pi/2}\approx 0.2079$.
    This is a \emph{real} number!
  \item $i^{1/2} = e^{\frac{1}{2}\Log i} = e^{i\pi/4}
    = \frac{\sqrt{2}}{2}(1+i)$.
  \item More generally, $i^\alpha = e^{i\alpha\pi/2}$ for any $\alpha\in\C$
    (principal value).
\end{itemize}
Using other branches $\log i = i\pi/2 + 2k\pi i$, we get the other values:
$i^i = e^{-\pi/2-2k\pi}$ for $k\in\Z$.
\end{example}

\begin{example}[The function $z^{1/2}$]\label{ex:sqrt-function}
The principal square root $z^{1/2} = e^{\frac{1}{2}\Log z}$ satisfies:
\begin{itemize}
  \item For $z = re^{i\theta}$ with $r>0$ and $\theta\in(-\pi,\pi)$:
    $z^{1/2} = \sqrt{r}\,e^{i\theta/2}$.
  \item The image of $z^{1/2}$ is the right half-plane $\{\re w > 0\}$
    together with the positive imaginary axis.
  \item On the negative real axis, $z^{1/2}$ is discontinuous: approaching
    $x<0$ from above gives $\sqrt{\abs{x}}\,i$, from below gives
    $-\sqrt{\abs{x}}\,i$.
\end{itemize}
\end{example}

\begin{center}
\begin{tikzpicture}[scale=0.9, >=Stealth]
  %% Left: z-plane with branch cut
  \begin{scope}[shift={(-4.5,0)}]
    \draw[->] (-2.8,0) -- (2.8,0) node[right]{$\re z$};
    \draw[->] (0,-2.3) -- (0,2.5) node[above]{$\im z$};
    \draw[red, line width=2.5pt] (-2.5,0) -- (0,0);
    \node[red, font=\small, below] at (-1.3,-0.15) {branch cut};
    \fill[red] (0,0) circle (2.5pt);
    %% Colored regions
    \fill[blue!10, opacity=0.5] (0.05,0.05) -- (2.5,0.05) -- (2.5,2.2) -- (0.05,2.2) -- cycle;
    \fill[green!10, opacity=0.5] (0.05,-0.05) -- (2.5,-0.05) -- (2.5,-2.2) -- (0.05,-2.2) -- cycle;
    \node[font=\small\bfseries] at (0,-2.8) {$z$-plane};
  \end{scope}
  %% Arrow
  \draw[->, line width=1.5pt] (-1.2,0) -- (0.8,0)
    node[midway, above, font=\small]{$z^{1/2}$};
  %% Right: w-plane
  \begin{scope}[shift={(4,0)}]
    \draw[->] (-2.3,0) -- (2.8,0) node[right]{$\re w$};
    \draw[->] (0,-2.3) -- (0,2.5) node[above]{$\im w$};
    \fill[blue!10, opacity=0.5] (0.05,0.05) -- (2.5,0.05) -- (2.5,2.2) -- (0.05,2.2) -- cycle;
    \fill[green!10, opacity=0.5] (0.05,-0.05) -- (2.5,-0.05) -- (2.5,-2.2) -- (0.05,-2.2) -- cycle;
    \node[font=\small] at (1.5,1.2) {image};
    \node[font=\small\bfseries] at (0,-2.8) {$w$-plane};
    \draw[dashed, gray] (0,-2.2) -- (0,2.2);
  \end{scope}
\end{tikzpicture}
\captionof{figure}{The principal square root maps $\C\setminus(-\infty,0]$
bijectively onto the open right half-plane $\{\re w>0\}$.}
\label{fig:sqrt-mapping}
\end{center}

%% ════════════════════════════════════════════════════════════
\section{Trigonometric and Hyperbolic Functions}
\label{sec:trig-hyperbolic}
%% ════════════════════════════════════════════════════════════

\subsection{Definitions via the exponential}

\begin{definition}[Complex trigonometric functions]\label{def:complex-trig}
For $z\in\C$, the \textbf{complex cosine}\index{cosine!complex} and
\textbf{complex sine}\index{sine!complex} are defined by
\[
  \cos z \;=\; \frac{e^{iz}+e^{-iz}}{2}, \qquad
  \sin z \;=\; \frac{e^{iz}-e^{-iz}}{2i}.
\]
The other trigonometric functions are defined as usual:
\[
  \tan z = \frac{\sin z}{\cos z}, \qquad
  \cot z = \frac{\cos z}{\sin z}, \qquad
  \sec z = \frac{1}{\cos z}, \qquad
  \csc z = \frac{1}{\sin z},
\]
wherever the denominators are nonzero.
\end{definition}

\begin{definition}[Complex hyperbolic functions]\label{def:complex-hyp}
For $z\in\C$, the \textbf{hyperbolic cosine}\index{hyperbolic cosine} and
\textbf{hyperbolic sine}\index{hyperbolic sine} are defined by
\[
  \cosh z \;=\; \frac{e^z+e^{-z}}{2}, \qquad
  \sinh z \;=\; \frac{e^z-e^{-z}}{2}.
\]
Similarly, $\tanh z = \sinh z/\cosh z$, etc.
\end{definition}

\subsection{Properties}

\begin{proposition}[Properties of complex trigonometric functions]%
\label{prop:trig-properties}
\leavevmode
\begin{enumerate}[label=\textup{(\roman*)}]
  \item \textbf{Entirety.} $\cos$ and $\sin$ are entire functions with
    derivatives $\cos' z = -\sin z$ and $\sin' z = \cos z$.
  \item \textbf{Euler identity.} $e^{iz} = \cos z + i\sin z$ for all $z\in\C$.
  \item \textbf{Pythagorean identity.} $\cos^2 z + \sin^2 z = 1$.
  \item \textbf{Periodicity.} $\cos$ and $\sin$ have period $2\pi$:
    $\cos(z+2\pi)=\cos z$, $\sin(z+2\pi)=\sin z$.
  \item \textbf{Zeros.} $\sin z = 0 \iff z = n\pi\;(n\in\Z)$;
    $\cos z = 0 \iff z = (n+\tfrac{1}{2})\pi\;(n\in\Z)$.
  \item \textbf{Unboundedness.} Unlike the real case, $\cos$ and $\sin$ are
    unbounded on~$\C$.  For instance,
    $\abs{\sin(iy)} = \abs{\sinh y}\to\infty$ as $\abs{y}\to\infty$.
  \item \textbf{Addition formulas.} All real addition formulas extend to $\C$:
    \begin{align*}
      \cos(z+w) &= \cos z\cos w - \sin z\sin w, \\
      \sin(z+w) &= \sin z\cos w + \cos z\sin w.
    \end{align*}
\end{enumerate}
\end{proposition}

\begin{proof}
All these follow from the definitions and properties of $\exp$.

\textbf{(i)} is immediate since $e^{iz}$ and $e^{-iz}$ are entire.

\textbf{(ii)} $\cos z + i\sin z = \frac{e^{iz}+e^{-iz}}{2}
+ i\cdot\frac{e^{iz}-e^{-iz}}{2i} = \frac{e^{iz}+e^{-iz}}{2}
+ \frac{e^{iz}-e^{-iz}}{2} = e^{iz}$.

\textbf{(iii)} $\cos^2 z + \sin^2 z
= \left(\frac{e^{iz}+e^{-iz}}{2}\right)^{\!2}
+ \left(\frac{e^{iz}-e^{-iz}}{2i}\right)^{\!2}$.
Expanding: $= \frac{e^{2iz}+2+e^{-2iz}}{4} - \frac{e^{2iz}-2+e^{-2iz}}{4}
= \frac{4}{4} = 1$.

\textbf{(v)} $\sin z = 0 \iff e^{iz} = e^{-iz} \iff e^{2iz} = 1
\iff 2iz \in 2\pi i\Z \iff z\in\pi\Z$.  The argument for $\cos$ is similar.

\textbf{(vi)} $\sin(iy) = \frac{e^{-y}-e^{y}}{2i} = i\frac{e^{y}-e^{-y}}{2}
= i\sinh y$, so $\abs{\sin(iy)}=\abs{\sinh y}\to\infty$.
\end{proof}

\begin{proposition}[Relations between trigonometric and hyperbolic functions]%
\label{prop:trig-hyp-relation}
For all $z\in\C$:
\[
  \cos(iz) = \cosh z, \qquad \sin(iz) = i\sinh z, \qquad
  \cosh(iz) = \cos z, \qquad \sinh(iz) = i\sin z.
\]
\end{proposition}

\begin{proof}
Direct computation: $\cos(iz) = \frac{e^{i(iz)}+e^{-i(iz)}}{2}
= \frac{e^{-z}+e^{z}}{2} = \cosh z$.  The other identities are proved
similarly.
\end{proof}

\begin{proposition}[Decomposition of $\sin$ and $\cos$ into real and imaginary parts]%
\label{prop:trig-real-imag}
Writing $z = x+iy$ with $x,y\in\R$:
\begin{align*}
  \sin z &= \sin x\cosh y + i\cos x\sinh y, \\
  \cos z &= \cos x\cosh y - i\sin x\sinh y.
\end{align*}
In particular,
\[
  \abs{\sin z}^2 = \sin^2 x + \sinh^2 y, \qquad
  \abs{\cos z}^2 = \cos^2 x + \sinh^2 y.
\]
\end{proposition}

\begin{proof}
By the addition formula,
$\sin(x+iy) = \sin x\cos(iy) + \cos x\sin(iy) = \sin x\cosh y + i\cos x\sinh y$.
Then $\abs{\sin z}^2 = \sin^2 x\cosh^2 y + \cos^2 x\sinh^2 y
= \sin^2 x(1+\sinh^2 y) + (1-\sin^2 x)\sinh^2 y = \sin^2 x + \sinh^2 y$.
The computation for $\cos z$ is analogous.
\end{proof}

\subsection{Mapping properties of $\sin$ and $\cos$}

\begin{example}[Image of a vertical strip under $\sin$]\label{ex:sin-vertical-strip}
Consider $\sin z$ restricted to the strip $S = \{z : -\pi/2 < \re z < \pi/2\}$.
The function $\sin$ maps $S$ conformally onto $\C$ minus the two rays
$(-\infty,-1]$ and $[1,+\infty)$ on the real axis.  This is because:
\begin{itemize}
  \item On the boundary $\re z = \pm\pi/2$, $\sin z$ takes values on the
    real axis with $\abs{\sin z}\ge 1$ (since $\sin(\pm\pi/2+iy)=\pm\cosh y$).
  \item On the real segment $(-\pi/2,\pi/2)$, $\sin$ maps onto $(-1,1)$.
  \item The imaginary axis is mapped onto itself (since $\sin(iy) = i\sinh y$).
\end{itemize}
\end{example}

%% ════════════════════════════════════════════════════════════
\section{Inverse Trigonometric Functions}
\label{sec:inverse-trig}
%% ════════════════════════════════════════════════════════════

\begin{proposition}[Complex arcsine]\label{prop:arcsin}
The equation $\sin w = z$ can be solved for $w$ in terms of $z$:
\[
  \arcsin z \;=\; -i\Log\bigl(iz + (1-z^2)^{1/2}\bigr),
\]
where $(1-z^2)^{1/2}$ denotes a branch of the square root.  This defines a
multi-valued function; the principal branch is obtained by using the principal
$\Log$ and the principal square root.
\end{proposition}

\begin{proof}
Let $w$ satisfy $\sin w = z$, i.e., $\frac{e^{iw}-e^{-iw}}{2i} = z$.
Set $\zeta = e^{iw}$.  Then $\zeta - \zeta^{-1} = 2iz$, so
$\zeta^2 - 2iz\zeta - 1 = 0$.  By the quadratic formula,
\[
  \zeta = iz + \sqrt{1-z^2},
\]
(choosing a square root).  Then $iw = \log\zeta$, so $w = -i\log\zeta$.
\end{proof}

\begin{remark}\label{rem:arccos}
Similarly, $\arccos z = -i\Log\bigl(z + i(1-z^2)^{1/2}\bigr)$ and
$\arctan z = \frac{1}{2i}\Log\!\left(\frac{1+iz}{1-iz}\right)$.
\end{remark}

%% ════════════════════════════════════════════════════════════
\section{Worked Examples}
\label{sec:ch3-examples}
%% ════════════════════════════════════════════════════════════

\begin{example}[Computing $\Log$ for various values]\label{ex:log-values}
\leavevmode
\begin{enumerate}[label=\textup{(\alph*)}]
  \item $\Log(-1) = \ln 1 + i\pi = i\pi$.  (Since $\Arg(-1) = \pi$\ldots\
    but wait: we defined $\Arg$ to lie in $(-\pi,\pi)$.  This means $-1$ lies
    on the branch cut, so $\Log(-1)$ is \emph{not defined} under our convention.
    If we adopt the convention $\Arg\in(-\pi,\pi]$, then $\Log(-1)=i\pi$.)
  \item $\Log(-i) = \ln 1 + i(-\pi/2) = -i\pi/2$.
  \item $\Log(1+i) = \ln\sqrt{2} + i\pi/4 = \tfrac{1}{2}\ln 2 + i\pi/4$.
  \item $\Log(e) = \ln e + i\cdot 0 = 1$.
  \item $\Log(-e^2) = \ln(e^2) + i\pi = 2 + i\pi$
    (if we allow $\Arg(-e^2) = \pi$).
\end{enumerate}
\end{example}

\begin{example}[Computing complex powers]\label{ex:complex-powers-computed}
\leavevmode
\begin{enumerate}[label=\textup{(\alph*)}]
  \item $(-1)^i = e^{i\Log(-1)}$.  Using $\Log(-1)=i\pi$:
    $(-1)^i = e^{i\cdot i\pi} = e^{-\pi}\approx 0.0432$.
  \item $2^i = e^{i\Log 2} = e^{i\ln 2} = \cos(\ln 2)+i\sin(\ln 2)
    \approx 0.7692 + 0.6390\,i$.
  \item $(1+i)^{1+i}$: We have $\Log(1+i) = \frac{1}{2}\ln 2 + i\frac{\pi}{4}$.
    Then
    \begin{align*}
      (1+i)^{1+i} &= e^{(1+i)(\frac{1}{2}\ln 2+i\frac{\pi}{4})}
      = e^{\frac{1}{2}\ln 2-\frac{\pi}{4}+i(\frac{1}{2}\ln 2+\frac{\pi}{4})} \\
      &= e^{\frac{1}{2}\ln 2-\frac{\pi}{4}}\bigl[\cos\bigl(\tfrac{1}{2}\ln 2
        +\tfrac{\pi}{4}\bigr)+i\sin\bigl(\tfrac{1}{2}\ln 2+\tfrac{\pi}{4}\bigr)\bigr].
    \end{align*}
    Numerically, $\abs{(1+i)^{1+i}} = e^{(\ln 2)/2-\pi/4}\approx 0.6440$.
\end{enumerate}
\end{example}

\begin{example}[Solving $e^z = -2$]\label{ex:solve-exp}
We need $z$ such that $e^z = -2$.  Since $-2 = 2e^{i\pi}$, the general
solution is $z = \ln 2 + i(\pi + 2k\pi) = \ln 2 + i(2k+1)\pi$ for $k\in\Z$.
\end{example}

\begin{example}[Solving $\cos z = 3$]\label{ex:solve-cos}
This equation has \emph{no real solution}, since $\abs{\cos x}\le 1$ for
$x\in\R$.  But it has complex solutions:
$\cos z = 3$ means $\frac{e^{iz}+e^{-iz}}{2}=3$, so $e^{2iz}-6e^{iz}+1=0$.
Setting $\zeta = e^{iz}$: $\zeta = 3\pm\sqrt{8} = 3\pm 2\sqrt{2}$.
Both values are positive reals, so $iz = \ln(3\pm 2\sqrt{2})+2k\pi i$,
hence $z = 2k\pi - i\ln(3\pm 2\sqrt{2})$ for $k\in\Z$.

Note: $\ln(3-2\sqrt{2}) = -\ln(3+2\sqrt{2})$ since
$(3-2\sqrt{2})(3+2\sqrt{2})=1$.  So the solutions are
$z = 2k\pi \pm i\ln(3+2\sqrt{2})$ for $k\in\Z$.
\end{example}

%% ════════════════════════════════════════════════════════════
\section{Exercises}
\label{sec:ch3-exercises}
%% ════════════════════════════════════════════════════════════

\begin{exercise}[Exponential equations]\label{exo:exp-equations}
Solve the following equations for $z\in\C$.
\begin{enumerate}[label=\textup{(\alph*)}]
  \item $e^z = 1+i$.
  \item $e^z = -e^3$.
  \item $e^{2z} - 3e^z + 2 = 0$.
  \item $e^z + e^{-z} = 1$.
\end{enumerate}
\end{exercise}

\begin{exercise}[Logarithm computations]\label{exo:log-compute}
Compute all values of $\log z$ and the principal value $\Log z$ for:
\begin{enumerate}[label=\textup{(\alph*)}]
  \item $z = -3$.
  \item $z = 1-i$.
  \item $z = -1-i\sqrt{3}$.
  \item $z = e^{2+i\pi/3}$.
\end{enumerate}
\end{exercise}

\begin{exercise}[Complex powers]\label{exo:complex-powers}
Compute the principal value and all values of:
\begin{enumerate}[label=\textup{(\alph*)}]
  \item $(-1)^{1/3}$.
  \item $(1-i)^{2i}$.
  \item $3^{3i}$.
  \item $i^{1/i}$.
\end{enumerate}
\end{exercise}

\begin{exercise}[Trigonometric equations]\label{exo:trig-equations}
Find all $z\in\C$ satisfying:
\begin{enumerate}[label=\textup{(\alph*)}]
  \item $\sin z = 2$.
  \item $\cos z = i$.
  \item $\tan z = 2i$.
\end{enumerate}
\end{exercise}

\begin{exercise}[Mapping properties]\label{exo:exp-mapping}
\leavevmode
\begin{enumerate}[label=\textup{(\alph*)}]
  \item Describe the image of the strip
    $\{z : 0<\im z<\pi/4\}$ under $\exp$.
  \item Describe the image of the strip
    $\{z : -\pi<\im z<0\}$ under $\exp$.
  \item Show that $w = e^z$ maps the rectangle
    $\{a<\re z<b,\;\alpha<\im z<\beta\}$ (with $\beta-\alpha<2\pi$)
    onto the region $\{e^a<\abs{w}<e^b,\;\alpha<\arg w<\beta\}$.
\end{enumerate}
\end{exercise}

\begin{exercise}[Identities]\label{exo:trig-identities}
Prove the following for all $z\in\C$:
\begin{enumerate}[label=\textup{(\alph*)}]
  \item $\abs{\cos z}^2+\abs{\sin z}^2 = \cosh(2y)+\cos(2x)$
    does \emph{not} hold.  Find the correct identity.
  \item $\cos(2z) = \cos^2 z - \sin^2 z$.
  \item $\sin(z+\pi/2) = \cos z$.
  \item $\cosh^2 z - \sinh^2 z = 1$.
\end{enumerate}
\end{exercise}

\begin{exercise}[Branches of the logarithm]\label{exo:log-branches}
Let $\Omega = \C\setminus[0,+\infty)$ (the plane with the non-negative real
axis removed).
\begin{enumerate}[label=\textup{(\alph*)}]
  \item Show that a branch of $\log$ exists on~$\Omega$ with
    $\im(\log z)\in(0,2\pi)$.
  \item Compute this branch at $z=-1$, $z=i$, $z=-i$.
  \item Is this branch continuous on $\Omega$?  Holomorphic?
\end{enumerate}
\end{exercise}

\begin{exercise}[Continuity across the branch cut]\label{exo:branch-cut-continuity}
Show that $\Log z$ is discontinuous along $(-\infty,0)$: for $x<0$,
\[
  \lim_{\varepsilon\to 0^+}\Log(x+i\varepsilon) -
  \lim_{\varepsilon\to 0^+}\Log(x-i\varepsilon) = 2\pi i.
\]
This ``jump'' of $2\pi i$ is characteristic of the branch cut.
\end{exercise}

\begin{exercise}[The function $z^{1/3}$]\label{exo:cube-root}
Using the principal branch of $\Log$:
\begin{enumerate}[label=\textup{(\alph*)}]
  \item Show that $z^{1/3} = \abs{z}^{1/3}e^{i\Arg(z)/3}$ on
    $\C\setminus(-\infty,0]$.
  \item Determine the image of $\C\setminus(-\infty,0]$ under $z\mapsto z^{1/3}$.
  \item Find all three cube roots of $-8$ and identify which one is
    the principal value.
\end{enumerate}
\end{exercise}

\begin{exercise}[The Joukowski function]\label{exo:joukowski}
The \textbf{Joukowski function} $J(z) = \frac{1}{2}(z+1/z)$ appears in
aerodynamics.
\begin{enumerate}[label=\textup{(\alph*)}]
  \item Show that $J$ maps circles $\abs{z}=r>1$ to ellipses, and circles
    $\abs{z}=1$ to the segment $[-1,1]$.
  \item Show that $J$ is conformal on $\{\abs{z}>1\}$.
  \item Express $\cos z$ in terms of $J$ and $\exp$.
\end{enumerate}
\end{exercise}


%%%%%%%%%%%%%%%%%%%%%%%%%%%%%%%%%%%%%%%%%%%%%%%%%%%%%%%%%%%%
%%
%%  CHAPTER 4: Complex Integration — Cauchy's Theorem
%%
%%%%%%%%%%%%%%%%%%%%%%%%%%%%%%%%%%%%%%%%%%%%%%%%%%%%%%%%%%%%
\chapter{Complex Integration --- Cauchy's Theorem}
\label{chap:complex-integration}

\section*{Introduction}
\addcontentsline{toc}{section}{Introduction}

Integration in the complex plane is a far richer theory than its real
counterpart.  The integral $\int_\gamma f(z)\,\dd z$ of a complex-valued
function along a curve~$\gamma$ depends, in general, on the \emph{path} taken
--- not just on the endpoints.  The spectacular discovery of Cauchy is that for
\emph{holomorphic} functions, this path-dependence disappears under topological
conditions on the domain: the integral around any closed curve that can be
continuously shrunk to a point is \emph{zero}.  This is Cauchy's theorem, the
cornerstone of complex analysis, from which an enormous number of consequences
will flow in the chapters ahead.

%% ════════════════════════════════════════════════════════════
\section{Curves and Contours}
\label{sec:curves-contours}
%% ════════════════════════════════════════════════════════════

\begin{definition}[Parametrized curve]\label{def:parametrized-curve}
A \textbf{parametrized curve}\index{curve!parametrized} (or \textbf{path}) in
$\C$ is a continuous function $\gamma\colon[a,b]\to\C$ where $[a,b]\subseteq\R$
is a closed interval.  We call $\gamma(a)$ the \textbf{initial point} and
$\gamma(b)$ the \textbf{terminal point}.  The \textbf{trace} (or \textbf{image})
of~$\gamma$ is the set $\gamma^* = \gamma([a,b])\subseteq\C$.
\end{definition}

\begin{definition}[Smooth and piecewise smooth curves]\label{def:smooth-curve}
A parametrized curve $\gamma\colon[a,b]\to\C$ is \textbf{smooth}\index{curve!smooth}
if $\gamma$ is continuously differentiable on $[a,b]$ (with one-sided derivatives
at the endpoints) and $\gamma'(t)\ne 0$ for all $t\in[a,b]$.

A curve is \textbf{piecewise smooth}\index{curve!piecewise smooth} (or a
\textbf{contour}\index{contour}) if there exists a partition
$a = t_0 < t_1 < \cdots < t_n = b$ such that $\gamma|_{[t_{k-1},t_k]}$ is
smooth for each $k=1,\dots,n$.
\end{definition}

\begin{definition}[Closed curve and simple curve]\label{def:closed-simple-curve}
A curve $\gamma\colon[a,b]\to\C$ is \textbf{closed}\index{curve!closed} if
$\gamma(a)=\gamma(b)$.  It is \textbf{simple}\index{curve!simple} (or a
\textbf{Jordan curve}\index{Jordan curve}) if $\gamma$ is injective on
$[a,b)$, i.e., $\gamma$ does not cross itself except possibly at the endpoints.
A simple closed curve is called a \textbf{Jordan curve}.
\end{definition}

\begin{definition}[Opposite curve]\label{def:opposite-curve}
Given a curve $\gamma\colon[a,b]\to\C$, the \textbf{opposite curve}
(or \textbf{reverse}) is $\gamma^-\colon[a,b]\to\C$ defined by
$\gamma^-(t) = \gamma(a+b-t)$.  It traverses the same trace in the
opposite direction.
\end{definition}

\begin{definition}[Concatenation]\label{def:concatenation}
If $\gamma_1\colon[a,b]\to\C$ and $\gamma_2\colon[b,c]\to\C$ satisfy
$\gamma_1(b)=\gamma_2(b)$, their \textbf{concatenation} is the curve
$\gamma_1+\gamma_2\colon[a,c]\to\C$ defined by
\[
  (\gamma_1+\gamma_2)(t) = \begin{cases}
    \gamma_1(t) & \text{if } t\in[a,b], \\
    \gamma_2(t) & \text{if } t\in[b,c].
  \end{cases}
\]
\end{definition}

\begin{example}[Standard parametrizations]\label{ex:standard-param}
\leavevmode
\begin{enumerate}[label=\textup{(\alph*)}]
  \item \textbf{Line segment.} The segment from $z_1$ to $z_2$:
    $\gamma(t)=(1-t)z_1+tz_2$, $t\in[0,1]$.  Then $\gamma'(t)=z_2-z_1$.
  \item \textbf{Circle.} The circle of center $z_0$ and radius $r>0$,
    traversed counterclockwise:
    $\gamma(t) = z_0 + re^{it}$, $t\in[0,2\pi]$.
    Then $\gamma'(t)=ire^{it}$.
  \item \textbf{Circular arc.} The arc from angle $\alpha$ to $\beta$:
    $\gamma(t) = z_0 + re^{it}$, $t\in[\alpha,\beta]$.
\end{enumerate}
\end{example}

\begin{center}
\begin{tikzpicture}[scale=1.0, >=Stealth]
  %% Example 1: line segment
  \begin{scope}[shift={(-5,0)}]
    \draw[->] (-0.5,0) -- (3,0) node[right, font=\tiny]{$\re$};
    \draw[->] (0,-0.5) -- (0,2.5) node[above, font=\tiny]{$\im$};
    \fill (0.3,0.3) circle (1.5pt) node[below left, font=\small]{$z_1$};
    \fill (2.2,1.8) circle (1.5pt) node[above right, font=\small]{$z_2$};
    \draw[thick, blue!70!black, decoration={markings,
      mark=at position 0.5 with {\arrow{Stealth}}}, postaction={decorate}]
      (0.3,0.3) -- (2.2,1.8);
    \node[below, font=\small\bfseries] at (1.2,-0.8) {(a) Segment};
  \end{scope}
  %% Example 2: circle
  \begin{scope}[shift={(0,0.8)}]
    \draw[->] (-1.8,0) -- (1.8,0) node[right, font=\tiny]{$\re$};
    \draw[->] (0,-1.8) -- (0,1.8) node[above, font=\tiny]{$\im$};
    \fill (0,0) circle (1.2pt) node[below right, font=\small]{$z_0$};
    \draw[thick, blue!70!black, decoration={markings,
      mark=at position 0.12 with {\arrow{Stealth}},
      mark=at position 0.37 with {\arrow{Stealth}},
      mark=at position 0.62 with {\arrow{Stealth}},
      mark=at position 0.87 with {\arrow{Stealth}}}, postaction={decorate}]
      (0,0) circle (1.2);
    \draw[dashed, gray] (0,0) -- node[above, font=\tiny]{$r$} (0.85,0.85);
    \node[below, font=\small\bfseries] at (0,-2.1) {(b) Circle};
  \end{scope}
  %% Example 3: piecewise
  \begin{scope}[shift={(5,0)}]
    \draw[->] (-0.5,0) -- (3,0) node[right, font=\tiny]{$\re$};
    \draw[->] (0,-0.5) -- (0,2.5) node[above, font=\tiny]{$\im$};
    \fill (0.2,0.2) circle (1.5pt) node[below, font=\small]{$z_1$};
    \fill (1.5,2) circle (1.5pt) node[above, font=\small]{$z_2$};
    \fill (2.5,0.5) circle (1.5pt) node[right, font=\small]{$z_3$};
    \draw[thick, blue!70!black, decoration={markings,
      mark=at position 0.3 with {\arrow{Stealth}},
      mark=at position 0.7 with {\arrow{Stealth}}}, postaction={decorate}]
      (0.2,0.2) -- (1.5,2) -- (2.5,0.5);
    \node[below, font=\small\bfseries] at (1.3,-0.8) {(c) Piecewise};
  \end{scope}
\end{tikzpicture}
\captionof{figure}{Standard types of contours: (a) a line segment,
(b) a circle traversed counterclockwise, (c) a piecewise linear path.}
\label{fig:standard-contours}
\end{center}

%% ════════════════════════════════════════════════════════════
\section{The Complex Line Integral}
\label{sec:complex-line-integral}
%% ════════════════════════════════════════════════════════════

\begin{definition}[Complex line integral]\label{def:complex-line-integral}
Let $\gamma\colon[a,b]\to\C$ be a piecewise smooth curve, and let
$f\colon\gamma^*\to\C$ be continuous.  The \textbf{complex line integral}%
\index{line integral!complex} of $f$ along~$\gamma$ is
\[
  \int_\gamma f(z)\,\dd z
  \;=\; \int_a^b f\bigl(\gamma(t)\bigr)\,\gamma'(t)\,\dd t.
\]
\end{definition}

\begin{remark}\label{rem:integral-well-defined}
The right-hand side is a standard Riemann integral of a complex-valued function
of a real variable: if $g\colon[a,b]\to\C$ is continuous, then
$\int_a^b g(t)\,\dd t = \int_a^b \re g(t)\,\dd t + i\int_a^b \im g(t)\,\dd t$.
\end{remark}

\begin{proposition}[Elementary properties]\label{prop:integral-properties}
Let $f,g$ be continuous on the traces of the indicated curves, and let
$\alpha\in\C$.
\begin{enumerate}[label=\textup{(\roman*)}]
  \item \textbf{Linearity.}
    $\displaystyle\int_\gamma(\alpha f+g)\,\dd z
    = \alpha\int_\gamma f\,\dd z + \int_\gamma g\,\dd z$.
  \item \textbf{Reversal.}
    $\displaystyle\int_{\gamma^-} f\,\dd z = -\int_\gamma f\,\dd z$.
  \item \textbf{Concatenation.} If $\gamma=\gamma_1+\gamma_2$, then
    $\displaystyle\int_\gamma f\,\dd z
    = \int_{\gamma_1}f\,\dd z + \int_{\gamma_2}f\,\dd z$.
  \item \textbf{Reparametrization invariance.} The integral depends only on the
    oriented curve, not on the particular parametrization.
\end{enumerate}
\end{proposition}

\begin{proof}
All parts follow from the corresponding properties of the real Riemann integral.

\textbf{(ii)} Let $\gamma\colon[a,b]\to\C$.  Setting $s = a+b-t$:
\[
  \int_{\gamma^-}f\,\dd z
  = \int_a^b f(\gamma(a+b-t))\cdot(-\gamma'(a+b-t))\,\dd t
  = -\int_a^b f(\gamma(s))\gamma'(s)\,\dd s
  = -\int_\gamma f\,\dd z.
\]
The other parts are left as exercises.
\end{proof}

\begin{theorem}[ML-inequality]\label{thm:ML-inequality}
If $\gamma$ is a piecewise smooth curve of length $L(\gamma)$ and
$\abs{f(z)}\le M$ for all $z\in\gamma^*$, then
\[
  \abs[\bigg]{\int_\gamma f(z)\,\dd z} \;\le\; M\cdot L(\gamma),
\]
where $L(\gamma) = \int_a^b \abs{\gamma'(t)}\,\dd t$ is the
\textbf{arc length}\index{arc length} of $\gamma$.
\end{theorem}

\begin{proof}
Using the standard estimate for integrals of complex-valued functions:
\[
  \abs[\bigg]{\int_\gamma f(z)\,\dd z}
  = \abs[\bigg]{\int_a^b f(\gamma(t))\gamma'(t)\,\dd t}
  \le \int_a^b \abs{f(\gamma(t))}\,\abs{\gamma'(t)}\,\dd t
  \le M\int_a^b\abs{\gamma'(t)}\,\dd t
  = ML(\gamma). \qedhere
\]
\end{proof}

\subsection{First computations}

\begin{example}[Integral of $z^n$ around a circle]\label{ex:integral-zn-circle}
Let $\gamma(t) = z_0+re^{it}$, $t\in[0,2\pi]$, and let $n\in\Z$.
We compute $\int_\gamma (z-z_0)^n\,\dd z$.

We have $z - z_0 = re^{it}$ and $\dd z = ire^{it}\,\dd t$, so
\[
  \int_\gamma (z-z_0)^n\,\dd z
  = \int_0^{2\pi}r^n e^{int}\cdot ire^{it}\,\dd t
  = ir^{n+1}\int_0^{2\pi}e^{i(n+1)t}\,\dd t.
\]

\textbf{Case $n\ne -1$:}
\[
  \int_0^{2\pi}e^{i(n+1)t}\,\dd t
  = \left[\frac{e^{i(n+1)t}}{i(n+1)}\right]_0^{2\pi}
  = \frac{e^{2\pi i(n+1)}-1}{i(n+1)} = 0,
\]
since $e^{2\pi i(n+1)}=1$.

\textbf{Case $n=-1$:}
\[
  \int_\gamma\frac{\dd z}{z-z_0}
  = ir^0\int_0^{2\pi}e^{0}\,\dd t = i\cdot 2\pi = 2\pi i.
\]

Therefore:
\begin{equation}\label{eq:integral-zn-circle}
  \boxed{\int_\gamma(z-z_0)^n\,\dd z
  = \begin{cases} 2\pi i & \text{if } n=-1, \\ 0 & \text{if } n\in\Z,\;n\ne -1. \end{cases}}
\end{equation}
This result is fundamental and will be used repeatedly.
\end{example}

\begin{example}[Integral of $\conj{z}$ around the unit circle]\label{ex:integral-zbar}
Let $\gamma(t) = e^{it}$, $t\in[0,2\pi]$.  Note that $\conj{z}$ is
\emph{not} holomorphic.  We compute:
\[
  \int_\gamma\conj{z}\,\dd z
  = \int_0^{2\pi}\conj{e^{it}}\cdot ie^{it}\,\dd t
  = \int_0^{2\pi}e^{-it}\cdot ie^{it}\,\dd t
  = \int_0^{2\pi}i\,\dd t = 2\pi i.
\]
\end{example}

\begin{example}[Integral along different paths]\label{ex:path-dependence}
Compute $\int_\gamma z\,\dd z$ where $\gamma$ goes from $0$ to $1+i$.

\textbf{Path 1:} the straight line $\gamma_1(t) = (1+i)t$, $t\in[0,1]$.
$\gamma_1'(t)=1+i$, so
\[
  \int_{\gamma_1}z\,\dd z = \int_0^1(1+i)t\cdot(1+i)\,\dd t
  = (1+i)^2\int_0^1 t\,\dd t = 2i\cdot\tfrac{1}{2} = i.
\]

\textbf{Path 2:} first along the real axis from $0$ to $1$, then vertically
from $1$ to $1+i$.  Let $\gamma_2 = \alpha+\beta$ with $\alpha(t)=t$,
$t\in[0,1]$, and $\beta(t)=1+it$, $t\in[0,1]$.
\begin{align*}
  \int_\alpha z\,\dd z &= \int_0^1 t\cdot 1\,\dd t = \tfrac{1}{2}, \\
  \int_\beta z\,\dd z &= \int_0^1(1+it)\cdot i\,\dd t
  = i\int_0^1\dd t - \int_0^1 t\,\dd t = i - \tfrac{1}{2}.
\end{align*}
So $\int_{\gamma_2}z\,\dd z = \frac{1}{2}+i-\frac{1}{2} = i$.

Both paths give the same answer.  This is no accident: $f(z)=z$ is entire and
has a primitive $F(z)=z^2/2$.
\end{example}

%% ════════════════════════════════════════════════════════════
\section{Primitives and Path Independence}
\label{sec:primitives}
%% ════════════════════════════════════════════════════════════

\begin{definition}[Primitive]\label{def:primitive}
Let $\Omega\subseteq\C$ be open and $f\colon\Omega\to\C$ continuous.
A \textbf{primitive}\index{primitive} (or \textbf{antiderivative}) of $f$
on~$\Omega$ is a holomorphic function $F\colon\Omega\to\C$ such that
$F'(z)=f(z)$ for all $z\in\Omega$.
\end{definition}

\begin{theorem}[Fundamental theorem of contour integration]%
\label{thm:fundamental-contour}
Let $\Omega\subseteq\C$ be open, $f\colon\Omega\to\C$ continuous, and suppose
$f$ has a primitive $F$ on~$\Omega$.  Then for every piecewise smooth curve
$\gamma\colon[a,b]\to\Omega$,
\begin{equation}\label{eq:FTC-complex}
  \int_\gamma f(z)\,\dd z \;=\; F(\gamma(b)) - F(\gamma(a)).
\end{equation}
In particular, if $\gamma$ is closed then $\int_\gamma f(z)\,\dd z = 0$.
\end{theorem}

\begin{proof}
By the chain rule (for functions of a real variable composed with holomorphic
functions),
\[
  \frac{\dd}{\dd t}\bigl[F(\gamma(t))\bigr]
  = F'(\gamma(t))\cdot\gamma'(t)
  = f(\gamma(t))\cdot\gamma'(t).
\]
Hence
\[
  \int_\gamma f(z)\,\dd z
  = \int_a^b f(\gamma(t))\gamma'(t)\,\dd t
  = \int_a^b \frac{\dd}{\dd t}\bigl[F(\gamma(t))\bigr]\,\dd t
  = F(\gamma(b))-F(\gamma(a)). \qedhere
\]
\end{proof}

\begin{corollary}\label{cor:primitive-path-independence}
If $f$ has a primitive on a domain $\Omega$, then the integral
$\int_\gamma f\,\dd z$ depends only on the endpoints of~$\gamma$, not on
the particular path within~$\Omega$.
\end{corollary}

\begin{example}[Using the fundamental theorem]\label{ex:FTC-application}
\leavevmode
\begin{enumerate}[label=\textup{(\alph*)}]
  \item $\displaystyle\int_\gamma z^3\,\dd z = \Bigl[\frac{z^4}{4}\Bigr]_0^{1+i}
    = \frac{(1+i)^4}{4} = \frac{-4}{4} = -1$
    for any path from $0$ to $1+i$, since $z^4/4$ is a primitive.
  \item $\displaystyle\int_\gamma e^z\,\dd z = \bigl[e^z\bigr]_0^{i\pi}
    = e^{i\pi}-e^0 = -1-1 = -2$
    for any path from $0$ to $i\pi$.
  \item $\displaystyle\int_\gamma\frac{\dd z}{z}$ from $1$ to $-1$
    \textbf{depends on the path}: through the upper half-plane gives $i\pi$,
    through the lower gives $-i\pi$.  There is no primitive of $1/z$ on
    all of $\C^*$.
\end{enumerate}
\end{example}

\begin{theorem}[Characterization of existence of primitives]%
\label{thm:primitive-characterization}
Let $\Omega\subseteq\C$ be a domain and $f\colon\Omega\to\C$ continuous.
The following are equivalent:
\begin{enumerate}[label=\textup{(\roman*)}]
  \item $f$ has a primitive on~$\Omega$.
  \item $\int_\gamma f\,\dd z = 0$ for every closed piecewise smooth curve
    $\gamma$ in~$\Omega$.
  \item $\int_\gamma f\,\dd z$ depends only on the endpoints for every
    piecewise smooth curve $\gamma$ in~$\Omega$.
\end{enumerate}
\end{theorem}

\begin{proof}
(i)$\Rightarrow$(ii): Theorem~\ref{thm:fundamental-contour}.

(ii)$\Rightarrow$(iii): If $\gamma_1$ and $\gamma_2$ have the same endpoints,
then $\gamma_1+\gamma_2^-$ is closed, so $0 = \int_{\gamma_1+\gamma_2^-}f\,\dd z
= \int_{\gamma_1}f\,\dd z - \int_{\gamma_2}f\,\dd z$.

(iii)$\Rightarrow$(i): Fix $z_0\in\Omega$ and define
$F(z) = \int_{\gamma_z}f(\zeta)\,\dd\zeta$
where $\gamma_z$ is any piecewise smooth path from $z_0$ to $z$ in~$\Omega$.
By (iii), $F$ is well-defined.  For $h$ small enough,
\[
  F(z+h)-F(z) = \int_{[z,z+h]}f(\zeta)\,\dd\zeta
  = f(z)\cdot h + \int_{[z,z+h]}\bigl(f(\zeta)-f(z)\bigr)\,\dd\zeta.
\]
By continuity of $f$, the last integral is $o(\abs{h})$, so
$F'(z) = \lim_{h\to 0}\frac{F(z+h)-F(z)}{h} = f(z)$.
\end{proof}

%% ════════════════════════════════════════════════════════════
\section{Cauchy's Theorem --- Statement and Proof}
\label{sec:cauchy-theorem}
%% ════════════════════════════════════════════════════════════

Cauchy's theorem asserts that the integral of a holomorphic function around a
closed contour is zero, under appropriate topological conditions on the domain
or the contour.  We present several versions, starting with the simplest.

\subsection{Simply connected domains}

\begin{definition}[Simply connected domain]\label{def:simply-connected}
A domain $\Omega\subseteq\C$ is \textbf{simply connected}%
\index{simply connected domain} if every closed curve in~$\Omega$ can be
continuously deformed (within~$\Omega$) to a point.  Equivalently,
$\Omega$ has ``no holes'': $\C\setminus\Omega$ is connected in
$\C\cup\{\infty\}$ (the Riemann sphere).
\end{definition}

\begin{center}
\begin{tikzpicture}[scale=0.8, >=Stealth]
  %% Simply connected
  \begin{scope}[shift={(-4,0)}]
    \draw[thick, blue!70!black, fill=blue!8]
      plot[smooth cycle, tension=0.7] coordinates
      {(-2,0) (-1.2,1.5) (0.8,1.8) (2,0.8) (2.2,-0.5) (1,-1.5) (-0.5,-1.6) (-2,-0.8)};
    \node[blue!70!black, font=\bfseries] at (0,2.3) {Simply connected};
    \node[font=\small] at (0,0) {no holes};
  \end{scope}
  %% Multiply connected
  \begin{scope}[shift={(4,0)}]
    \fill[blue!8]
      plot[smooth cycle, tension=0.7] coordinates
      {(-2,0) (-1.2,1.5) (0.8,1.8) (2,0.8) (2.2,-0.5) (1,-1.5) (-0.5,-1.6) (-2,-0.8)};
    \fill[white] (0,0) circle (0.6);
    \draw[thick, blue!70!black]
      plot[smooth cycle, tension=0.7] coordinates
      {(-2,0) (-1.2,1.5) (0.8,1.8) (2,0.8) (2.2,-0.5) (1,-1.5) (-0.5,-1.6) (-2,-0.8)};
    \draw[thick, red!70!black] (0,0) circle (0.6);
    \node[red!70!black, font=\small] at (0,0) {hole};
    \node[blue!70!black, font=\bfseries] at (0,2.3) {Not simply connected};
  \end{scope}
\end{tikzpicture}
\captionof{figure}{A simply connected domain (left) versus a multiply connected
domain with a hole (right).}
\label{fig:simply-multiply-connected}
\end{center}

\begin{example}[Simply connected domains]\label{ex:simply-connected}
The following are simply connected:
\begin{itemize}
  \item $\C$ itself, and any open disk $D(z_0,R)$.
  \item Any convex open set (e.g., a half-plane, a strip).
  \item Any star-shaped domain: a domain $\Omega$ such that there exists
    $z_0\in\Omega$ with $[z_0,z]\subseteq\Omega$ for all $z\in\Omega$.
  \item $\C\setminus(-\infty,0]$ (the slit plane).
\end{itemize}
The following are \emph{not} simply connected:
\begin{itemize}
  \item $\C^* = \C\setminus\{0\}$ (a loop around the origin cannot be shrunk).
  \item An annulus $\{z : r<\abs{z}<R\}$.
\end{itemize}
\end{example}

\subsection{Cauchy's theorem for triangles (Goursat's lemma)}

We begin with a version that requires only holomorphicity --- no assumption on
the continuity of the derivative.

\begin{theorem}[Goursat's lemma]\label{thm:goursat}
Let $\Omega\subseteq\C$ be open and $f\colon\Omega\to\C$ holomorphic.  If
$\Delta$ is a closed triangle (together with its interior) contained
in~$\Omega$, and $\partial\Delta$ denotes its boundary traversed
counterclockwise, then
\[
  \int_{\partial\Delta} f(z)\,\dd z \;=\; 0.
\]
\end{theorem}

\begin{proof}
We use a bisection argument.  Let $\Delta^{(0)} = \Delta$ and let
$I_0 = \int_{\partial\Delta^{(0)}}f(z)\,\dd z$.  Join the midpoints of the
three sides of $\Delta^{(0)}$ to form four congruent triangles
$\Delta_1,\Delta_2,\Delta_3,\Delta_4$.  Orienting each sub-triangle
counterclockwise, the integrals over the internal edges cancel in pairs, so
\[
  I_0 = \sum_{j=1}^{4}\int_{\partial\Delta_j}f(z)\,\dd z.
\]
By the triangle inequality, at least one of the four sub-triangles, call it
$\Delta^{(1)}$, satisfies
\[
  \abs[\bigg]{\int_{\partial\Delta^{(1)}}f(z)\,\dd z} \;\ge\; \frac{\abs{I_0}}{4}.
\]

\begin{center}
\begin{tikzpicture}[scale=1.2]
  %% Original triangle
  \coordinate (A) at (0,0);
  \coordinate (B) at (4,0);
  \coordinate (C) at (1.5,3);
  %% Midpoints
  \coordinate (MAB) at ($(A)!0.5!(B)$);
  \coordinate (MBC) at ($(B)!0.5!(C)$);
  \coordinate (MAC) at ($(A)!0.5!(C)$);
  %% Outer triangle
  \draw[thick, blue!70!black] (A) -- (B) -- (C) -- cycle;
  %% Inner lines
  \draw[thick, red!60!black, dashed] (MAB) -- (MBC) -- (MAC) -- cycle;
  %% Labels
  \node[below left, font=\small] at (A) {$A$};
  \node[below right, font=\small] at (B) {$B$};
  \node[above, font=\small] at (C) {$C$};
  %% Highlight one sub-triangle
  \fill[red!15, opacity=0.5] (MAC) -- (MAB) -- (MBC) -- cycle;
  \node[font=\small, red!60!black] at ($(MAC)!0.33!(MAB)!0.33!(MBC)$) {$\Delta^{(1)}$};
\end{tikzpicture}
\captionof{figure}{Bisection in Goursat's proof: the triangle $\Delta^{(0)}$
is subdivided into four; we select the sub-triangle with the largest integral.}
\label{fig:goursat-bisection}
\end{center}

Repeating inductively, we obtain a nested sequence of triangles
$\Delta^{(0)}\supseteq\Delta^{(1)}\supseteq\Delta^{(2)}\supseteq\cdots$ with
\[
  \abs[\bigg]{\int_{\partial\Delta^{(n)}}f(z)\,\dd z} \ge \frac{\abs{I_0}}{4^n},
  \quad
  L(\partial\Delta^{(n)}) = \frac{L}{2^n},
  \quad
  \diam(\Delta^{(n)}) = \frac{d}{2^n},
\]
where $L = L(\partial\Delta^{(0)})$ and $d = \diam(\Delta^{(0)})$.

By the nested compact sets theorem, $\bigcap_n\Delta^{(n)} = \{z_*\}$ for
some $z_*\in\Omega$.  Since $f$ is holomorphic at $z_*$,
\[
  f(z) = f(z_*) + f'(z_*)(z-z_*) + \varphi(z)(z-z_*),
\]
where $\varphi(z)\to 0$ as $z\to z_*$.  The function $g(z)=f(z_*)+f'(z_*)(z-z_*)$
is a polynomial and thus has a primitive, so
$\int_{\partial\Delta^{(n)}}g(z)\,\dd z=0$.  Therefore
\[
  \int_{\partial\Delta^{(n)}}f(z)\,\dd z
  = \int_{\partial\Delta^{(n)}}\varphi(z)(z-z_*)\,\dd z.
\]
For any $\varepsilon>0$, there exists $N$ such that for $n\ge N$,
$\abs{\varphi(z)}<\varepsilon$ on $\Delta^{(n)}$.
Then by the ML-inequality:
\[
  \abs[\bigg]{\int_{\partial\Delta^{(n)}}f(z)\,\dd z}
  \le \varepsilon\cdot\frac{d}{2^n}\cdot\frac{L}{2^n}
  = \frac{\varepsilon\,dL}{4^n}.
\]
Combined with the lower bound:
\[
  \frac{\abs{I_0}}{4^n}
  \le \frac{\varepsilon\,dL}{4^n},
\]
so $\abs{I_0}\le\varepsilon\,dL$.  Since $\varepsilon>0$ is arbitrary,
$I_0 = 0$.
\end{proof}

\subsection{Cauchy's theorem for convex and star-shaped domains}

\begin{theorem}[Cauchy's theorem for star-shaped domains]%
\label{thm:cauchy-star-shaped}
Let $\Omega\subseteq\C$ be a star-shaped domain, and let $f\colon\Omega\to\C$
be holomorphic.  Then:
\begin{enumerate}[label=\textup{(\roman*)}]
  \item $f$ has a primitive $F$ on~$\Omega$.
  \item $\displaystyle\int_\gamma f(z)\,\dd z = 0$ for every closed
    piecewise smooth curve $\gamma$ in~$\Omega$.
\end{enumerate}
\end{theorem}

\begin{proof}
Let $z_0$ be a star center of $\Omega$, i.e., $[z_0,z]\subseteq\Omega$ for
all $z\in\Omega$.  Define
\[
  F(z) = \int_{[z_0,z]} f(\zeta)\,\dd\zeta
  = \int_0^1 f\bigl(z_0+t(z-z_0)\bigr)(z-z_0)\,\dd t.
\]
We show $F'(z)=f(z)$.  For $h$ small, consider the triangle with vertices
$z_0, z, z+h$.  By Goursat's lemma (Theorem~\ref{thm:goursat}),
\[
  \int_{[z_0,z]}f\,\dd\zeta + \int_{[z,z+h]}f\,\dd\zeta
  + \int_{[z+h,z_0]}f\,\dd\zeta = 0.
\]
Hence $F(z+h)-F(z) = \int_{[z,z+h]}f(\zeta)\,\dd\zeta$.  Since
$\int_{[z,z+h]}\dd\zeta = h$, we can write
\[
  F(z+h)-F(z)-f(z)h = \int_{[z,z+h]}\bigl(f(\zeta)-f(z)\bigr)\,\dd\zeta.
\]
By continuity of $f$ and the ML-inequality,
$\abs{F(z+h)-F(z)-f(z)h}\le\sup_{\zeta\in[z,z+h]}\abs{f(\zeta)-f(z)}\cdot\abs{h}
= o(\abs{h})$.  Thus $F'(z)=f(z)$.  Part (ii) follows from
Theorem~\ref{thm:fundamental-contour}.
\end{proof}

\subsection{Cauchy's theorem via Green's theorem}

When $f'$ is continuous (not just existent), one can give a short proof using
Green's theorem from multivariable calculus.  Although this assumption is
slightly stronger, the argument is instructive.

\begin{theorem}[Green's theorem --- recalled]\label{thm:green-recalled}
Let $\Omega\subseteq\R^2$ be a bounded domain with piecewise smooth boundary
$\partial\Omega$ oriented counterclockwise.  If $P,Q\colon\overline{\Omega}\to\R$
are $C^1$, then
\[
  \oint_{\partial\Omega}(P\,\dd x + Q\,\dd y)
  = \iint_\Omega\Bigl(\frac{\partial Q}{\partial x}
    - \frac{\partial P}{\partial y}\Bigr)\,\dd x\,\dd y.
\]
\end{theorem}

\begin{theorem}[Cauchy's theorem via Green's theorem]%
\label{thm:cauchy-via-green}
Let $\Omega\subseteq\C$ be a bounded domain with piecewise smooth boundary
$\partial\Omega$ oriented positively, and let $f$ be holomorphic on an open set
containing $\overline{\Omega}$, with $f'$ continuous.  Then
\[
  \oint_{\partial\Omega} f(z)\,\dd z \;=\; 0.
\]
\end{theorem}

\begin{proof}
Write $f = u+iv$ and $z = x+iy$, so $\dd z = \dd x + i\,\dd y$.  Then
\[
  f(z)\,\dd z = (u+iv)(\dd x+i\,\dd y)
  = (u\,\dd x - v\,\dd y) + i(v\,\dd x + u\,\dd y).
\]
By Green's theorem applied to each real part:
\begin{align*}
  \re\oint_{\partial\Omega}f\,\dd z
  &= \oint_{\partial\Omega}(u\,\dd x-v\,\dd y)
  = \iint_\Omega\Bigl(-\frac{\partial v}{\partial x}
    -\frac{\partial u}{\partial y}\Bigr)\dd x\,\dd y, \\
  \im\oint_{\partial\Omega}f\,\dd z
  &= \oint_{\partial\Omega}(v\,\dd x+u\,\dd y)
  = \iint_\Omega\Bigl(\frac{\partial u}{\partial x}
    -\frac{\partial v}{\partial y}\Bigr)\dd x\,\dd y.
\end{align*}
By the Cauchy--Riemann equations ($u_x=v_y$ and $u_y=-v_x$), both integrands
vanish identically.
\end{proof}

\subsection{The homotopy version}

The most general version of Cauchy's theorem involves the topological notion
of homotopy.

\begin{definition}[Homotopy of curves]\label{def:homotopy}
Let $\Omega\subseteq\C$ be open, and let $\gamma_0,\gamma_1\colon[a,b]\to\Omega$
be two closed curves (with the same base point, or free homotopy).  We say
$\gamma_0$ and $\gamma_1$ are \textbf{homotopic}\index{homotopy} in~$\Omega$
if there exists a continuous function $H\colon[a,b]\times[0,1]\to\Omega$
such that
\[
  H(t,0) = \gamma_0(t), \quad
  H(t,1) = \gamma_1(t), \quad
  H(a,s) = H(b,s) \quad\text{for all } t\in[a,b],\;s\in[0,1].
\]
The map $H$ is called a \textbf{homotopy} between $\gamma_0$ and $\gamma_1$.
A closed curve homotopic to a constant (a point) is called
\textbf{null-homotopic}\index{null-homotopic}.
\end{definition}

\begin{theorem}[Cauchy's theorem --- homotopy version]\label{thm:cauchy-homotopy}
Let $\Omega\subseteq\C$ be open and $f\colon\Omega\to\C$ holomorphic.
If $\gamma_0$ and $\gamma_1$ are homotopic closed piecewise smooth curves
in~$\Omega$, then
\[
  \int_{\gamma_0}f(z)\,\dd z \;=\; \int_{\gamma_1}f(z)\,\dd z.
\]
In particular, if $\gamma$ is null-homotopic in~$\Omega$, then
$\int_\gamma f(z)\,\dd z = 0$.
\end{theorem}

\begin{proof}[Proof sketch]
The continuous homotopy $H$ maps the rectangle $[a,b]\times[0,1]$ into $\Omega$.
By compactness, the image $H([a,b]\times[0,1])$ has a positive distance to
$\C\setminus\Omega$. We can therefore approximate $H$ by a ``piecewise linear''
homotopy through a fine grid.  Each elementary deformation (moving the curve
across one small rectangle in the grid) changes the integral by the integral
around a small rectangle contained in~$\Omega$.  By Goursat's lemma
(applied to triangulate each rectangle), each such integral is zero.
Hence the total integral is unchanged by the homotopy.

A complete proof can be found in standard references; it relies on a careful
approximation argument and the compactness of $[a,b]\times[0,1]$.
\end{proof}

\begin{corollary}[Cauchy's theorem for simply connected domains]%
\label{cor:cauchy-simply-connected}
Let $\Omega\subseteq\C$ be a simply connected domain and $f\colon\Omega\to\C$
holomorphic.  Then:
\begin{enumerate}[label=\textup{(\roman*)}]
  \item $\displaystyle\oint_\gamma f(z)\,\dd z = 0$ for every closed piecewise smooth
    curve $\gamma$ in $\Omega$.
  \item $f$ has a primitive on $\Omega$.
\end{enumerate}
\end{corollary}

\begin{proof}
\textbf{(i)} In a simply connected domain, every closed curve is null-homotopic.
By Theorem~\ref{thm:cauchy-homotopy}, the integral is zero.

\textbf{(ii)} By Theorem~\ref{thm:primitive-characterization},
(i) implies the existence of a primitive.
\end{proof}

%% ════════════════════════════════════════════════════════════
\section{The Winding Number}
\label{sec:winding-number}
%% ════════════════════════════════════════════════════════════

\begin{definition}[Winding number]\label{def:winding-number}
Let $\gamma\colon[a,b]\to\C$ be a closed piecewise smooth curve and let
$w\notin\gamma^*$.  The \textbf{winding number}\index{winding number}
(or \textbf{index}) of $\gamma$ with respect to $w$ is
\[
  \ind(\gamma,w) \;=\; n(\gamma,w) \;=\;
  \frac{1}{2\pi i}\oint_\gamma\frac{\dd z}{z-w}.
\]
\end{definition}

\begin{theorem}[Properties of the winding number]\label{thm:winding-number-properties}
\leavevmode
\begin{enumerate}[label=\textup{(\roman*)}]
  \item $\ind(\gamma,w)\in\Z$ for every $w\notin\gamma^*$.
  \item The function $w\mapsto\ind(\gamma,w)$ is continuous on
    $\C\setminus\gamma^*$, hence locally constant (since it is integer-valued
    and continuous).
  \item $\ind(\gamma,w) = 0$ for all $w$ in the unbounded component of
    $\C\setminus\gamma^*$.
  \item If $\gamma(t) = z_0+re^{int}$, $t\in[0,2\pi]$ (the circle of center
    $z_0$ and radius $r$, traversed $n$ times), then
    $\ind(\gamma,z_0) = n$.
\end{enumerate}
\end{theorem}

\begin{proof}
\textbf{(i)} Define $\varphi(t) = \int_a^t\frac{\gamma'(s)}{\gamma(s)-w}\,\dd s$
for $t\in[a,b]$.  Then $\varphi'(t) = \frac{\gamma'(t)}{\gamma(t)-w}$
(at points of smoothness).  Consider
\[
  h(t) \;=\; (\gamma(t)-w)\,e^{-\varphi(t)}.
\]
Then $h'(t) = \gamma'(t)e^{-\varphi(t)} - (\gamma(t)-w)\varphi'(t)e^{-\varphi(t)}
= \gamma'(t)e^{-\varphi(t)} - \gamma'(t)e^{-\varphi(t)} = 0$,
so $h$ is constant: $h(t)=h(a)=(\gamma(a)-w)e^{0}=\gamma(a)-w$.
Therefore $(\gamma(t)-w)e^{-\varphi(t)}=\gamma(a)-w$ for all $t$.
At $t=b$, using $\gamma(b)=\gamma(a)$:
$(\gamma(a)-w)e^{-\varphi(b)} = \gamma(a)-w$, so $e^{-\varphi(b)}=1$
(since $\gamma(a)\ne w$).  This means $\varphi(b)\in 2\pi i\Z$, i.e.,
$\frac{1}{2\pi i}\varphi(b) = \ind(\gamma,w)\in\Z$.

\textbf{(ii)} For $w_1$ close to $w_0$ (both outside $\gamma^*$),
\[
  \abs{\ind(\gamma,w_1)-\ind(\gamma,w_0)}
  = \frac{1}{2\pi}\abs[\bigg]{\int_\gamma\Bigl(\frac{1}{z-w_1}-\frac{1}{z-w_0}\Bigr)\dd z}
  = \frac{\abs{w_1-w_0}}{2\pi}\abs[\bigg]{\int_\gamma\frac{\dd z}{(z-w_0)(z-w_1)}},
\]
which tends to $0$ as $w_1\to w_0$ by the ML-inequality.

\textbf{(iii)} For $\abs{w}$ sufficiently large (larger than $\max_{t}\abs{\gamma(t)}$),
we can bound $\abs{1/(z-w)}\le 1/(\abs{w}-\max\abs{\gamma})$ and the integral
tends to $0$ as $\abs{w}\to\infty$.  Since $\ind(\gamma,w)$ is integer-valued
and continuous on the unbounded component, it must be $0$ everywhere on that
component.

\textbf{(iv)} We already computed this in Example~\ref{ex:integral-zn-circle}
(for $n=1$; the general case is similar).
\end{proof}

\begin{center}
\begin{tikzpicture}[scale=0.85, >=Stealth]
  %% A closed curve
  \draw[thick, blue!70!black, decoration={markings,
    mark=at position 0.1 with {\arrow{Stealth}},
    mark=at position 0.35 with {\arrow{Stealth}},
    mark=at position 0.6 with {\arrow{Stealth}},
    mark=at position 0.85 with {\arrow{Stealth}}},
    postaction={decorate}]
    plot[smooth cycle, tension=0.8] coordinates
    {(0,0) (2,1.5) (3.5,0.5) (3,-1) (1.5,-1.8) (-0.5,-1)};
  %% Points
  \fill[red!80!black] (1.5,0) circle (2pt) node[above right, font=\small]{$w_1$};
  \fill[green!50!black] (5,0) circle (2pt) node[right, font=\small]{$w_2$};
  \fill[orange!80!black] (-1.5,0.5) circle (2pt) node[left, font=\small]{$w_3$};
  %% Labels
  \node[blue!70!black, font=\small] at (3.8,1.2) {$\gamma$};
  \node[red!80!black, font=\small] at (1.5,-0.5) {$\ind=1$};
  \node[green!50!black, font=\small] at (5,-0.5) {$\ind=0$};
  \node[orange!80!black, font=\small] at (-1.5,-0.2) {$\ind=0$};
\end{tikzpicture}
\captionof{figure}{Winding numbers: the curve $\gamma$ winds once around $w_1$
(inside the curve) and zero times around $w_2$ and $w_3$ (outside).}
\label{fig:winding-number}
\end{center}

\begin{remark}[Geometric interpretation]\label{rem:winding-geometric}
The winding number $\ind(\gamma,w)$ counts the number of times $\gamma$ wraps
around $w$ in the counterclockwise direction.  As $z$ traverses $\gamma$, the
argument of $z-w$ changes by $2\pi\cdot\ind(\gamma,w)$ in total.
\end{remark}

%% ════════════════════════════════════════════════════════════
\section{Cauchy's Theorem --- General Formulation}
\label{sec:cauchy-general}
%% ════════════════════════════════════════════════════════════

\begin{theorem}[Cauchy's theorem --- general version]\label{thm:cauchy-general}
Let $f$ be holomorphic on an open set $\Omega\subseteq\C$, and let
$\gamma$ be a closed piecewise smooth curve in~$\Omega$ such that
$\ind(\gamma,w)=0$ for every $w\in\C\setminus\Omega$.  Then
\[
  \oint_\gamma f(z)\,\dd z \;=\; 0.
\]
\end{theorem}

This version subsumes all previous versions: in a simply connected domain,
every closed curve satisfies $\ind(\gamma,w)=0$ for $w\notin\Omega$.

\begin{corollary}[Deformation of contours]\label{cor:deformation-contours}
Let $f$ be holomorphic on $\Omega$, and let $\gamma_1,\gamma_2$ be two closed
piecewise smooth curves in $\Omega$ such that $\ind(\gamma_1,w) = \ind(\gamma_2,w)$
for all $w\notin\Omega$.  Then
\[
  \oint_{\gamma_1}f(z)\,\dd z \;=\; \oint_{\gamma_2}f(z)\,\dd z.
\]
\end{corollary}

\begin{proof}
Apply Theorem~\ref{thm:cauchy-general} to the curve $\gamma = \gamma_1 + \gamma_2^-$,
noting that $\ind(\gamma,w)=\ind(\gamma_1,w)-\ind(\gamma_2,w)=0$ for $w\notin\Omega$.
\end{proof}

\begin{example}[Deforming a contour around a pole]\label{ex:deformation-around-pole}
Let $\Omega = \C\setminus\{0\}$ and $f(z)=1/z$.  Let $\gamma_1$ be any simple
closed curve in $\Omega$ that winds once counterclockwise around $0$, and let
$\gamma_2(t)=e^{it}$, $t\in[0,2\pi]$, be the unit circle.  Since both have
winding number $1$ around $0$ (the only point outside $\Omega$), we get
\[
  \oint_{\gamma_1}\frac{\dd z}{z} = \oint_{\gamma_2}\frac{\dd z}{z} = 2\pi i.
\]
This is extremely useful: to evaluate the integral of $1/z$ around any curve,
we need only know how many times it winds around the origin.
\end{example}

%% ════════════════════════════════════════════════════════════
\section{Applications and Consequences}
\label{sec:cauchy-applications}
%% ════════════════════════════════════════════════════════════

\begin{theorem}[Existence of primitives on simply connected domains]%
\label{thm:primitive-simply-connected}
Let $\Omega$ be a simply connected domain and $f\colon\Omega\to\C$ holomorphic.
Then $f$ has a primitive on $\Omega$.  In particular:
\begin{itemize}
  \item A branch of the logarithm exists on every simply connected subdomain
    of $\C^*$.
  \item For any $\alpha\in\C$, the function $z^\alpha$ can be defined as a
    single-valued holomorphic function on any simply connected subdomain
    of $\C^*$.
\end{itemize}
\end{theorem}

\begin{proof}
By Corollary~\ref{cor:cauchy-simply-connected} and
Theorem~\ref{thm:primitive-characterization}.  For the logarithm: $1/z$ has a
primitive $L$ on $\Omega$; then $ze^{-L(z)}$ is constant (as shown in
Proposition~\ref{prop:log-branch-existence}), and adjusting the constant of
integration gives $e^{L(z)}=z$.
\end{proof}

\begin{proposition}[Integral of $1/z$ around a closed curve]%
\label{prop:integral-1-over-z}
Let $\gamma$ be a closed piecewise smooth curve in $\C^*$.  Then
\[
  \oint_\gamma\frac{\dd z}{z} \;=\; 2\pi i\cdot\ind(\gamma,0).
\]
\end{proposition}

\begin{proof}
This is the definition of the winding number.
\end{proof}

%% ════════════════════════════════════════════════════════════
\section{Worked Examples}
\label{sec:ch4-examples}
%% ════════════════════════════════════════════════════════════

\begin{example}[Direct computation on a circle]\label{ex:circle-computations}
Let $\gamma(t) = 2e^{it}$, $t\in[0,2\pi]$ (circle of radius $2$, centered
at the origin).
\begin{enumerate}[label=\textup{(\alph*)}]
  \item $\displaystyle\oint_\gamma z^2\,\dd z = 0$
    (since $z^2$ has primitive $z^3/3$ on $\C$).
  \item $\displaystyle\oint_\gamma\frac{\dd z}{z-1} = 2\pi i$
    (since $\ind(\gamma,1)=1$: the point $1$ is inside the circle of radius $2$).
  \item $\displaystyle\oint_\gamma\frac{\dd z}{z-3} = 0$
    (since $\ind(\gamma,3)=0$: the point $3$ is outside the circle).
  \item $\displaystyle\oint_\gamma\frac{\dd z}{z^2+1}
    = \oint_\gamma\frac{\dd z}{(z-i)(z+i)}$.
    Both $i$ and $-i$ lie inside $\gamma$.  By partial fractions:
    $\frac{1}{z^2+1} = \frac{1}{2i}\bigl(\frac{1}{z-i}-\frac{1}{z+i}\bigr)$.
    So the integral is $\frac{1}{2i}(2\pi i - 2\pi i) = 0$.
\end{enumerate}
\end{example}

\begin{example}[Integral of $\re z$ around the unit circle]\label{ex:re-z-integral}
Let $\gamma(t)=e^{it}$, $t\in[0,2\pi]$.  Then $\re z = \frac{z+\conj{z}}{2}$
and on the unit circle $\conj{z} = 1/z$, so $\re z = \frac{z+1/z}{2}$.
Hence
\[
  \oint_\gamma\re z\,\dd z
  = \frac{1}{2}\oint_\gamma z\,\dd z + \frac{1}{2}\oint_\gamma\frac{\dd z}{z}
  = 0 + \frac{1}{2}\cdot 2\pi i = \pi i.
\]
\end{example}

\begin{example}[A non-trivial winding number computation]\label{ex:winding-figure-eight}
Consider a figure-eight curve $\gamma$ that winds once counterclockwise
around $z=1$ and once clockwise around $z=-1$.  For any function $f$
holomorphic on $\C\setminus\{-1,1\}$, the winding numbers are
$\ind(\gamma,1) = 1$ and $\ind(\gamma,-1)=-1$.  Hence
\[
  \oint_\gamma\frac{\dd z}{z-1} = 2\pi i, \quad
  \oint_\gamma\frac{\dd z}{z+1} = -2\pi i, \quad
  \oint_\gamma\frac{\dd z}{z^2-1} = \frac{1}{2}(2\pi i - (-2\pi i)) = 2\pi i.
\]
\end{example}

\begin{example}[Integral of $e^{1/z}$ around the unit circle]\label{ex:exp-1-over-z}
Let $\gamma(t) = e^{it}$, $t\in[0,2\pi]$.  The function $e^{1/z}$ is
holomorphic on $\C^*$ but has an essential singularity at $0$.  Using the
Laurent expansion:
\[
  e^{1/z} = \sum_{n=0}^{\infty}\frac{1}{n!\,z^n}
  = 1 + \frac{1}{z} + \frac{1}{2z^2} + \frac{1}{6z^3} + \cdots
\]
By \eqref{eq:integral-zn-circle}, only the $1/z$ term contributes to
the integral:
\[
  \oint_\gamma e^{1/z}\,\dd z = 2\pi i.
\]
(This anticipates the \emph{residue theorem} of later chapters.)
\end{example}

%% ════════════════════════════════════════════════════════════
\section{Exercises}
\label{sec:ch4-exercises}
%% ════════════════════════════════════════════════════════════

\begin{exercise}[Direct computations]\label{exo:direct-integrals}
Compute the following integrals, where $\gamma(t)=Re^{it}$, $t\in[0,2\pi]$,
is the circle of radius $R$ centered at the origin.
\begin{enumerate}[label=\textup{(\alph*)}]
  \item $\displaystyle\oint_\gamma\conj{z}\,\dd z$ (with $R=1$).
  \item $\displaystyle\oint_\gamma\abs{z}^2\,\dd z$ (with $R=1$).
  \item $\displaystyle\oint_\gamma z^n\,\dd z$ for $n\in\Z$ and arbitrary $R>0$.
  \item $\displaystyle\oint_\gamma\frac{z+1}{z-1}\,\dd z$ with $R=2$.
\end{enumerate}
\end{exercise}

\begin{exercise}[Path integrals]\label{exo:path-integrals}
\leavevmode
\begin{enumerate}[label=\textup{(\alph*)}]
  \item Compute $\int_\gamma z^2\,\dd z$ where $\gamma$ is the segment from
    $0$ to $1+i$.
  \item Compute $\int_\gamma\re z\,\dd z$ along (i) the segment from $0$ to $1$,
    and (ii) the upper semicircle from $0$ to $1$.  Are the results equal?  Why or
    why not?
  \item Compute $\int_\gamma\frac{\dd z}{z}$ where $\gamma$ goes from $-1$ to $1$
    along (i) the upper semicircle, (ii) the lower semicircle.
\end{enumerate}
\end{exercise}

\begin{exercise}[ML-inequality applications]\label{exo:ML-inequality}
\leavevmode
\begin{enumerate}[label=\textup{(\alph*)}]
  \item Show that $\abs[\big]{\oint_{\abs{z}=R}\frac{e^z}{z^2}\,\dd z}
    \le \frac{2\pi e^R}{R}$ for $R>0$.
  \item Show that $\abs[\big]{\oint_{\abs{z}=R}\frac{\dd z}{z^n+1}}
    \to 0$ as $R\to\infty$ for $n\ge 2$.
  \item Let $\gamma_R$ be the upper semicircle $Re^{it}$, $t\in[0,\pi]$.
    Show that $\abs[\big]{\int_{\gamma_R}\frac{e^{iz}}{z}\,\dd z}\to 0$
    as $R\to\infty$.
\end{enumerate}
\end{exercise}

\begin{exercise}[Cauchy's theorem --- applications]\label{exo:cauchy-applications}
\leavevmode
\begin{enumerate}[label=\textup{(\alph*)}]
  \item Let $f$ be entire.  Show that $\oint_\gamma f(z)\,\dd z = 0$ for
    every closed contour $\gamma$.
  \item Let $f$ be holomorphic on $\C\setminus\{0\}$ and suppose $f$ has a
    primitive on $\C\setminus\{0\}$.  Show that
    $\oint_{\abs{z}=1}f(z)\,\dd z = 0$.
  \item Show that $1/z$ does \emph{not} have a primitive on $\C\setminus\{0\}$.
\end{enumerate}
\end{exercise}

\begin{exercise}[Winding numbers]\label{exo:winding-numbers}
\leavevmode
\begin{enumerate}[label=\textup{(\alph*)}]
  \item Compute $\ind(\gamma,0)$ where $\gamma(t)=e^{2it}$, $t\in[0,2\pi]$.
  \item Let $\gamma(t) = 2\cos t + i\sin t$, $t\in[0,2\pi]$ (an ellipse).
    What is $\ind(\gamma,0)$?  What is $\ind(\gamma,3)$?
  \item Let $\gamma$ be the boundary of the square with vertices $\pm 2\pm 2i$,
    traversed counterclockwise.  Compute $\ind(\gamma,0)$, $\ind(\gamma,1+i)$,
    and $\ind(\gamma,5)$.
\end{enumerate}
\end{exercise}

\begin{exercise}[Primitives on domains]\label{exo:primitives}
\leavevmode
\begin{enumerate}[label=\textup{(\alph*)}]
  \item Show that $z^2 e^z$ has a primitive on $\C$ and find it explicitly.
  \item Does $1/z^2$ have a primitive on $\C\setminus\{0\}$?  If so, find it.
    Does this contradict anything?
  \item Show that $1/(z^2+1)$ has a primitive on the right half-plane
    $\{\re z > 0\}$ but not on $\C\setminus\{\pm i\}$.
\end{enumerate}
\end{exercise}

\begin{exercise}[Cauchy's theorem for multiply connected domains]\label{exo:multiply-connected}
Let $\gamma_1(t) = 3e^{it}$ and $\gamma_2(t) = e^{it}$, $t\in[0,2\pi]$.
Let $f$ be holomorphic on the annulus $\Omega = \{1<\abs{z}<3\}$.
\begin{enumerate}[label=\textup{(\alph*)}]
  \item Explain why we cannot directly conclude $\oint_{\gamma_1}f\,\dd z = 0$.
  \item Show that $\oint_{\gamma_1}f\,\dd z = \oint_{\gamma_2}f\,\dd z$.
    \textit{Hint:} connect the two circles by cuts and apply Cauchy's theorem
    to a simply connected subdomain.
  \item Verify this when $f(z) = 1/z$.
\end{enumerate}
\end{exercise}

\begin{exercise}[The integral $\int_0^{2\pi}\frac{\dd\theta}{a+\cos\theta}$]%
\label{exo:trig-integral}
Let $a>1$.  Use the substitution $z=e^{i\theta}$ to show that
\[
  \int_0^{2\pi}\frac{\dd\theta}{a+\cos\theta}
  = \frac{2\pi}{\sqrt{a^2-1}}.
\]
\textit{Hint:} On $\abs{z}=1$, $\cos\theta = (z+z^{-1})/2$ and
$\dd\theta = \dd z/(iz)$.
\end{exercise}

\begin{exercise}[Green's function approach]\label{exo:green-approach}
Let $\Omega$ be a bounded domain with piecewise smooth boundary, and let
$f,g$ be holomorphic on an open set containing $\overline{\Omega}$.
Show that
\[
  \oint_{\partial\Omega}f(z)g(z)\,\dd z = 0.
\]
\end{exercise}

\begin{exercise}[An integral estimate]\label{exo:integral-estimate}
Let $f$ be holomorphic on $\{z:\abs{z}>R_0\}$ and suppose
$\abs{f(z)}\le M/\abs{z}^2$ for $\abs{z}>R_0$.  Show that
$\oint_{\abs{z}=R}f(z)\,\dd z\to 0$ as $R\to\infty$.
\end{exercise}

\begin{exercise}[Deformation principle]\label{exo:deformation-principle}
Let $f$ be holomorphic on $\C\setminus\{i,-i\}$.  Let $\gamma$ be a simple
closed curve enclosing both $i$ and $-i$, and let $\gamma_1$ (resp.\ $\gamma_2$)
be a small circle around $i$ (resp.\ $-i$), both positively oriented.
Show that
\[
  \oint_\gamma f(z)\,\dd z
  = \oint_{\gamma_1}f(z)\,\dd z + \oint_{\gamma_2}f(z)\,\dd z.
\]
Apply this to $f(z) = 1/(z^2+1)$ to recover the result of
Example~\ref{ex:circle-computations}(d).
\end{exercise}
%%%%%%%%%%%%%%%%%%%%%%%%%%%%%%%%%%%%%%%%%%%%%%%%%%%%%%%%%%%%
%%
%%  CHAPTER 5: Cauchy Integral Formula and Consequences
%%
%%%%%%%%%%%%%%%%%%%%%%%%%%%%%%%%%%%%%%%%%%%%%%%%%%%%%%%%%%%%
\chapter{Cauchy Integral Formula and Consequences}
\label{chap:cauchy-formula}

\section*{Introduction}
\addcontentsline{toc}{section}{Introduction}

The Cauchy integral theorem, established in the preceding chapter, tells us that
the integral of a holomorphic function around a closed contour is zero. In this
chapter we explore the remarkable \emph{converse} direction: the values of a
holomorphic function inside a domain are completely determined by its values on
the boundary. This is the content of the \textbf{Cauchy integral formula}, one
of the most powerful results in all of analysis. Its consequences are sweeping:
every holomorphic function is infinitely differentiable, every holomorphic
function is analytic (representable by a convergent power series), bounded entire
functions are constant (Liouville's theorem), and the modulus of a holomorphic
function cannot attain a local maximum in the interior of its domain.

%% ════════════════════════════════════════════════════════════
\section{The Cauchy Integral Formula}
\label{sec:cauchy-formula}
%% ════════════════════════════════════════════════════════════

\subsection{Statement for a disk}

We begin with the simplest and most fundamental version of the formula.

\begin{theorem}[Cauchy Integral Formula for a Disk]\label{thm:cauchy-integral-formula-disk}
Let $f$ be holomorphic on an open set $\Omega\subseteq\C$, and let
$\overline{D}(z_0,R) = \{z\in\C : \abs{z-z_0}\le R\}\subseteq\Omega$ be a
closed disk contained in~$\Omega$. Denote by $\gamma$ the positively oriented
circle $\gamma(t)=z_0+Re^{it}$, $t\in[0,2\pi]$. Then for every $z$ with
$\abs{z-z_0}<R$,
\begin{equation}\label{eq:cauchy-formula}
  f(z) \;=\; \frac{1}{2\pi i}\oint_\gamma \frac{f(\zeta)}{\zeta - z}\,\dd\zeta.
\end{equation}
\end{theorem}

\begin{center}
\begin{tikzpicture}[scale=1.4]
  %% Disk
  \fill[blue!6] (0,0) circle (2);
  \draw[thick, blue!70!black, decoration={markings, mark=at position 0.15 with {\arrow{Stealth}},
    mark=at position 0.4 with {\arrow{Stealth}},
    mark=at position 0.65 with {\arrow{Stealth}},
    mark=at position 0.9 with {\arrow{Stealth}}}, postaction={decorate}]
    (0,0) circle (2);
  %% Center
  \fill (0,0) circle (1.5pt) node[below left]{$z_0$};
  %% Interior point
  \fill (0.6,0.4) circle (1.5pt) node[below right]{$z$};
  %% Radius
  \draw[dashed, gray] (0,0) -- node[above,font=\small]{$R$} (1.414,1.414);
  %% Labels
  \node[blue!70!black] at (2.4,1.2) {$\gamma$};
  \node at (0,-2.6) {\small Positively oriented circle};
\end{tikzpicture}
\end{center}

\begin{proof}
Fix $z$ with $\abs{z-z_0}<R$. The key idea is to deform the contour: since the
integrand $\zeta\mapsto f(\zeta)/(\zeta-z)$ is holomorphic on $\Omega\setminus\{z\}$,
we may replace $\gamma$ by a small circle around~$z$ without changing the integral.

More precisely, let $\varepsilon>0$ be small enough so that $\overline{D}(z,\varepsilon)\subset D(z_0,R)$, and let $\gamma_\varepsilon(t)=z+\varepsilon e^{it}$, $t\in[0,2\pi]$.

\medskip
\noindent\textbf{Step 1.} \emph{Reduction to the small circle.}
The function $g(\zeta)=f(\zeta)/(\zeta-z)$ is holomorphic on the annular region
between $\gamma$ and $\gamma_\varepsilon$. By the Cauchy integral theorem for
multiply connected domains (or by introducing a slit connecting the two circles),
\[
  \oint_\gamma g(\zeta)\,\dd\zeta \;=\; \oint_{\gamma_\varepsilon} g(\zeta)\,\dd\zeta.
\]

\begin{center}
\begin{tikzpicture}[scale=1.2]
  %% Outer disk
  \fill[blue!6] (0,0) circle (2);
  \fill[white] (0.6,0.4) circle (0.45);
  %% Outer circle
  \draw[thick, blue!70!black, decoration={markings,
    mark=at position 0.12 with {\arrow{Stealth}},
    mark=at position 0.62 with {\arrow{Stealth}}}, postaction={decorate}]
    (0,0) circle (2);
  %% Inner circle
  \draw[thick, red!70!black, decoration={markings,
    mark=at position 0.25 with {\arrow{Stealth}},
    mark=at position 0.75 with {\arrow{Stealth}}}, postaction={decorate}]
    (0.6,0.4) circle (0.45);
  %% Points
  \fill (0,0) circle (1.3pt) node[below left]{$z_0$};
  \fill (0.6,0.4) circle (1.3pt) node[below,font=\small,yshift=-5pt]{$z$};
  %% Labels
  \node[blue!70!black] at (2.3,0.8) {$\gamma$};
  \node[red!70!black] at (1.4,0.7) {$\gamma_\varepsilon$};
\end{tikzpicture}
\end{center}

\noindent\textbf{Step 2.} \emph{Evaluation on the small circle.}
On $\gamma_\varepsilon$ we write $\zeta = z + \varepsilon e^{it}$, so
$\dd\zeta = i\varepsilon e^{it}\,\dd t$ and $\zeta - z = \varepsilon e^{it}$.
Then
\begin{align*}
  \oint_{\gamma_\varepsilon}\frac{f(\zeta)}{\zeta-z}\,\dd\zeta
  &= \int_0^{2\pi}\frac{f(z+\varepsilon e^{it})}{\varepsilon e^{it}}\cdot
     i\varepsilon e^{it}\,\dd t \\
  &= i\int_0^{2\pi}f(z+\varepsilon e^{it})\,\dd t.
\end{align*}
Since $f$ is continuous at $z$, as $\varepsilon\to 0$ the integrand converges
uniformly to $f(z)$, so
\[
  \oint_{\gamma_\varepsilon}\frac{f(\zeta)}{\zeta-z}\,\dd\zeta
  \;=\; i\int_0^{2\pi}f(z+\varepsilon e^{it})\,\dd t
  \;\xrightarrow{\varepsilon\to 0}\; 2\pi i\,f(z).
\]
But the integral on the left equals $\oint_\gamma f(\zeta)/(\zeta-z)\,\dd\zeta$
for all sufficiently small $\varepsilon$ (by Step~1), hence this value is
independent of $\varepsilon$. Therefore
\[
  \oint_\gamma\frac{f(\zeta)}{\zeta-z}\,\dd\zeta = 2\pi i\,f(z),
\]
which gives \eqref{eq:cauchy-formula}.
\end{proof}

\begin{remark}[Representation Formula]\label{rem:representation}
Equation~\eqref{eq:cauchy-formula} is extraordinary: the values of $f$ in the
\emph{interior} of the disk are entirely determined by the values of $f$ on
the \emph{boundary}. No analogous result exists for smooth functions of a real
variable. This rigidity is one of the hallmarks of complex analysis.
\end{remark}

\subsection{General version}

\begin{theorem}[Cauchy Integral Formula --- General Version]\label{thm:cauchy-integral-formula-general}
Let $\Omega\subseteq\C$ be a simply connected open set, and let $\gamma$ be a
positively oriented, simple, piecewise smooth closed curve in $\Omega$. If $f$
is holomorphic on $\Omega$ and $z$ lies in the interior of $\gamma$, then
\[
  f(z) = \frac{1}{2\pi i}\oint_\gamma \frac{f(\zeta)}{\zeta-z}\,\dd\zeta.
\]
\end{theorem}

\begin{proof}
The proof follows the same deformation argument as for the disk: introduce a
small circle $\gamma_\varepsilon$ around $z$ contained in the interior of
$\gamma$. Since $f(\zeta)/(\zeta-z)$ is holomorphic in the region between
$\gamma$ and $\gamma_\varepsilon$, the Cauchy integral theorem for multiply
connected domains gives
\[
  \oint_\gamma \frac{f(\zeta)}{\zeta-z}\,\dd\zeta
  = \oint_{\gamma_\varepsilon}\frac{f(\zeta)}{\zeta-z}\,\dd\zeta
  = 2\pi i\,f(z),
\]
where the last equality is the computation from the proof of
Theorem~\ref{thm:cauchy-integral-formula-disk}.
\end{proof}

\begin{example}[A First Application]\label{ex:cauchy-formula-basic}
Let us compute
\[
  I = \oint_{\abs{z}=2}\frac{e^z}{z-1}\,\dd z.
\]
Here $f(z)=e^z$ is entire, and $z_0=1$ lies inside the circle $\abs{z}=2$. By
the Cauchy integral formula,
\[
  I = 2\pi i\,f(1) = 2\pi i\,e.
\]
\end{example}

\begin{example}[Distinguishing Inside and Outside]\label{ex:cauchy-inside-outside}
Compute
\[
  J = \oint_{\abs{z}=1}\frac{\cos z}{z - 3}\,\dd z.
\]
Since $z_0=3$ lies \emph{outside} the unit circle, the integrand
$\zeta\mapsto\cos\zeta/(\zeta-3)$ is holomorphic on and inside $\abs{z}=1$.
By the Cauchy integral theorem, $J=0$.
\end{example}

%% ════════════════════════════════════════════════════════════
\section{Cauchy Integral Formula for Higher Derivatives}
\label{sec:higher-derivatives}
%% ════════════════════════════════════════════════════════════

One of the most striking consequences of the Cauchy integral formula is that
holomorphic functions are automatically infinitely differentiable, and all
derivatives admit integral representations.

\begin{theorem}[Cauchy Formula for Derivatives]\label{thm:cauchy-formula-derivatives}
Let $f$ be holomorphic on an open set $\Omega$, and let $\gamma$ be a positively
oriented, simple, piecewise smooth closed curve in $\Omega$ whose interior is
contained in $\Omega$. Then $f$ is infinitely differentiable in the interior
of $\gamma$, and for every $n\ge 0$ and every $z$ in the interior of $\gamma$,
\begin{equation}\label{eq:cauchy-derivative}
  f^{(n)}(z) \;=\; \frac{n!}{2\pi i}\oint_\gamma
    \frac{f(\zeta)}{(\zeta-z)^{n+1}}\,\dd\zeta.
\end{equation}
\end{theorem}

\begin{proof}
We prove the formula by induction on $n$. The base case $n=0$ is the Cauchy
integral formula (Theorem~\ref{thm:cauchy-integral-formula-general}).

\medskip\noindent
\textbf{Step 1: The case $n=1$.}
Let $z$ be in the interior of $\gamma$. For $h\ne 0$ small enough so that $z+h$ also
lies in the interior, the Cauchy integral formula gives
\[
  \frac{f(z+h)-f(z)}{h}
  = \frac{1}{2\pi i}\oint_\gamma f(\zeta)
    \cdot\frac{1}{h}\left(\frac{1}{\zeta-z-h}-\frac{1}{\zeta-z}\right)\dd\zeta.
\]
A direct computation shows
\[
  \frac{1}{h}\left(\frac{1}{\zeta-z-h}-\frac{1}{\zeta-z}\right)
  = \frac{1}{(\zeta-z-h)(\zeta-z)}.
\]
Let $d = \min_{\zeta\in\gamma}\abs{\zeta-z}>0$. For $\abs{h}<d/2$ we have
$\abs{\zeta-z-h}\ge d/2$, so
\[
  \frac{1}{(\zeta-z-h)(\zeta-z)} \;\xrightarrow{h\to 0}\;
  \frac{1}{(\zeta-z)^2}
\]
uniformly for $\zeta\in\gamma$. Since $f$ is continuous (hence bounded) on $\gamma$,
we may pass the limit under the integral sign to obtain
\[
  f'(z) = \frac{1}{2\pi i}\oint_\gamma \frac{f(\zeta)}{(\zeta-z)^2}\,\dd\zeta,
\]
which is \eqref{eq:cauchy-derivative} for $n=1$.

\medskip\noindent
\textbf{Step 2: Inductive step.}
Assume \eqref{eq:cauchy-derivative} holds for some $n\ge 1$. We differentiate
under the integral sign. Set
\[
  g_n(\zeta,z) = \frac{f(\zeta)}{(\zeta-z)^{n+1}}.
\]
The difference quotient argument analogous to Step~1 shows that for $\abs{h}$
small,
\[
  \frac{g_n(\zeta,z+h)-g_n(\zeta,z)}{h}
  \;\to\; \frac{(n+1)f(\zeta)}{(\zeta-z)^{n+2}}
\]
uniformly in $\zeta\in\gamma$. Passing the limit under the integral,
\[
  f^{(n+1)}(z) = \frac{n!}{2\pi i}\cdot(n+1)\oint_\gamma
    \frac{f(\zeta)}{(\zeta-z)^{n+2}}\,\dd\zeta
  = \frac{(n+1)!}{2\pi i}\oint_\gamma
    \frac{f(\zeta)}{(\zeta-z)^{n+2}}\,\dd\zeta,
\]
completing the induction.
\end{proof}

\begin{corollary}[Infinite Differentiability]\label{cor:holomorphic-smooth}
Every holomorphic function on an open set $\Omega\subseteq\C$ is infinitely
(complex) differentiable on $\Omega$. In particular, $f'$ is itself holomorphic.
\end{corollary}

\begin{remark}\label{rem:real-vs-complex}
This result has no analogue in real analysis: a real-valued function can be
differentiable everywhere yet have a discontinuous second derivative. In complex
analysis, a single complex derivative implies all higher derivatives exist and are
continuous.
\end{remark}

\begin{example}[Computing with the Derivative Formula]\label{ex:derivative-formula}
Compute
\[
  \oint_{\abs{z}=2}\frac{e^z}{(z-1)^3}\,\dd z.
\]
By Theorem~\ref{thm:cauchy-formula-derivatives} with $f(\zeta)=e^\zeta$, $z_0=1$,
and $n=2$,
\[
  \oint_{\abs{z}=2}\frac{e^z}{(z-1)^3}\,\dd z
  = \frac{2\pi i}{2!}\,f''(1) = \pi i\cdot e.
\]
\end{example}

\begin{example}\label{ex:derivative-formula2}
Compute
\[
  \oint_{\abs{z}=1}\frac{\sin z}{z^4}\,\dd z.
\]
Here $f(z)=\sin z$ and we evaluate at $z_0=0$ with $n=3$:
\[
  \oint_{\abs{z}=1}\frac{\sin z}{z^4}\,\dd z
  = \frac{2\pi i}{3!}\,f'''(0)
  = \frac{2\pi i}{6}\cdot(-\cos 0)
  = -\frac{\pi i}{3}.
\]
\end{example}

%% ════════════════════════════════════════════════════════════
\section{Morera's Theorem}
\label{sec:morera}
%% ════════════════════════════════════════════════════════════

\begin{theorem}[Morera's Theorem]\label{thm:morera}
Let $\Omega\subseteq\C$ be an open set, and let $f\colon\Omega\to\C$ be continuous.
If
\[
  \oint_T f(z)\,\dd z = 0
\]
for every triangle $T$ whose closure is contained in~$\Omega$, then $f$ is
holomorphic on~$\Omega$.
\end{theorem}

\begin{proof}
It suffices to prove that $f$ is holomorphic on each open disk $D(z_0,r)\subseteq\Omega$.
Fix such a disk and define
\[
  F(z) = \int_{z_0}^z f(\zeta)\,\dd\zeta, \qquad z\in D(z_0,r),
\]
where the integral is taken along the line segment from $z_0$ to $z$. The
hypothesis that $\oint_T f = 0$ for every triangle ensures that $F$ is
well-defined and independent of path within $D(z_0,r)$.

We claim that $F'(z)=f(z)$. Indeed, for $h$ small,
\[
  F(z+h)-F(z) = \int_z^{z+h}f(\zeta)\,\dd\zeta,
\]
where the integral is along the segment $[z,z+h]$. Since $f$ is continuous
at $z$, for every $\varepsilon>0$ there exists $\delta>0$ such that
$\abs{f(\zeta)-f(z)}<\varepsilon$ whenever $\abs{\zeta-z}<\delta$. For
$\abs{h}<\delta$,
\[
  \abs{F(z+h)-F(z)-hf(z)}
  = \abs*{\int_z^{z+h}\bigl(f(\zeta)-f(z)\bigr)\,\dd\zeta}
  \le \varepsilon\abs{h}.
\]
Hence $F'(z)=f(z)$. Since $F$ is holomorphic on $D(z_0,r)$,
Corollary~\ref{cor:holomorphic-smooth} implies $F' = f$ is holomorphic as well.
\end{proof}

\begin{remark}\label{rem:morera-converse}
Morera's theorem is a partial converse to the Cauchy integral theorem: the Cauchy
theorem says that if $f$ is holomorphic then $\oint_\gamma f = 0$ for closed
curves, while Morera's theorem says that if $\oint_T f = 0$ for all triangles
then $f$ is holomorphic. The two together characterize holomorphy among
continuous functions.
\end{remark}

%% ════════════════════════════════════════════════════════════
\section{Cauchy Inequalities}
\label{sec:cauchy-inequalities}
%% ════════════════════════════════════════════════════════════

\begin{theorem}[Cauchy Inequalities]\label{thm:cauchy-inequalities}
Let $f$ be holomorphic on an open set containing $\overline{D}(z_0,R)$, and let
\[
  M(R) = \max_{\abs{z-z_0}=R}\abs{f(z)}.
\]
Then for all $n\ge 0$,
\begin{equation}\label{eq:cauchy-inequality}
  \abs{f^{(n)}(z_0)} \;\le\; \frac{n!\,M(R)}{R^n}.
\end{equation}
\end{theorem}

\begin{proof}
By Theorem~\ref{thm:cauchy-formula-derivatives} with $\gamma$ the circle of
radius $R$ centered at $z_0$,
\[
  f^{(n)}(z_0) = \frac{n!}{2\pi i}\oint_{\abs{\zeta-z_0}=R}
    \frac{f(\zeta)}{(\zeta-z_0)^{n+1}}\,\dd\zeta.
\]
By the standard estimate for contour integrals,
\[
  \abs{f^{(n)}(z_0)}
  \le \frac{n!}{2\pi}\cdot\frac{M(R)}{R^{n+1}}\cdot 2\pi R
  = \frac{n!\,M(R)}{R^n}.\qedhere
\]
\end{proof}

%% ════════════════════════════════════════════════════════════
\section{Liouville's Theorem and the Fundamental Theorem of Algebra}
\label{sec:liouville}
%% ════════════════════════════════════════════════════════════

\begin{definition}[Entire Function]\label{def:entire}
A function $f\colon\C\to\C$ is called \textbf{entire} if it is holomorphic on
all of~$\C$.
\end{definition}

\begin{theorem}[Liouville's Theorem]\label{thm:liouville}
Every bounded entire function is constant.
\end{theorem}

\begin{proof}
Let $f$ be entire with $\abs{f(z)}\le M$ for all $z\in\C$. Fix any $z_0\in\C$.
By the Cauchy inequality \eqref{eq:cauchy-inequality} with $n=1$,
\[
  \abs{f'(z_0)} \le \frac{M}{R}
\]
for every $R>0$. Letting $R\to\infty$, we obtain $f'(z_0)=0$. Since $z_0$ was
arbitrary, $f'\equiv 0$ on $\C$, so $f$ is constant.
\end{proof}

\begin{remark}\label{rem:liouville-generalizations}
Liouville's theorem admits several generalizations. For instance, if $f$ is
entire and $\abs{f(z)}\le C(1+\abs{z})^n$ for some $C>0$ and some non-negative
integer $n$, then $f$ is a polynomial of degree at most $n$. The proof uses
the Cauchy inequality for the $(n+1)$-st derivative.
\end{remark}

\begin{theorem}[Fundamental Theorem of Algebra]\label{thm:fta}
Every non-constant polynomial $p(z)=a_nz^n+\cdots+a_1z+a_0$ with complex
coefficients and $n\ge 1$ has at least one root in $\C$.
\end{theorem}

\begin{proof}
Suppose for contradiction that $p(z)\ne 0$ for all $z\in\C$. Then $g(z)=1/p(z)$
is entire. We show that $g$ is bounded.

Since $a_n\ne 0$ and $n\ge 1$, we have
\[
  \abs{p(z)} = \abs{a_n}\abs{z}^n
  \abs*{1 + \frac{a_{n-1}}{a_n z}+\cdots+\frac{a_0}{a_n z^n}}.
\]
For $\abs{z}\ge R_0$ sufficiently large, the sum in the second factor has modulus
at least $1/2$, so
\[
  \abs{p(z)} \ge \frac{\abs{a_n}}{2}\abs{z}^n \ge \frac{\abs{a_n}}{2}R_0^n
  \qquad\text{for } \abs{z}\ge R_0.
\]
Hence $\abs{g(z)}\le 2/(\abs{a_n}R_0^n)$ for $\abs{z}\ge R_0$.

On the compact disk $\overline{D}(0,R_0)$, the continuous function $g$ is
bounded, say $\abs{g(z)}\le M_0$.

Taking $M=\max\{M_0, \, 2/(\abs{a_n}R_0^n)\}$, we see that $\abs{g(z)}\le M$
for all $z\in\C$. By Liouville's theorem, $g$ is constant, hence $p$ is
constant, contradicting $n\ge 1$.
\end{proof}

\begin{corollary}\label{cor:fta-factors}
Every polynomial of degree $n\ge 1$ with complex coefficients factors completely
as
\[
  p(z) = a_n(z-\alpha_1)(z-\alpha_2)\cdots(z-\alpha_n),
\]
where $\alpha_1,\ldots,\alpha_n\in\C$ (not necessarily distinct) are its roots.
\end{corollary}

\begin{proof}
By Theorem~\ref{thm:fta}, $p$ has a root $\alpha_1$. By polynomial division,
$p(z)=(z-\alpha_1)q(z)$ where $\deg q = n-1$. The result follows by induction.
\end{proof}

%% ════════════════════════════════════════════════════════════
\section{Mean Value Property}
\label{sec:mean-value}
%% ════════════════════════════════════════════════════════════

\begin{theorem}[Mean Value Property]\label{thm:mean-value}
Let $f$ be holomorphic on an open set containing $\overline{D}(z_0,R)$. Then
\begin{equation}\label{eq:mean-value}
  f(z_0) = \frac{1}{2\pi}\int_0^{2\pi}f(z_0+Re^{it})\,\dd t.
\end{equation}
That is, the value of $f$ at the center of a disk equals the average of $f$ over
the boundary circle.
\end{theorem}

\begin{proof}
By the Cauchy integral formula with $\gamma(t)=z_0+Re^{it}$,
\[
  f(z_0) = \frac{1}{2\pi i}\oint_\gamma\frac{f(\zeta)}{\zeta-z_0}\,\dd\zeta
  = \frac{1}{2\pi i}\int_0^{2\pi}\frac{f(z_0+Re^{it})}{Re^{it}}\cdot iRe^{it}\,\dd t
  = \frac{1}{2\pi}\int_0^{2\pi}f(z_0+Re^{it})\,\dd t.\qedhere
\]
\end{proof}

\begin{center}
\begin{tikzpicture}[scale=1.3]
  \fill[blue!6] (0,0) circle (1.8);
  \draw[thick, blue!70!black] (0,0) circle (1.8);
  \fill (0,0) circle (1.5pt) node[below]{$z_0$};
  %% Sample points on circle
  \foreach \angle in {0,30,...,330} {
    \fill[red!70!black] ({1.8*cos(\angle)},{1.8*sin(\angle)}) circle (1pt);
  }
  \draw[dashed, gray] (0,0) -- node[above,font=\small]{$R$} (1.8,0);
  \node[align=center,font=\small] at (0,-2.5) {$f(z_0)$ equals the average of $f$\\
    over the red sample points (in the limit)};
\end{tikzpicture}
\end{center}

%% ════════════════════════════════════════════════════════════
\section{Maximum and Minimum Modulus Principles}
\label{sec:max-modulus}
%% ════════════════════════════════════════════════════════════

\begin{theorem}[Maximum Modulus Principle]\label{thm:max-modulus}
Let $\Omega\subseteq\C$ be an open connected set, and let $f$ be holomorphic on
$\Omega$. If $\abs{f}$ attains a local maximum at some point $z_0\in\Omega$,
then $f$ is constant on $\Omega$.
\end{theorem}

\begin{proof}
Suppose $\abs{f(z_0)}\ge\abs{f(z)}$ for all $z$ in some disk $D(z_0,r)\subseteq\Omega$.
We first show $f$ is constant on $D(z_0,r)$.

If $f(z_0)=0$, then $\abs{f}\equiv 0$ on $D(z_0,r)$ and the claim is clear.
Assume $f(z_0)\ne 0$.

By the mean value property, for any $0<\rho<r$,
\[
  f(z_0) = \frac{1}{2\pi}\int_0^{2\pi}f(z_0+\rho e^{it})\,\dd t.
\]
Taking moduli and using $\abs{f(z_0+\rho e^{it})}\le\abs{f(z_0)}$,
\[
  \abs{f(z_0)} = \abs*{\frac{1}{2\pi}\int_0^{2\pi}f(z_0+\rho e^{it})\,\dd t}
  \le \frac{1}{2\pi}\int_0^{2\pi}\abs{f(z_0+\rho e^{it})}\,\dd t
  \le \abs{f(z_0)}.
\]
Hence equality holds throughout. In particular,
\[
  \frac{1}{2\pi}\int_0^{2\pi}\bigl[\abs{f(z_0)}-\abs{f(z_0+\rho e^{it})}\bigr]\,\dd t = 0.
\]
Since the integrand is non-negative and continuous, it must vanish identically:
$\abs{f(z_0+\rho e^{it})}=\abs{f(z_0)}$ for all $t$ and all $0<\rho<r$. Thus
$\abs{f}$ is constant on $D(z_0,r)$.

Now, $\abs{f}$ constant on an open disk implies $f$ is constant there (since a
holomorphic function with constant modulus on an open set is constant---this
follows from the Cauchy--Riemann equations or from the open mapping theorem).

The set $S=\{z\in\Omega:f(z)=f(z_0)\}$ is closed in $\Omega$ (by continuity)
and open (by the argument above). Since $\Omega$ is connected and $S$ is
non-empty, $S=\Omega$.
\end{proof}

\begin{corollary}[Maximum Modulus on Compact Sets]\label{cor:max-modulus-boundary}
Let $\Omega$ be a bounded open connected set, and let $f$ be holomorphic on
$\Omega$ and continuous on $\overline{\Omega}$. Then
\[
  \max_{z\in\overline{\Omega}}\abs{f(z)} = \max_{z\in\partial\Omega}\abs{f(z)}.
\]
That is, $\abs{f}$ attains its maximum on the boundary $\partial\Omega$.
\end{corollary}

\begin{proof}
Since $\overline{\Omega}$ is compact and $\abs{f}$ is continuous, the maximum is
attained at some $z_0\in\overline{\Omega}$. If $z_0\in\partial\Omega$, we are
done. If $z_0\in\Omega$, Theorem~\ref{thm:max-modulus} implies $f$ is constant,
and the result is trivially true.
\end{proof}

\begin{theorem}[Minimum Modulus Principle]\label{thm:min-modulus}
Let $\Omega\subseteq\C$ be an open connected set, and let $f$ be holomorphic on
$\Omega$ with $f(z)\ne 0$ for all $z\in\Omega$. If $\abs{f}$ attains a local
minimum at some $z_0\in\Omega$, then $f$ is constant.
\end{theorem}

\begin{proof}
Since $f$ is never zero on $\Omega$, the function $g=1/f$ is holomorphic on
$\Omega$. A local minimum of $\abs{f}$ at $z_0$ is a local maximum of $\abs{g}$
at $z_0$. By Theorem~\ref{thm:max-modulus}, $g$ is constant, hence $f$ is constant.
\end{proof}

\begin{remark}\label{rem:min-modulus-zero}
The hypothesis $f\ne 0$ is essential. For instance, $f(z)=z$ on $D(0,1)$ has
$\abs{f(0)}=0$, which is a minimum of $\abs{f}$, yet $f$ is not constant. The
zeros of $f$ are precisely the points where the minimum modulus principle can fail.
\end{remark}

\begin{center}
\begin{tikzpicture}[scale=1.0]
  %% Domain
  \draw[thick, blue!70!black, fill=blue!5]
    plot[smooth cycle, tension=0.7] coordinates
    {(-2.5,0) (-1.5,1.8) (0.8,2) (2.5,1.2) (2.8,-0.5) (1.5,-1.8) (-0.5,-2) (-2.2,-1.2)};
  \node[blue!70!black] at (0,2.4) {$\Omega$};
  %% Max point on boundary
  \fill[red!80!black] (2.65,0.35) circle (2pt);
  \node[red!80!black, right, font=\small] at (2.65,0.35) {$\max\abs{f}$};
  %% Interior point
  \fill (0,0) circle (1.5pt);
  \node[below, font=\small] at (0,-0.15) {interior: no max};
  %% Boundary label
  \node[gray, font=\footnotesize] at (-2.6,1.2) {$\partial\Omega$};
\end{tikzpicture}
\end{center}

%% ════════════════════════════════════════════════════════════
\section{Taylor Series of Holomorphic Functions}
\label{sec:taylor-series}
%% ════════════════════════════════════════════════════════════

\begin{theorem}[Taylor Series Expansion]\label{thm:taylor}
Let $f$ be holomorphic on an open set $\Omega\subseteq\C$, and let $z_0\in\Omega$.
Let $R$ be the largest radius such that $D(z_0,R)\subseteq\Omega$ (possibly $R=\infty$).
Then $f$ has a power series expansion
\begin{equation}\label{eq:taylor}
  f(z) = \sum_{n=0}^{\infty}a_n(z-z_0)^n, \qquad \abs{z-z_0}<R,
\end{equation}
where
\begin{equation}\label{eq:taylor-coefficients}
  a_n = \frac{f^{(n)}(z_0)}{n!}
  = \frac{1}{2\pi i}\oint_{\abs{\zeta-z_0}=\rho}\frac{f(\zeta)}{(\zeta-z_0)^{n+1}}\,\dd\zeta
\end{equation}
for any $0<\rho<R$. The convergence is absolute and uniform on each compact subset
of $D(z_0,R)$.
\end{theorem}

\begin{center}
\begin{tikzpicture}[scale=0.9]
  %% Omega region
  \draw[thick, blue!70!black, fill=blue!5]
    plot[smooth cycle, tension=0.6] coordinates
    {(-3.5,0.5) (-2,3) (1,3.5) (3.5,2) (4,-0.5) (3,-3) (-0.5,-3.5) (-3,-2)};
  \node[blue!70!black] at (-3,3) {$\Omega$};
  %% Center
  \fill (0,0) circle (1.8pt) node[below right]{$z_0$};
  %% Largest disk
  \draw[thick, red!70!black, dashed] (0,0) circle (2.8);
  \draw[<->, red!70!black] (0,0) -- node[above,font=\small]{$R$} (2.8,0);
  \node[red!70!black, font=\small] at (2.2,-2.5) {largest disk in $\Omega$};
  %% Boundary touching
  \draw[->, gray, thick] (3.3,-2.3) to[out=120,in=-30] (2.4,-0.8);
  \node[gray, font=\footnotesize, align=center] at (4.2,-2.5) {disk touches\\$\partial\Omega$};
\end{tikzpicture}
\end{center}

\begin{proof}
Fix $z\in D(z_0,R)$ and choose $\rho$ with $\abs{z-z_0}<\rho<R$. Let
$\gamma$ be the circle $\abs{\zeta-z_0}=\rho$. By the Cauchy integral formula,
\[
  f(z) = \frac{1}{2\pi i}\oint_\gamma\frac{f(\zeta)}{\zeta-z}\,\dd\zeta.
\]
We expand the Cauchy kernel:
\[
  \frac{1}{\zeta-z}
  = \frac{1}{(\zeta-z_0)-(z-z_0)}
  = \frac{1}{\zeta-z_0}\cdot\frac{1}{1-\frac{z-z_0}{\zeta-z_0}}.
\]
Since $\abs{z-z_0}<\rho=\abs{\zeta-z_0}$ for $\zeta\in\gamma$, we have
$\abs{(z-z_0)/(\zeta-z_0)}<1$, and the geometric series converges:
\[
  \frac{1}{\zeta-z}
  = \sum_{n=0}^{\infty}\frac{(z-z_0)^n}{(\zeta-z_0)^{n+1}}.
\]
The convergence is uniform in $\zeta\in\gamma$ (with ratio $\abs{z-z_0}/\rho<1$),
so we may interchange summation and integration:
\begin{align*}
  f(z)
  &= \frac{1}{2\pi i}\oint_\gamma f(\zeta)\sum_{n=0}^{\infty}
     \frac{(z-z_0)^n}{(\zeta-z_0)^{n+1}}\,\dd\zeta \\
  &= \sum_{n=0}^{\infty}\left[\frac{1}{2\pi i}\oint_\gamma
     \frac{f(\zeta)}{(\zeta-z_0)^{n+1}}\,\dd\zeta\right](z-z_0)^n \\
  &= \sum_{n=0}^{\infty}\frac{f^{(n)}(z_0)}{n!}(z-z_0)^n,
\end{align*}
where the last equality uses the Cauchy formula for derivatives \eqref{eq:cauchy-derivative}.

\medskip\noindent
\textbf{Uniform convergence.}
Let $K\subset D(z_0,R)$ be compact. Choose $r$ with $\sup_{z\in K}\abs{z-z_0}<r<R$
and set $\rho=(r+R)/2$. For $z\in K$ and $\zeta$ on $\abs{\zeta-z_0}=\rho$,
\[
  \abs*{\frac{(z-z_0)^n}{(\zeta-z_0)^{n+1}}}
  \le \frac{r^n}{\rho^{n+1}}.
\]
Since $M=\max_{\abs{\zeta-z_0}=\rho}\abs{f(\zeta)}<\infty$ and
$\sum_n(r/\rho)^n<\infty$, the Weierstrass $M$-test gives uniform convergence on $K$.
\end{proof}

\begin{corollary}[Holomorphic Implies Analytic]\label{cor:holomorphic-analytic}
A function is holomorphic on an open set $\Omega$ if and only if it is complex
analytic on $\Omega$ (i.e., locally representable by a convergent power series).
\end{corollary}

\begin{remark}\label{rem:radius-convergence}
The radius of convergence of the Taylor series \eqref{eq:taylor} is
\emph{at least} $R=\operatorname{dist}(z_0,\partial\Omega)$. In fact, the
radius of convergence equals the distance from $z_0$ to the nearest singularity
of $f$ (which may lie outside $\Omega$ if $f$ can be analytically continued).
\end{remark}

\begin{example}[Taylor Series of Standard Functions]\label{ex:taylor-standard}
\begin{enumerate}[label=(\alph*)]
\item $\displaystyle e^z = \sum_{n=0}^{\infty}\frac{z^n}{n!}$, valid for all $z\in\C$
  (radius of convergence $R=\infty$).
\item $\displaystyle\frac{1}{1-z}=\sum_{n=0}^{\infty}z^n$, valid for $\abs{z}<1$
  (singularity at $z=1$, $R=1$).
\item $\displaystyle\Log(1+z)=\sum_{n=1}^{\infty}\frac{(-1)^{n-1}}{n}z^n$, valid for $\abs{z}<1$
  (branch point at $z=-1$, $R=1$).
\item $\displaystyle\frac{1}{z-2}=\frac{-1}{2}\cdot\frac{1}{1-z/2}
  = -\sum_{n=0}^{\infty}\frac{z^n}{2^{n+1}}$, centered at $z_0=0$, valid for $\abs{z}<2$.
\end{enumerate}
\end{example}

\begin{example}[Taylor Expansion at a Non-Zero Center]\label{ex:taylor-nonzero-center}
Expand $f(z)=1/(1-z)$ in a Taylor series centered at $z_0=i$.

We write
\[
  \frac{1}{1-z} = \frac{1}{(1-i)-(z-i)}
  = \frac{1}{1-i}\cdot\frac{1}{1-\frac{z-i}{1-i}}
  = \frac{1}{1-i}\sum_{n=0}^{\infty}\frac{(z-i)^n}{(1-i)^n}
  = \sum_{n=0}^{\infty}\frac{(z-i)^n}{(1-i)^{n+1}},
\]
valid for $\abs{z-i}<\abs{1-i}=\sqrt{2}$, which is the distance from $i$ to the
singularity at $z=1$.
\end{example}

%% ════════════════════════════════════════════════════════════
\section{Further Applications}
\label{sec:further-applications}
%% ════════════════════════════════════════════════════════════

\begin{proposition}[Schwarz Lemma]\label{prop:schwarz}
Let $f\colon D(0,1)\to D(0,1)$ be holomorphic with $f(0)=0$. Then:
\begin{enumerate}[label=(\roman*)]
  \item $\abs{f(z)}\le\abs{z}$ for all $z\in D(0,1)$.
  \item $\abs{f'(0)}\le 1$.
  \item If equality holds in (i) for some $z\ne 0$, or in (ii), then $f(z)=e^{i\theta}z$
    for some $\theta\in\R$.
\end{enumerate}
\end{proposition}

\begin{proof}
Since $f(0)=0$, the function $g(z)=f(z)/z$ has a removable singularity at $0$
(set $g(0)=f'(0)$) and is holomorphic on $D(0,1)$.

For $0<r<1$, the maximum modulus principle applied to $g$ on $\overline{D}(0,r)$ gives
\[
  \abs{g(z)} \le \max_{\abs{\zeta}=r}\abs{g(\zeta)}
  = \max_{\abs{\zeta}=r}\frac{\abs{f(\zeta)}}{r}
  \le \frac{1}{r}.
\]
Letting $r\to 1^-$, we get $\abs{g(z)}\le 1$ for all $z\in D(0,1)$. This gives
(i) and (ii).

If equality holds (i.e., $\abs{g(z_0)}=1$ for some $z_0\in D(0,1)$), then
$\abs{g}$ attains its maximum in the interior. By the maximum modulus principle,
$g$ is constant, say $g(z)=e^{i\theta}$, so $f(z)=e^{i\theta}z$.
\end{proof}

%% ════════════════════════════════════════════════════════════
\section{Exercises}
\label{sec:exercises-ch5}
%% ════════════════════════════════════════════════════════════

\begin{exercise}[$\bigstar$]\label{exo:cauchy-formula-1}
Use the Cauchy integral formula to evaluate:
\begin{enumerate}[label=(\alph*)]
  \item $\displaystyle\oint_{\abs{z}=3}\frac{z^2+1}{z-2}\,\dd z$
  \item $\displaystyle\oint_{\abs{z}=1}\frac{e^z}{z^2}\,\dd z$
  \item $\displaystyle\oint_{\abs{z-i}=2}\frac{\cos z}{z^2+1}\,\dd z$
\end{enumerate}
\textit{Hint:} For (c), factor the denominator and determine which poles lie inside the contour.
\end{exercise}

\begin{exercise}[$\bigstar$]\label{exo:morera-application}
Let $(f_n)$ be a sequence of holomorphic functions on an open set $\Omega$ that
converges uniformly on compact subsets of $\Omega$ to a function $f$. Use
Morera's theorem to prove that $f$ is holomorphic on $\Omega$.
\end{exercise}

\begin{exercise}[$\bigstar\!\bigstar$]\label{exo:cauchy-inequality-application}
Let $f$ be entire and suppose there exist constants $C>0$ and $k\in\N$ such that
$\abs{f(z)}\le C\abs{z}^k$ for all $\abs{z}\ge 1$. Prove that $f$ is a
polynomial of degree at most $k$.
\end{exercise}

\begin{exercise}[$\bigstar\!\bigstar$]\label{exo:max-modulus-application}
Let $f$ be holomorphic on $D(0,1)$ and continuous on $\overline{D}(0,1)$, with
$\abs{f(z)}=1$ for all $\abs{z}=1$. Prove that $f$ has at least one zero in
$D(0,1)$.

\textit{Hint:} Apply the minimum modulus principle.
\end{exercise}

\begin{exercise}[$\bigstar\!\bigstar$]\label{exo:mean-value-harmonic}
Let $u\colon\Omega\to\R$ be harmonic (i.e., $u=\re f$ for some holomorphic $f$).
Prove that $u$ satisfies the mean value property:
\[
  u(z_0) = \frac{1}{2\pi}\int_0^{2\pi}u(z_0+re^{it})\,\dd t
\]
for every $\overline{D}(z_0,r)\subseteq\Omega$.
\end{exercise}

\begin{exercise}[$\bigstar\!\bigstar$]\label{exo:liouville-harmonic}
Let $u\colon\C\to\R$ be a harmonic function that is bounded above. Prove that $u$
is constant.

\textit{Hint:} Consider $g(z)=e^{u(z)+iv(z)}$ where $v$ is a harmonic conjugate.
\end{exercise}

\begin{exercise}[$\bigstar\!\bigstar\!\bigstar$]\label{exo:schwarz-pick}
\textbf{(Schwarz--Pick Lemma.)}
Let $f\colon D(0,1)\to D(0,1)$ be holomorphic. Prove that for all $z,w\in D(0,1)$,
\[
  \abs*{\frac{f(z)-f(w)}{1-\overline{f(w)}f(z)}}
  \le \abs*{\frac{z-w}{1-\conj{w}z}},
\]
and $\abs{f'(z)}\le\dfrac{1-\abs{f(z)}^2}{1-\abs{z}^2}$.

\textit{Hint:} Compose $f$ with suitable M\"obius transformations and apply the Schwarz lemma.
\end{exercise}

\begin{exercise}[$\bigstar\!\bigstar\!\bigstar$]\label{exo:hadamard-three-circles}
\textbf{(Hadamard Three-Circles Theorem.)}
Let $f$ be holomorphic on the annulus $\{z : r_1 < \abs{z} < r_2\}$ and continuous
on its closure. Define $M(r)=\max_{\abs{z}=r}\abs{f(z)}$ for $r_1\le r\le r_2$.
Prove that $\log M(r)$ is a convex function of $\log r$. More precisely, for
$r_1\le r\le r_2$,
\[
  \log M(r) \le \frac{\log(r_2/r)}{\log(r_2/r_1)}\log M(r_1)
  + \frac{\log(r/r_1)}{\log(r_2/r_1)}\log M(r_2).
\]
\textit{Hint:} Apply the maximum modulus principle to $z^{\alpha}f(z)$ for a suitable
real number $\alpha$.
\end{exercise}

\begin{exercise}[$\bigstar\!\bigstar\!\bigstar$]\label{exo:open-mapping}
\textbf{(Open Mapping Theorem.)}
Let $f$ be a non-constant holomorphic function on a connected open set $\Omega$.
Prove that $f(\Omega)$ is open.

\textit{Hint:} Given $w_0=f(z_0)$, use the argument principle (or a direct
approach with Rouch\'e's theorem) to show that $f$ takes all values in a
neighborhood of $w_0$.
\end{exercise}

\begin{exercise}[$\bigstar\!\bigstar$]\label{exo:cauchy-formula-partial-fractions}
Compute the following integrals using partial fractions and the Cauchy integral
formula:
\begin{enumerate}[label=(\alph*)]
  \item $\displaystyle\oint_{\abs{z}=3}\frac{z}{(z-1)(z-2)}\,\dd z$
  \item $\displaystyle\oint_{\abs{z}=2}\frac{1}{z^2(z-3)}\,\dd z$
  \item $\displaystyle\oint_{\abs{z}=4}\frac{e^z}{z^2-1}\,\dd z$
\end{enumerate}
\end{exercise}

\begin{exercise}[$\bigstar\!\bigstar$]\label{exo:taylor-computation}
Find the Taylor series centered at $z_0=0$ for:
\begin{enumerate}[label=(\alph*)]
  \item $f(z)=\dfrac{z}{(1-z)^2}$
  \item $f(z)=\dfrac{1}{z^2+4}$
  \item $f(z)=z\cos(z^2)$
\end{enumerate}
Determine the radius of convergence in each case.
\end{exercise}


%%%%%%%%%%%%%%%%%%%%%%%%%%%%%%%%%%%%%%%%%%%%%%%%%%%%%%%%%%%%
%%
%%  CHAPTER 6: Laurent Series and Isolated Singularities
%%
%%%%%%%%%%%%%%%%%%%%%%%%%%%%%%%%%%%%%%%%%%%%%%%%%%%%%%%%%%%%
\chapter{Laurent Series and Isolated Singularities}
\label{chap:laurent}

\section*{Introduction}
\addcontentsline{toc}{section}{Introduction}

The Taylor series expansion developed in Chapter~\ref{chap:cauchy-formula}
represents a holomorphic function as a power series $\sum_{n=0}^\infty a_n(z-z_0)^n$
in a disk centered at $z_0$. However, many functions of interest are not defined
at certain isolated points---for instance, $1/z$ at the origin or $e^{1/z}$ at
the origin. To study such functions, we need an expansion that allows
\emph{negative} powers of $(z-z_0)$. This is the \textbf{Laurent series}, which
converges in an \emph{annulus} rather than a disk. The nature of the negative-power
part (the \emph{principal part}) leads to a classification of singularities that
is fundamental to the residue calculus of the next chapter.

%% ════════════════════════════════════════════════════════════
\section{Laurent Series}
\label{sec:laurent-series}
%% ════════════════════════════════════════════════════════════

\subsection{Statement and uniqueness}

\begin{definition}[Annulus]\label{def:annulus}
For $0\le r_1 < r_2 \le +\infty$ and $z_0\in\C$, the \textbf{open annulus}
centered at $z_0$ with radii $r_1$ and $r_2$ is
\[
  A(z_0;\,r_1,r_2) = \{z\in\C : r_1 < \abs{z-z_0} < r_2\}.
\]
When $r_1=0$ this is a \emph{punctured disk} $D(z_0,r_2)\setminus\{z_0\}$.
When $r_2=+\infty$ this is the complement of a closed disk.
\end{definition}

\begin{center}
\begin{tikzpicture}[scale=1.0]
  %% Annulus
  \fill[blue!8] (0,0) circle (2.5);
  \fill[white] (0,0) circle (1);
  \draw[thick, blue!70!black, dashed] (0,0) circle (2.5);
  \draw[thick, blue!70!black, dashed] (0,0) circle (1);
  %% Center
  \fill (0,0) circle (1.5pt) node[below left]{$z_0$};
  %% Radii
  \draw[<->, gray] (0,0) -- node[above,font=\small]{$r_1$} (1,0);
  \draw[<->, gray] (0,0) -- node[above,font=\small,pos=0.7]{$r_2$} (0,2.5);
  %% Label
  \node[blue!70!black] at (1.9,1.5) {$A(z_0;r_1,r_2)$};
\end{tikzpicture}
\end{center}

\begin{theorem}[Laurent Series Expansion]\label{thm:laurent}
Let $f$ be holomorphic on the annulus $A(z_0;\,r_1,r_2)$. Then $f$ has a unique
expansion
\begin{equation}\label{eq:laurent}
  f(z) = \sum_{n=-\infty}^{+\infty}a_n(z-z_0)^n
  = \underbrace{\sum_{n=0}^{\infty}a_n(z-z_0)^n}_{\text{analytic part}}
  \;+\; \underbrace{\sum_{n=1}^{\infty}a_{-n}(z-z_0)^{-n}}_{\text{principal part}},
\end{equation}
valid for all $z\in A(z_0;\,r_1,r_2)$. The convergence is absolute and uniform
on compact subsets of the annulus. The coefficients are given by
\begin{equation}\label{eq:laurent-coefficients}
  a_n = \frac{1}{2\pi i}\oint_{\gamma_\rho}\frac{f(\zeta)}{(\zeta-z_0)^{n+1}}\,\dd\zeta,
  \qquad n\in\Z,
\end{equation}
where $\gamma_\rho$ is the positively oriented circle $\abs{\zeta-z_0}=\rho$ for
any $r_1<\rho<r_2$.
\end{theorem}

\begin{remark}[Uniqueness]\label{rem:laurent-uniqueness}
The Laurent expansion is \emph{unique}: if $f(z)=\sum_{n=-\infty}^{\infty}b_n(z-z_0)^n$
converges on $A(z_0;\,r_1,r_2)$, then $b_n=a_n$ for all $n$. This follows from
the integral formula \eqref{eq:laurent-coefficients} and the fact that
\[
  \frac{1}{2\pi i}\oint_{\abs{\zeta-z_0}=\rho}(\zeta-z_0)^{m}\,\dd\zeta
  = \begin{cases} 1 & \text{if } m=-1, \\ 0 & \text{if } m\ne -1, \end{cases}
\]
for any integer $m$.
\end{remark}

\subsection{Proof via the Cauchy formula}

\begin{proof}[Proof of Theorem~\ref{thm:laurent}]
Fix $z\in A(z_0;\,r_1,r_2)$. Choose radii $\rho_1,\rho_2$ with
$r_1<\rho_1<\abs{z-z_0}<\rho_2<r_2$. Let $C_1$ be the circle
$\abs{\zeta-z_0}=\rho_1$ (positively oriented) and $C_2$ the circle
$\abs{\zeta-z_0}=\rho_2$ (positively oriented).

\begin{center}
\begin{tikzpicture}[scale=1.1]
  %% Annulus shading
  \fill[blue!6] (0,0) circle (2.8);
  \fill[white] (0,0) circle (0.8);
  %% Outer boundary (dashed, boundary of annulus)
  \draw[gray, dashed] (0,0) circle (2.8);
  \draw[gray, dashed] (0,0) circle (0.8);
  %% Integration circles
  \draw[thick, red!70!black, decoration={markings,
    mark=at position 0.15 with {\arrow{Stealth}},
    mark=at position 0.65 with {\arrow{Stealth}}}, postaction={decorate}]
    (0,0) circle (2.3);
  \draw[thick, orange!80!black, decoration={markings,
    mark=at position 0.15 with {\arrow{Stealth}},
    mark=at position 0.65 with {\arrow{Stealth}}}, postaction={decorate}]
    (0,0) circle (1.2);
  %% Center
  \fill (0,0) circle (1.3pt) node[below left]{$z_0$};
  %% Point z
  \fill (1.6,0.6) circle (1.5pt) node[right]{$z$};
  %% Labels
  \node[red!70!black,font=\small] at (2.3,2.0) {$C_2$ ($\rho_2$)};
  \node[orange!80!black,font=\small] at (-1.5,0.8) {$C_1$ ($\rho_1$)};
  \node[gray,font=\footnotesize] at (3.2,1.5) {$r_2$};
  \node[gray,font=\footnotesize] at (0.5,0.9) {$r_1$};
\end{tikzpicture}
\end{center}

\medskip\noindent
\textbf{Step 1: Cauchy formula with two circles.}
By the Cauchy integral formula for multiply connected domains, the function
$\zeta\mapsto f(\zeta)/(\zeta-z)$ is holomorphic in the annular region between
$C_1$ and $C_2$ (except at $\zeta=z$, which lies in this region). We introduce a
keyhole contour (or use the standard deformation argument) to obtain:
\begin{equation}\label{eq:two-circle}
  f(z) = \frac{1}{2\pi i}\oint_{C_2}\frac{f(\zeta)}{\zeta-z}\,\dd\zeta
       - \frac{1}{2\pi i}\oint_{C_1}\frac{f(\zeta)}{\zeta-z}\,\dd\zeta.
\end{equation}
The minus sign arises because $C_1$ must be traversed clockwise (negatively) in
the boundary of the region between the two circles, but we have written it with
positive orientation.

\medskip\noindent
\textbf{Step 2: Expanding the outer integral.}
For $\zeta\in C_2$ we have $\abs{z-z_0}<\rho_2=\abs{\zeta-z_0}$, so exactly
as in the proof of the Taylor theorem,
\[
  \frac{1}{\zeta-z} = \sum_{n=0}^{\infty}\frac{(z-z_0)^n}{(\zeta-z_0)^{n+1}},
\]
uniformly for $\zeta\in C_2$. Substituting into the first integral in
\eqref{eq:two-circle} and interchanging sum and integral:
\begin{equation}\label{eq:outer-sum}
  \frac{1}{2\pi i}\oint_{C_2}\frac{f(\zeta)}{\zeta-z}\,\dd\zeta
  = \sum_{n=0}^{\infty}\underbrace{\left[\frac{1}{2\pi i}\oint_{C_2}
    \frac{f(\zeta)}{(\zeta-z_0)^{n+1}}\,\dd\zeta\right]}_{=\,a_n}(z-z_0)^n.
\end{equation}

\medskip\noindent
\textbf{Step 3: Expanding the inner integral.}
For $\zeta\in C_1$ we have $\abs{\zeta-z_0}=\rho_1<\abs{z-z_0}$. We write
\[
  \frac{1}{\zeta-z} = \frac{-1}{(z-z_0)-({\zeta-z_0})}
  = \frac{-1}{z-z_0}\cdot\frac{1}{1-\frac{\zeta-z_0}{z-z_0}}
  = -\sum_{m=0}^{\infty}\frac{(\zeta-z_0)^m}{(z-z_0)^{m+1}},
\]
uniformly for $\zeta\in C_1$ (since the ratio $\rho_1/\abs{z-z_0}<1$).
Substituting:
\begin{align}
  -\frac{1}{2\pi i}\oint_{C_1}\frac{f(\zeta)}{\zeta-z}\,\dd\zeta
  &= \sum_{m=0}^{\infty}\left[\frac{1}{2\pi i}\oint_{C_1}
     f(\zeta)(\zeta-z_0)^m\,\dd\zeta\right]\frac{1}{(z-z_0)^{m+1}} \notag\\
  &= \sum_{m=0}^{\infty}\underbrace{\left[\frac{1}{2\pi i}\oint_{C_1}
     \frac{f(\zeta)}{(\zeta-z_0)^{-m}}\,\dd\zeta\right]}_{}\frac{1}{(z-z_0)^{m+1}}. \label{eq:inner-sum}
\end{align}
Setting $n=-(m+1)$ (so $m=-n-1$ and $n$ runs over $-1,-2,-3,\ldots$), the
generic term becomes
\[
  a_n(z-z_0)^n \quad\text{where}\quad
  a_n = \frac{1}{2\pi i}\oint_{C_1}\frac{f(\zeta)}{(\zeta-z_0)^{n+1}}\,\dd\zeta.
\]

\medskip\noindent
\textbf{Step 4: Combining.}
Adding \eqref{eq:outer-sum} and \eqref{eq:inner-sum},
\[
  f(z) = \sum_{n=0}^{\infty}a_n(z-z_0)^n + \sum_{n=-\infty}^{-1}a_n(z-z_0)^n
  = \sum_{n=-\infty}^{+\infty}a_n(z-z_0)^n.
\]

\medskip\noindent
\textbf{Step 5: Independence of $\rho$.}
Since $f(\zeta)/(\zeta-z_0)^{n+1}$ is holomorphic on the annulus (for each fixed $n$),
the Cauchy theorem implies that the integral in \eqref{eq:laurent-coefficients}
is independent of the choice of $\rho\in(r_1,r_2)$.

\medskip\noindent
\textbf{Step 6: Uniqueness.}
If $f(z)=\sum_{n=-\infty}^{\infty}b_n(z-z_0)^n$ converges in $A(z_0;\,r_1,r_2)$,
multiply both sides by $(z-z_0)^{-k-1}$ and integrate over $\abs{z-z_0}=\rho$:
\[
  \frac{1}{2\pi i}\oint_{\abs{z-z_0}=\rho}\frac{f(z)}{(z-z_0)^{k+1}}\,\dd z
  = \sum_{n=-\infty}^{\infty}b_n\cdot\frac{1}{2\pi i}\oint_{\abs{z-z_0}=\rho}
    (z-z_0)^{n-k-1}\,\dd z = b_k.
\]
Hence $b_k=a_k$ for all $k$.
\end{proof}

%% ════════════════════════════════════════════════════════════
\section{Classification of Isolated Singularities}
\label{sec:singularities}
%% ════════════════════════════════════════════════════════════

\begin{definition}[Isolated Singularity]\label{def:isolated-singularity}
A point $z_0\in\C$ is an \textbf{isolated singularity} of a function $f$ if
$f$ is holomorphic on some punctured disk $D(z_0,R)\setminus\{z_0\}$ for some
$R>0$, but $f$ is not holomorphic at $z_0$ itself (either $f$ is not defined at
$z_0$ or not differentiable there).
\end{definition}

At an isolated singularity $z_0$, the function has a Laurent expansion
$f(z)=\sum_{n=-\infty}^{\infty}a_n(z-z_0)^n$ valid in $0<\abs{z-z_0}<R$. The
nature of the principal part $\sum_{n=1}^{\infty}a_{-n}(z-z_0)^{-n}$ determines
the type of singularity.

\begin{definition}[Classification of Singularities]\label{def:singularity-types}
Let $z_0$ be an isolated singularity of $f$ with Laurent expansion
$\sum_{n=-\infty}^{\infty}a_n(z-z_0)^n$ in $0<\abs{z-z_0}<R$.
\begin{enumerate}[label=(\roman*)]
  \item $z_0$ is a \textbf{removable singularity} if $a_n=0$ for all $n<0$.
    In this case the Laurent series is just a Taylor series, and $f$ can be
    extended to a holomorphic function at $z_0$ by setting $f(z_0)=a_0$.
  \item $z_0$ is a \textbf{pole} if there exists $m\ge 1$ such that $a_{-m}\ne 0$
    and $a_n=0$ for all $n<-m$. The integer $m$ is called the \textbf{order}
    (or \textbf{multiplicity}) of the pole.
  \item $z_0$ is an \textbf{essential singularity} if $a_n\ne 0$ for infinitely
    many $n<0$.
\end{enumerate}
\end{definition}

\begin{definition}[Principal Part]\label{def:principal-part}
The \textbf{principal part} of $f$ at an isolated singularity $z_0$ is
\[
  \sum_{n=1}^{\infty}a_{-n}(z-z_0)^{-n}.
\]
At a removable singularity, the principal part is zero. At a pole of order $m$,
it is a finite sum $\sum_{n=1}^{m}a_{-n}(z-z_0)^{-n}$. At an essential
singularity, it is an infinite series.
\end{definition}

\begin{center}
\begin{tikzpicture}[
  box/.style={draw, rounded corners, minimum width=3.2cm, minimum height=1cm,
    align=center, font=\small},
  >=Stealth
  ]
  \node[box, fill=blue!10] (iso) at (0,2.5) {Isolated singularity\\at $z_0$};
  \node[box, fill=green!15] (rem) at (-4.5,0) {Removable\\$a_n=0$ for $n<0$};
  \node[box, fill=yellow!20] (pole) at (0,0) {Pole of order $m$\\$a_{-m}\ne 0$, $a_n=0$ for $n<-m$};
  \node[box, fill=red!15] (ess) at (4.5,0) {Essential\\infinitely many $a_n\ne 0$, $n<0$};
  \draw[->] (iso) -- (rem);
  \draw[->] (iso) -- (pole);
  \draw[->] (iso) -- (ess);
\end{tikzpicture}
\end{center}

%% ════════════════════════════════════════════════════════════
\section{Removable Singularities}
\label{sec:removable}
%% ════════════════════════════════════════════════════════════

\begin{theorem}[Riemann's Criterion for Removable Singularities]\label{thm:riemann-removable}
Let $z_0$ be an isolated singularity of $f$. The following are equivalent:
\begin{enumerate}[label=(\roman*)]
  \item $z_0$ is a removable singularity.
  \item $f$ is bounded in some punctured neighborhood of $z_0$.
  \item $\lim_{z\to z_0}(z-z_0)f(z) = 0$.
\end{enumerate}
\end{theorem}

\begin{proof}\
\begin{itemize}[leftmargin=2em]
\item[\textbf{(i)$\Rightarrow$(ii)}:]
If $z_0$ is removable, $f$ extends to a holomorphic (hence continuous) function
at $z_0$. A continuous function is bounded near $z_0$.

\item[\textbf{(ii)$\Rightarrow$(iii)}:]
If $\abs{f(z)}\le M$ for $0<\abs{z-z_0}<r$, then
$\abs{(z-z_0)f(z)}\le M\abs{z-z_0}\to 0$ as $z\to z_0$.

\item[\textbf{(iii)$\Rightarrow$(i)}:]
Define $g(z)=(z-z_0)^2 f(z)$ for $z\ne z_0$ and $g(z_0)=0$. The hypothesis
$\lim_{z\to z_0}(z-z_0)f(z)=0$ implies $g$ is continuous at $z_0$, and in fact
\[
  g'(z_0) = \lim_{z\to z_0}\frac{g(z)-g(z_0)}{z-z_0}
  = \lim_{z\to z_0}(z-z_0)f(z) = 0.
\]
So $g$ is holomorphic at $z_0$ (and on the whole disk). The Taylor expansion of
$g$ at $z_0$ starts as $g(z)=g(z_0)+g'(z_0)(z-z_0)+\cdots = 0 + 0\cdot(z-z_0)+a_2(z-z_0)^2+\cdots$
Hence $g(z)=(z-z_0)^2 h(z)$ where $h$ is holomorphic at $z_0$. For $z\ne z_0$
we have $f(z)=h(z)$, so $h$ is the holomorphic extension of $f$.

Alternatively, using the Laurent series: for $n<0$, the coefficient is
\[
  a_n = \frac{1}{2\pi i}\oint_{\abs{\zeta-z_0}=\rho}
    \frac{f(\zeta)}{(\zeta-z_0)^{n+1}}\,\dd\zeta.
\]
The hypothesis (ii) gives $\abs{f(\zeta)}\le M$ on the circle. For $n\le -1$
(so $n+1\le 0$),
\[
  \abs{a_n} \le \frac{1}{2\pi}\cdot\frac{M}{\rho^{n+1}}\cdot 2\pi\rho
  = M\rho^{-n}.
\]
Letting $\rho\to 0^+$, and since $-n\ge 1>0$, we get $a_n=0$ for all $n<0$.
\end{itemize}
\end{proof}

\begin{example}[A Classic Removable Singularity]\label{ex:removable-sinc}
Consider $f(z)=\dfrac{\sin z}{z}$ for $z\ne 0$. Since $\sin z = z - z^3/6 + \cdots$,
we have
\[
  f(z) = 1 - \frac{z^2}{6} + \frac{z^4}{120} - \cdots
\]
The Laurent series has no negative powers, so $z=0$ is a removable singularity.
Setting $f(0)=1$ gives a holomorphic extension to all of $\C$.
\end{example}

%% ════════════════════════════════════════════════════════════
\section{Poles}
\label{sec:poles}
%% ════════════════════════════════════════════════════════════

\begin{proposition}[Characterization of Poles]\label{prop:poles-characterization}
Let $z_0$ be an isolated singularity of $f$. The following are equivalent:
\begin{enumerate}[label=(\roman*)]
  \item $z_0$ is a pole of order $m$.
  \item $\lim_{z\to z_0}\abs{f(z)}=+\infty$.
  \item There exist a positive integer $m$ and a holomorphic function $h$ with
    $h(z_0)\ne 0$ such that
    \[
      f(z) = \frac{h(z)}{(z-z_0)^m} \qquad\text{for } z\ne z_0.
    \]
  \item The function $g(z)=(z-z_0)^m f(z)$ extends to a holomorphic function at
    $z_0$ with $g(z_0)\ne 0$, and $m$ is the smallest such positive integer.
\end{enumerate}
\end{proposition}

\begin{proof}\
\begin{itemize}[leftmargin=2em]
\item[\textbf{(i)$\Rightarrow$(iii)}:]
If $z_0$ is a pole of order $m$, the Laurent expansion is
\[
  f(z) = \frac{a_{-m}}{(z-z_0)^m} + \cdots + \frac{a_{-1}}{z-z_0}
  + a_0 + a_1(z-z_0) + \cdots
\]
with $a_{-m}\ne 0$. Setting $h(z)=(z-z_0)^m f(z)=a_{-m}+a_{-m+1}(z-z_0)+\cdots$,
we see $h$ is holomorphic at $z_0$ with $h(z_0)=a_{-m}\ne 0$.

\item[\textbf{(iii)$\Rightarrow$(ii)}:]
$\abs{f(z)} = \abs{h(z)}/\abs{z-z_0}^m \to +\infty$ as $z\to z_0$, since
$h(z_0)\ne 0$.

\item[\textbf{(ii)$\Rightarrow$(i)}:]
Since $\abs{f(z)}\to\infty$, for $\abs{z-z_0}$ small we have $f(z)\ne 0$, so
$g(z)=1/f(z)$ is holomorphic near $z_0$ with $g(z)\to 0$. By the Riemann criterion,
$g$ has a removable singularity at $z_0$ with $g(z_0)=0$. Since $g\not\equiv 0$,
the zero of $g$ at $z_0$ has some finite order $m\ge 1$: $g(z)=(z-z_0)^m\phi(z)$
with $\phi(z_0)\ne 0$. Then $f(z)=1/g(z)=(z-z_0)^{-m}/\phi(z)$, and $h(z)=1/\phi(z)$
is holomorphic with $h(z_0)\ne 0$.

The equivalence with (iv) is immediate from (iii).
\end{itemize}
\end{proof}

\begin{example}[Poles of Various Orders]\label{ex:poles-orders}
\begin{enumerate}[label=(\alph*)]
  \item $f(z)=\dfrac{1}{z-1}$ has a simple pole (pole of order 1) at $z=1$.
  \item $f(z)=\dfrac{1}{(z+i)^3}$ has a pole of order 3 at $z=-i$.
  \item $f(z)=\dfrac{e^z-1}{z^2}$ at $z=0$: since $e^z-1=z+z^2/2+\cdots$,
    we have $f(z)=1/z+1/2+\cdots$, so $z=0$ is a simple pole.
  \item $f(z)=\dfrac{\sin z}{z^3}$ at $z=0$: since $\sin z = z - z^3/6+\cdots$,
    we have $f(z) = 1/z^2 - 1/6 + \cdots$, a pole of order~2.
\end{enumerate}
\end{example}

%% ════════════════════════════════════════════════════════════
\section{Essential Singularities and the Casorati--Weierstrass Theorem}
\label{sec:essential}
%% ════════════════════════════════════════════════════════════

Essential singularities exhibit dramatically different behavior from removable
singularities and poles. While the function stays bounded near a removable
singularity and blows up near a pole, near an essential singularity the function
takes values arbitrarily close to every complex number.

\begin{theorem}[Casorati--Weierstrass Theorem]\label{thm:casorati-weierstrass}
Let $z_0$ be an essential singularity of $f$. Then for every $\delta>0$, the
image $f\bigl(D(z_0,\delta)\setminus\{z_0\}\bigr)$ is dense in $\C$. That is,
for every $w\in\C$ and every $\varepsilon>0$, there exists $z$ with
$0<\abs{z-z_0}<\delta$ and $\abs{f(z)-w}<\varepsilon$.
\end{theorem}

\begin{proof}
Suppose for contradiction that the image is not dense. Then there exist $w_0\in\C$,
$\varepsilon>0$, and $\delta>0$ such that $\abs{f(z)-w_0}\ge\varepsilon$ for all
$z$ with $0<\abs{z-z_0}<\delta$.

Define $g(z)=\dfrac{1}{f(z)-w_0}$ for $0<\abs{z-z_0}<\delta$. Then $g$ is
holomorphic and $\abs{g(z)}\le 1/\varepsilon$, so $g$ is bounded. By the Riemann
criterion (Theorem~\ref{thm:riemann-removable}), $z_0$ is a removable singularity
of $g$, so $g$ extends holomorphically to $z_0$.

\medskip\noindent
\textbf{Case 1:} $g(z_0)\ne 0$. Then $f(z)=w_0+1/g(z)$ is holomorphic at $z_0$,
contradicting the assumption that $z_0$ is a singularity of $f$ (let alone an
essential one).

\medskip\noindent
\textbf{Case 2:} $g(z_0)=0$. Then $g$ has a zero of some order $m\ge 1$:
$g(z)=(z-z_0)^m\phi(z)$ with $\phi(z_0)\ne 0$. Hence
$f(z)=w_0+1/g(z)=w_0+(z-z_0)^{-m}/\phi(z)$, which shows $z_0$ is a pole of
order $m$, again a contradiction.

In both cases we reach a contradiction, so the image must be dense.
\end{proof}

\begin{remark}[The Great Picard Theorem]\label{rem:picard}
The Casorati--Weierstrass theorem is a ``weak'' version of a much deeper result.
The \textbf{Great Picard Theorem} states that in any punctured neighborhood of an
essential singularity, $f$ assumes every complex value infinitely often, with at
most one exception. For instance, $e^{1/z}$ near $z=0$ takes every value except
$0$ infinitely often. The proof of Picard's theorem requires more advanced tools
(modular functions or Bloch--Schottky type theorems) and is beyond the scope of
this course.
\end{remark}

%% ════════════════════════════════════════════════════════════
\section{Worked Examples of Laurent Expansions}
\label{sec:laurent-examples}
%% ════════════════════════════════════════════════════════════

\begin{example}[Laurent Expansion of $1/(z(z-1))$]\label{ex:laurent-1}
We find the Laurent series of $f(z)=\dfrac{1}{z(z-1)}$ in different annular regions
centered at $z_0=0$.

\medskip\noindent
\textbf{Partial fractions:}
\[
  \frac{1}{z(z-1)} = \frac{-1}{z} + \frac{1}{z-1}.
\]

\medskip\noindent
\textbf{Region I: $0<\abs{z}<1$.}
We need to expand $1/(z-1)$ in powers of $z$. Since $\abs{z}<1$:
\[
  \frac{1}{z-1} = \frac{-1}{1-z} = -\sum_{n=0}^{\infty}z^n.
\]
Therefore
\[
  f(z) = -\frac{1}{z} - \sum_{n=0}^{\infty}z^n
  = -\frac{1}{z} - 1 - z - z^2 - \cdots
\]
This is the Laurent series in $0<\abs{z}<1$. The singularity at $z=0$ is a
simple pole.

\begin{center}
\begin{tikzpicture}[scale=1.0]
  \fill[blue!8] (0,0) circle (1.5);
  \fill[white] (0,0) circle (0.08);
  \draw[thick, blue!70!black, dashed] (0,0) circle (1.5);
  \fill (0,0) circle (1.5pt) node[below left]{$0$};
  \fill (1.5,0) circle (1.5pt) node[right]{$1$};
  \node at (0,-2) {\small Region I: $0<|z|<1$};
\end{tikzpicture}
\end{center}

\medskip\noindent
\textbf{Region II: $\abs{z}>1$.}
We need to expand $1/(z-1)$ in negative powers of $z$. Since $\abs{z}>1$,
$\abs{1/z}<1$:
\[
  \frac{1}{z-1} = \frac{1}{z}\cdot\frac{1}{1-1/z}
  = \frac{1}{z}\sum_{n=0}^{\infty}\frac{1}{z^n}
  = \sum_{n=0}^{\infty}\frac{1}{z^{n+1}}
  = \sum_{k=1}^{\infty}\frac{1}{z^k}.
\]
Therefore
\[
  f(z) = -\frac{1}{z} + \sum_{k=1}^{\infty}\frac{1}{z^k}
  = \sum_{k=2}^{\infty}\frac{1}{z^k}
  = \frac{1}{z^2}+\frac{1}{z^3}+\frac{1}{z^4}+\cdots
\]
This Laurent series converges for $\abs{z}>1$.
\end{example}

\begin{example}[Laurent Expansion of $e^{1/z}$]\label{ex:laurent-exp}
The function $f(z)=e^{1/z}$ is holomorphic on $\C\setminus\{0\}$, with an
isolated singularity at $z=0$. Replacing $z$ by $1/z$ in the exponential series:
\[
  e^{1/z} = \sum_{n=0}^{\infty}\frac{1}{n!\,z^n}
  = 1 + \frac{1}{z} + \frac{1}{2!\,z^2} + \frac{1}{3!\,z^3} + \cdots
\]
This is the Laurent series in $0<\abs{z}<\infty$. Since there are infinitely many
non-zero coefficients $a_{-n}=1/n!$ for $n\ge 1$, the singularity at $z=0$ is
\textbf{essential}.

Note that $e^{1/z}$ never vanishes (consistent with the Great Picard Theorem: the
value $0$ is the single exceptional value).
\end{example}

\begin{example}[Laurent Series Centered at a Pole]\label{ex:laurent-centered-pole}
Find the Laurent expansion of $f(z)=\dfrac{z}{(z-1)^2(z-2)}$ centered at $z_0=1$.

\medskip\noindent
\textbf{Partial fractions:} Let $w=z-1$, so $z=w+1$ and $z-2=w-1$.
\[
  f(z) = \frac{w+1}{w^2(w-1)}.
\]
We decompose:
\[
  \frac{w+1}{w^2(w-1)} = \frac{A}{w} + \frac{B}{w^2} + \frac{C}{w-1}.
\]
Multiplying through: $w+1 = Aw(w-1)+B(w-1)+Cw^2$.
\begin{itemize}
  \item $w=0$: $1=-B$, so $B=-1$.
  \item $w=1$: $2=C$.
  \item Comparing $w^2$ coefficients: $0=A+C$, so $A=-2$.
\end{itemize}
Thus
\[
  f(z) = \frac{-2}{w} + \frac{-1}{w^2} + \frac{2}{w-1}, \qquad w=z-1.
\]

\textbf{For $0<\abs{w}<1$:}
\[
  \frac{2}{w-1} = \frac{-2}{1-w} = -2\sum_{n=0}^{\infty}w^n.
\]
So
\[
  f(z) = -\frac{1}{(z-1)^2} - \frac{2}{z-1} - 2\sum_{n=0}^{\infty}(z-1)^n.
\]
This exhibits a pole of order 2 at $z_0=1$.
\end{example}

\begin{example}[Laurent Expansion of $\cos(1/z)$]\label{ex:laurent-cos}
\[
  \cos\left(\frac{1}{z}\right)
  = \sum_{n=0}^{\infty}\frac{(-1)^n}{(2n)!\,z^{2n}}
  = 1 - \frac{1}{2!\,z^2} + \frac{1}{4!\,z^4} - \cdots
\]
This converges for $0<\abs{z}<\infty$. Since infinitely many negative-power
coefficients are non-zero, $z=0$ is an essential singularity of $\cos(1/z)$.
\end{example}

\begin{example}[Laurent Series in Multiple Annuli]\label{ex:laurent-multiple-annuli}
Consider $f(z)=\dfrac{1}{(z-1)(z-3)}$, centered at $z_0=0$. The singularities are
at $z=1$ and $z=3$, giving three regions.

By partial fractions: $f(z)=\dfrac{1}{2}\left(\dfrac{1}{z-3}-\dfrac{1}{z-1}\right)$.

\medskip\noindent
\textbf{Region I: $\abs{z}<1$.} Both $1/(z-1)$ and $1/(z-3)$ expand in non-negative
powers of $z$:
\begin{align*}
  \frac{-1}{z-1} &= \frac{1}{1-z} = \sum_{n=0}^{\infty}z^n, \\
  \frac{1}{z-3} &= \frac{-1}{3}\cdot\frac{1}{1-z/3}
  = -\frac{1}{3}\sum_{n=0}^{\infty}\frac{z^n}{3^n}
  = -\sum_{n=0}^{\infty}\frac{z^n}{3^{n+1}}.
\end{align*}
So $f(z)=\dfrac{1}{2}\sum_{n=0}^{\infty}\left(-\dfrac{1}{3^{n+1}}-1\right)z^n$
(a Taylor series; this region has no singularity).

\medskip\noindent
\textbf{Region II: $1<\abs{z}<3$.} We expand $1/(z-1)$ in negative powers and
$1/(z-3)$ in non-negative powers:
\begin{align*}
  \frac{-1}{z-1} &= \frac{-1}{z}\cdot\frac{1}{1-1/z}
  = -\sum_{n=0}^{\infty}\frac{1}{z^{n+1}}
  = -\sum_{k=1}^{\infty}\frac{1}{z^k}, \\
  \frac{1}{z-3} &= -\sum_{n=0}^{\infty}\frac{z^n}{3^{n+1}}.
\end{align*}

\medskip\noindent
\textbf{Region III: $\abs{z}>3$.} Both terms expand in negative powers:
\begin{align*}
  \frac{-1}{z-1} &= -\sum_{k=1}^{\infty}\frac{1}{z^k}, \\
  \frac{1}{z-3} &= \frac{1}{z}\cdot\frac{1}{1-3/z}
  = \sum_{n=0}^{\infty}\frac{3^n}{z^{n+1}}
  = \sum_{k=1}^{\infty}\frac{3^{k-1}}{z^k}.
\end{align*}
So $f(z) = \dfrac{1}{2}\sum_{k=1}^{\infty}\dfrac{3^{k-1}-1}{z^k}$.
\end{example}

\begin{center}
\begin{tikzpicture}[scale=0.75]
  %% Three regions
  \fill[green!8] (0,0) circle (1.3);
  \fill[blue!6] (0,0) circle (3.9);
  \fill[white] (0,0) circle (1.3);
  \fill[green!8] (0,0) circle (1.3);
  %% Outer region (outside of r=3) - indicated by label
  \draw[thick, dashed, green!60!black] (0,0) circle (1.3);
  \draw[thick, dashed, blue!60!black] (0,0) circle (3.9);
  %% Singularities
  \fill[red!80!black] (1.3,0) circle (2pt) node[above right,font=\small]{$1$};
  \fill[red!80!black] (3.9,0) circle (2pt) node[above right,font=\small]{$3$};
  \fill (0,0) circle (1.5pt) node[below left]{$0$};
  %% Region labels
  \node[green!50!black,font=\small] at (0,0.5) {I};
  \node[blue!50!black,font=\small] at (0,2.5) {II};
  \node[font=\small] at (0,4.5) {III};
\end{tikzpicture}
\end{center}

%% ════════════════════════════════════════════════════════════
\section{Singularity at Infinity}
\label{sec:singularity-infinity}
%% ════════════════════════════════════════════════════════════

\begin{definition}[Singularity at Infinity]\label{def:singularity-infinity}
A function $f$ is said to have an \textbf{isolated singularity at $\infty$} if
$f$ is holomorphic on $\{z:\abs{z}>R\}$ for some $R>0$. The nature of this
singularity is defined by studying $g(w)=f(1/w)$ near $w=0$.
\begin{enumerate}[label=(\roman*)]
  \item $f$ has a \textbf{removable singularity at $\infty$} if $g$ has a
    removable singularity at $0$ (equivalently, $f$ is bounded for $\abs{z}$
    large).
  \item $f$ has a \textbf{pole of order $m$ at $\infty$} if $g$ has a pole of
    order $m$ at $0$ (equivalently, $f(z)/z^m$ has a finite non-zero limit as
    $\abs{z}\to\infty$).
  \item $f$ has an \textbf{essential singularity at $\infty$} if $g$ has an
    essential singularity at $0$.
\end{enumerate}
\end{definition}

\begin{remark}\label{rem:infinity-laurent}
If $f$ is holomorphic for $\abs{z}>R$, it has a Laurent expansion
$f(z)=\sum_{n=-\infty}^{\infty}a_n z^n$ converging for $\abs{z}>R$. The
singularity at $\infty$ is:
\begin{itemize}
  \item removable if $a_n=0$ for all $n\ge 1$;
  \item a pole of order $m$ if $a_m\ne 0$ and $a_n=0$ for $n>m$;
  \item essential if $a_n\ne 0$ for infinitely many $n>0$.
\end{itemize}
Note the reversal: at $\infty$, it is the \emph{positive} powers that constitute
the ``singular part.''
\end{remark}

\begin{example}[Singularities at Infinity]\label{ex:singularity-infinity}
\begin{enumerate}[label=(\alph*)]
  \item $f(z)=\dfrac{1}{z}$: $g(w)=f(1/w)=w$ has a zero at $0$, so $f$ has a
    removable singularity at $\infty$ (with $f(\infty)=0$).
  \item $f(z)=z^3+2z$: $g(w)=1/w^3+2/w$ has a pole of order 3 at $0$. So $f$
    has a pole of order 3 at $\infty$.
  \item $f(z)=e^z$: $g(w)=e^{1/w}$ has an essential singularity at $0$. So
    $e^z$ has an essential singularity at $\infty$.
  \item A rational function $p(z)/q(z)$ with $\deg p > \deg q$ has a pole at
    $\infty$ of order $\deg p - \deg q$.
\end{enumerate}
\end{example}

\begin{proposition}\label{prop:entire-singularity-infinity}
An entire function $f$ has a pole of order $m$ at $\infty$ if and only if $f$ is
a polynomial of degree $m$. An entire function has a removable singularity at
$\infty$ if and only if it is constant.
\end{proposition}

\begin{proof}
An entire function has Taylor series $f(z)=\sum_{n=0}^{\infty}a_n z^n$ converging
for all $z$.

If $f$ has a pole of order $m$ at $\infty$, then $a_n=0$ for all $n>m$
(with $a_m\ne 0$), so $f$ is a polynomial of degree $m$.

If $f$ has a removable singularity at $\infty$, then $a_n=0$ for all $n\ge 1$,
so $f=a_0$ is constant. (Alternatively, a removable singularity at $\infty$ means
$f$ is bounded, and the conclusion follows from Liouville's theorem.)
\end{proof}

%% ════════════════════════════════════════════════════════════
\section{Meromorphic Functions}
\label{sec:meromorphic}
%% ════════════════════════════════════════════════════════════

\begin{definition}[Meromorphic Function]\label{def:meromorphic}
A function $f$ is \textbf{meromorphic} on an open set $\Omega\subseteq\C$ if
$f$ is holomorphic on $\Omega$ except for isolated singularities, all of which
are poles.
\end{definition}

\begin{example}\label{ex:meromorphic}
\begin{enumerate}[label=(\alph*)]
  \item Every rational function $p(z)/q(z)$ is meromorphic on $\C$, with poles at
    the zeros of $q$.
  \item $\tan z = \sin z / \cos z$ is meromorphic on $\C$, with simple poles at
    $z=(n+1/2)\pi$, $n\in\Z$.
  \item $e^{1/z}$ is \emph{not} meromorphic on $\C$ (the singularity at $0$ is essential).
\end{enumerate}
\end{example}

\begin{remark}\label{rem:meromorphic-field}
The set of meromorphic functions on a connected open set $\Omega$ forms a
\emph{field} under pointwise addition and multiplication (with the convention
that poles are handled in the natural way). This field contains the field of
holomorphic functions as a subring.
\end{remark}

%% ════════════════════════════════════════════════════════════
\section{Further Results on Singularities}
\label{sec:further-singularities}
%% ════════════════════════════════════════════════════════════

\begin{proposition}[Behavior Near Different Singularities --- Summary]\label{prop:singularity-summary}
Let $z_0$ be an isolated singularity of $f$. Then:
\begin{enumerate}[label=(\roman*)]
  \item $z_0$ is removable $\iff$ $f$ is bounded near $z_0$ $\iff$ $\lim_{z\to z_0}f(z)$ exists in $\C$.
  \item $z_0$ is a pole $\iff$ $\lim_{z\to z_0}\abs{f(z)}=+\infty$.
  \item $z_0$ is essential $\iff$ $\lim_{z\to z_0}f(z)$ does not exist in $\C\cup\{\infty\}$
    $\iff$ $f(D(z_0,\delta)\setminus\{z_0\})$ is dense in $\C$ for all $\delta>0$.
\end{enumerate}
\end{proposition}

\begin{proof}
Parts (i) and (ii) are Theorem~\ref{thm:riemann-removable} and
Proposition~\ref{prop:poles-characterization}. Part~(iii): if $z_0$ is essential,
it is neither removable nor a pole, so $f$ neither has a limit in $\C$ nor tends
to $\infty$. The density statement is the Casorati--Weierstrass theorem
(Theorem~\ref{thm:casorati-weierstrass}).
\end{proof}

\begin{proposition}[Counting Negative Laurent Coefficients]\label{prop:pole-order-residue}
If $z_0$ is a pole of order $m$ of $f$, then the Laurent expansion at $z_0$ is
\[
  f(z) = \frac{a_{-m}}{(z-z_0)^m} + \frac{a_{-m+1}}{(z-z_0)^{m-1}}
  + \cdots + \frac{a_{-1}}{z-z_0} + a_0 + a_1(z-z_0) + \cdots
\]
The coefficient $a_{-1}$ is called the \textbf{residue} of $f$ at $z_0$, denoted
$\Res(f,z_0)$ or $\Res_{z=z_0}f(z)$. For a simple pole ($m=1$),
\[
  \Res(f,z_0) = \lim_{z\to z_0}(z-z_0)f(z).
\]
For a pole of order $m$,
\[
  \Res(f,z_0) = \frac{1}{(m-1)!}\lim_{z\to z_0}\frac{\dd^{m-1}}{\dd z^{m-1}}
    \bigl[(z-z_0)^m f(z)\bigr].
\]
\end{proposition}

\begin{proof}
The Laurent expansion with finitely many negative terms is immediate from the
definition of a pole of order $m$. For the residue formula: $(z-z_0)^m f(z)=
a_{-m}+a_{-m+1}(z-z_0)+\cdots+a_{-1}(z-z_0)^{m-1}+\cdots$ is holomorphic at
$z_0$. Taking the $(m-1)$-st derivative and evaluating at $z_0$ extracts the
coefficient of $(z-z_0)^{m-1}$, which is $(m-1)!\,a_{-1}$.
\end{proof}

\begin{example}[Computing a Residue]\label{ex:residue-computation}
Find the residue of $f(z)=\dfrac{e^z}{(z-1)^2}$ at $z_0=1$.

Here $z_0=1$ is a pole of order 2. Using the formula:
\[
  \Res(f,1) = \lim_{z\to 1}\frac{\dd}{\dd z}\bigl[(z-1)^2 f(z)\bigr]
  = \lim_{z\to 1}\frac{\dd}{\dd z}(e^z) = e.
\]
\end{example}

%% ════════════════════════════════════════════════════════════
\section{Exercises}
\label{sec:exercises-ch6}
%% ════════════════════════════════════════════════════════════

\begin{exercise}[$\bigstar$]\label{exo:laurent-basic}
Find the Laurent series of $f(z)=\dfrac{1}{z^2-z}$ centered at $z_0=0$ in the
regions:
\begin{enumerate}[label=(\alph*)]
  \item $0<\abs{z}<1$
  \item $\abs{z}>1$
\end{enumerate}
Identify the type of singularity at $z=0$.
\end{exercise}

\begin{exercise}[$\bigstar$]\label{exo:classify-singularities}
Classify the singularities (removable, pole with order, or essential) of:
\begin{enumerate}[label=(\alph*)]
  \item $\dfrac{z^2-1}{z-1}$ at $z=1$
  \item $\dfrac{1-\cos z}{z^2}$ at $z=0$
  \item $z^2 e^{1/z}$ at $z=0$
  \item $\dfrac{1}{\sin z}$ at $z=0$
  \item $\dfrac{z}{\sin^2 z}$ at $z=0$
\end{enumerate}
\end{exercise}

\begin{exercise}[$\bigstar$]\label{exo:laurent-exp}
Find the Laurent expansion of $f(z)=e^{z+1/z}$ centered at $z_0=0$, valid for
$0<\abs{z}<\infty$. Express the coefficient $a_n$ as a series.

\textit{Hint:} Write $e^{z+1/z}=e^z\cdot e^{1/z}$ and multiply the two series.
\end{exercise}

\begin{exercise}[$\bigstar\!\bigstar$]\label{exo:casorati-example}
Let $f(z)=e^{1/z}$.
\begin{enumerate}[label=(\alph*)]
  \item Show directly (without invoking the Casorati--Weierstrass theorem) that
    for every $w\in\C\setminus\{0\}$ and every $\delta>0$, there exists $z$ with
    $0<\abs{z}<\delta$ and $f(z)=w$.
  \item Why is $w=0$ not attained?
\end{enumerate}
\end{exercise}

\begin{exercise}[$\bigstar\!\bigstar$]\label{exo:laurent-multiple-centers}
Find the Laurent expansions of $f(z)=\dfrac{1}{z(z-1)(z-2)}$ in the annuli:
\begin{enumerate}[label=(\alph*)]
  \item $0<\abs{z}<1$
  \item $1<\abs{z}<2$
  \item $\abs{z}>2$
\end{enumerate}
\end{exercise}

\begin{exercise}[$\bigstar\!\bigstar$]\label{exo:removable-criterion}
Let $f$ be holomorphic on $D(z_0,R)\setminus\{z_0\}$ and suppose that
$\lim_{z\to z_0}(z-z_0)^{\alpha}f(z)=0$ for some $\alpha\in(0,1)$. Prove that
$z_0$ is a removable singularity.

\textit{Hint:} Show that the hypothesis implies $\lim_{z\to z_0}(z-z_0)f(z)=0$,
then apply Riemann's criterion.
\end{exercise}

\begin{exercise}[$\bigstar\!\bigstar$]\label{exo:residue-computations}
Compute the residues at all singularities of:
\begin{enumerate}[label=(\alph*)]
  \item $\dfrac{z^3}{z^2+1}$
  \item $\dfrac{e^z}{z^2(z^2+\pi^2)}$
  \item $\dfrac{1}{z^2\sin z}$
\end{enumerate}
\end{exercise}

\begin{exercise}[$\bigstar\!\bigstar$]\label{exo:singularity-infinity-classify}
Classify the singularity at $\infty$ for each of the following:
\begin{enumerate}[label=(\alph*)]
  \item $f(z)=\dfrac{z^3+1}{z^2-1}$
  \item $f(z)=e^{1/z}$
  \item $f(z)=\sin z$
  \item $f(z)=\dfrac{z}{z+1}$
\end{enumerate}
\end{exercise}

\begin{exercise}[$\bigstar\!\bigstar\!\bigstar$]\label{exo:essential-dense}
Let $f$ have an essential singularity at $z_0$. Prove that for every $n\ge 1$,
the function $f^n$ (i.e., $f$ raised to the $n$-th power) also has an essential
singularity at $z_0$.
\end{exercise}

\begin{exercise}[$\bigstar\!\bigstar\!\bigstar$]\label{exo:isolated-zeros}
\textbf{(Identity Theorem.)}
Let $f$ be holomorphic on a connected open set $\Omega$, and suppose $f$ is not
identically zero. Prove that the zeros of $f$ are isolated: for every $z_0\in\Omega$
with $f(z_0)=0$, there exists $\delta>0$ such that $f(z)\ne 0$ for
$0<\abs{z-z_0}<\delta$.

\textit{Hint:} Use the Taylor expansion of $f$ at $z_0$.
\end{exercise}

\begin{exercise}[$\bigstar\!\bigstar\!\bigstar$]\label{exo:meromorphic-rational}
Let $f$ be meromorphic on $\C$ with only finitely many poles, and suppose $f$ has
a pole (or removable singularity) at $\infty$. Prove that $f$ is a rational
function.

\textit{Hint:} Subtract the principal parts at each pole to reduce to an entire
function with a pole at $\infty$.
\end{exercise}

\begin{exercise}[$\bigstar\!\bigstar\!\bigstar$]\label{exo:weierstrass-no-essential}
Prove that if $f$ is meromorphic on $D(z_0,R)\setminus\{z_0\}$ (i.e., $f$ is
holomorphic except for poles in the punctured disk), and $z_0$ is not an
accumulation point of the poles, then $f$ has an isolated singularity at $z_0$
(which may be removable, a pole, or essential).
\end{exercise}
%%%%%%%%%%%%%%%%%%%%%%%%%%%%%%%%%%%%%%%%%%%%%%%%%%%%%%%%%%%%%%%%%%%%
%  CHAPTER 7 — Residue Theorem and Applications
%%%%%%%%%%%%%%%%%%%%%%%%%%%%%%%%%%%%%%%%%%%%%%%%%%%%%%%%%%%%%%%%%%%%
\chapter{Residue Theorem and Applications}\label{ch:residues}

\section{Residues: Definition and Computation}\label{sec:residue-def}

Throughout this chapter, we let $\Omega\subset\C$ be an open set and
$f\colon\Omega\setminus\{z_0\}\to\C$ a function holomorphic on the
punctured neighbourhood of an isolated singularity $z_0$.

\begin{definition}[Residue]\label{def:residue}
Let $z_0$ be an isolated singularity of $f$.  The \emph{residue} of $f$
at $z_0$ is the coefficient $a_{-1}$ in the Laurent expansion
\[
  f(z)=\sum_{n=-\infty}^{\infty}a_n(z-z_0)^n,
\]
i.e.\
\[
  \Res_{z=z_0}f(z)\;=\;a_{-1}
  \;=\;\frac{1}{2\pi i}\oint_{\gamma}f(z)\,\dd z,
\]
where $\gamma$ is any positively oriented simple closed curve in $\Omega$
enclosing $z_0$ and no other singularity of $f$.
\end{definition}

\begin{remark}[Residue is intrinsic]\label{rem:residue-intrinsic}
The residue does not depend on the choice of the curve $\gamma$: by the
deformation theorem for holomorphic functions, any two such curves yield
the same integral.
\end{remark}

\subsection{Residue at a simple pole}

\begin{proposition}[Residue at a simple pole]\label{prop:res-simple}
If $f$ has a simple pole at $z_0$, then
\[
  \Res_{z=z_0}f(z)=\lim_{z\to z_0}(z-z_0)f(z).
\]
\end{proposition}

\begin{proof}
Since $z_0$ is a simple pole, the Laurent series has the form
\[
  f(z)=\frac{a_{-1}}{z-z_0}+a_0+a_1(z-z_0)+\cdots
\]
Multiplying both sides by $(z-z_0)$ and letting $z\to z_0$ gives
$(z-z_0)f(z)\to a_{-1}$.
\end{proof}

\begin{proposition}[Quotient rule for simple poles]\label{prop:res-quotient}
If $f(z)=g(z)/h(z)$ where $g$ and $h$ are holomorphic near $z_0$,
$g(z_0)\neq 0$, $h(z_0)=0$, and $h'(z_0)\neq 0$ (so $z_0$ is a simple
zero of $h$), then
\[
  \Res_{z=z_0}\frac{g(z)}{h(z)}=\frac{g(z_0)}{h'(z_0)}.
\]
\end{proposition}

\begin{proof}
We have
\[
  \Res_{z=z_0}\frac{g(z)}{h(z)}
  =\lim_{z\to z_0}(z-z_0)\frac{g(z)}{h(z)}
  =\lim_{z\to z_0}g(z)\cdot\frac{z-z_0}{h(z)-h(z_0)}
  =\frac{g(z_0)}{h'(z_0)}.  \qedhere
\]
\end{proof}

\subsection{Residue at a pole of order \texorpdfstring{$n$}{n}}

\begin{proposition}[Residue at a pole of order $n$]\label{prop:res-order-n}
If $f$ has a pole of order $n\geq 1$ at $z_0$, then
\[
  \Res_{z=z_0}f(z)
  =\frac{1}{(n-1)!}\lim_{z\to z_0}
    \frac{\dd^{n-1}}{\dd z^{n-1}}\bigl[(z-z_0)^n f(z)\bigr].
\]
\end{proposition}

\begin{proof}
Write the Laurent expansion
\[
  f(z)=\sum_{k=-n}^{\infty}a_k(z-z_0)^k.
\]
Then $(z-z_0)^n f(z)=\sum_{k=-n}^{\infty}a_k(z-z_0)^{k+n}
=\sum_{j=0}^{\infty}a_{j-n}(z-z_0)^j$, which is holomorphic at $z_0$.
Its Taylor coefficients satisfy
\[
  a_{j-n}=\frac{1}{j!}\frac{\dd^{j}}{\dd z^{j}}
    \bigl[(z-z_0)^n f(z)\bigr]\bigg|_{z=z_0}.
\]
Setting $j=n-1$ gives $a_{-1}=\frac{1}{(n-1)!}
\frac{\dd^{n-1}}{\dd z^{n-1}}\bigl[(z-z_0)^n f(z)\bigr]\big|_{z=z_0}$.
\end{proof}

\begin{example}[Computing residues]\label{ex:res-compute-1}
Compute the residues of $f(z)=\dfrac{e^z}{(z-1)^2 z}$.

\medskip\noindent\textbf{Solution.}
The singularities are $z=0$ (simple pole) and $z=1$ (pole of order $2$).

\emph{At $z=0$:} Using Proposition~\ref{prop:res-simple},
\[
  \Res_{z=0}f(z)=\lim_{z\to 0}z\cdot\frac{e^z}{(z-1)^2 z}
  =\frac{e^0}{(0-1)^2}=1.
\]

\emph{At $z=1$:} Using Proposition~\ref{prop:res-order-n} with $n=2$,
\[
  \Res_{z=1}f(z)=\lim_{z\to 1}\frac{\dd}{\dd z}\left[(z-1)^2\cdot
  \frac{e^z}{(z-1)^2 z}\right]
  =\lim_{z\to 1}\frac{\dd}{\dd z}\frac{e^z}{z}
  =\lim_{z\to 1}\frac{e^z(z-1)}{z^2}=0.
\]
\end{example}

\begin{example}[Residue via Laurent series]\label{ex:res-laurent}
Find $\Res_{z=0}\dfrac{1}{z^2\sin z}$.

\medskip\noindent\textbf{Solution.}
We know $\sin z=z-z^3/6+z^5/120-\cdots$, so
\[
  \frac{1}{z^2\sin z}=\frac{1}{z^3(1-z^2/6+\cdots)}
  =\frac{1}{z^3}\left(1+\frac{z^2}{6}+\cdots\right)
  =\frac{1}{z^3}+\frac{1}{6z}+\cdots
\]
Therefore $\Res_{z=0}\dfrac{1}{z^2\sin z}=\dfrac{1}{6}$.
\end{example}

%-------------------------------------------------------------------
\section{The Residue Theorem}\label{sec:residue-thm}
%-------------------------------------------------------------------

\begin{theorem}[Residue Theorem]\label{thm:residue}
Let $\Omega\subset\C$ be a simply connected open set and let
$z_1,\dots,z_N$ be finitely many distinct points in $\Omega$.  If
$f\colon\Omega\setminus\{z_1,\dots,z_N\}\to\C$ is holomorphic and
$\gamma$ is a positively oriented simple closed contour in $\Omega$
enclosing exactly the singularities $z_1,\dots,z_N$, then
\[
  \boxed{\frac{1}{2\pi i}\oint_{\gamma}f(z)\,\dd z
  =\sum_{k=1}^{N}\Res_{z=z_k}f(z).}
\]
\end{theorem}

\begin{proof}
For each $k=1,\dots,N$ let $\gamma_k$ be a small positively oriented
circle of radius $\varepsilon_k$ centred at $z_k$, where the radii are
chosen small enough so that the disks $\overline{D(z_k,\varepsilon_k)}$
are pairwise disjoint and contained in the interior of $\gamma$.

Define the region
$G=(\text{interior of }\gamma)\setminus\bigcup_{k=1}^N
\overline{D(z_k,\varepsilon_k)}$.
Then $f$ is holomorphic on $\overline{G}$, and $\partial G$ consists of
$\gamma$ (traversed positively) together with the circles $\gamma_k$
(each traversed negatively, as seen from $G$).

By Cauchy's theorem applied to the multiply connected region $G$,
\[
  \oint_{\gamma}f(z)\,\dd z
  -\sum_{k=1}^{N}\oint_{\gamma_k}f(z)\,\dd z=0.
\]
Hence
\[
  \oint_{\gamma}f(z)\,\dd z
  =\sum_{k=1}^{N}\oint_{\gamma_k}f(z)\,\dd z
  =\sum_{k=1}^{N}2\pi i\,\Res_{z=z_k}f(z),
\]
where the last equality follows from the definition of the residue
(Definition~\ref{def:residue}).
\end{proof}

\begin{remark}[Winding number version]\label{rem:winding-residue}
More generally, if $\gamma$ is a closed curve (not necessarily simple)
that does not pass through any $z_k$, the residue theorem takes the form
\[
  \frac{1}{2\pi i}\oint_{\gamma}f(z)\,\dd z
  =\sum_{k=1}^{N}\ind(\gamma,z_k)\,\Res_{z=z_k}f(z),
\]
where $\ind(\gamma,z_k)=\frac{1}{2\pi i}\oint_{\gamma}\frac{\dd z}{z-z_k}$
is the winding number (index) of $\gamma$ about $z_k$.
\end{remark}

%-------------------------------------------------------------------
\section{Applications to Real Integrals}\label{sec:real-integrals}
%-------------------------------------------------------------------

One of the most striking applications of the residue theorem is the
evaluation of definite real integrals that are difficult or impossible
to compute by real-variable methods alone.  The general strategy is:
\begin{enumerate}[label=(\roman*)]
  \item embed the real integral as part of a complex contour integral,
  \item close the contour so that it encloses certain poles,
  \item apply the residue theorem,
  \item show that the contributions from the added arcs vanish
        (or can be evaluated) in a suitable limit.
\end{enumerate}

%-------------------------------------------------------------------
\subsection{Rational integrals on the real line}
\label{ssec:rational-integrals}
%-------------------------------------------------------------------

\begin{proposition}[Rational integrals]\label{prop:rational-integral}
Let $R(x)=P(x)/Q(x)$ be a rational function where
$\deg Q\geq\deg P+2$ and $Q$ has no real zeros.  Then
\[
  \int_{-\infty}^{+\infty}R(x)\,\dd x
  =2\pi i\sum_{\im z_k>0}\Res_{z=z_k}R(z),
\]
where the sum is over all poles of $R$ in the upper half-plane
$\{z:\im z>0\}$.
\end{proposition}

\begin{proof}
Consider the semicircular contour $\Gamma_R$ consisting of the segment
$[-R,R]$ on the real axis and the upper semicircle
$C_R=\{Re^{i\theta}:0\leq\theta\leq\pi\}$, traversed counter-clockwise.

\begin{center}
\begin{tikzpicture}[scale=1.1,>=Stealth]
  % Real axis
  \draw[->] (-3.8,0) -- (3.8,0) node[right]{$\re$};
  \draw[->] (0,-0.5) -- (0,3.3) node[above]{$\im$};
  % Semicircle
  \draw[blue,thick,->,postaction={decorate,decoration={markings,
    mark=at position 0.25 with {\arrow{>}},
    mark=at position 0.75 with {\arrow{>}}}}]
    (3,0) arc(0:180:3);
  \draw[blue,thick,->,postaction={decorate,decoration={markings,
    mark=at position 0.5 with {\arrow{>}}}}]
    (-3,0) -- (3,0);
  % Labels
  \node[below] at (3,0) {$R$};
  \node[below] at (-3,0) {$-R$};
  \node[blue] at (2.5,2.2) {$C_R$};
  % Poles
  \node[red] at (0.8,1.5) {$\times$};
  \node[red,right] at (0.9,1.5) {$z_1$};
  \node[red] at (-1,0.8) {$\times$};
  \node[red,left] at (-1.1,0.8) {$z_2$};
\end{tikzpicture}
\end{center}

For $R$ large enough, $\Gamma_R$ encloses all poles of $R(z)$ in the
upper half-plane.  By the residue theorem,
\[
  \oint_{\Gamma_R}R(z)\,\dd z
  =2\pi i\sum_{\im z_k>0}\Res_{z=z_k}R(z).
\]
Now we split:
\[
  \oint_{\Gamma_R}R(z)\,\dd z
  =\int_{-R}^{R}R(x)\,\dd x+\int_{C_R}R(z)\,\dd z.
\]
On $C_R$ with $z=Re^{i\theta}$, since $\deg Q\geq\deg P+2$, we have
$\abs{R(z)}\leq M/R^2$ for large $R$.  Therefore
\[
  \left|\int_{C_R}R(z)\,\dd z\right|\leq\frac{M}{R^2}\cdot\pi R
  =\frac{M\pi}{R}\to 0\quad\text{as }R\to\infty.
\]
Taking $R\to\infty$ yields the result.
\end{proof}

\begin{example}[A basic rational integral]\label{ex:rational-1}
Evaluate $\displaystyle I=\int_{-\infty}^{+\infty}\frac{\dd x}{1+x^2}$.

\medskip\noindent\textbf{Solution.}
The function $f(z)=1/(1+z^2)=1/((z+i)(z-i))$ has simple poles at
$z=\pm i$.  Only $z=i$ lies in the upper half-plane.
\[
  \Res_{z=i}\frac{1}{1+z^2}=\frac{1}{2i}.
\]
Therefore $I=2\pi i\cdot\dfrac{1}{2i}=\pi$.
\end{example}

\begin{example}[Degree condition is essential]\label{ex:rational-2}
Evaluate $\displaystyle I=\int_{-\infty}^{+\infty}\frac{\dd x}{(1+x^2)^2}$.

\medskip\noindent\textbf{Solution.}
Here $f(z)=1/(1+z^2)^2$ has a pole of order $2$ at $z=i$ in the upper
half-plane.  We compute
\[
  \Res_{z=i}\frac{1}{(z^2+1)^2}
  =\lim_{z\to i}\frac{\dd}{\dd z}\frac{1}{(z+i)^2}
  =\lim_{z\to i}\frac{-2}{(z+i)^3}
  =\frac{-2}{(2i)^3}=\frac{-2}{-8i}=\frac{1}{4i}.
\]
Therefore $I=2\pi i\cdot\dfrac{1}{4i}=\dfrac{\pi}{2}$.
\end{example}

%-------------------------------------------------------------------
\subsection{Trigonometric integrals}
\label{ssec:trig-integrals}
%-------------------------------------------------------------------

\begin{proposition}[Trigonometric integrals]\label{prop:trig-integral}
Let $R(\cos\theta,\sin\theta)$ be a rational function of $\cos\theta$
and $\sin\theta$ with no singularities on the unit circle.  Then
\[
  \int_0^{2\pi}R(\cos\theta,\sin\theta)\,\dd\theta
  =\oint_{\abs{z}=1}R\!\left(\frac{z+z^{-1}}{2},
  \frac{z-z^{-1}}{2i}\right)\frac{\dd z}{iz}.
\]
The right-hand side is evaluated by the residue theorem, summing over
poles inside the unit disk $\D$.
\end{proposition}

\begin{proof}
Set $z=e^{i\theta}$, so $\dd z=ie^{i\theta}\dd\theta=iz\,\dd\theta$,
hence $\dd\theta=\dd z/(iz)$.  Also
\[
  \cos\theta=\frac{e^{i\theta}+e^{-i\theta}}{2}=\frac{z+z^{-1}}{2},
  \qquad
  \sin\theta=\frac{e^{i\theta}-e^{-i\theta}}{2i}=\frac{z-z^{-1}}{2i}.
\]
As $\theta$ ranges over $[0,2\pi)$, $z$ traverses the unit circle
$\abs{z}=1$ once counter-clockwise.
\end{proof}

\begin{example}[Trigonometric integral]\label{ex:trig-1}
Evaluate $\displaystyle I=\int_0^{2\pi}\frac{\dd\theta}{2+\cos\theta}$.

\medskip\noindent\textbf{Solution.}
Substitute $z=e^{i\theta}$:
\[
  I=\oint_{\abs{z}=1}\frac{1}{2+\frac{z+z^{-1}}{2}}\,\frac{\dd z}{iz}
  =\oint_{\abs{z}=1}\frac{1}{\frac{4+z+z^{-1}}{2}}\,\frac{\dd z}{iz}
  =\oint_{\abs{z}=1}\frac{2\,\dd z}{iz(4+z+z^{-1})}.
\]
Multiply numerator and denominator by $z$:
\[
  I=\oint_{\abs{z}=1}\frac{2\,\dd z}{i(z^2+4z+1)}.
\]
The roots of $z^2+4z+1=0$ are $z=-2\pm\sqrt{3}$.
Only $z_0=-2+\sqrt{3}$ lies inside $\D$ (since $\abs{-2+\sqrt{3}}
=2-\sqrt{3}<1$).

\[
  \Res_{z=z_0}\frac{2}{i(z^2+4z+1)}
  =\frac{2}{i\cdot 2z_0+i\cdot 4}\bigg|_{z=z_0}
  =\frac{2}{i(2(-2+\sqrt{3})+4)}
  =\frac{2}{i\cdot 2\sqrt{3}}=\frac{1}{i\sqrt{3}}.
\]
Therefore
\[
  I=2\pi i\cdot\frac{1}{i\sqrt{3}}=\frac{2\pi}{\sqrt{3}}.
\]
\end{example}

\begin{example}[Another trigonometric integral]\label{ex:trig-2}
Evaluate $\displaystyle I=\int_0^{2\pi}\frac{\dd\theta}{5-4\cos\theta}$.

\medskip\noindent\textbf{Solution.}
With $z=e^{i\theta}$:
\[
  I=\oint_{\abs{z}=1}\frac{\dd z}{iz(5-2(z+z^{-1}))}
  =\oint_{\abs{z}=1}\frac{\dd z}{i(5z-2z^2-2)}
  =\oint_{\abs{z}=1}\frac{\dd z}{-i(2z^2-5z+2)}.
\]
Factor: $2z^2-5z+2=(2z-1)(z-2)$.  The roots are $z=1/2$ and $z=2$.
Only $z=1/2$ lies inside $\D$.
\[
  \Res_{z=1/2}\frac{1}{-i(2z-1)(z-2)}
  =\frac{1}{-i\cdot 2\cdot(1/2-2)}
  =\frac{1}{-i\cdot 2\cdot(-3/2)}
  =\frac{1}{3i}.
\]
Therefore $I=2\pi i\cdot\dfrac{1}{3i}=\dfrac{2\pi}{3}$.
\end{example}

%-------------------------------------------------------------------
\subsection{Jordan's lemma and Fourier-type integrals}
\label{ssec:jordan}
%-------------------------------------------------------------------

\begin{lemma}[Jordan's Lemma]\label{lem:jordan}
Let $f$ be continuous on $\{z:\im z\geq 0,\;\abs{z}\geq R_0\}$ for some
$R_0>0$, and suppose $\abs{f(Re^{i\theta})}\to 0$ uniformly in
$\theta\in[0,\pi]$ as $R\to\infty$.  Then for every $a>0$,
\[
  \lim_{R\to\infty}\int_{C_R}f(z)e^{iaz}\,\dd z=0,
\]
where $C_R=\{Re^{i\theta}:0\leq\theta\leq\pi\}$ is the upper semicircle.
\end{lemma}

\begin{proof}
Set $M(R)=\max_{0\leq\theta\leq\pi}\abs{f(Re^{i\theta})}$.  On $C_R$,
$z=Re^{i\theta}$, so $e^{iaz}=e^{iaR\cos\theta-aR\sin\theta}$ and
$\abs{e^{iaz}}=e^{-aR\sin\theta}$.  Therefore
\begin{align*}
  \left|\int_{C_R}f(z)e^{iaz}\,\dd z\right|
  &\leq M(R)\int_0^{\pi}e^{-aR\sin\theta}\,R\,\dd\theta \\
  &= 2M(R)R\int_0^{\pi/2}e^{-aR\sin\theta}\,\dd\theta.
\end{align*}
By Jordan's inequality, $\sin\theta\geq 2\theta/\pi$ for
$\theta\in[0,\pi/2]$, so
\[
  \int_0^{\pi/2}e^{-aR\sin\theta}\,\dd\theta
  \leq\int_0^{\pi/2}e^{-2aR\theta/\pi}\,\dd\theta
  =\frac{\pi}{2aR}(1-e^{-aR})\leq\frac{\pi}{2aR}.
\]
Hence
\[
  \left|\int_{C_R}f(z)e^{iaz}\,\dd z\right|
  \leq 2M(R)R\cdot\frac{\pi}{2aR}=\frac{\pi M(R)}{a}\to 0.
  \qedhere
\]
\end{proof}

\begin{center}
\begin{tikzpicture}[scale=1.0,>=Stealth]
  \draw[->] (-3.5,0) -- (3.5,0) node[right]{$\re$};
  \draw[->] (0,-0.5) -- (0,3.2) node[above]{$\im$};
  % Semicircular contour
  \draw[blue,thick,->] (-2.8,0) -- (0,0);
  \draw[blue,thick,->] (0,0) -- (2.8,0);
  \draw[blue,thick,->] (2.8,0) arc(0:90:2.8);
  \draw[blue,thick,->] (0,2.8) arc(90:180:2.8);
  \node[below] at (2.8,0) {$R$};
  \node[below] at (-2.8,0) {$-R$};
  \node[blue] at (2.3,2.0) {$C_R$};
  % Shading to indicate decay
  \fill[blue!8] (0,0) -- (2.8,0) arc(0:180:2.8) -- cycle;
  \node at (0,1.3) {$e^{-aR\sin\theta}\to 0$};
\end{tikzpicture}
\end{center}

\begin{example}[Fourier-type integral]\label{ex:fourier-1}
Evaluate $\displaystyle I=\int_{-\infty}^{+\infty}
\frac{\cos x}{x^2+1}\,\dd x$.

\medskip\noindent\textbf{Solution.}
Since $\cos x=\re(e^{ix})$, we compute
\[
  J=\int_{-\infty}^{+\infty}\frac{e^{ix}}{x^2+1}\,\dd x
\]
and take the real part.  Using the upper semicircular contour with
$f(z)=1/(z^2+1)$, Jordan's lemma gives $\int_{C_R}f(z)e^{iz}\dd z\to 0$.
The only pole in the upper half-plane is $z=i$:
\[
  \Res_{z=i}\frac{e^{iz}}{z^2+1}
  =\frac{e^{i\cdot i}}{2i}=\frac{e^{-1}}{2i}.
\]
Therefore $J=2\pi i\cdot\dfrac{e^{-1}}{2i}=\dfrac{\pi}{e}$, and
$I=\re(J)=\dfrac{\pi}{e}$.
\end{example}

\begin{example}[Integral with $\sin x$ in the numerator]\label{ex:fourier-2}
Evaluate $\displaystyle I=\int_{0}^{+\infty}\frac{\sin x}{x}\,\dd x$.

\medskip\noindent\textbf{Solution.}
The integrand $(\sin x)/x$ extends continuously to $x=0$ (with value~$1$),
so there is no issue at the origin in the real sense.  However, the
complex function $e^{iz}/z$ has a simple pole at $z=0$.  We use an
\emph{indented contour}: let $\Gamma$ consist of
\begin{enumerate}[label=(\roman*)]
  \item the segment $[\varepsilon,R]$ on the real axis,
  \item the upper semicircle $C_R$ from $R$ to $-R$,
  \item the segment $[-R,-\varepsilon]$ on the real axis,
  \item the small semicircle $c_\varepsilon$ from $-\varepsilon$ to
        $\varepsilon$, traversed clockwise (i.e.\ in the upper half-plane).
\end{enumerate}

\begin{center}
\begin{tikzpicture}[scale=1.0,>=Stealth]
  \draw[->] (-3.5,0) -- (3.5,0) node[right]{$\re$};
  \draw[->] (0,-0.5) -- (0,3.2) node[above]{$\im$};
  % Large semicircle
  \draw[blue,thick,->] (3,0) arc(0:90:3);
  \draw[blue,thick,->] (0,3) arc(90:180:3);
  % Real segments
  \draw[blue,thick,->] (0.4,0) -- (1.7,0);
  \draw[blue,thick,->] (1.7,0) -- (3,0);
  \draw[blue,thick,->] (-3,0) -- (-1.7,0);
  \draw[blue,thick,->] (-1.7,0) -- (-0.4,0);
  % Small semicircle (clockwise = going through upper half-plane from -eps to +eps)
  \draw[red,thick,->] (-0.4,0) arc(180:90:0.4);
  \draw[red,thick,->] (0,0.4) arc(90:0:0.4);
  % Labels
  \node[below] at (3,0) {$R$};
  \node[below] at (-3,0) {$-R$};
  \node[below] at (0.4,0) {$\varepsilon$};
  \node[below] at (-0.4,0) {$-\varepsilon$};
  \node[blue] at (2.4,2.0) {$C_R$};
  \node[red] at (0.1,0.75) {$c_\varepsilon$};
\end{tikzpicture}
\end{center}

Since $e^{iz}/z$ is holomorphic inside $\Gamma$ (the pole at $0$ has been
excluded), the residue theorem gives
\[
  \oint_{\Gamma}\frac{e^{iz}}{z}\,\dd z=0.
\]
By Jordan's lemma, $\int_{C_R}\frac{e^{iz}}{z}\dd z\to 0$ as
$R\to\infty$.  For the small semicircle, parametrise
$z=\varepsilon e^{i\theta}$, $\theta:\pi\to 0$ (clockwise):
\[
  \int_{c_\varepsilon}\frac{e^{iz}}{z}\,\dd z
  =\int_{\pi}^{0}\frac{e^{i\varepsilon e^{i\theta}}}
  {\varepsilon e^{i\theta}}\,i\varepsilon e^{i\theta}\,\dd\theta
  =i\int_{\pi}^{0}e^{i\varepsilon e^{i\theta}}\,\dd\theta
  \xrightarrow{\varepsilon\to 0}i\int_{\pi}^{0}1\,\dd\theta=-i\pi.
\]
The real-axis contributions give
\[
  \int_{\varepsilon}^{R}\frac{e^{ix}}{x}\,\dd x
  +\int_{-R}^{-\varepsilon}\frac{e^{ix}}{x}\,\dd x
  =\int_{\varepsilon}^{R}\frac{e^{ix}-e^{-ix}}{x}\,\dd x
  =2i\int_{\varepsilon}^{R}\frac{\sin x}{x}\,\dd x.
\]
Combining and letting $R\to\infty$, $\varepsilon\to 0$:
\[
  2i\int_0^{\infty}\frac{\sin x}{x}\,\dd x-i\pi=0
  \implies
  \int_0^{\infty}\frac{\sin x}{x}\,\dd x=\frac{\pi}{2}.
\]
\end{example}

%-------------------------------------------------------------------
\subsection{Integrals with a logarithmic branch cut}
\label{ssec:log-branch-cut}
%-------------------------------------------------------------------

A powerful technique for evaluating integrals of the form
$\int_0^{\infty}x^{\alpha-1}R(x)\,\dd x$ (where $0<\alpha<1$ and $R$
is rational with no poles on $[0,\infty)$) uses the \emph{keyhole contour}.

\begin{center}
\begin{tikzpicture}[scale=1.0,>=Stealth]
  \draw[->] (-3.5,0) -- (3.5,0) node[right]{$\re$};
  \draw[->] (0,-3.2) -- (0,3.2) node[above]{$\im$};
  % Large circle
  \draw[blue,thick] (3,0.15) arc(3:357:3);
  \draw[blue,thick,->] (3,0.15) arc(3:180:3);
  \draw[blue,thick,->] (-3,-0.06) arc(181:357:3);
  % Small circle
  \draw[red,thick] (0.4,-0.15) arc(-20:20:0.43);
  \draw[red,thick,->] (0.4,-0.15) arc(-20:160:0.43);
  \draw[red,thick] (0.4,-0.15) arc(-20:340:0.43);
  % Upper lip
  \draw[blue,thick,->] (0.4,0.15) -- (1.7,0.15);
  \draw[blue,thick] (1.7,0.15) -- (3,0.15);
  % Lower lip
  \draw[blue,thick,->] (3,-0.15) -- (1.7,-0.15);
  \draw[blue,thick] (1.7,-0.15) -- (0.4,-0.15);
  % Labels
  \node[above right] at (2.5,1.8) {$C_R$};
  \node[above] at (0.6,0.4) {$c_\varepsilon$};
  \node[above] at (1.7,0.15) {upper lip};
  \node[below] at (1.7,-0.15) {lower lip};
  % Branch cut
  \draw[thick,red!60!black,decorate,decoration={zigzag,segment length=4,amplitude=1}]
    (0,0) -- (3.4,0);
  \node[red!60!black,below right] at (2,-0.4) {branch cut};
  % Poles
  \node[red] at (-1.5,1.2) {$\times$};
  \node[red] at (-1.5,-1.2) {$\times$};
  \node[red,left] at (-1.6,1.2) {$z_k$};
\end{tikzpicture}
\end{center}

\begin{proposition}[Keyhole contour method]\label{prop:keyhole}
Let $0<\alpha<1$ and let $R(z)=P(z)/Q(z)$ be a rational function with
$\deg Q\geq\deg P+1$, no poles on $[0,\infty)$.
Define $z^{\alpha-1}=e^{(\alpha-1)\Log z}$ using the principal branch
with the branch cut on $[0,\infty)$ (i.e.\ $0<\Arg z<2\pi$).  Then
\[
  \int_0^{\infty}x^{\alpha-1}R(x)\,\dd x
  =\frac{2\pi i}{1-e^{2\pi i\alpha}}
  \sum_{k}\Res_{z=z_k}\bigl[z^{\alpha-1}R(z)\bigr],
\]
where the sum is over all poles $z_k$ of $R$.
\end{proposition}

\begin{proof}
Integrate $f(z)=z^{\alpha-1}R(z)$ over the keyhole contour.
On the upper lip of the cut, $\Arg z=0^+$, so $z^{\alpha-1}
=x^{\alpha-1}$.
On the lower lip, $\Arg z=2\pi^-$, so $z^{\alpha-1}
=x^{\alpha-1}e^{2\pi i(\alpha-1)}$.

By the residue theorem:
\[
  \oint_{\text{keyhole}}f(z)\,\dd z
  =2\pi i\sum_k\Res_{z=z_k}f(z).
\]
The contributions from the large circle $C_R$ and small circle
$c_\varepsilon$ vanish as $R\to\infty$ and $\varepsilon\to 0$
(using $\deg Q\geq\deg P+1$ and $0<\alpha<1$).
The remaining contributions are:
\begin{align*}
  &\text{Upper lip: }\int_0^{\infty}x^{\alpha-1}R(x)\,\dd x,\\
  &\text{Lower lip: }-\int_0^{\infty}x^{\alpha-1}
    e^{2\pi i(\alpha-1)}R(x)\,\dd x
  =-e^{2\pi i(\alpha-1)}\int_0^{\infty}x^{\alpha-1}R(x)\,\dd x.
\end{align*}
(The minus sign comes from the reversed direction of traversal.)
Therefore
\[
  (1-e^{2\pi i(\alpha-1)})\int_0^{\infty}x^{\alpha-1}R(x)\,\dd x
  =2\pi i\sum_k\Res_{z=z_k}f(z).
\]
Since $e^{2\pi i(\alpha-1)}=e^{2\pi i\alpha}\cdot e^{-2\pi i}
=e^{2\pi i\alpha}$, the result follows.
\end{proof}

\begin{example}[Integral via keyhole contour]\label{ex:keyhole-1}
Evaluate $\displaystyle I=\int_0^{\infty}\frac{x^{-1/2}}{1+x}\,\dd x$.

\medskip\noindent\textbf{Solution.}
Here $\alpha=1/2$, $R(x)=1/(1+x)$.  The only pole of $R$ is at $z=-1$
(not on $[0,\infty)$).  With $z^{-1/2}=e^{-\frac{1}{2}\Log z}$ where
$0<\Arg z<2\pi$, at $z=-1$ we have $\Arg(-1)=\pi$, so
$(-1)^{-1/2}=e^{-i\pi/2}=-i$.
\[
  \Res_{z=-1}\frac{z^{-1/2}}{1+z}=(-1)^{-1/2}\cdot 1=-i.
\]
Therefore
\[
  I=\frac{2\pi i}{1-e^{i\pi}}\cdot(-i)
  =\frac{2\pi i}{1-(-1)}\cdot(-i)
  =\frac{2\pi i}{2}\cdot(-i)
  =\pi i\cdot(-i)=\pi.
\]
\end{example}

%-------------------------------------------------------------------
\subsection{Rectangular contour technique}
\label{ssec:rectangular}
%-------------------------------------------------------------------

\begin{example}[Gaussian-type integral with rectangular contour]\label{ex:rect-contour}
Evaluate $\displaystyle I=\int_{-\infty}^{+\infty}
\frac{e^{ax}}{1+e^x}\,\dd x$ for $0<a<1$.

\medskip\noindent\textbf{Solution.}
Consider the rectangular contour $\Gamma$ with vertices at
$-R$, $R$, $R+2\pi i$, $-R+2\pi i$, traversed counter-clockwise.

\begin{center}
\begin{tikzpicture}[scale=0.65,>=Stealth]
  \draw[->] (-4.5,0) -- (4.5,0) node[right]{$\re$};
  \draw[->] (0,-0.8) -- (0,4.5) node[above]{$\im$};
  % Rectangle
  \draw[blue,thick,->] (-3.5,0) -- (0,0);
  \draw[blue,thick,->] (0,0) -- (3.5,0);
  \draw[blue,thick,->] (3.5,0) -- (3.5,1.8);
  \draw[blue,thick,->] (3.5,3.5) -- (0,3.5);
  \draw[blue,thick,->] (0,3.5) -- (-3.5,3.5);
  \draw[blue,thick,->] (-3.5,3.5) -- (-3.5,1.8);
  \draw[blue,thick] (-3.5,1.8) -- (-3.5,0);
  \draw[blue,thick] (3.5,1.8) -- (3.5,3.5);
  % Labels
  \node[below] at (3.5,0) {$R$};
  \node[below] at (-3.5,0) {$-R$};
  \node[left] at (-3.7,3.5) {$-R+2\pi i$};
  \node[right] at (3.7,3.5) {$R+2\pi i$};
  % Pole
  \node[red] at (0,1.75) {$\times$};
  \node[red,right] at (0.2,1.75) {$\pi i$};
  \node at (0,3.9) {$\im z=2\pi$};
\end{tikzpicture}
\end{center}

The function $f(z)=e^{az}/(1+e^z)$ has poles where $e^z=-1$, i.e.\
$z=(2k+1)\pi i$ for $k\in\Z$.  Inside $\Gamma$, the only pole is
$z=\pi i$.
\[
  \Res_{z=\pi i}\frac{e^{az}}{1+e^z}
  =\frac{e^{a\pi i}}{e^{\pi i}}=-e^{a\pi i}.
\]

On the bottom side: $\int_{-R}^{R}\frac{e^{ax}}{1+e^x}\dd x$.

On the top side ($z=x+2\pi i$, $x$ from $R$ to $-R$):
$\int_{R}^{-R}\frac{e^{a(x+2\pi i)}}{1+e^{x+2\pi i}}\dd x
=-e^{2\pi ia}\int_{-R}^{R}\frac{e^{ax}}{1+e^x}\dd x$.

The integrals on the vertical sides tend to $0$ as $R\to\infty$
(using $0<a<1$).

Therefore:
\[
  (1-e^{2\pi ia})\int_{-\infty}^{\infty}\frac{e^{ax}}{1+e^x}\dd x
  =2\pi i\cdot(-e^{a\pi i})
\]
\[
  I=\frac{-2\pi i\,e^{a\pi i}}{1-e^{2\pi ia}}
  =\frac{-2\pi i\,e^{a\pi i}}{e^{a\pi i}(e^{-a\pi i}-e^{a\pi i})}
  =\frac{-2\pi i}{-2i\sin(a\pi)}
  =\frac{\pi}{\sin(a\pi)}.
\]
\end{example}

%-------------------------------------------------------------------
\section{Further Worked Examples}\label{sec:more-examples}
%-------------------------------------------------------------------

\begin{example}[Integral involving $x^4$]\label{ex:extra-1}
Evaluate $\displaystyle I=\int_{-\infty}^{+\infty}
\frac{\dd x}{x^4+1}$.

\medskip\noindent\textbf{Solution.}
The poles of $f(z)=1/(z^4+1)$ are the fourth roots of $-1$:
$z_k=e^{i(\pi+2k\pi)/4}$ for $k=0,1,2,3$.  Those in the upper
half-plane are
\[
  z_0=e^{i\pi/4}=\frac{1+i}{\sqrt{2}},\qquad
  z_1=e^{3i\pi/4}=\frac{-1+i}{\sqrt{2}}.
\]
Since each is a simple pole of $1/(z^4+1)$ and $\frac{\dd}{\dd z}(z^4+1)=4z^3$:
\[
  \Res_{z=z_k}\frac{1}{z^4+1}=\frac{1}{4z_k^3}=\frac{z_k}{4z_k^4}
  =\frac{z_k}{4(-1)}=-\frac{z_k}{4}.
\]
Therefore
\begin{align*}
  I &= 2\pi i\left(-\frac{z_0}{4}-\frac{z_1}{4}\right)
  = 2\pi i\cdot\left(-\frac{1}{4}\right)(z_0+z_1) \\
  &= 2\pi i\cdot\left(-\frac{1}{4}\right)
  \left(\frac{1+i}{\sqrt{2}}+\frac{-1+i}{\sqrt{2}}\right)
  = 2\pi i\cdot\left(-\frac{1}{4}\right)\cdot\frac{2i}{\sqrt{2}} \\
  &= 2\pi i\cdot\frac{-i}{2\sqrt{2}}
  = 2\pi\cdot\frac{1}{2\sqrt{2}}
  = \frac{\pi}{\sqrt{2}}.
\end{align*}
\end{example}

\begin{example}[Integral with exponential]\label{ex:extra-2}
Evaluate $\displaystyle I=\int_{-\infty}^{+\infty}
\frac{x\sin x}{x^2+a^2}\,\dd x$ for $a>0$.

\medskip\noindent\textbf{Solution.}
We write $\sin x=\im(e^{ix})$ and compute
\[
  J=\int_{-\infty}^{+\infty}\frac{xe^{ix}}{x^2+a^2}\,\dd x.
\]
Then $I=\im(J)$.  The function $g(z)=ze^{iz}/(z^2+a^2)$ has poles at
$z=\pm ia$.  Only $z=ia$ is in the upper half-plane.

By Jordan's lemma, the semicircular arc contribution vanishes.
\[
  \Res_{z=ia}\frac{ze^{iz}}{z^2+a^2}
  =\frac{ia\,e^{i(ia)}}{2ia}=\frac{e^{-a}}{2}.
\]
Therefore $J=2\pi i\cdot\dfrac{e^{-a}}{2}=i\pi e^{-a}$, and
$I=\im(J)=\pi e^{-a}$.
\end{example}

\begin{example}[A more challenging integral]\label{ex:extra-3}
Evaluate $\displaystyle I=\int_0^{\infty}\frac{\ln x}{(1+x^2)^2}\,\dd x$.

\medskip\noindent\textbf{Solution.}
Consider $f(z)=(\Log z)/(1+z^2)^2$ with $\Log z$ taken with branch cut
on $(-\infty,0]$ (principal branch, $-\pi<\Arg z\leq\pi$).

We integrate over the contour consisting of $[\varepsilon,R]$ on the
real axis, upper semicircle $C_R$, $[-R,-\varepsilon]$, and indented
semicircle $c_\varepsilon$ above the origin.

The only pole in the upper half-plane is $z=i$ (order~$2$).  With
$\Log(i)=i\pi/2$:
\begin{align*}
  \Res_{z=i}\frac{\Log z}{(1+z^2)^2}
  &=\lim_{z\to i}\frac{\dd}{\dd z}\frac{\Log z}{(z+i)^2} \\
  &=\lim_{z\to i}\frac{\frac{1}{z}(z+i)^2-2(z+i)\Log z}{(z+i)^4} \\
  &=\frac{\frac{1}{i}(2i)^2-2(2i)\cdot\frac{i\pi}{2}}{(2i)^4} \\
  &=\frac{\frac{-4}{i}+2\pi}{16}
  =\frac{4i+2\pi}{16}
  =\frac{\pi+2i}{8}.
\end{align*}

On $[\varepsilon,R]$: $\Log x=\ln x$ (real).
On $[-R,-\varepsilon]$ (write $x=-t$, $t>0$): $\Log(-t)=\ln t+i\pi$.

As $R\to\infty$, $\varepsilon\to 0$, the arc contributions vanish, and
\begin{align*}
  \int_0^{\infty}\frac{\ln x}{(1+x^2)^2}\dd x
  +\int_0^{\infty}\frac{\ln t+i\pi}{(1+t^2)^2}\dd t
  &=2\pi i\cdot\frac{\pi+2i}{8}.
\end{align*}
So $2\int_0^{\infty}\frac{\ln x}{(1+x^2)^2}\dd x
+i\pi\int_0^{\infty}\frac{\dd t}{(1+t^2)^2}
=\frac{2\pi i(\pi+2i)}{8}=\frac{\pi(\pi i-2)}{4}$.

We already know (Example~\ref{ex:rational-2})
$\int_0^{\infty}\frac{\dd t}{(1+t^2)^2}=\pi/4$.
Equating real parts:
\[
  2I = -\frac{\pi}{2} \implies
  I = -\frac{\pi}{4}.
\]
\end{example}

%-------------------------------------------------------------------
\section{Argument Principle and Rouch\'e's Theorem}\label{sec:argument-principle}
%-------------------------------------------------------------------

\begin{theorem}[Argument Principle]\label{thm:argument-principle}
Let $f$ be meromorphic inside and on a simple closed contour $\gamma$,
with zeros $a_1,\dots,a_p$ (counted with multiplicity) and poles
$b_1,\dots,b_q$ (counted with multiplicity) inside $\gamma$, and no
zeros or poles on $\gamma$.  Then
\[
  \frac{1}{2\pi i}\oint_{\gamma}\frac{f'(z)}{f(z)}\,\dd z
  = p - q.
\]
\end{theorem}

\begin{proof}
Near a zero $a_j$ of order $m_j$, write $f(z)=(z-a_j)^{m_j}g(z)$ with
$g(a_j)\neq 0$.  Then
\[
  \frac{f'(z)}{f(z)}=\frac{m_j}{z-a_j}+\frac{g'(z)}{g(z)},
\]
so $\Res_{z=a_j}\frac{f'}{f}=m_j$.  Similarly, near a pole $b_k$ of
order $n_k$, $f(z)=(z-b_k)^{-n_k}h(z)$ with $h(b_k)\neq 0$, giving
$\Res_{z=b_k}\frac{f'}{f}=-n_k$.  By the Residue Theorem,
\[
  \frac{1}{2\pi i}\oint_{\gamma}\frac{f'(z)}{f(z)}\,\dd z
  =\sum_j m_j - \sum_k n_k = p - q. \qedhere
\]
\end{proof}

\begin{theorem}[Rouch\'e's Theorem]\label{thm:rouche}
Let $f$ and $g$ be holomorphic inside and on a simple closed contour
$\gamma$.  If $\abs{g(z)}<\abs{f(z)}$ for all $z\in\gamma$, then $f$
and $f+g$ have the same number of zeros (counted with multiplicity)
inside $\gamma$.
\end{theorem}

\begin{proof}
Set $h_t(z)=f(z)+tg(z)$ for $t\in[0,1]$.  On $\gamma$,
$\abs{h_t(z)}\geq\abs{f(z)}-t\abs{g(z)}>0$, so $h_t$ has no zeros on
$\gamma$.  The function
\[
  N(t)=\frac{1}{2\pi i}\oint_{\gamma}\frac{h_t'(z)}{h_t(z)}\,\dd z
\]
is continuous in $t$ and integer-valued by the Argument Principle.
Hence $N(t)$ is constant, so $N(0)=N(1)$, i.e.\ the number of zeros
of $f$ equals the number of zeros of $f+g$ inside $\gamma$.
\end{proof}

\begin{example}[Application of Rouch\'e's theorem]\label{ex:rouche}
Show that the polynomial $p(z)=z^5+3z+1$ has exactly one zero in the
disk $\abs{z}<1$.

\medskip\noindent\textbf{Solution.}
On $\abs{z}=1$, let $f(z)=3z$ and $g(z)=z^5+1$.  Then $\abs{f(z)}=3$
and $\abs{g(z)}\leq\abs{z}^5+1=2<3=\abs{f(z)}$.  By Rouch\'e's theorem,
$f+g=p$ has the same number of zeros as $f(z)=3z$ inside $\abs{z}<1$,
namely one zero.
\end{example}

%-------------------------------------------------------------------
\section{Exercises}\label{sec:ex-ch7}
%-------------------------------------------------------------------

\begin{exercise}[Residue computations]\label{exer:res-compute}
Compute the residues of the following functions at all their isolated
singularities:
\begin{enumerate}[label=(\alph*)]
  \item $\dfrac{z^2}{(z-1)(z+2)^3}$
  \item $\dfrac{e^{1/z}}{z^2}$
  \item $\dfrac{\cos z}{z^2(z-\pi)}$
  \item $\dfrac{1}{z^2\sin z}$
  \item $\dfrac{z}{(e^z-1)^2}$
\end{enumerate}
\end{exercise}

\begin{exercise}[Real integrals via residues I]\label{exer:real-int-1}
Evaluate the following integrals using the residue theorem:
\begin{enumerate}[label=(\alph*)]
  \item $\displaystyle\int_{-\infty}^{+\infty}\frac{\dd x}{x^4+4}$
  \item $\displaystyle\int_{-\infty}^{+\infty}
        \frac{x^2\,\dd x}{(x^2+1)(x^2+4)}$
  \item $\displaystyle\int_{-\infty}^{+\infty}
        \frac{\dd x}{(x^2+a^2)^3}$ \quad ($a>0$)
  \item $\displaystyle\int_0^{2\pi}
        \frac{\dd\theta}{3+2\cos\theta+\sin\theta}$
\end{enumerate}
\end{exercise}

\begin{exercise}[Fourier-type integrals]\label{exer:fourier}
Evaluate:
\begin{enumerate}[label=(\alph*)]
  \item $\displaystyle\int_{-\infty}^{+\infty}
        \frac{\cos(3x)}{x^2+4}\,\dd x$
  \item $\displaystyle\int_0^{\infty}
        \frac{x\sin(2x)}{x^2+9}\,\dd x$
  \item $\displaystyle\int_{-\infty}^{+\infty}
        \frac{x\cos x}{(x^2+1)^2}\,\dd x$
\end{enumerate}
\end{exercise}

\begin{exercise}[Keyhole contour]\label{exer:keyhole}
Using a keyhole contour, evaluate:
\begin{enumerate}[label=(\alph*)]
  \item $\displaystyle\int_0^{\infty}\frac{x^{1/3}}{1+x^2}\,\dd x$
  \item $\displaystyle\int_0^{\infty}\frac{\ln x}{x^2+1}\,\dd x$
  \item $\displaystyle\int_0^{\infty}\frac{x^{a-1}}{1+x}\,\dd x$
        \quad ($0<a<1$) \quad
        [\emph{Euler's reflection formula:
        $\Gamma(a)\Gamma(1-a)=\pi/\sin(\pi a)$}]
\end{enumerate}
\end{exercise}

\begin{exercise}[Argument principle and Rouch\'e]\label{exer:rouche}
\begin{enumerate}[label=(\alph*)]
  \item Show that $z^4-5z+1$ has exactly one zero in $\abs{z}<1$.
  \item Show that $z^7-5z^3+12$ has no zeros in $\abs{z}<1$ and
        exactly three zeros in $1<\abs{z}<2$.
  \item Let $p(z)=z^n+a_{n-1}z^{n-1}+\cdots+a_0$ with
        $\abs{a_0}>\abs{a_1}+\cdots+\abs{a_{n-1}}+1$.  Show that $p$
        has no zeros in $\abs{z}\leq 1$.
\end{enumerate}
\end{exercise}

\begin{exercise}[Summation by residues]\label{exer:summation}
Use the residue theorem with the function $\pi\cot(\pi z)/g(z)$ to show:
\begin{enumerate}[label=(\alph*)]
  \item $\displaystyle\sum_{n=1}^{\infty}\frac{1}{n^2}=\frac{\pi^2}{6}$
  \item $\displaystyle\sum_{n=1}^{\infty}\frac{1}{n^2+a^2}
        =\frac{\pi}{2a}\coth(\pi a)-\frac{1}{2a^2}$ \quad ($a>0$)
  \item $\displaystyle\sum_{n=0}^{\infty}\frac{(-1)^n}{(2n+1)^2}
        =G$ \quad (Catalan's constant; express as a contour integral)
\end{enumerate}
[\emph{Hint:} For (a), integrate $\pi\cot(\pi z)/z^2$ over a square
contour with vertices at $\pm(N+\tfrac{1}{2})\pm(N+\tfrac{1}{2})i$.]
\end{exercise}

\begin{exercise}[Integral challenge]\label{exer:challenge}
Evaluate, using an appropriate contour:
\begin{enumerate}[label=(\alph*)]
  \item $\displaystyle\int_0^{\infty}
        \frac{\dd x}{(1+x^2)(1+x^3)}$
  \item $\displaystyle\int_0^{\infty}
        \frac{(\ln x)^2}{1+x^2}\,\dd x$
  \item $\displaystyle\int_0^{\infty}
        \frac{\cos(x^2)\,\dd x}{\phantom{x}}$ \quad
        (Fresnel integral; use the contour consisting of $[0,R]$,
        an arc of angle $\pi/4$, and the line $te^{i\pi/4}$,
        $t\in[R,0]$).
\end{enumerate}
\end{exercise}


%%%%%%%%%%%%%%%%%%%%%%%%%%%%%%%%%%%%%%%%%%%%%%%%%%%%%%%%%%%%%%%%%%%%
%  CHAPTER 8 — Conformal Mappings
%%%%%%%%%%%%%%%%%%%%%%%%%%%%%%%%%%%%%%%%%%%%%%%%%%%%%%%%%%%%%%%%%%%%
\chapter{Conformal Mappings}\label{ch:conformal}

\section{Conformal Maps: Definition and First Properties}
\label{sec:conformal-def}

Throughout this chapter, we study holomorphic maps that preserve angles
and local geometry.  Conformal mappings are indispensable in potential
theory, fluid dynamics, and electrostatics, where they allow one to
transfer solutions of Laplace's equation between domains.

\begin{definition}[Conformal map]\label{def:conformal}
Let $\Omega\subset\C$ be an open set.  A holomorphic function
$f\colon\Omega\to\C$ is said to be \emph{conformal} at a point
$z_0\in\Omega$ if $f'(z_0)\neq 0$.  We say $f$ is \emph{conformal on
$\Omega$} if $f'(z)\neq 0$ for all $z\in\Omega$.

A \emph{conformal equivalence} (or \emph{biholomorphism}) between open
sets $\Omega_1$ and $\Omega_2$ is a bijective holomorphic map
$f\colon\Omega_1\to\Omega_2$ whose inverse $f^{-1}\colon\Omega_2\to\Omega_1$
is also holomorphic.
\end{definition}

\begin{theorem}[Angle preservation]\label{thm:angle-preservation}
If $f$ is conformal at $z_0$, then $f$ preserves angles between curves
at $z_0$: if two smooth curves $\gamma_1,\gamma_2$ meet at $z_0$ with
angle $\alpha$, then their images $f\circ\gamma_1$ and $f\circ\gamma_2$
meet at $f(z_0)$ with the same angle $\alpha$, both in magnitude and
orientation.
\end{theorem}

\begin{proof}
Let $\gamma_j(t)$ be smooth curves with $\gamma_j(0)=z_0$ and
$\gamma_j'(0)\neq 0$ for $j=1,2$.  The tangent direction of
$f\circ\gamma_j$ at $t=0$ is
\[
  (f\circ\gamma_j)'(0)=f'(z_0)\,\gamma_j'(0).
\]
Since $f'(z_0)\neq 0$, write $f'(z_0)=re^{i\phi}$ with $r>0$.
Then
\[
  \arg(f\circ\gamma_j)'(0)=\phi+\arg\gamma_j'(0).
\]
The angle between the image curves is
\[
  \arg(f\circ\gamma_2)'(0)-\arg(f\circ\gamma_1)'(0)
  =\arg\gamma_2'(0)-\arg\gamma_1'(0),
\]
which equals the angle between the original curves.
\end{proof}

\begin{remark}[Conformal $\Leftrightarrow$ locally bijective]\label{rem:conf-local-bij}
By the inverse function theorem for holomorphic functions, $f$ is
conformal at $z_0$ if and only if $f$ is locally biholomorphic near
$z_0$.  A holomorphic bijection between open sets is automatically
conformal (its inverse is holomorphic).
\end{remark}

\begin{proposition}[Anti-conformal maps]\label{prop:anti-conformal}
The map $z\mapsto\conj{z}$ preserves angle magnitudes but reverses
orientation.  Composing a conformal map with conjugation yields an
\emph{anti-conformal} map.
\end{proposition}

%-------------------------------------------------------------------
\section{Basic Examples of Conformal Maps}\label{sec:basic-conformal}
%-------------------------------------------------------------------

\begin{example}[Translation]\label{ex:translation}
$f(z)=z+b$ for $b\in\C$.  This is conformal everywhere ($f'(z)=1$) and
maps every region to a translated copy.
\end{example}

\begin{example}[Rotation and dilation]\label{ex:rotation-dilation}
$f(z)=az$ for $a\in\C\setminus\{0\}$.  Writing $a=re^{i\theta}$, this
is a rotation by $\theta$ and dilation by factor $r$.  It is conformal
everywhere ($f'(z)=a\neq 0$).
\end{example}

\begin{example}[The power map]\label{ex:power-map}
$f(z)=z^n$ for $n\geq 2$.  This is conformal except at $z=0$ (where
$f'(0)=0$) and maps a sector of opening angle $\alpha$ to a sector
of opening angle $n\alpha$.  In particular, $z^2$ maps the first quadrant
$\{z:\re z>0,\,\im z>0\}$ conformally onto the upper half-plane.

\begin{center}
\begin{tikzpicture}[scale=0.95,>=Stealth]
  % Left: quarter plane
  \begin{scope}[shift={(-3.5,0)}]
    \fill[blue!10] (0,0) -- (2.3,0) -- (2.3,2.3) -- (0,2.3) -- cycle;
    \draw[->] (-0.5,0) -- (2.5,0) node[right]{$\re$};
    \draw[->] (0,-0.5) -- (0,2.5) node[above]{$\im$};
    \draw[blue!50,thick] (0.4,0) -- (0.4,2.1);
    \draw[blue!50,thick] (0.8,0) -- (0.8,2.1);
    \draw[blue!50,thick] (1.2,0) -- (1.2,2.1);
    \draw[blue!50,thick] (1.6,0) -- (1.6,2.1);
    \draw[red!50,thick] (0,0.4) -- (2.1,0.4);
    \draw[red!50,thick] (0,0.8) -- (2.1,0.8);
    \draw[red!50,thick] (0,1.2) -- (2.1,1.2);
    \draw[red!50,thick] (0,1.6) -- (2.1,1.6);
    \node[below] at (1.2,-0.7) {first quadrant};
  \end{scope}
  % Arrow
  \draw[->,thick] (-0.5,1) -- (0.5,1) node[midway,above]{$z^2$};
  % Right: upper half-plane
  \begin{scope}[shift={(3.5,0)}]
    \fill[red!8] (-2.3,0) rectangle (2.3,2.3);
    \draw[->] (-2.5,0) -- (2.5,0) node[right]{$\re$};
    \draw[->] (0,-0.5) -- (0,2.5) node[above]{$\im$};
    % Images of vertical lines: right-opening parabolas
    \draw[blue!50,thick,samples=40,domain=0:2.1]
      plot({0.16-(\x)^2/0.64},{2*0.4*\x/0.8});
    \draw[blue!50,thick,samples=40,domain=0:1.8]
      plot({0.64-(\x)^2/1.28},{2*0.8*\x/1.6});
    \draw[blue!50,thick,samples=40,domain=0:1.5]
      plot({1.44-(\x)^2/1.92},{2*1.2*\x/2.4});
    % Images of horizontal lines: left-opening parabolas
    \draw[red!50,thick,samples=40,domain=0:2.1]
      plot({(\x)^2/0.64-0.16},{2*0.4*\x/0.8});
    \draw[red!50,thick,samples=40,domain=0:1.8]
      plot({(\x)^2/1.28-0.64},{2*0.8*\x/1.6});
    \node[below] at (0,-0.7) {upper half-plane};
  \end{scope}
\end{tikzpicture}
\end{center}
\end{example}

\begin{example}[The exponential map]\label{ex:exp-map}
$f(z)=e^z$ is conformal everywhere ($f'(z)=e^z\neq 0$).  It maps the
horizontal strip $S=\{z:0<\im z<\pi\}$ conformally onto the upper
half-plane $\{w:\im w>0\}$.  Vertical lines $\re z=c$ map to
semicircles $\abs{w}=e^c$ in the upper half-plane, and horizontal lines
$\im z=d$ map to rays $\arg w=d$.

\begin{center}
\begin{tikzpicture}[scale=0.85,>=Stealth]
  % Left: strip
  \begin{scope}[shift={(-3.5,0)}]
    \fill[blue!10] (-2,0) rectangle (2,2);
    \draw[->] (-2.3,0) -- (2.3,0) node[right]{$\re$};
    \draw[->] (0,-0.5) -- (0,2.7) node[above]{$\im$};
    \draw[blue,thick,dashed] (-2,0) -- (2,0);
    \draw[blue,thick,dashed] (-2,2) -- (2,2);
    \node[left] at (-2,2) {$\im z=\pi$};
    \node[left] at (-2,0) {$\im z=0$};
    % Grid
    \draw[green!50!black,thin] (-1.2,0) -- (-1.2,2);
    \draw[green!50!black,thin] (0,0) -- (0,2);
    \draw[green!50!black,thin] (1.2,0) -- (1.2,2);
    \draw[orange,thin] (-2,0.7) -- (2,0.7);
    \draw[orange,thin] (-2,1.4) -- (2,1.4);
    \node[below] at (0,-0.7) {strip $0<\im z<\pi$};
  \end{scope}
  \draw[->,thick] (-0.7,1) -- (0.7,1) node[midway,above]{$e^z$};
  % Right: upper half-plane
  \begin{scope}[shift={(3.5,0)}]
    \fill[red!8] (-2.5,0) rectangle (2.5,2.5);
    \draw[->] (-2.7,0) -- (2.7,0) node[right]{$\re$};
    \draw[->] (0,-0.5) -- (0,2.7) node[above]{$\im$};
    % Semicircles (images of vertical lines)
    \begin{scope}
      \clip (-2.7,0) rectangle (2.7,2.7);
      \draw[green!50!black,thin] (0,0) circle(0.3);
      \draw[green!50!black,thin] (0,0) circle(1.0);
      \draw[green!50!black,thin] (0,0) circle(2.3);
    \end{scope}
    % Rays (images of horizontal lines)
    \draw[orange,thin] (0,0) -- ({2.5*cos(40)},{2.5*sin(40)});
    \draw[orange,thin] (0,0) -- ({2.5*cos(80)},{2.5*sin(80)});
    \node[below] at (0,-0.7) {upper half-plane};
  \end{scope}
\end{tikzpicture}
\end{center}
\end{example}

\begin{example}[Inversion]\label{ex:inversion}
$f(z)=1/z$ is conformal on $\C\setminus\{0\}$.  It maps circles and
lines to circles and lines (treating lines as circles through $\infty$).
The unit disk $\D$ maps to its exterior
$\{\abs{z}>1\}\cup\{\infty\}$ and the unit circle $\T$ maps to itself.
\end{example}

%-------------------------------------------------------------------
\section{M\"obius Transformations}\label{sec:mobius}
%-------------------------------------------------------------------

\begin{definition}[M\"obius transformation]\label{def:mobius}
A \emph{M\"obius transformation} (or \emph{fractional linear
transformation}) is a map of the form
\[
  T(z)=\frac{az+b}{cz+d},\qquad a,b,c,d\in\C,\quad ad-bc\neq 0.
\]
The condition $ad-bc\neq 0$ ensures that $T$ is not constant.
We extend $T$ to the Riemann sphere $\hat{\C}=\C\cup\{\infty\}$ by
setting $T(-d/c)=\infty$ and $T(\infty)=a/c$ (when $c\neq 0$; if $c=0$,
$T(\infty)=\infty$).
\end{definition}

\begin{proposition}[Properties of M\"obius transformations]\label{prop:mobius-props}
\begin{enumerate}[label=(\roman*)]
  \item Every M\"obius transformation is a conformal bijection
        $\hat{\C}\to\hat{\C}$.
  \item The set of all M\"obius transformations forms a group under
        composition, isomorphic to $\mathrm{PGL}(2,\C)$.
  \item The inverse of $T(z)=\frac{az+b}{cz+d}$ is
        $T^{-1}(z)=\frac{dz-b}{-cz+a}$.
  \item Composition corresponds to matrix multiplication:
        \[
          \begin{pmatrix}a_2&b_2\\c_2&d_2\end{pmatrix}
          \begin{pmatrix}a_1&b_1\\c_1&d_1\end{pmatrix}
          =\begin{pmatrix}a_2a_1+b_2c_1 & a_2b_1+b_2d_1 \\
                          c_2a_1+d_2c_1 & c_2b_1+d_2d_1\end{pmatrix}.
        \]
  \item Every M\"obius transformation is a composition of translations,
        rotations, dilations, and inversions.
\end{enumerate}
\end{proposition}

\begin{proof}
We prove~(v).  If $c=0$, then
$T(z)=(a/d)z+b/d$, a dilation-rotation followed by a translation.
If $c\neq 0$, perform polynomial long division:
\[
  T(z)=\frac{az+b}{cz+d}
  =\frac{a}{c}-\frac{ad-bc}{c^2}\cdot\frac{1}{z+d/c}.
\]
This is the composition of four elementary maps:
\begin{enumerate}[label=(\arabic*)]
  \item $z\mapsto z+d/c$ \quad (translation),
  \item $z\mapsto 1/z$ \quad (inversion),
  \item $z\mapsto -(ad-bc)/c^2\cdot z$ \quad (rotation-dilation),
  \item $z\mapsto z+a/c$ \quad (translation).
\end{enumerate}
Each is conformal ($ad-bc\neq 0$ ensures step~(3) is non-degenerate).
\end{proof}

\begin{theorem}[Circles to circles]\label{thm:mobius-circles}
M\"obius transformations map generalised circles (circles and lines in
$\C$, viewed as circles on $\hat{\C}$) to generalised circles.
\end{theorem}

\begin{proof}
By Proposition~\ref{prop:mobius-props}(v), it suffices to check each
elementary type.  Translations, rotations, and dilations clearly preserve
the class of generalised circles.

For the inversion $w=1/z$: a generalised circle has equation
$A\abs{z}^2+\re(\conj{B}z)+D=0$ with $A,D\in\R$ and $B\in\C$
(a line if $A=0$, a circle if $A\neq 0$).  Under $z=1/w$,
$\abs{z}^2=1/\abs{w}^2$ and $\re(\conj{B}z)=\re(\conj{B}/w)
=\re(B\conj{w})/\abs{w}^2$.
Multiplying through by $\abs{w}^2$:
\[
  A+\re(\conj{\conj{B}}\,\conj{w})+D\abs{w}^2=0,
\]
which is again a generalised circle equation (with $A$ and $D$ swapped,
and $B$ replaced by $\conj{B}$).
\end{proof}

%-------------------------------------------------------------------
\subsection{Fixed points and classification}\label{ssec:fixed-pts}
%-------------------------------------------------------------------

\begin{definition}[Fixed points]\label{def:fixed-pts}
A \emph{fixed point} of a M\"obius transformation $T$ is a point
$z_0\in\hat{\C}$ with $T(z_0)=z_0$.
\end{definition}

\begin{proposition}[Number of fixed points]\label{prop:fixed-pts-count}
A M\"obius transformation $T\neq\mathrm{id}$ has exactly $1$ or $2$
fixed points in $\hat{\C}$.
\end{proposition}

\begin{proof}
If $c\neq 0$, the fixed points in $\C$ satisfy $cz^2+(d-a)z-b=0$.
Since $T\neq\mathrm{id}$, this quadratic is non-trivial, yielding
$0$, $1$, or $2$ finite roots.  Including $z=\infty$ (which is fixed
iff $c=0$), the total is at most~$2$.  It is at least~$1$ because a
quadratic over $\C$ always has a root.
\end{proof}

\begin{definition}[Classification of M\"obius transformations]\label{def:mobius-class}
Let $T\neq\mathrm{id}$ with matrix $M=\bigl(\begin{smallmatrix}a&b\\c&d
\end{smallmatrix}\bigr)$ normalised so that $ad-bc=1$.  The \emph{trace}
$\tau=a+d$ determines the type (up to conjugation in the M\"obius group):
\begin{itemize}
  \item \textbf{Parabolic:} $\tau^2=4$ (one fixed point).
        Conjugate to $z\mapsto z+1$.
  \item \textbf{Elliptic:} $\tau^2\in[0,4)$ with $\tau\in\R$
        (two fixed points).  Conjugate to $z\mapsto e^{i\theta}z$ (rotation).
  \item \textbf{Hyperbolic:} $\tau^2\in(4,+\infty)$ with $\tau\in\R$
        (two fixed points).  Conjugate to $z\mapsto\lambda z$, $\lambda>0$,
        $\lambda\neq 1$ (dilation).
  \item \textbf{Loxodromic:} $\tau\notin\R$ (two fixed points).
        Conjugate to $z\mapsto\lambda z$ with $\abs{\lambda}\neq 1$,
        $\lambda\notin\R$ (spiral).
\end{itemize}
\end{definition}

\begin{example}[Classifying a M\"obius transformation]\label{ex:mobius-classify}
Classify $T(z)=\dfrac{3z-4}{z-1}$.

\medskip\noindent\textbf{Solution.}
Here $a=3$, $b=-4$, $c=1$, $d=-1$, so $ad-bc=-3+4=1$.  The trace is
$\tau=a+d=3+(-1)=2$, and $\tau^2=4$.  Therefore $T$ is \textbf{parabolic}.

The unique fixed point: $z^2-4z+4=0$, i.e.\ $(z-2)^2=0$, so $z=2$.
\end{example}

%-------------------------------------------------------------------
\subsection{Cross-ratio}\label{ssec:cross-ratio}
%-------------------------------------------------------------------

\begin{definition}[Cross-ratio]\label{def:cross-ratio}
Given four distinct points $z_1,z_2,z_3,z_4\in\hat{\C}$, the
\emph{cross-ratio} is
\[
  (z_1,z_2;z_3,z_4)
  =\frac{(z_1-z_3)(z_2-z_4)}{(z_1-z_4)(z_2-z_3)}.
\]
If one point is $\infty$, the definition is obtained by taking the
appropriate limit (e.g.\ $(z_1,z_2;z_3,\infty)=(z_1-z_3)/(z_2-z_3)$).
\end{definition}

\begin{theorem}[Cross-ratio is a M\"obius invariant]\label{thm:cross-ratio-invariant}
If $T$ is a M\"obius transformation, then
\[
  (T(z_1),T(z_2);T(z_3),T(z_4))=(z_1,z_2;z_3,z_4).
\]
\end{theorem}

\begin{proof}
Define $S(z)=(z,z_2;z_3,z_4)$, a M\"obius transformation with
$S(z_2)=1$, $S(z_3)=0$, $S(z_4)=\infty$.  Similarly,
$S'(w)=(w,Tz_2;Tz_3,Tz_4)$ sends $Tz_2\mapsto 1$, $Tz_3\mapsto 0$,
$Tz_4\mapsto\infty$.  Then $S'\circ T$ and $S$ are M\"obius
transformations agreeing at three points $z_2,z_3,z_4$, hence they are
identical.  Evaluating at $z_1$ gives the result.
\end{proof}

\begin{proposition}[Three-point determination]\label{prop:three-point}
Given two ordered triples of distinct points $(z_1,z_2,z_3)$ and
$(w_1,w_2,w_3)$ in $\hat{\C}$, there exists a unique M\"obius
transformation $T$ with $T(z_j)=w_j$ for $j=1,2,3$.  It is determined
implicitly by
\[
  (w,w_1;w_2,w_3)=(z,z_1;z_2,z_3).
\]
\end{proposition}

\begin{example}[Finding a M\"obius transformation]\label{ex:mobius-find}
Find the M\"obius transformation $T$ sending $0\mapsto 1$,
$1\mapsto i$, $\infty\mapsto -1$.

\medskip\noindent\textbf{Solution.}
Write $T(z)=(az+b)/(cz+d)$.  From the conditions:
\begin{align*}
  T(\infty)=a/c=-1 &\implies a=-c, \\
  T(0)=b/d=1 &\implies b=d, \\
  T(1)=(a+b)/(c+d)=i &\implies (-c+d)/(c+d)=i.
\end{align*}
From the last equation: $-c+d=ic+id$, so $d(1-i)=c(1+i)$, giving
$c=d(1-i)/(1+i)=d\cdot(-i)=-id$.  Take $d=1$: then $c=-i$,
$a=i$, $b=1$.
\[
  \boxed{T(z)=\frac{iz+1}{-iz+1}=\frac{iz+1}{1-iz}.}
\]
Verification: $T(0)=1/1=1$; \; $T(1)=(i+1)/(1-i)=(1+i)^2/2=2i/2=i$; \;
$T(\infty)=i/(-i)=-1$. \checkmark
\end{example}

%-------------------------------------------------------------------
\section{The Cayley Transform}\label{sec:cayley}
%-------------------------------------------------------------------

\begin{definition}[Cayley transform]\label{def:cayley}
The \emph{Cayley transform} is the M\"obius transformation
\[
  \varphi(z)=\frac{z-i}{z+i}.
\]
Its inverse is $\varphi^{-1}(w)=i\,\dfrac{1+w}{1-w}$.
\end{definition}

\begin{theorem}[Cayley maps upper half-plane to disk]\label{thm:cayley}
The Cayley transform $\varphi$ is a conformal equivalence between the
upper half-plane $\mathbb{H}=\{z\in\C:\im z>0\}$ and the unit disk
$\D=\{w\in\C:\abs{w}<1\}$.  It maps the real line
$\R\cup\{\infty\}$ bijectively onto the unit circle $\T$.
\end{theorem}

\begin{proof}
For $z=x+iy$ with $z\neq -i$:
\[
  \abs{\varphi(z)}^2=\frac{\abs{z-i}^2}{\abs{z+i}^2}
  =\frac{x^2+(y-1)^2}{x^2+(y+1)^2}.
\]
This quantity is less than $1$ iff $(y-1)^2<(y+1)^2$, i.e.\ $y>0$.
Thus $\abs{\varphi(z)}<1\Leftrightarrow\im z>0$.  Similarly,
$\abs{\varphi(z)}=1\Leftrightarrow y=0$.

Since $\varphi$ is a M\"obius transformation, it is a bijection
$\hat{\C}\to\hat{\C}$, restricting to a bijection $\mathbb{H}\to\D$.
It is holomorphic with $\varphi'(z)=2i/(z+i)^2\neq 0$, hence conformal.
\end{proof}

\begin{center}
\begin{tikzpicture}[scale=1.05,>=Stealth]
  % Left: upper half-plane
  \begin{scope}[shift={(-3.5,0)}]
    \fill[blue!8] (-2.3,0) rectangle (2.3,2.3);
    \draw[->] (-2.5,0) -- (2.5,0) node[right]{$\re$};
    \draw[->] (0,-0.8) -- (0,2.5) node[above]{$\im$};
    \draw[thick,blue] (-2.3,0) -- (2.3,0);
    \node[below] at (0,-1) {$\mathbb{H}$};
    \fill (0,1) circle(1.5pt) node[right]{\small$i$};
    \fill (0,0) circle(1.5pt) node[below right]{\small$0$};
    \fill (1,0) circle(1.5pt) node[below]{\small$1$};
    \fill (-1,0) circle(1.5pt) node[below]{\small$-1$};
  \end{scope}
  % Arrow
  \draw[->,thick] (-0.5,1) -- (0.5,1)
    node[midway,above]{\small$\frac{z-i}{z+i}$};
  % Right: unit disk
  \begin{scope}[shift={(3.5,0)}]
    \fill[red!8] (0,0.7) circle(1.5);
    \draw[->] (-2.2,0) -- (2.2,0) node[right]{$\re$};
    \draw[->] (0,-1.1) -- (0,2.5) node[above]{$\im$};
    \draw[thick,red] (0,0.7) circle(1.5);
    \node[below] at (0,-1.3) {$\D$};
    \fill (0,0.7) circle(1.5pt) node[above right]{\small$0$};
    \fill (0,2.2) circle(1.5pt) node[above right]{\small$1$};
    \fill (0,-0.8) circle(1.5pt) node[below right]{\small$-1$};
    \fill (1.5,0.7) circle(1.5pt) node[right]{\small$-i$};
    \fill (-1.5,0.7) circle(1.5pt) node[left]{\small$i$};
  \end{scope}
\end{tikzpicture}
\end{center}

\begin{remark}[Key correspondences under the Cayley transform]\label{rem:cayley-correspondences}
\[
\begin{array}{c|c}
  \mathbb{H} & \D \\\hline
  i & 0 \\
  0 & -1 \\
  1 & -i \\
  -1 & i \\
  \infty & 1
\end{array}
\]
The upper half of the imaginary axis maps to $(-1,0)$ on the real axis
in $\D$, and the real axis maps to $\T$.
\end{remark}

\begin{remark}[Automorphisms of $\D$]\label{rem:aut-disk}
Every conformal automorphism of the unit disk has the form
\[
  T(z)=e^{i\theta}\frac{z-a}{1-\conj{a}z},
  \qquad \theta\in\R,\;\abs{a}<1.
\]
These are \emph{Blaschke factors}.  The group $\mathrm{Aut}(\D)$ is
isomorphic to $\mathrm{PSU}(1,1)$.
\end{remark}

\begin{remark}[Automorphisms of $\mathbb{H}$]\label{rem:aut-H}
Every conformal automorphism of the upper half-plane has the form
\[
  T(z)=\frac{az+b}{cz+d},\qquad a,b,c,d\in\R,\quad ad-bc>0.
\]
The group $\mathrm{Aut}(\mathbb{H})\cong\mathrm{PSL}(2,\R)$.
\end{remark}

%-------------------------------------------------------------------
\section{The Riemann Mapping Theorem}\label{sec:riemann-mapping}
%-------------------------------------------------------------------

\begin{theorem}[Riemann Mapping Theorem]\label{thm:riemann-mapping}
Let $\Omega\subsetneq\C$ be a simply connected open set with
$\Omega\neq\C$.  Then there exists a conformal equivalence
$f\colon\Omega\to\D$.

Moreover, for any $z_0\in\Omega$, there is a unique such $f$
satisfying $f(z_0)=0$ and $f'(z_0)>0$.
\end{theorem}

\begin{remark}[Sketch of proof]\label{rem:riemann-proof}
The proof is one of the great achievements of complex analysis.  We
outline the main steps.

\begin{enumerate}[label=\textbf{Step~\arabic*:}]
  \item \textbf{Family.}
  Let $\mathcal{F}$ be the family of all injective holomorphic functions
  $g\colon\Omega\to\D$ with $g(z_0)=0$.  One shows $\mathcal{F}\neq
  \emptyset$ using the hypothesis $\Omega\neq\C$ and a square-root
  construction.

  \item \textbf{Extremal problem.}
  Consider $\sup\{\abs{g'(z_0)}:g\in\mathcal{F}\}$.  This is finite by
  the Schwarz--Pick lemma.

  \item \textbf{Normal families.}
  By Montel's theorem, $\mathcal{F}$ is a normal family.  A subsequence
  converges uniformly on compact subsets to some $f\in\mathcal{F}$
  achieving the supremum.

  \item \textbf{Surjectivity.}
  If $f(\Omega)\neq\D$, one constructs $\tilde{f}\in\mathcal{F}$ with
  $\abs{\tilde{f}'(z_0)}>\abs{f'(z_0)}$, contradicting maximality.
\end{enumerate}

Uniqueness follows from the Schwarz lemma: if $f_1,f_2$ are two such
maps, then $f_1\circ f_2^{-1}$ is an automorphism of $\D$ fixing $0$
with positive derivative, so it must be the identity.
\end{remark}

\begin{corollary}[Equivalence of simply connected domains]\label{cor:sc-equiv}
Any two simply connected proper open subsets of $\C$ are conformally
equivalent.
\end{corollary}

\begin{proof}
If $\Omega_1,\Omega_2\subsetneq\C$ are simply connected, let
$f_j\colon\Omega_j\to\D$ be Riemann maps.  Then
$f_2^{-1}\circ f_1\colon\Omega_1\to\Omega_2$ is a biholomorphism.
\end{proof}

%-------------------------------------------------------------------
\section{Classical Conformal Maps}\label{sec:classical-maps}
%-------------------------------------------------------------------

We catalogue the most important conformal maps between standard domains.
These are the building blocks for solving boundary value problems.

\subsection{Half-plane, disk, and strips}

\begin{example}[Upper half-plane $\to$ unit disk]\label{ex:H-to-D}
The Cayley transform (Definition~\ref{def:cayley}):
$w=\dfrac{z-i}{z+i}:\;\mathbb{H}\xrightarrow{\;\sim\;}\D$.
\end{example}

\begin{example}[Strip $\to$ upper half-plane]\label{ex:strip-to-H}
The horizontal strip $S=\{z:0<\im z<\pi\}$ is mapped to $\mathbb{H}$
by $w=e^z$.  If $z=x+iy$ with $0<y<\pi$, then $w=e^x e^{iy}$ has
$\im w=e^x\sin y>0$.
\end{example}

\begin{example}[Strip $\to$ unit disk]\label{ex:strip-to-D}
Composing, $w=\dfrac{e^z-i}{e^z+i}$ maps $S=\{0<\im z<\pi\}$ to $\D$.
More generally, $w=\tanh\!\bigl(\frac{\pi z}{4a}\bigr)$ maps
$\{-a<\im z<a\}$ to $\D$.
\end{example}

\begin{example}[Right half-plane $\to$ unit disk]\label{ex:right-hp-to-D}
$w=\dfrac{z-1}{z+1}$ maps $\{\re z>0\}\to\D$.  This follows from the
Cayley transform by a $90^\circ$ rotation: if $\re z>0$, then $\im(iz)>0$,
so $\frac{iz-i}{iz+i}=\frac{z-1}{z+1}\in\D$.
\end{example}

\subsection{Sectors and wedges}

\begin{example}[Sector $\to$ upper half-plane]\label{ex:sector-to-H}
The sector $S_\alpha=\{z:0<\arg z<\alpha\}$, $0<\alpha\leq 2\pi$, is
mapped to $\mathbb{H}$ by $w=z^{\pi/\alpha}$ (principal branch).

\begin{center}
\begin{tikzpicture}[scale=0.9,>=Stealth]
  % Left: sector
  \begin{scope}[shift={(-3.5,0)}]
    \fill[blue!10] (0,0) -- (2.5,0) arc(0:70:2.5) -- cycle;
    \draw[->] (-0.5,0) -- (2.8,0) node[right]{$\re$};
    \draw[->] (0,-0.5) -- (0,2.8) node[above]{$\im$};
    \draw[blue,thick] (0,0) -- (2.5,0);
    \draw[blue,thick] (0,0) -- ({2.5*cos(70)},{2.5*sin(70)});
    \draw[<->,blue] (1,0) arc(0:70:1);
    \node[blue] at (0.6,0.5) {$\alpha$};
    \node[below] at (1.2,-0.7) {sector $S_\alpha$};
  \end{scope}
  \draw[->,thick] (-0.5,1) -- (0.5,1) node[midway,above]{$z^{\pi/\alpha}$};
  % Right: upper half-plane
  \begin{scope}[shift={(3.5,0)}]
    \fill[red!8] (-2.3,0) rectangle (2.3,2.3);
    \draw[->] (-2.5,0) -- (2.5,0) node[right]{$\re$};
    \draw[->] (0,-0.5) -- (0,2.5) node[above]{$\im$};
    \draw[red,thick] (-2.3,0) -- (2.3,0);
    \node[below] at (0,-0.7) {$\mathbb{H}$};
  \end{scope}
\end{tikzpicture}
\end{center}

Special cases:
\begin{itemize}
  \item $\alpha=\pi/2$ (first quadrant): $w=z^2$ maps to $\mathbb{H}$.
  \item $\alpha=\pi$ (upper half-plane to itself): $w=z$ (identity).
  \item $\alpha=2\pi$ (slit plane): $w=z^{1/2}$ maps to $\mathbb{H}$.
\end{itemize}
\end{example}

\subsection{The Joukowski map}

\begin{example}[Joukowski (Zhukovsky) map]\label{ex:joukowski}
The map $J(z)=z+1/z$ is conformal on $\{\abs{z}>1\}$ (since
$J'(z)=1-1/z^2\neq 0$ there).  It maps:
\begin{itemize}
  \item the exterior of the unit disk $\{\abs{z}>1\}$ conformally onto
        $\C\setminus[-2,2]$;
  \item the unit circle to the segment $[-2,2]$, traversed twice;
  \item circles $\abs{z}=r>1$ to ellipses with semi-axes $r+1/r$ and
        $r-1/r$;
  \item rays $\arg z=\theta$ to branches of hyperbolas.
\end{itemize}

\begin{center}
\begin{tikzpicture}[scale=0.75,>=Stealth]
  \begin{scope}[shift={(-3.5,0)}]
    \draw[->] (-2.5,0) -- (2.5,0) node[right]{$\re$};
    \draw[->] (0,-2.3) -- (0,2.3) node[above]{$\im$};
    \draw[blue,thick] (0,0) circle(1.3);
    \draw[blue!50] (0,0) circle(1.7);
    \draw[blue!50] (0,0) circle(2.1);
    \draw[red!50] (0,0) -- (2.3,0);
    \draw[red!50] (0,0) -- ({2.3*cos(45)},{2.3*sin(45)});
    \draw[red!50] (0,0) -- (0,2.3);
    \draw[red!50] (0,0) -- ({2.3*cos(135)},{2.3*sin(135)});
    \node[below] at (0,-2.6) {$z$-plane};
  \end{scope}
  \draw[->,thick] (-0.7,0) -- (0.7,0) node[midway,above]{$z+\frac{1}{z}$};
  \begin{scope}[shift={(3.5,0)}]
    \draw[->] (-3,0) -- (3,0) node[right]{$\re$};
    \draw[->] (0,-2.3) -- (0,2.3) node[above]{$\im$};
    \draw[blue,thick] (0,0) ellipse(2.07 and 0.53);
    \draw[blue!50] (0,0) ellipse(2.29 and 1.11);
    \draw[blue!50] (0,0) ellipse(2.58 and 1.62);
    \draw[red,thick] (-2,0) -- (2,0);
    \node[below left] at (-2,0) {\small$-2$};
    \node[below right] at (2,0) {\small$2$};
    \node[below] at (0,-2.6) {$w$-plane};
  \end{scope}
\end{tikzpicture}
\end{center}

This map is fundamental in aerodynamics: by applying $J$ to a circle
passing close to $z=-1$, one obtains airfoil profiles (Joukowski airfoils).
\end{example}

\subsection{Lens domains}

\begin{example}[Lens to half-plane]\label{ex:lens}
A \emph{lens} (or \emph{lune}) is a region bounded by two circular arcs
meeting at angles $\alpha$ at their intersection points $p$ and $q$.
To map a lens with vertices at $\pm 1$ and interior angle $\alpha$ to
the upper half-plane:
\begin{enumerate}[label=(\arabic*)]
  \item Apply $T(z)=(z-1)/(z+1)$, sending $1\mapsto 0$,
        $-1\mapsto\infty$.  The lens becomes a sector of angle $\alpha$.
  \item Apply $w\mapsto w^{\pi/\alpha}$ to open the sector to a
        half-plane.
\end{enumerate}
The composition
$w=\left(\dfrac{z-1}{z+1}\right)^{\!\pi/\alpha}$
maps the lens conformally to $\mathbb{H}$.
\end{example}

\subsection{The logarithm and related maps}

\begin{example}[Slit plane $\to$ strip]\label{ex:log-map}
The principal logarithm $\Log z=\ln\abs{z}+i\Arg z$ (with $-\pi<\Arg z
<\pi$) maps $\C\setminus(-\infty,0]$ conformally onto the horizontal
strip $\{w:-\pi<\im w<\pi\}$.
\end{example}

\begin{example}[Upper half-plane $\to$ strip via $\Log$]\label{ex:log-H-to-strip}
The map $w=\Log z$ sends $\mathbb{H}$ to the strip $\{0<\im w<\pi\}$.
This is the inverse of $e^w:\{0<\im w<\pi\}\to\mathbb{H}$.
\end{example}

\subsection{Summary table of standard conformal maps}

\begin{center}
\renewcommand{\arraystretch}{1.5}
\begin{tabular}{lll}
\toprule
\textbf{Domain} & \textbf{Map} & \textbf{Image} \\
\midrule
$\mathbb{H}$ & $\dfrac{z-i}{z+i}$ & $\D$ \\[6pt]
$\{0<\im z<\pi\}$ & $e^z$ & $\mathbb{H}$ \\[4pt]
$\{0<\arg z<\alpha\}$ & $z^{\pi/\alpha}$ & $\mathbb{H}$ \\[4pt]
$\{\re z>0\}$ & $\dfrac{z-1}{z+1}$ & $\D$ \\[6pt]
$\C\setminus(-\infty,0]$ & $\Log z$ & $\{-\pi<\im w<\pi\}$ \\[4pt]
$\C\setminus[-2,2]$ & $\frac{1}{2}(z+\sqrt{z^2-4})$ & $\{\abs{w}>1\}$ \\[4pt]
$\D$ & $\dfrac{z}{(1-z)^2}$ (Koebe) & $\C\setminus(-\infty,-\tfrac{1}{4}]$ \\
\bottomrule
\end{tabular}
\end{center}

%-------------------------------------------------------------------
\section{Images of Grids Under Conformal Maps}\label{sec:grid-images}
%-------------------------------------------------------------------

A powerful way to visualise a conformal map $f$ is to draw the images
of a Cartesian grid.  Since $f$ is conformal, the image curves intersect
at right angles wherever $f'\neq 0$.

\begin{example}[Grid image under $f(z)=z^2$]\label{ex:grid-z2}
Under $w=z^2=(x+iy)^2=(x^2-y^2)+2ixy$:
\begin{itemize}
  \item Vertical lines $x=c$ map to the parabolas $u=c^2-v^2/(4c^2)$,
        opening to the left.
  \item Horizontal lines $y=c$ map to the parabolas $u=v^2/(4c^2)-c^2$,
        opening to the right.
\end{itemize}
These two families of parabolas are orthogonal, confirming conformality.
\end{example}

\begin{example}[Grid image under $f(z)=e^z$]\label{ex:grid-exp}
Under $w=e^z=e^x(\cos y+i\sin y)$:
\begin{itemize}
  \item Vertical lines $x=c$ map to circles $\abs{w}=e^c$.
  \item Horizontal lines $y=c$ map to rays $\arg w=c$.
\end{itemize}
The Cartesian grid becomes a polar grid.
\end{example}

\begin{example}[Grid image under $f(z)=1/z$]\label{ex:grid-inv}
Under $w=1/z=\conj{z}/\abs{z}^2$, with $z=x+iy$:
\[
  u=\frac{x}{x^2+y^2},\quad v=\frac{-y}{x^2+y^2}.
\]
\begin{itemize}
  \item Vertical lines $x=c\neq 0$ map to circles
        $\left(u-\frac{1}{2c}\right)^2+v^2=\frac{1}{4c^2}$
        (passing through the origin).
  \item Horizontal lines $y=c\neq 0$ map to circles
        $u^2+\left(v+\frac{1}{2c}\right)^2=\frac{1}{4c^2}$.
\end{itemize}

\begin{center}
\begin{tikzpicture}[scale=0.8,>=Stealth]
  % Left: grid
  \begin{scope}[shift={(-3.5,0)}]
    \draw[->] (-2.5,0) -- (2.5,0) node[right]{$\re$};
    \draw[->] (0,-2.2) -- (0,2.2) node[above]{$\im$};
    \foreach \c in {0.5,1.0,1.5,2.0} {
      \draw[blue!60,thin] (\c,-2) -- (\c,2);
      \draw[blue!60,thin] (-\c,-2) -- (-\c,2);
    }
    \foreach \c in {0.5,1.0,1.5} {
      \draw[red!60,thin] (-2.2,\c) -- (2.2,\c);
      \draw[red!60,thin] (-2.2,-\c) -- (2.2,-\c);
    }
    \node[below] at (0,-2.5) {$z$-plane};
  \end{scope}
  \draw[->,thick] (-0.7,0) -- (0.7,0) node[midway,above]{$1/z$};
  % Right: image (circles through origin)
  \begin{scope}[shift={(3.5,0)}]
    \draw[->] (-2.5,0) -- (2.5,0) node[right]{$\re$};
    \draw[->] (0,-2.2) -- (0,2.2) node[above]{$\im$};
    % Images of x=c: circles centered at (1/(2c),0)
    \foreach \c in {0.5,1.0,1.5,2.0} {
      \draw[blue!60,thin] ({1/(2*\c)},0) circle({1/(2*\c)});
      \draw[blue!60,thin] ({-1/(2*\c)},0) circle({1/(2*\c)});
    }
    % Images of y=c: circles centered at (0,-1/(2c))
    \foreach \c in {0.5,1.0,1.5} {
      \draw[red!60,thin] (0,{-1/(2*\c)}) circle({1/(2*\c)});
      \draw[red!60,thin] (0,{1/(2*\c)}) circle({1/(2*\c)});
    }
    \node[below] at (0,-2.5) {$w$-plane};
  \end{scope}
\end{tikzpicture}
\end{center}
\end{example}

\begin{example}[M\"obius transformation action on circles]\label{ex:mobius-grid}
Consider $T(z)=\dfrac{z}{1-z}$, which sends $0\mapsto 0$,
$\infty\mapsto -1$, $1\mapsto\infty$.  It maps the unit circle
$\abs{z}=1$ to the vertical line $\re w=-1/2$ (since $T$ maps
the boundary of $\D$ to a generalised circle through $T(1)=\infty$).

The family of circles $\abs{z}=r$ for $0<r<1$ maps to a family of
circles in $\{\re w>-1/2\}$ that shrink to the point $0$ as $r\to 0$.

\begin{center}
\begin{tikzpicture}[scale=0.85,>=Stealth]
  \begin{scope}[shift={(-3.5,0)}]
    \draw[->] (-2.2,0) -- (2.2,0) node[right]{$\re$};
    \draw[->] (0,-2) -- (0,2) node[above]{$\im$};
    \draw[blue,thick] (0,0) circle(1.5);
    \draw[blue!60] (0,0) circle(0.5);
    \draw[blue!60] (0,0) circle(1.0);
    \draw[red!60] (0,0) -- (2,0);
    \draw[red!60] (0,0) -- ({2*cos(60)},{2*sin(60)});
    \draw[red!60] (0,0) -- ({2*cos(120)},{2*sin(120)});
    \draw[red!60] (0,0) -- (0,2);
    \node[below] at (0,-2.3) {$z$-plane};
  \end{scope}
  \draw[->,thick] (-0.7,0) -- (0.7,0) node[midway,above]{$\frac{z}{1-z}$};
  \begin{scope}[shift={(3.5,0)}]
    \draw[->] (-2.2,0) -- (2.2,0) node[right]{$\re$};
    \draw[->] (0,-2) -- (0,2) node[above]{$\im$};
    \draw[blue,thick] (-0.75,0) -- (-0.75,2);
    \draw[blue,thick] (-0.75,0) -- (-0.75,-2);
    \draw[blue!60] (-0.167,0) circle(0.167);
    \draw[blue!60] (-0.5,0) circle(0.5);
    \node[below] at (0,-2.3) {$w$-plane};
  \end{scope}
\end{tikzpicture}
\end{center}
\end{example}

%-------------------------------------------------------------------
\section{The Schwarz--Christoffel Formula}\label{sec:schwarz-christoffel}
%-------------------------------------------------------------------

We briefly state the Schwarz--Christoffel formula, which gives an
explicit conformal map from the upper half-plane to any polygon.

\begin{theorem}[Schwarz--Christoffel formula]\label{thm:schwarz-christoffel}
Let $P$ be a polygon with vertices $w_1,\dots,w_n$ (in order)
and interior angles $\alpha_1\pi,\dots,\alpha_n\pi$ (where
$0<\alpha_k<2$ and $\sum\alpha_k=n-2$).  Then a conformal map
$f\colon\mathbb{H}\to\mathrm{int}(P)$ has the form
\[
  f(z)=A\int_{z_0}^{z}\prod_{k=1}^{n}(\zeta-x_k)^{\alpha_k-1}\,
  \dd\zeta+B,
\]
where $x_1<x_2<\cdots<x_n$ are points on $\R$ with $f(x_k)=w_k$,
and $A,B\in\C$ are constants.
\end{theorem}

\begin{example}[Map to a half-strip]\label{ex:sc-half-strip}
The infinite vertical half-strip
$\{w:\re w\in(-\pi/2,\pi/2),\,\im w>0\}$ is the image of
$\mathbb{H}$ under $f(z)=\arcsin z$.  This is a Schwarz--Christoffel
map with two finite vertices at $w=\pm\pi/2$ (angles $\pi/2$ each)
and one vertex at $\infty$.
\end{example}

%-------------------------------------------------------------------
\section{Exercises}\label{sec:ex-ch8}
%-------------------------------------------------------------------

\begin{exercise}[Basic conformal maps]\label{exer:basic-conformal}
\begin{enumerate}[label=(\alph*)]
  \item Show that $f(z)=z+1/z$ maps $\D\setminus\{0\}$ to
        $\C\setminus[-2,2]$ (but is $2$-to-$1$ on $\D$).
  \item Find a conformal map from the first quadrant
        $\{\re z>0,\,\im z>0\}$ to the unit disk $\D$.
  \item Find a conformal map from $\{z:\abs{z}<1,\,\im z>0\}$
        (upper half of the disk) to $\D$.
  \item Find a conformal map from the vertical strip
        $\{z:0<\re z<1\}$ to $\mathbb{H}$.
\end{enumerate}
\end{exercise}

\begin{exercise}[M\"obius transformations]\label{exer:mobius}
\begin{enumerate}[label=(\alph*)]
  \item Find the M\"obius transformation sending
        $1\mapsto 0$, $i\mapsto 1$, $-1\mapsto\infty$.
  \item Show that the M\"obius transformations preserving $\D$ are
        exactly those of the form $e^{i\theta}\frac{z-a}{1-\conj{a}z}$.
  \item Prove that a M\"obius transformation mapping $\R\cup\{\infty\}$
        to itself has the form $(az+b)/(cz+d)$ with $a,b,c,d\in\R$.
  \item Let $T(z)=(2z+i)/(iz+2)$.  Find the image of $\D$ under $T$.
\end{enumerate}
\end{exercise}

\begin{exercise}[Cross-ratio]\label{exer:cross-ratio}
\begin{enumerate}[label=(\alph*)]
  \item Compute $(1,i;-1,-i)$.
  \item Show that four points lie on a generalised circle if and only if
        their cross-ratio is real.
  \item Show that the cross-ratio takes only six values under
        permutations of the four points.  If $\lambda=(z_1,z_2;z_3,z_4)$,
        list all six values in terms of $\lambda$.
\end{enumerate}
\end{exercise}

\begin{exercise}[Conformal maps to standard domains]\label{exer:conf-standard}
Find explicit conformal maps from the following regions to the unit disk
$\D$:
\begin{enumerate}[label=(\alph*)]
  \item The sector $\{z:0<\arg z<\pi/3\}$.
  \item The region $\{z:\abs{z}>1,\,\im z>0\}$
        (upper half of the complement of $\D$).
  \item The lens-shaped region bounded by $\abs{z-i}=\sqrt{2}$ and
        $\abs{z+i}=\sqrt{2}$.
  \item The region $\{z:\re z>0\}\setminus[0,1]$
        (right half-plane with a slit).
\end{enumerate}
\end{exercise}

\begin{exercise}[Schwarz--Christoffel]\label{exer:sc}
\begin{enumerate}[label=(\alph*)]
  \item Verify that $f(z)=\sqrt{z}$ (principal branch) is the
        Schwarz--Christoffel map from $\mathbb{H}$ to the first
        quadrant.
  \item Show that $f(z)=\cosh^{-1}(z)$ maps $\mathbb{H}$ to the
        semi-infinite strip $\{w:0<\im w<\pi,\,\re w>0\}$.
  \item Using the Schwarz--Christoffel formula, find the conformal map
        from $\mathbb{H}$ to an equilateral triangle.  Express the
        answer as an integral.
\end{enumerate}
\end{exercise}

\begin{exercise}[Geometric properties]\label{exer:geometric}
\begin{enumerate}[label=(\alph*)]
  \item Let $f\colon\D\to\D$ be holomorphic with $f(0)=0$.
        Show that $\abs{f(z)}\leq\abs{z}$ for all $z\in\D$ (Schwarz
        lemma) and deduce that $\abs{f'(0)}\leq 1$.
  \item Prove that $\mathrm{Aut}(\C)=\{z\mapsto az+b:a\neq 0\}$.
  \item Prove that $\mathrm{Aut}(\C\setminus\{0\})$ consists of
        $z\mapsto az$ and $z\mapsto a/z$ for $a\neq 0$.
  \item Show that there is no conformal equivalence between
        $\C$ and $\D$.  (\emph{Hint:} Liouville's theorem.)
\end{enumerate}
\end{exercise}

\begin{exercise}[Visualisation]\label{exer:visualisation}
Describe (or sketch) the image of the Cartesian grid
$\{x+iy:x\in\Z\text{ or }y\in\Z\}$ under the following maps:
\begin{enumerate}[label=(\alph*)]
  \item $f(z)=z^2$
  \item $f(z)=e^z$
  \item $f(z)=\sin z$
  \item $f(z)=(z-1)/(z+1)$
\end{enumerate}
For (c), note that $\sin(x+iy)=\sin x\cosh y+i\cos x\sinh y$, so
vertical lines map to hyperbolas and horizontal lines map to ellipses
(all confocal with foci at $\pm 1$).
\end{exercise}
%%%%%%%%%%%%%%%%%%%%%%%%%%%%%%%%%%%%%%%%%%%%%%%%%%%%%%%%%%%%%%%%%%%%
%  Complex Analysis — Chapters 9–10 + Back Matter
%  Course Notes — English Edition
%%%%%%%%%%%%%%%%%%%%%%%%%%%%%%%%%%%%%%%%%%%%%%%%%%%%%%%%%%%%%%%%%%%%

%% ═══════════════════════════════════════════════════════════════════
%%  CHAPTER 9 : Rouché's Theorem and the Argument Principle
%% ═══════════════════════════════════════════════════════════════════
\chapter{Rouch\'{e}'s Theorem and the Argument Principle}\label{ch:rouche-argument}

\section{Zeros of Holomorphic Functions}\label{sec:zeros-holo}

We begin by establishing the fundamental algebraic structure
of zeros of holomorphic functions, which is considerably more
rigid than in the real-variable setting.

\begin{definition}[Zero of order $m$]\label{def:zero-order}
Let $f$ be holomorphic on a domain $\Omega \subseteq \C$ and let $z_0 \in \Omega$.
We say that $f$ has a \textbf{zero of order}\index{zero!order of}
(or \textbf{multiplicity}\index{multiplicity}) $m \geq 1$ at $z_0$ if
\[
  f(z_0) = f'(z_0) = \cdots = f^{(m-1)}(z_0) = 0, \qquad f^{(m)}(z_0) \neq 0.
\]
Equivalently, $f$ can be written as
\[
  f(z) = (z - z_0)^m \, g(z),
\]
where $g$ is holomorphic on $\Omega$ and $g(z_0) \neq 0$.
\end{definition}

\begin{proposition}[Characterisation of zero order]\label{prop:zero-order-char}
Let $f$ be holomorphic on a domain $\Omega$ with $f \not\equiv 0$, and
suppose $f(z_0) = 0$ for some $z_0 \in \Omega$. Then there exists a
unique integer $m \geq 1$ such that $f(z) = (z-z_0)^m g(z)$ with
$g$ holomorphic on $\Omega$ and $g(z_0) \neq 0$.
\end{proposition}

\begin{proof}
Since $f \not\equiv 0$, the Taylor series of $f$ about $z_0$,
\[
  f(z) = \sum_{n=0}^{\infty} a_n (z - z_0)^n,
  \qquad a_n = \frac{f^{(n)}(z_0)}{n!},
\]
is not identically zero. Let $m$ be the smallest index with $a_m \neq 0$.
Then $m \geq 1$ (since $a_0 = f(z_0) = 0$) and
\[
  f(z) = (z - z_0)^m \sum_{n=0}^{\infty} a_{m+n} (z-z_0)^n
  = (z-z_0)^m \, g(z).
\]
The function $g(z) = \sum_{n=0}^{\infty} a_{m+n}(z-z_0)^n$ is
holomorphic on the disc of convergence and satisfies
$g(z_0) = a_m \neq 0$. Extending $g$ to all of $\Omega$ by
$g(z) = f(z)/(z-z_0)^m$ for $z \neq z_0$ gives the result.
Uniqueness of $m$ follows from the identity theorem.
\end{proof}

\begin{theorem}[Isolation of zeros]\label{thm:isolated-zeros}
\index{zeros!isolation of}
Let $f$ be holomorphic on a domain $\Omega$ with $f \not\equiv 0$.
Then the zeros of $f$ are isolated: for every zero $z_0$,
there exists $r > 0$ such that $f(z) \neq 0$ for
$0 < \abs{z - z_0} < r$.
\end{theorem}

\begin{proof}
Write $f(z) = (z-z_0)^m g(z)$ as in Proposition~\ref{prop:zero-order-char}.
Since $g(z_0) \neq 0$ and $g$ is continuous, there exists $r > 0$
such that $g(z) \neq 0$ for $\abs{z-z_0} < r$. For $0 < \abs{z-z_0} < r$,
we then have $f(z) = (z-z_0)^m g(z) \neq 0$.
\end{proof}

\begin{corollary}[Identity theorem]\label{cor:identity-thm}
\index{identity theorem}
If $f$ and $g$ are holomorphic on a domain $\Omega$ and
$f(z_n) = g(z_n)$ for a sequence $(z_n)$ with a limit point
in $\Omega$, then $f \equiv g$ on $\Omega$.
\end{corollary}

\begin{proof}
Apply Theorem~\ref{thm:isolated-zeros} to $h = f - g$. If
$h \not\equiv 0$, then its zeros are isolated, contradicting
the existence of an accumulation point of zeros in $\Omega$.
\end{proof}

\begin{definition}[Meromorphic function — zeros and poles]\label{def:zeros-poles-count}
\index{meromorphic function}
Let $f$ be meromorphic on $\Omega$. For $z_0 \in \Omega$, define
the \textbf{order} of $f$ at $z_0$, denoted $\ord_{z_0}(f)$, by:
\begin{itemize}
  \item $\ord_{z_0}(f) = m > 0$ if $f$ has a zero of order $m$ at $z_0$;
  \item $\ord_{z_0}(f) = -m < 0$ if $f$ has a pole of order $m$ at $z_0$;
  \item $\ord_{z_0}(f) = 0$ if $f$ is holomorphic and nonzero at $z_0$.
\end{itemize}
\end{definition}

\begin{remark}[Logarithmic derivative]\label{rem:log-derivative}
\index{logarithmic derivative}
If $f$ is meromorphic and not identically zero, then the
\textbf{logarithmic derivative} $f'/f$ is meromorphic, and at
each point $z_0$ where $\ord_{z_0}(f) = m$ (possibly negative), we have
\[
  \frac{f'(z)}{f(z)} = \frac{m}{z - z_0} + \text{holomorphic terms near } z_0.
\]
Hence $\Res_{z_0}\!\bigl(\frac{f'}{f}\bigr) = m = \ord_{z_0}(f)$.
This observation is the key to the argument principle.
\end{remark}


\section{The Argument Principle}\label{sec:argument-principle}

\begin{theorem}[Argument Principle]\label{thm:argument-principle}
\index{argument principle}
Let $f$ be meromorphic on a domain containing $\overline{D}$, where
$D$ is a bounded domain with piecewise-smooth boundary $\partial D$,
oriented positively. Suppose $f$ has no zeros or poles on $\partial D$.
Let $Z$ denote the number of zeros and $P$ the number of poles of $f$
inside $D$, counted with multiplicity. Then
\begin{equation}\label{eq:argument-principle}
  \frac{1}{2\pi i} \oint_{\partial D} \frac{f'(z)}{f(z)} \, \dd z
  = Z - P.
\end{equation}
\end{theorem}

\begin{proof}
Since $f$ is meromorphic on a domain containing $\overline{D}$ and
has no zeros or poles on $\partial D$, the function $f'/f$ is
meromorphic on a neighbourhood of $\overline{D}$ with poles only
at the zeros and poles of $f$.

Let $z_1, \ldots, z_N$ be the zeros of $f$ in $D$, with respective orders
$m_1, \ldots, m_N$, and let $w_1, \ldots, w_M$ be the poles of $f$ in $D$,
with respective orders $p_1, \ldots, p_M$. Near a zero $z_k$ of order $m_k$,
we can write $f(z) = (z - z_k)^{m_k} g_k(z)$ with $g_k$ holomorphic and nonvanishing, so
\[
  \frac{f'(z)}{f(z)} = \frac{m_k}{z - z_k} + \frac{g_k'(z)}{g_k(z)}.
\]
The second term is holomorphic near $z_k$, so
$\Res_{z_k}\!\bigl(\frac{f'}{f}\bigr) = m_k$.

Similarly, near a pole $w_j$ of order $p_j$, writing
$f(z) = (z - w_j)^{-p_j} h_j(z)$ with $h_j$ holomorphic and nonvanishing:
\[
  \frac{f'(z)}{f(z)} = \frac{-p_j}{z - w_j} + \frac{h_j'(z)}{h_j(z)},
\]
so $\Res_{w_j}\!\bigl(\frac{f'}{f}\bigr) = -p_j$.

By the Residue Theorem,
\[
  \frac{1}{2\pi i} \oint_{\partial D} \frac{f'(z)}{f(z)} \, \dd z
  = \sum_{k=1}^{N} m_k + \sum_{j=1}^{M} (-p_j) = Z - P. \qedhere
\]
\end{proof}

\begin{remark}[Topological interpretation]\label{rem:argument-principle-winding}
The integral on the left-hand side of \eqref{eq:argument-principle} can be
interpreted as the \textbf{winding number}\index{winding number} of the
image curve $f \circ \gamma$ around the origin:
\[
  \frac{1}{2\pi i} \oint_{\gamma} \frac{f'(z)}{f(z)} \, \dd z
  = \ind(f \circ \gamma, 0),
\]
where $\gamma$ parametrises $\partial D$. Indeed, writing
$w = f(z)$ gives $\dd w = f'(z)\, \dd z$, so the integral becomes
$\frac{1}{2\pi i}\oint_{f \circ \gamma} \frac{\dd w}{w}$.
Thus the argument principle states that the number of times the image
$f(\gamma)$ winds around the origin equals $Z - P$.
\end{remark}

\medskip

\begin{center}
\begin{tikzpicture}[scale=1.0, >=Stealth]
  % Left: domain with zeros and poles
  \begin{scope}[xshift=-4cm]
    \draw[thick, blue!70!black, postaction={decorate,
      decoration={markings, mark=at position 0.15 with {\arrow{>}},
        mark=at position 0.45 with {\arrow{>}},
        mark=at position 0.75 with {\arrow{>}}}}]
      (0,0) circle (2.2);
    \node at (0, -2.7) {$\gamma = \partial D$};
    % zeros
    \filldraw[red!70!black] (0.5, 0.8) circle (2pt) node[above right]{\small $z_1$};
    \filldraw[red!70!black] (-0.7, -0.5) circle (2pt) node[below left]{\small $z_2$};
    \filldraw[red!70!black] (0.9, -0.6) circle (2pt) node[below right]{\small $z_3$};
    % pole
    \draw[thick, blue] (-0.3, 0.3) node{$\times$};
    \node at (-0.3, 0.65) {\small $w_1$};
    % labels
    \node at (0, -3.3) {\small $Z = 3,\; P = 1$};
    \node[red!70!black] at (1.5, 1.5) {\small zeros: $\bullet$};
    \node[blue] at (1.5, 1.0) {\small poles: $\times$};
  \end{scope}
  % Arrow
  \draw[->, thick, line width=1.5pt] (-0.8, 0) -- (0.8, 0)
    node[midway, above]{$f$};
  % Right: image curve winding around origin
  \begin{scope}[xshift=4cm]
    \draw[thick, red!60!black, postaction={decorate,
      decoration={markings, mark=at position 0.1 with {\arrow{>}},
        mark=at position 0.4 with {\arrow{>}},
        mark=at position 0.7 with {\arrow{>}},
        mark=at position 0.95 with {\arrow{>}}}}]
      plot[smooth cycle, tension=0.7] coordinates
      {(1.8,0.2) (1.0,1.6) (-0.8,1.4) (-1.6,0.2)
       (-1.4,-1.0) (-0.2,-1.6) (1.2,-1.2) (1.9,-0.3)};
    \filldraw (0,0) circle (2pt) node[below right]{\small $0$};
    \node at (0, -2.7) {$f \circ \gamma$};
    \node at (0, -3.3) {\small $\ind(f \circ \gamma, 0) = 2$};
  \end{scope}
\end{tikzpicture}
\captionof{figure}{The argument principle: $f$ maps the boundary $\gamma$
to a curve winding $Z - P = 3 - 1 = 2$ times around the origin.}
\label{fig:argument-principle}
\end{center}

\bigskip

\begin{example}[Counting zeros of a polynomial]\label{ex:argument-zeros-poly}
Let $f(z) = z^4 + 6z + 3$. We want to determine how many zeros $f$ has
in the disc $\abs{z} < 2$. On $\abs{z} = 2$:
\[
  \abs{z^4} = 16, \qquad \abs{6z + 3} \leq 6 \cdot 2 + 3 = 15 < 16.
\]
Since $f$ is a polynomial (no poles), the argument principle gives
$Z = \frac{1}{2\pi i}\oint_{\abs{z}=2} \frac{f'(z)}{f(z)}\,\dd z$.
By the estimate above and Rouch\'{e}'s theorem (which we shall prove shortly),
$f$ has the same number of zeros in $\abs{z} < 2$ as $z^4$, namely $4$.
Since $\deg f = 4$, all zeros of $f$ lie in $\abs{z} < 2$.
\end{example}

\begin{proposition}[Argument principle for $f - a$]\label{prop:argument-principle-a}
Under the hypotheses of Theorem~\ref{thm:argument-principle}, if
$a \in \C$ is not in the image $f(\partial D)$, then
\[
  \frac{1}{2\pi i}\oint_{\partial D} \frac{f'(z)}{f(z) - a}\,\dd z
  = N(f, a, D) - P,
\]
where $N(f, a, D)$ denotes the number of solutions of $f(z) = a$ in $D$,
counted with multiplicity, and $P$ is the number of poles.
\end{proposition}

\begin{proof}
Apply Theorem~\ref{thm:argument-principle} to $g(z) = f(z) - a$,
noting that $g$ has the same poles as $f$ and zeros exactly where $f(z) = a$.
\end{proof}

\section{Rouch\'{e}'s Theorem}\label{sec:rouche}

Rouch\'{e}'s theorem is one of the most powerful tools for locating zeros
of holomorphic functions. It allows one to compare the zeros of a
``complicated'' function with those of a simpler one.

\begin{theorem}[Rouch\'{e}'s Theorem]\label{thm:rouche}
\index{Rouch\'{e}'s theorem}
Let $f$ and $g$ be holomorphic on a domain containing $\overline{D}$,
where $D$ is a bounded domain with piecewise-smooth positively oriented
boundary $\partial D$. If
\begin{equation}\label{eq:rouche-condition}
  \abs{g(z)} < \abs{f(z)} \quad \text{for all } z \in \partial D,
\end{equation}
then $f$ and $f + g$ have the same number of zeros (counted with multiplicity)
inside $D$.
\end{theorem}

\begin{proof}
The condition $\abs{g(z)} < \abs{f(z)}$ on $\partial D$ implies in
particular that $f(z) \neq 0$ on $\partial D$, and also that
$(f+g)(z) \neq 0$ on $\partial D$ (since $\abs{f+g} \geq \abs{f} - \abs{g} > 0$).

Consider the family of functions
\[
  f_t(z) = f(z) + t\, g(z), \qquad t \in [0,1].
\]
For each $t \in [0,1]$ and $z \in \partial D$:
\[
  \abs{f_t(z)} = \abs{f(z) + tg(z)} \geq \abs{f(z)} - t\abs{g(z)}
  \geq \abs{f(z)} - \abs{g(z)} > 0.
\]
Thus $f_t$ has no zeros on $\partial D$ for any $t$. By the argument
principle, the number of zeros of $f_t$ inside $D$ is
\[
  N(t) = \frac{1}{2\pi i} \oint_{\partial D}
  \frac{f_t'(z)}{f_t(z)}\,\dd z
  = \frac{1}{2\pi i} \oint_{\partial D}
  \frac{f'(z) + t\, g'(z)}{f(z) + t\, g(z)}\,\dd z.
\]

The integrand depends continuously on $t$ (uniformly on $\partial D$,
since the denominator is bounded away from zero). Therefore $N(t)$ is
a continuous function of $t$. But $N(t)$ is always a non-negative
integer, so it must be constant. Hence
$N(0) = N(1)$, i.e., $f$ and $f + g$ have the same number of zeros
in $D$.
\end{proof}

\begin{remark}[Symmetric form of Rouch\'{e}'s theorem]\label{rem:rouche-symmetric}
A stronger, symmetric version states: if $\abs{f(z) - g(z)} < \abs{f(z)} + \abs{g(z)}$
on $\partial D$, then $f$ and $g$ have the same number of zeros in $D$.
The classical form follows by setting $g_{\text{sym}} = f + g$ and noting
that $\abs{(f+g) - f} = \abs{g} < \abs{f} \leq \abs{f} + \abs{f+g}$ whenever
$\abs{g} < \abs{f}$.
\end{remark}

\subsection{Applications of Rouch\'{e}'s Theorem}

\begin{theorem}[Fundamental Theorem of Algebra via Rouch\'{e}]\label{thm:FTA-rouche}
\index{Fundamental Theorem of Algebra}
Every non-constant polynomial $p(z) = a_n z^n + a_{n-1}z^{n-1} + \cdots + a_0$
with $a_n \neq 0$ has exactly $n$ zeros in $\C$ (counted with multiplicity).
\end{theorem}

\begin{proof}
Set $f(z) = a_n z^n$ and $g(z) = a_{n-1}z^{n-1} + \cdots + a_0$. On
$\abs{z} = R$:
\[
  \abs{g(z)} \leq \abs{a_{n-1}} R^{n-1} + \cdots + \abs{a_0}
  \leq C\, R^{n-1},
\]
where $C = \abs{a_{n-1}} + \cdots + \abs{a_0}$. Also
$\abs{f(z)} = \abs{a_n}\, R^n$. For $R > C/\abs{a_n}$ we have
$\abs{g(z)} < \abs{f(z)}$ on $\abs{z} = R$. By Rouch\'{e}'s theorem,
$p = f + g$ has the same number of zeros in $\abs{z} < R$ as $f(z) = a_n z^n$,
which has exactly $n$ zeros (all at the origin). Hence $p$ has $n$ zeros
in $\C$.
\end{proof}

\begin{example}[Counting zeros in a specific region]\label{ex:rouche-counting}
Determine the number of roots of $p(z) = z^5 + 3z^3 + 7$ in the
annulus $1 < \abs{z} < 2$.

\textbf{Step 1: Zeros in $\abs{z} < 1$.}
On $\abs{z} = 1$, set $f(z) = 7$ and $g(z) = z^5 + 3z^3$.
Then $\abs{g(z)} \leq 1 + 3 = 4 < 7 = \abs{f(z)}$.
By Rouch\'{e}, $p$ has the same number of zeros in $\abs{z} < 1$ as
$f(z) = 7$, which is $0$.

\textbf{Step 2: Zeros in $\abs{z} < 2$.}
On $\abs{z} = 2$, set $f(z) = z^5$ and $g(z) = 3z^3 + 7$.
Then $\abs{f(z)} = 32$ and $\abs{g(z)} \leq 3 \cdot 8 + 7 = 31 < 32$.
By Rouch\'{e}, $p$ has $5$ zeros in $\abs{z} < 2$.

\textbf{Conclusion.} There are $5 - 0 = 5$ zeros in $1 \leq \abs{z} < 2$.
Since $\abs{p(z)}$ on $\abs{z} = 1$ satisfies $\abs{p(z)} \geq 7 - 4 = 3 > 0$,
there are no zeros on $\abs{z} = 1$, so all $5$ zeros lie in $1 < \abs{z} < 2$.
\end{example}

\begin{example}[Perturbation of zeros]\label{ex:rouche-perturbation}
Show that for small $\varepsilon > 0$, the equation $e^z = 1 + z + \varepsilon z^2$
has exactly one root in $\abs{z} < 1$.

Set $f(z) = 1 + z$ and $g(z) = e^z - (1 + z) - \varepsilon z^2$. Note that
$f + g = e^z - \varepsilon z^2$. We use a different decomposition: let
$F(z) = e^z$ and $G(z) = -1 - z - \varepsilon z^2$, so $F + G = e^z - 1 - z - \varepsilon z^2$.

Actually, we apply Rouch\'{e} differently. Consider $h(z) = e^z - 1 - z$ and
$k(z) = \varepsilon z^2$. On $\abs{z} = 1$, we have
$\abs{h(z)} = \abs{\sum_{n=2}^{\infty} z^n/n!} \leq \sum_{n=2}^{\infty} 1/n! = e - 2 \approx 0.718$.
Also $\abs{k(z)} = \varepsilon$.

We want to compare $e^z$ and $1 + z + \varepsilon z^2$. Let $f(z) = e^z$
and $g(z) = -(1+z+\varepsilon z^2)$. On $\abs{z}=1$: $\abs{f(z)} = e^{\re z} \geq e^{-1} \approx 0.368$
and $\abs{g(z)} \leq 1 + 1 + \varepsilon = 2 + \varepsilon$. This does not directly work.

Instead, set $f(z) = 1 + z$ and $g(z) = \varepsilon z^2$, so $f + g = 1 + z + \varepsilon z^2$.
On $\abs{z} = r$ for suitable $r$: we need $\abs{f(z)} > \abs{g(z)}$.
On $\abs{z} = 1/2$: $\abs{f(z)} \geq 1 - 1/2 = 1/2$ and $\abs{g(z)} = \varepsilon/4$.
For $\varepsilon < 2$, Rouch\'{e} gives that $1 + z + \varepsilon z^2$ has
$1$ zero in $\abs{z} < 1/2$ (same as $1 + z$, which has its zero at $z = -1$,
outside the disc). This gives $0$ zeros.

Let us reconsider the problem. The equation is $e^z = 1 + z + \varepsilon z^2$.
Set $h(z) = e^z - 1 - z = \frac{z^2}{2} + \frac{z^3}{6} + \cdots$, so
$h$ has a zero of order $2$ at $z=0$. Now $e^z - (1+z+\varepsilon z^2) = h(z) - \varepsilon z^2
= (\frac{1}{2}-\varepsilon)z^2 + \frac{z^3}{6}+\cdots$. On $\abs{z}=\delta$ small:
set $f(z) = (\frac{1}{2}-\varepsilon)z^2$ and $g(z) = \frac{z^3}{6}+\cdots$, so
$\abs{f(z)} = \abs{\frac{1}{2}-\varepsilon}\delta^2$ and $\abs{g(z)} \leq C\delta^3$.
For $\delta$ small enough, $\abs{g} < \abs{f}$, giving $2$ zeros of $e^z - 1 - z - \varepsilon z^2$ near the origin.

On $\abs{z} = 1$: let $f(z) = e^z - 1 - z$ and $g(z) = -\varepsilon z^2$.
Then $\abs{g(z)} = \varepsilon$. We need $\abs{f(z)} > \varepsilon$ on $\abs{z}=1$.
Since $f(z) = z^2/2 + z^3/6 + \cdots$ and $\abs{f(z)}$ achieves a minimum on
$\abs{z}=1$ that is strictly positive (as $f$ has no zeros on $\abs{z}=1$;
its only zero is at $z=0$), for $\varepsilon$ sufficiently small, Rouch\'{e}
applies and $e^z - 1 - z - \varepsilon z^2$ has $2$ zeros in $\abs{z}<1$
(same as $e^z - 1 - z$).

But the original equation asks for roots of $e^z - 1 - z - \varepsilon z^2 = 0$.
Since $e^z - 1 - z$ has a double zero at $0$, and the perturbation
$\varepsilon z^2$ splits it, we get exactly $2$ roots near the origin
(and hence in $\abs{z}<1$) for small $\varepsilon$. The claim of ``exactly one root''
holds only under specific conditions on the region or the
perturbation parameter.

For a cleaner application: the equation $z^7 - 5z^3 + 12 = 0$.
On $\abs{z}=1$: $\abs{12} > \abs{z^7 - 5z^3} \leq 1 + 5 = 6$, so $0$~zeros
in $\abs{z}<1$. On $\abs{z}=2$: $\abs{-5z^3}=40$ and $\abs{z^7+12} \leq 128+12=140$.
This is not sharp. Instead: $\abs{z^7}=128$ vs $\abs{-5z^3+12}\leq 40+12=52<128$,
giving $7$~zeros in $\abs{z}<2$. So $7$ zeros in $1\leq\abs{z}<2$.
\end{example}

\section{Hurwitz's Theorem}\label{sec:hurwitz}

\begin{theorem}[Hurwitz's Theorem]\label{thm:hurwitz}
\index{Hurwitz's theorem}
Let $(f_n)$ be a sequence of holomorphic functions on a domain $\Omega$
converging uniformly on compact subsets to $f$. If $f \not\equiv 0$
and $f(z_0) = 0$ for some $z_0 \in \Omega$, then for every sufficiently
small $r > 0$ such that $\overline{D(z_0, r)} \subseteq \Omega$ and $f$
has no other zeros in $\overline{D(z_0, r)}$, there exists $N$ such
that for all $n \geq N$, $f_n$ has exactly as many zeros in $D(z_0, r)$
as $f$ does (i.e., the order of the zero of $f$ at $z_0$).

More precisely, if $f \not\equiv 0$, then either:
\begin{enumerate}[label=(\roman*)]
  \item every $z_0 \in \Omega$ with $f(z_0) = 0$ attracts zeros of $f_n$
    with the correct multiplicity, or
  \item $f$ has no zeros in $\Omega$ and, for $n$ large enough, neither
    does $f_n$.
\end{enumerate}
\end{theorem}

\begin{proof}
Choose $r > 0$ small enough that $\overline{D(z_0,r)} \subseteq \Omega$
and $f(z) \neq 0$ for $0 < \abs{z - z_0} \leq r$. This is possible by
the isolation of zeros. Let
\[
  \delta = \min_{\abs{z-z_0}=r} \abs{f(z)} > 0.
\]
Since $f_n \to f$ uniformly on the circle $\abs{z-z_0} = r$, for $n$
large enough we have $\abs{f_n(z) - f(z)} < \delta \leq \abs{f(z)}$ on
this circle. By Rouch\'{e}'s theorem (Theorem~\ref{thm:rouche}), $f_n = f + (f_n - f)$
has the same number of zeros in $D(z_0, r)$ as $f$, namely the order of
the zero of $f$ at $z_0$.
\end{proof}

\begin{corollary}[Hurwitz — injective limit]\label{cor:hurwitz-injective}
\index{Hurwitz's theorem!injective limit}
If $(f_n)$ is a sequence of injective holomorphic functions on a domain $\Omega$
converging uniformly on compact subsets to $f$, then $f$ is either injective
or constant.
\end{corollary}

\begin{proof}
Suppose $f$ is not constant and $f(a) = f(b)$ for distinct $a, b \in \Omega$.
Set $c = f(a) = f(b)$ and apply Hurwitz's theorem to $g_n(z) = f_n(z) - c$
converging to $g(z) = f(z) - c$. Since $g(a) = 0$, for large $n$,
$g_n$ has a zero $a_n$ near $a$ (with $a_n \to a$). Similarly, $g_n$ has
a zero $b_n$ near $b$ (with $b_n \to b$). For large $n$, $a_n \neq b_n$,
so $f_n(a_n) = f_n(b_n) = c$, contradicting injectivity of $f_n$.
\end{proof}

\section{The Open Mapping Theorem}\label{sec:open-mapping}

\begin{theorem}[Open Mapping Theorem]\label{thm:open-mapping}
\index{open mapping theorem}
If $f$ is a non-constant holomorphic function on a domain $\Omega$,
then $f$ is an \textbf{open map}: for every open set $U \subseteq \Omega$,
the image $f(U)$ is open in $\C$.
\end{theorem}

\begin{proof}
Let $w_0 = f(z_0)$ for some $z_0 \in U \subseteq \Omega$. We must show
that $w_0$ is an interior point of $f(U)$, i.e., some disc centred at
$w_0$ is contained in $f(U)$.

Since $f$ is non-constant, the function $f(z) - w_0$ has an isolated zero
at $z_0$ (of some order $m \geq 1$). Choose $r > 0$ such that
$\overline{D(z_0, r)} \subseteq U$ and $f(z) - w_0 \neq 0$ for
$0 < \abs{z - z_0} \leq r$. Set
\[
  \delta = \min_{\abs{z - z_0} = r} \abs{f(z) - w_0} > 0.
\]

Now let $w$ be any point with $\abs{w - w_0} < \delta$. On the circle
$\abs{z - z_0} = r$:
\[
  \abs{(f(z) - w) - (f(z) - w_0)} = \abs{w - w_0} < \delta
  \leq \abs{f(z) - w_0}.
\]
By Rouch\'{e}'s theorem, $f(z) - w$ has the same number of zeros in
$D(z_0, r)$ as $f(z) - w_0$, namely $m \geq 1$. In particular,
$f(z) = w$ for some $z \in D(z_0, r) \subseteq U$, so $w \in f(U)$.

This shows $D(w_0, \delta) \subseteq f(U)$, hence $f(U)$ is open.
\end{proof}

\begin{corollary}[Maximum modulus principle — alternate proof]\label{cor:max-mod-open-map}
\index{maximum modulus principle}
If $f$ is holomorphic and non-constant on a domain $\Omega$, then
$\abs{f}$ has no local maximum in $\Omega$.
\end{corollary}

\begin{proof}
If $\abs{f}$ had a local maximum at $z_0 \in \Omega$, then
$f(z_0)$ would be in the boundary of $f(U)$ for any small neighbourhood
$U$ of $z_0$. But by the open mapping theorem, $f(U)$ is open, so
$f(z_0)$ is an interior point — a contradiction.
\end{proof}

\begin{corollary}[Inverse function theorem for holomorphic functions]\label{cor:inverse-function}
\index{inverse function theorem}
If $f$ is holomorphic on $\Omega$ with $f'(z_0) \neq 0$, then $f$ is
a local biholomorphism near $z_0$: there exist neighbourhoods $U$ of $z_0$
and $V$ of $f(z_0)$ such that $f\colon U \to V$ is bijective and
$f^{-1}\colon V \to U$ is holomorphic.
\end{corollary}

\begin{proof}
Since $f'(z_0) \neq 0$, the zero of $f(z) - f(z_0)$ at $z_0$ is simple
(order $m=1$). Choose $r>0$ as in the proof of the open mapping theorem.
Then for $w$ near $f(z_0)$, the equation $f(z) = w$ has exactly $1$
solution in $D(z_0,r)$, so $f$ is injective on some neighbourhood $U$
of $z_0$. By the open mapping theorem, $V = f(U)$ is open. The inverse
$f^{-1}: V \to U$ is continuous (as $f$ is open and bijective) and
holomorphic by the formula $(f^{-1})'(w) = 1/f'(f^{-1}(w))$.
\end{proof}

\medskip

\begin{center}
\begin{tikzpicture}[scale=1.1, >=Stealth]
  % Left: domain
  \begin{scope}[xshift=-3.5cm]
    \draw[thick, fill=blue!8] (0,0) circle (1.8);
    \draw[thick, fill=blue!15] (0,0) circle (0.8);
    \node at (0,0) {$z_0$};
    \filldraw (0,0) circle (1.5pt);
    \node at (0, -1.2) {\small $D(z_0,r)$};
    \node at (0, -2.3) {$\Omega$};
  \end{scope}
  % Arrow
  \draw[->, thick, line width=1.5pt] (-1.0, 0) -- (0.6, 0)
    node[midway, above]{$f$};
  % Right: image
  \begin{scope}[xshift=3.5cm]
    \draw[thick, fill=red!8, rotate=15]
      plot[smooth cycle, tension=0.8] coordinates
      {(1.8,0.3) (0.5,1.6) (-1.2,1.2) (-1.8,-0.2)
       (-0.8,-1.5) (0.8,-1.3)};
    \draw[thick, fill=red!20] (0,0) circle (0.7);
    \node at (0,0) {$w_0$};
    \filldraw (0,0) circle (1.5pt);
    \node at (0, -1.1) {\small $D(w_0,\delta) \subseteq f(\Omega)$};
    \node at (0, -2.3) {$f(\Omega)$ open};
  \end{scope}
\end{tikzpicture}
\captionof{figure}{The open mapping theorem: $f$ maps open sets to open sets.}
\label{fig:open-mapping}
\end{center}

\section{Further Applications}\label{sec:further-apps-rouche}

\begin{proposition}[Local mapping theorem]\label{prop:local-mapping}
\index{local mapping theorem}
If $f$ is holomorphic at $z_0$ with a zero of order $m$ at $z_0$ (i.e.,
$f(z) - f(z_0) = (z-z_0)^m h(z)$ with $h(z_0) \neq 0$), then $f$ is
locally $m$-to-$1$ near $z_0$: for $w$ close to $w_0 = f(z_0)$ (but
$w \neq w_0$), the equation $f(z) = w$ has exactly $m$ solutions near $z_0$.
\end{proposition}

\begin{proof}
This follows directly from Rouch\'{e}'s theorem, as in the proof of the
open mapping theorem: for $\abs{w - w_0} < \delta$, the function
$f(z) - w$ has exactly $m$ zeros in $D(z_0,r)$, counted with multiplicity.
When $w \neq w_0$ and $\delta$ is small enough, these $m$ zeros are
distinct (as the multiplicity phenomenon occurs only at $w_0$).
\end{proof}

\begin{theorem}[Counting zeros via the argument principle]\label{thm:counting-formula}
Under the hypotheses of Theorem~\ref{thm:argument-principle}, if $f$
is holomorphic (no poles) on $\overline{D}$, then the number of
zeros in $D$ (with multiplicity) is
\[
  Z = \frac{1}{2\pi i}\oint_{\partial D} \frac{f'(z)}{f(z)}\,\dd z
  = \frac{1}{2\pi}\Delta_{\partial D}\arg f(z),
\]
where $\Delta_{\partial D}\arg f(z)$ denotes the total change in argument
of $f(z)$ as $z$ traverses $\partial D$.
\end{theorem}

\begin{example}[Zeros of $e^z - 3z$ in the unit disc]\label{ex:ez-3z}
Let $f(z) = e^z - 3z$. On $\abs{z} = 1$, set $g(z) = -3z$ and
$h(z) = e^z$. We have $\abs{g(z)} = 3$ and $\abs{h(z)} = e^{\re z} \leq e$.
Since $e \approx 2.718 < 3$, we get $\abs{h(z)} < \abs{g(z)}$ on $\abs{z}=1$.
By Rouch\'{e}, $f = g + h = e^z - 3z$ has the same number of zeros in
$\abs{z}<1$ as $g(z) = -3z$, which has exactly $1$ zero. So $e^z - 3z$
has exactly $1$ zero in the unit disc.
\end{example}

\begin{example}[Zeros of $z^4 - 5z + 1$ in the annulus $1 < \abs{z} < 2$]\label{ex:annulus-zeros}
Let $p(z) = z^4 - 5z + 1$.

On $\abs{z} = 1$: $\abs{-5z} = 5$ and $\abs{z^4 + 1} \leq 2 < 5$.
By Rouch\'{e}, $p$ has $1$ zero in $\abs{z} < 1$ (same as $-5z$).

On $\abs{z} = 2$: $\abs{z^4} = 16$ and $\abs{-5z+1} \leq 11 < 16$.
By Rouch\'{e}, $p$ has $4$ zeros in $\abs{z} < 2$ (same as $z^4$).

Therefore $p$ has $4 - 1 = 3$ zeros in the annulus $1 < \abs{z} < 2$
(after verifying $p$ has no zeros on $\abs{z}=1$: $\abs{p(z)} \geq 5-2=3>0$).
\end{example}


\section{Exercises}\label{sec:exercises-ch9}

\begin{exercise}[Zeros of $z^6 + z^3 + 9$]\label{exer:rouche-1}
Show that all zeros of $p(z) = z^6 + z^3 + 9$ lie in the annulus
$1 < \abs{z} < 2$.

\emph{Hint:} On $\abs{z}=1$, compare $9$ with $z^6+z^3$.
On $\abs{z}=2$, compare $z^6$ with $z^3+9$.
\end{exercise}

\begin{exercise}[Zeros near a simple zero]\label{exer:hurwitz-app}
Let $f_n(z) = z + 1/n + z^2/n$. Show that for $n$ large enough,
$f_n$ has exactly one zero in $\abs{z} < 1/2$, and find the limit
of this zero as $n \to \infty$.
\end{exercise}

\begin{exercise}[Fixed points of holomorphic maps]\label{exer:fixed-point}
Let $f\colon \overline{\D} \to \D$ be continuous on $\overline{\D}$
and holomorphic on $\D$. Show that $f$ has a unique fixed point in $\D$.

\emph{Hint:} Apply Rouch\'{e}'s theorem to $g(z) = z - f(z)$ and $h(z) = z$.
\end{exercise}

\begin{exercise}[Argument principle computation]\label{exer:argument-comp}
Let $f(z) = \dfrac{z^3 - 1}{z^2 + 1}$. Compute
$\displaystyle\frac{1}{2\pi i}\oint_{\abs{z}=2} \frac{f'(z)}{f(z)}\,\dd z$.

\emph{Solution sketch:} In $\abs{z}<2$, $f$ has zeros at the three cube roots
of unity (each simple) and poles at $z=\pm i$ (each simple). So
$Z - P = 3 - 2 = 1$.
\end{exercise}

\begin{exercise}[Image of the boundary]\label{exer:image-boundary}
Let $f(z) = z^2 + z$ and $\gamma$ the circle $\abs{z} = 2$ traversed
once counterclockwise. Determine $\ind(f \circ \gamma, 0)$ and
$\ind(f \circ \gamma, -1/4)$.

\emph{Hint:} The zeros of $f$ in $\abs{z}<2$ are $z=0$ and $z=-1$ (both simple),
so $\ind(f\circ\gamma,0) = 2$. For $-1/4$: solve $z^2+z = -1/4$,
i.e., $(z+1/2)^2 = 0$, double root at $z=-1/2$.
\end{exercise}

\begin{exercise}[Rouch\'{e} for exponentials]\label{exer:rouche-exp}
Show that the equation $3e^{-z} = z - 1$ has exactly one root in the
half-plane $\re(z) > 0$.

\emph{Hint:} Consider the rectangle with vertices $\pm Ri$ and $M \pm Ri$
for large $R, M$ and apply the argument principle.
\end{exercise}

\begin{exercise}[Hurwitz and univalence]\label{exer:hurwitz-unival}
Let $f_n(z) = z + z^n / n^2$ for $n \geq 1$. Show that each $f_n$ is
injective on $\D$ for $n$ large enough, and that the limit $f(z) = z$
is injective on $\D$, consistent with Corollary~\ref{cor:hurwitz-injective}.
\end{exercise}

\begin{exercise}[Open mapping and boundary behaviour]\label{exer:open-map-boundary}
Let $f$ be holomorphic and non-constant on a bounded domain $\Omega$.
Show that if $\abs{f}$ extends continuously to $\overline{\Omega}$,
then $\max_{\overline{\Omega}}\abs{f}$ is achieved on $\partial \Omega$.
Deduce the maximum modulus principle from the open mapping theorem.
\end{exercise}

\begin{exercise}[Number of preimages]\label{exer:preimage-count}
Let $f(z) = z^3 - 3z$. Using the argument principle, find the number of
solutions of $f(z) = w$ in $\abs{z} < 2$ for: (a) $w = 0$; (b) $w = 2$;
(c) $w = 10$.
\end{exercise}

\begin{exercise}[Generalised Rouch\'{e}]\label{exer:gen-rouche}
Prove the symmetric version of Rouch\'{e}'s theorem
(Remark~\ref{rem:rouche-symmetric}):
if $f$ and $g$ are holomorphic on $\overline{D}$ with no zeros on $\partial D$
and $\abs{f(z) - g(z)} < \abs{f(z)} + \abs{g(z)}$ on $\partial D$,
then $f$ and $g$ have the same number of zeros in $D$.

\emph{Hint:} Consider $h_t(z) = (1-t)f(z) + tg(z)$ and show $h_t$
has no zeros on $\partial D$ for $t \in [0,1]$.
\end{exercise}


%% ═══════════════════════════════════════════════════════════════════
%%  CHAPTER 10 : Introduction to Entire and Meromorphic Functions
%% ═══════════════════════════════════════════════════════════════════
\chapter{Introduction to Entire and Meromorphic Functions}\label{ch:entire-meromorphic}

\section{Entire Functions}\label{sec:entire-functions}

\begin{definition}[Entire function]\label{def:entire-function}
\index{entire function}
A function $f\colon \C \to \C$ is called \textbf{entire} if it is
holomorphic on all of $\C$. Equivalently, $f$ has a power series
expansion $f(z) = \sum_{n=0}^{\infty} a_n z^n$ with infinite radius
of convergence.
\end{definition}

\begin{example}[Classical entire functions]\label{ex:classical-entire}
The following are entire functions:
\begin{enumerate}[label=(\roman*)]
  \item \textbf{Polynomials:} $p(z) = a_n z^n + \cdots + a_0$.
  \item \textbf{Exponential:} $e^z = \sum_{n=0}^{\infty} z^n/n!$.
  \item \textbf{Trigonometric:} $\sin z = \sum_{n=0}^{\infty} (-1)^n z^{2n+1}/(2n+1)!$
    and $\cos z = \sum_{n=0}^{\infty} (-1)^n z^{2n}/(2n)!$.
  \item \textbf{Hyperbolic:} $\sinh z$ and $\cosh z$.
  \item \textbf{Compositions:} $e^{e^z}$, $\sin(z^2)$, $e^{-z^2}$, etc.
\end{enumerate}
Note that $1/z$, $\tan z$, and $\Gamma(z)$ are \emph{not} entire
(they have singularities in $\C$).
\end{example}

\begin{definition}[Transcendental entire function]\label{def:transcendental-entire}
\index{transcendental entire function}
An entire function that is not a polynomial is called \textbf{transcendental}.
Equivalently, $f$ is transcendental entire if and only if its Taylor series
has infinitely many nonzero terms, i.e., $z = \infty$ is an essential singularity.
\end{definition}

\begin{remark}[Polynomial vs.\ transcendental]\label{rem:poly-vs-trans}
By the Fundamental Theorem of Algebra, a polynomial of degree $n$ has
exactly $n$ zeros (with multiplicity). A transcendental entire function
can have any number of zeros: finitely many (like $e^z + 1$, with
zeros at $z = (2k+1)\pi i$... actually infinitely many), or none
(like $e^z$), or infinitely many (like $\sin z$).
In fact, the only entire functions with no zeros are those of the form
$e^{g(z)}$ for some entire $g$ (since $\C$ is simply connected).
\end{remark}

\section{Liouville's Theorem Revisited}\label{sec:liouville-revisited}

We recall the fundamental theorem constraining bounded entire functions.

\begin{theorem}[Liouville's Theorem]\label{thm:liouville-revisited}
\index{Liouville's theorem}
Every bounded entire function is constant.
\end{theorem}

\begin{proof}
Let $\abs{f(z)} \leq M$ for all $z \in \C$. By Cauchy's inequality
applied to the disc $D(0, R)$:
\[
  \abs{f^{(n)}(0)} \leq \frac{n!\, M}{R^n}
  \quad \text{for all } R > 0.
\]
Letting $R \to \infty$ gives $f^{(n)}(0) = 0$ for all $n \geq 1$.
Hence $f(z) = f(0)$ is constant.
\end{proof}

\begin{theorem}[Generalised Liouville]\label{thm:gen-liouville}
\index{Liouville's theorem!generalised}
If $f$ is entire and $\abs{f(z)} \leq C(1 + \abs{z})^k$ for some
constants $C > 0$ and $k \geq 0$, then $f$ is a polynomial of degree
at most $\lfloor k \rfloor$.
\end{theorem}

\begin{proof}
By Cauchy's inequality for the $n$-th derivative on $D(0,R)$:
\[
  \abs{f^{(n)}(0)} \leq \frac{n!}{R^n}
  \max_{\abs{z}=R}\abs{f(z)}
  \leq \frac{n!\, C(1+R)^k}{R^n}.
\]
For $n > k$, as $R \to \infty$ the right-hand side tends to $0$,
so $f^{(n)}(0) = 0$ for all $n > k$. The Taylor series terminates:
$f$ is a polynomial of degree at most $\lfloor k \rfloor$.
\end{proof}

\begin{corollary}[Entire functions with polynomial growth]\label{cor:poly-growth}
If $f$ is an entire function satisfying $\abs{f(z)} = o(\abs{z}^n)$
as $\abs{z}\to\infty$ for some positive integer $n$, then $f$ is a
polynomial of degree at most $n-1$.
\end{corollary}

\section{Picard's Theorems}\label{sec:picard}

The little Picard theorem dramatically strengthens Liouville's theorem.

\begin{theorem}[Little Picard Theorem]\label{thm:little-picard}
\index{Picard's theorem!little}
Every non-constant entire function takes every complex value with at
most one exception. That is, if $f$ is entire and non-constant, then
$f(\C) \supseteq \C \setminus \{a\}$ for at most one point $a \in \C$.
\end{theorem}

The proof of Picard's theorem requires the theory of modular functions
and covering spaces, which is beyond the scope of this introductory
course. We state it without proof but illustrate its content.

\begin{example}[Picard's theorem in action]\label{ex:picard-examples}
\begin{enumerate}[label=(\roman*)]
  \item $e^z$ omits exactly one value: $0$. It takes every other value
    $w \neq 0$, since $e^z = w$ has solutions $z = \Log w + 2\pi i k$.
  \item $\sin z$ takes \emph{every} complex value (no exceptions).
    Indeed, for any $w \in \C$, $\sin z = w$ is equivalent to
    $e^{iz} - 2iw e^{-iz} = e^{-iz}$, which always has solutions.
  \item A polynomial of degree $n \geq 1$ takes every value (by the FTA),
    so there are zero exceptions.
  \item There is no entire function omitting two or more values.
    For instance, there is no entire function with range $\C \setminus \{0, 1\}$.
\end{enumerate}
\end{example}

\begin{center}
\begin{tikzpicture}[scale=0.9, >=Stealth]
  % Picard illustration for e^z
  \begin{scope}[xshift=-4.5cm]
    \draw[->] (-2.5, 0) -- (2.5, 0) node[right]{\small $\re$};
    \draw[->] (0, -2.2) -- (0, 2.5) node[above]{\small $\im$};
    % Horizontal strip
    \fill[blue!12] (-2.3, -1.8) rectangle (2.3, 1.8);
    \draw[dashed, blue!50] (-2.3, 1.8) -- (2.3, 1.8);
    \draw[dashed, blue!50] (-2.3, -1.8) -- (2.3, -1.8);
    \node at (0, -2.7) {\small Domain: strip $-\pi < \im(z) \leq \pi$};
    \node at (0, -3.3) {\small $z$-plane};
  \end{scope}
  % Arrow
  \draw[->, thick, line width=1.5pt] (-1.2, 0) -- (0.8, 0)
    node[midway, above]{\small $e^z$};
  % Image
  \begin{scope}[xshift=4.5cm]
    \draw[->] (-2.5, 0) -- (2.5, 0) node[right]{\small $\re$};
    \draw[->] (0, -2.2) -- (0, 2.5) node[above]{\small $\im$};
    \fill[red!12] (-2.3,-2.1) rectangle (2.3, 2.3);
    \fill[white] (0,0) circle (3pt);
    \draw (0,0) circle (3pt);
    \node[below right] at (0.15, -0.15) {\small $0$ (omitted)};
    \node at (0, -2.7) {\small Image: $\C \setminus \{0\}$};
    \node at (0, -3.3) {\small $w$-plane};
  \end{scope}
\end{tikzpicture}
\captionof{figure}{Little Picard: $e^z$ omits exactly one value ($0$).
Every other value is attained.}
\label{fig:picard-exp}
\end{center}

\begin{theorem}[Great Picard Theorem]\label{thm:great-picard}
\index{Picard's theorem!great}
If $f$ has an essential singularity at $z_0$, then in every
punctured neighbourhood of $z_0$, $f$ takes every complex value
with at most one exception.
\end{theorem}

This generalises the Casorati--Weierstrass theorem
(which only gives density of the image) and will not be proved here.

\section{Order of Growth}\label{sec:order-growth}

\begin{definition}[Order of an entire function]\label{def:order-entire}
\index{order of an entire function}
The \textbf{order} of an entire function $f$ is
\[
  \rho = \rho(f) = \limsup_{r \to \infty}
  \frac{\log \log M(r)}{\log r},
\]
where $M(r) = \max_{\abs{z}=r}\abs{f(z)}$ is the \textbf{maximum modulus}
\index{maximum modulus} of $f$.
\end{definition}

\begin{remark}[Interpretation of order]\label{rem:order-interpretation}
An entire function has order $\rho$ roughly if $M(r) \sim e^{r^\rho}$.
More precisely, $\rho$ is the infimum of $\alpha > 0$ such that
$M(r) \leq e^{r^\alpha}$ for all sufficiently large $r$.
\begin{itemize}
  \item Polynomials have order $\rho = 0$.
  \item $e^z$ has order $\rho = 1$ (since $M(r) = e^r$).
  \item $e^{z^2}$ has order $\rho = 2$.
  \item $e^{e^z}$ has infinite order.
  \item $\sin z$ and $\cos z$ have order $\rho = 1$
    (since $M(r) \sim e^r/2$).
\end{itemize}
\end{remark}

\begin{proposition}[Order from Taylor coefficients]\label{prop:order-coefficients}
\index{order!from Taylor coefficients}
If $f(z) = \sum_{n=0}^{\infty} a_n z^n$ is entire and of finite order, then
\[
  \rho = \limsup_{n \to \infty} \frac{n \log n}{-\log \abs{a_n}}.
\]
\end{proposition}

\begin{proof}
This follows from a careful comparison of the maximum modulus with the
coefficient estimates. For the upper bound, from $M(r) \leq \sum \abs{a_n}r^n$
and optimising in $r$, one obtains the formula. For the lower bound,
$\abs{a_n} \leq M(r)/r^n$ (Cauchy inequality), and choosing
$r = n^{1/\rho}$ leads to the reverse inequality. We refer to
Boas~\cite{boas1954} or Conway~\cite{conway1978} for the complete argument.
\end{proof}

\begin{definition}[Type of an entire function]\label{def:type-entire}
\index{type of an entire function}
If $f$ has finite order $\rho$, the \textbf{type} of $f$ is
\[
  \sigma = \limsup_{r \to \infty} \frac{\log M(r)}{r^\rho}.
\]
We say $f$ is of \textbf{minimal type} if $\sigma = 0$, \textbf{normal type}
if $0 < \sigma < \infty$, and \textbf{maximal type} if $\sigma = \infty$.
\end{definition}

\begin{example}[Orders and types]\label{ex:order-type}
\begin{enumerate}[label=(\roman*)]
  \item $e^z$: order $1$, type $1$ (normal).
  \item $e^{3z^2}$: order $2$, type $3$.
  \item $\sin z$: order $1$, type $1$.
  \item $e^{z^{1/2}}$ (defined on an appropriate domain, or
    its entire extension): order $1/2$, type $1$.
\end{enumerate}
\end{example}

\section{Weierstrass Products}\label{sec:weierstrass-products}

A fundamental question in function theory is: given a sequence of complex
numbers (the desired zero set), can we construct an entire function
having exactly those zeros? The Weierstrass factorisation theorem answers
this affirmatively.

\begin{definition}[Elementary factors]\label{def:elementary-factors}
\index{elementary factors}\index{Weierstrass!elementary factors}
The \textbf{elementary factors} are defined by
\[
  E_0(z) = 1 - z, \qquad
  E_p(z) = (1-z)\exp\!\Bigl(z + \frac{z^2}{2} + \cdots + \frac{z^p}{p}\Bigr)
  \quad \text{for } p \geq 1.
\]
\end{definition}

\begin{lemma}[Properties of elementary factors]\label{lem:elementary-factors-props}
For $\abs{z} \leq 1$ and $p \geq 0$:
\begin{equation}\label{eq:elementary-bound}
  \abs{1 - E_p(z)} \leq \abs{z}^{p+1}.
\end{equation}
In particular, $E_p(z)$ vanishes only at $z = 1$, with a simple zero.
\end{lemma}

\begin{proof}
Write $\log E_p(z) = \log(1-z) + z + z^2/2 + \cdots + z^p/p
= -\sum_{n=p+1}^{\infty} z^n/n$ for $\abs{z} < 1$. Then
$\abs{\log E_p(z)} \leq \sum_{n=p+1}^{\infty} \abs{z}^n/n
\leq \abs{z}^{p+1}\sum_{n=0}^{\infty}\abs{z}^n/(n+p+1)
\leq \abs{z}^{p+1}/(1-\abs{z})$ for $\abs{z}<1$.
The bound $\abs{1 - E_p(z)} \leq \abs{z}^{p+1}$ follows from
$\abs{e^w - 1} \leq \abs{w}e^{\abs{w}}$ and a more careful estimate
(see, e.g., Stein--Shakarchi \cite{stein2003}).
\end{proof}

\begin{theorem}[Weierstrass Factorisation Theorem]\label{thm:weierstrass-factorisation}
\index{Weierstrass!factorisation theorem}
Let $f$ be an entire function with $f(0) \neq 0$, and let
$a_1, a_2, \ldots$ be the nonzero zeros of $f$ (listed with multiplicity),
with $0 < \abs{a_1} \leq \abs{a_2} \leq \cdots$ and $\abs{a_n} \to \infty$.
Then there exist non-negative integers $p_n$ and an entire function $g$
such that
\begin{equation}\label{eq:weierstrass-product}
  f(z) = e^{g(z)} z^m \prod_{n=1}^{\infty} E_{p_n}\!\Bigl(\frac{z}{a_n}\Bigr),
\end{equation}
where $m = \ord_0(f)$ is the order of the zero at the origin (possibly $0$).

If the $p_n$ are chosen so that $\sum_n (r/\abs{a_n})^{p_n+1}$ converges
for all $r > 0$, the product converges uniformly on compact sets
and defines an entire function with the prescribed zeros.
\end{theorem}

The proof, which we omit, proceeds by showing that the product
$\prod E_{p_n}(z/a_n)$ converges, and then the entire function $g$
is determined by $f / \bigl(z^m \prod E_{p_n}(z/a_n)\bigr)$, which
is entire and zero-free, hence of the form $e^{g(z)}$.

\begin{definition}[Genus and canonical product]\label{def:genus}
\index{genus}
If the integers $p_n$ in \eqref{eq:weierstrass-product} can all be
taken equal to some fixed $p$, and this is the smallest such $p$, then
$p$ is called the \textbf{genus} of the canonical product. The
\textbf{genus of $f$} is defined as the maximum of $p$ and the degree of
$g$ (if $g$ is a polynomial).
\end{definition}

\subsection{Factorisation of $\sin z$}

\begin{theorem}[Product formula for sine]\label{thm:sin-product}
\index{sine!product formula}
\begin{equation}\label{eq:sin-product}
  \sin(\pi z) = \pi z \prod_{n=1}^{\infty}
  \Bigl(1 - \frac{z^2}{n^2}\Bigr).
\end{equation}
\end{theorem}

\begin{proof}
The function $\sin(\pi z)$ is entire with simple zeros at $z = n$ for
$n \in \Z$. We aim to write
\[
  \sin(\pi z) = \pi z \, e^{g(z)} \prod_{n=1}^{\infty}
  E_1\!\Bigl(\frac{z}{n}\Bigr) \cdot E_1\!\Bigl(\frac{z}{-n}\Bigr).
\]
Now, using $E_1(z) = (1-z)e^z$:
\[
  E_1\!\Bigl(\frac{z}{n}\Bigr) \cdot E_1\!\Bigl(\frac{z}{-n}\Bigr)
  = \Bigl(1 - \frac{z}{n}\Bigr)e^{z/n} \cdot
  \Bigl(1 + \frac{z}{n}\Bigr)e^{-z/n}
  = 1 - \frac{z^2}{n^2}.
\]
Since $\sum_{n=1}^{\infty} 1/n^2$ converges (equals $\pi^2/6$),
the product $\prod_{n=1}^{\infty}(1 - z^2/n^2)$ converges uniformly
on compact sets.

Denoting $P(z) = \pi z \prod_{n=1}^{\infty}(1-z^2/n^2)$, we must show
$g(z) = 0$. The quotient $h(z) = \sin(\pi z)/P(z)$ is entire and
nonvanishing. To show $h \equiv 1$, we use:

\textbf{Boundedness argument.} Consider the logarithmic derivative:
\[
  \frac{\pi\cos(\pi z)}{\sin(\pi z)}
  = \frac{1}{z} + \sum_{n=1}^{\infty}\Bigl(\frac{1}{z-n}+\frac{1}{z+n}\Bigr)
  = \frac{1}{z} + \sum_{n=1}^{\infty}\frac{2z}{z^2-n^2}.
\]
This is the partial fraction expansion of $\pi\cot(\pi z)$.
Similarly, the logarithmic derivative of $P(z)$ is
\[
  \frac{P'(z)}{P(z)} = \frac{1}{z} + \sum_{n=1}^{\infty}\frac{2z}{z^2-n^2}.
\]
Thus $h'(z)/h(z) = g'(z) = 0$, so $g$ is constant.
Evaluating at $z=0$:
$\lim_{z\to 0}\frac{\sin(\pi z)}{\pi z} = 1$ and
$\lim_{z\to 0} \prod(1-z^2/n^2) = 1$, so $e^{g(0)} = 1$,
giving $g \equiv 0$.
\end{proof}

\begin{center}
\begin{tikzpicture}[scale=0.7, >=Stealth]
  \draw[->] (-5.5, 0) -- (5.5, 0) node[right]{$\re$};
  \draw[->] (0, -3) -- (0, 3.5) node[above]{$\im$};
  % Zeros of sin(pi z) on the real axis
  \foreach \n in {-5,-4,-3,-2,-1,0,1,2,3,4,5} {
    \filldraw[red!70!black] (\n, 0) circle (3pt);
  }
  % Plot sin(pi x) for real x
  \draw[thick, blue!70!black, domain=-5.2:5.2, samples=200]
    plot (\x, {sin(180*\x)/1.2});
  \node at (0, -3.8) {Zeros of $\sin(\pi z)$: $z = n,\ n \in \Z$};
  \node[red!70!black] at (3.5, -1.5) {\small $\bullet$ = zeros};
  % Labels
  \node[below] at (1, -0.2) {\small $1$};
  \node[below] at (2, -0.2) {\small $2$};
  \node[below] at (-1, -0.2) {\small $-1$};
  \node[below] at (-2, -0.2) {\small $-2$};
\end{tikzpicture}
\captionof{figure}{The zeros of $\sin(\pi z)$ are precisely the integers.
The Weierstrass product reflects this by having factors $(1-z^2/n^2)$.}
\label{fig:zeros-sin}
\end{center}

\begin{corollary}[Euler's identity]\label{cor:euler-identity}
\[
  \sum_{n=1}^{\infty} \frac{1}{n^2} = \frac{\pi^2}{6}.
\]
\end{corollary}

\begin{proof}
Expand both sides of \eqref{eq:sin-product} in Taylor series about $z=0$.
The left-hand side gives
\[
  \sin(\pi z) = \pi z - \frac{(\pi z)^3}{3!} + \cdots
  = \pi z\Bigl(1 - \frac{\pi^2 z^2}{6} + \cdots\Bigr).
\]
The right-hand side gives
\[
  \pi z \prod_{n=1}^{\infty}\Bigl(1 - \frac{z^2}{n^2}\Bigr)
  = \pi z \Bigl(1 - z^2\sum_{n=1}^{\infty}\frac{1}{n^2} + \cdots\Bigr).
\]
Comparing the coefficients of $z^3$:
$-\pi^2/6 = -\sum_{n=1}^{\infty}1/n^2$, giving the result.
\end{proof}


\section{Meromorphic Functions and Partial Fractions}\label{sec:meromorphic-partial}

\begin{definition}[Meromorphic function on $\C$]\label{def:meromorphic-C}
\index{meromorphic function!on $\C$}
A function $f$ is \textbf{meromorphic on $\C$} if it is holomorphic
on $\C \setminus S$, where $S$ is a discrete subset of $\C$ consisting
entirely of poles of $f$.
\end{definition}

\begin{remark}[Field of meromorphic functions]\label{rem:meromorphic-field}
\index{meromorphic function!field of}
The set $\mathcal{M}(\C)$ of meromorphic functions on $\C$ forms a
\textbf{field} under the usual operations of addition and multiplication.
The entire functions $\mathcal{O}(\C)$ form an integral domain within
this field. Every meromorphic function on $\C$ can be written as a
quotient of two entire functions: $f = g/h$ with $g, h$ entire and
$h \not\equiv 0$.
\end{remark}

The Mittag-Leffler theorem is the analogue of the Weierstrass factorisation
theorem: whereas Weierstrass prescribes \emph{zeros}, Mittag-Leffler
prescribes \emph{principal parts} at poles.

\begin{theorem}[Mittag-Leffler Theorem]\label{thm:mittag-leffler}
\index{Mittag-Leffler theorem}
Let $(b_n)_{n \geq 1}$ be a sequence of distinct complex numbers with
$\abs{b_n} \to \infty$, and let $p_n(z)$ be polynomials in $1/(z-b_n)$
without constant term (\textbf{principal parts}). Then there exists a
meromorphic function $f$ on $\C$ whose poles are exactly the $b_n$
and whose principal part at each $b_n$ is $p_n(z)$.

Moreover, $f$ can be written as
\[
  f(z) = \sum_{n=1}^{\infty}\bigl[p_n(z) - q_n(z)\bigr] + g(z),
\]
where $q_n(z)$ are polynomials chosen to ensure convergence and $g$ is entire.
The function is unique up to addition of an entire function.
\end{theorem}

\begin{example}[Partial fraction of $\pi\cot(\pi z)$]\label{ex:cot-partial}
\index{cotangent!partial fraction}
The function $\pi\cot(\pi z)$ is meromorphic with simple poles at
$z = n$ ($n \in \Z$), each with residue $1$. The Mittag-Leffler expansion gives:
\begin{equation}\label{eq:cot-expansion}
  \pi\cot(\pi z)
  = \frac{1}{z} + \sum_{n=1}^{\infty}\Bigl(\frac{1}{z-n}+\frac{1}{z+n}\Bigr)
  = \frac{1}{z} + \sum_{n=1}^{\infty}\frac{2z}{z^2-n^2}.
\end{equation}
Here the ``correction polynomials'' $q_n$ are the constants $1/n$ and $-1/n$
combined in pairs to achieve convergence.
\end{example}

\begin{example}[Partial fraction of $\pi/\sin(\pi z)$]\label{ex:csc-partial}
Similarly, $\pi/\sin(\pi z)$ has simple poles at $z = n$ with residues
$(-1)^n$. Its Mittag-Leffler expansion is:
\[
  \frac{\pi}{\sin(\pi z)} = \frac{1}{z}
  + \sum_{n=1}^{\infty}(-1)^n\Bigl(\frac{1}{z-n}+\frac{1}{z+n}\Bigr)
  = \frac{1}{z}
  + \sum_{n=1}^{\infty}\frac{(-1)^n \cdot 2z}{z^2-n^2}.
\]
\end{example}

\begin{example}[Partial fraction of $1/\sin^2(\pi z)$]\label{ex:csc2-partial}
The function $\pi^2/\sin^2(\pi z)$ has double poles at each integer.
The principal part at $z = n$ is $1/(z-n)^2$. Its expansion is:
\[
  \frac{\pi^2}{\sin^2(\pi z)} = \sum_{n=-\infty}^{\infty} \frac{1}{(z-n)^2}.
\]
This series converges absolutely and uniformly on compact subsets
of $\C \setminus \Z$ (since $\sum 1/n^2 < \infty$, no correction
polynomials are needed).
\end{example}

\section{Introduction to Elliptic Functions}\label{sec:elliptic-intro}

We give a brief overview of elliptic functions, one of the
crowning achievements of 19th-century complex analysis.

\begin{definition}[Elliptic function]\label{def:elliptic-function}
\index{elliptic function}
An \textbf{elliptic function} is a meromorphic function
$f\colon \C \to \C \cup \{\infty\}$ that is doubly periodic:
there exist $\omega_1, \omega_2 \in \C$, linearly independent over $\R$, such that
\[
  f(z + \omega_1) = f(z), \qquad f(z + \omega_2) = f(z)
  \quad \text{for all } z \in \C.
\]
The \textbf{period lattice} is $\Lambda = \Z\omega_1 + \Z\omega_2$, and
the \textbf{fundamental period parallelogram} (or \textbf{fundamental domain})
is
\[
  \Pi = \{t_1 \omega_1 + t_2 \omega_2 : 0 \leq t_1, t_2 < 1\}.
\]
\end{definition}

\begin{center}
\begin{tikzpicture}[scale=1.0, >=Stealth]
  \def\wa{2.5} \def\wai{0.5}  % omega_1 = (2.5, 0.5)
  \def\wb{0.8} \def\wbi{2.0}  % omega_2 = (0.8, 2.0)

  % Draw lattice points and parallelogram
  \foreach \m in {-1,0,1,2} {
    \foreach \n in {-1,0,1} {
      \pgfmathsetmacro{\px}{\m*\wa + \n*\wb}
      \pgfmathsetmacro{\py}{\m*\wai + \n*\wbi}
      \filldraw[gray!60] (\px, \py) circle (2pt);
    }
  }

  % Fundamental domain
  \fill[blue!15, opacity=0.6]
    (0,0) -- (\wa, \wai) -- ({\wa+\wb}, {\wai+\wbi}) -- (\wb, \wbi) -- cycle;
  \draw[thick, blue!70!black]
    (0,0) -- (\wa, \wai) -- ({\wa+\wb}, {\wai+\wbi}) -- (\wb, \wbi) -- cycle;

  % Labels
  \filldraw (0,0) circle (2.5pt) node[below left]{$0$};
  \draw[->, thick, red!70!black] (0,0) -- (\wa, \wai) node[below right]{$\omega_1$};
  \draw[->, thick, red!70!black] (0,0) -- (\wb, \wbi) node[above left]{$\omega_2$};
  \node[blue!60!black] at ({(\wa+\wb)/2}, {(\wai+\wbi)/2}) {$\Pi$};

  \node at ({(\wa+\wb)/2}, -1.2) {Period lattice $\Lambda = \Z\omega_1 + \Z\omega_2$};
\end{tikzpicture}
\captionof{figure}{The period lattice and fundamental parallelogram
of an elliptic function.}
\label{fig:lattice}
\end{center}

\begin{theorem}[Basic properties of elliptic functions]\label{thm:elliptic-basic}
\index{elliptic function!basic properties}
Let $f$ be an elliptic function with period lattice $\Lambda$.
\begin{enumerate}[label=(\roman*)]
  \item If $f$ is entire (i.e., has no poles), then $f$ is constant.
  \item The sum of the residues of $f$ in any fundamental parallelogram
    (whose boundary contains no poles) is zero.
  \item The number of zeros equals the number of poles in any fundamental
    parallelogram (counted with multiplicity). This common number is
    called the \textbf{order}\index{order!of elliptic function} of $f$.
  \item A non-constant elliptic function has order at least $2$.
\end{enumerate}
\end{theorem}

\begin{proof}
\begin{enumerate}[label=(\roman*)]
  \item An entire elliptic function is bounded on $\Pi$ (continuous on
    compact $\overline{\Pi}$), hence bounded on all of $\C$ by periodicity.
    By Liouville's theorem, $f$ is constant.
  \item By the residue theorem, the sum of residues in $\Pi$ equals
    $\frac{1}{2\pi i}\oint_{\partial \Pi} f(z)\,\dd z$. By periodicity,
    the integrals over opposite sides cancel, giving $0$.
  \item Apply the argument principle to $f$ over $\partial \Pi$.
    Again, opposite sides cancel, giving $Z - P = 0$.
  \item If $f$ had order $1$, it would have a single simple pole in $\Pi$,
    with residue (by (ii)) equal to $0$ — contradicting simplicity. \qedhere
\end{enumerate}
\end{proof}

\begin{definition}[Weierstrass $\wp$-function]\label{def:weierstrass-pe}
\index{Weierstrass pe-function@Weierstrass $\wp$-function}
The \textbf{Weierstrass $\wp$-function} associated to a lattice
$\Lambda$ is defined by
\[
  \wp(z) = \wp(z; \Lambda)
  = \frac{1}{z^2} + \sum_{\substack{\omega \in \Lambda \\ \omega \neq 0}}
  \Bigl(\frac{1}{(z-\omega)^2} - \frac{1}{\omega^2}\Bigr).
\]
\end{definition}

\begin{proposition}[Properties of $\wp$]\label{prop:pe-properties}
\begin{enumerate}[label=(\roman*)]
  \item $\wp$ is an elliptic function of order $2$ with period lattice $\Lambda$.
  \item $\wp$ has a double pole at each lattice point and no other poles.
  \item $\wp$ is an even function: $\wp(-z) = \wp(z)$.
  \item $\wp$ satisfies the differential equation
    \begin{equation}\label{eq:pe-ode}
      (\wp')^2 = 4\wp^3 - g_2 \wp - g_3,
    \end{equation}
    where $g_2 = 60\sum_{\omega \neq 0}\omega^{-4}$ and
    $g_3 = 140\sum_{\omega \neq 0}\omega^{-6}$ are the
    \textbf{Eisenstein series}\index{Eisenstein series}.
\end{enumerate}
\end{proposition}

\begin{proof}[Proof sketch]
Convergence of the series defining $\wp$ follows from
$\sum_{\omega \neq 0}\abs{\omega}^{-3} < \infty$ (which holds for
any lattice in $\C$). Periodicity: $\wp(z+\omega_k) - \wp(z)$ is
entire and bounded, hence constant by Liouville; checking at $z = -\omega_k/2$
shows the constant is $0$. The differential equation is obtained by
computing the Laurent expansion of both sides about $z = 0$ and verifying
that the difference is an entire elliptic function vanishing at $z=0$,
hence identically zero.
\end{proof}

\begin{remark}[Connection to elliptic curves]\label{rem:elliptic-curves}
\index{elliptic curve}
The differential equation \eqref{eq:pe-ode} defines an \textbf{elliptic curve}
$y^2 = 4x^3 - g_2 x - g_3$ in the $(x,y)$-plane. The map
$z \mapsto (\wp(z), \wp'(z))$ establishes an isomorphism between
the complex torus $\C/\Lambda$ and (the projective completion of)
this elliptic curve. This profound connection between complex analysis,
algebraic geometry, and number theory is one of the most beautiful
themes in modern mathematics.
\end{remark}


\section{The Hadamard Factorisation Theorem}\label{sec:hadamard}

For entire functions of finite order, the Weierstrass product can be
made more precise.

\begin{theorem}[Hadamard Factorisation Theorem]\label{thm:hadamard}
\index{Hadamard factorisation theorem}
Let $f$ be an entire function of finite order $\rho$, with zeros
$a_1, a_2, \ldots$ (listed with multiplicity, $a_n \neq 0$) and a
zero of order $m$ at the origin. Set $p = \lfloor \rho \rfloor$. Then
\begin{equation}\label{eq:hadamard}
  f(z) = e^{Q(z)} z^m \prod_{n=1}^{\infty} E_p\!\Bigl(\frac{z}{a_n}\Bigr),
\end{equation}
where $Q$ is a polynomial of degree at most $p$.
\end{theorem}

The key improvement over the general Weierstrass theorem is that the
integers $p_n$ can all be taken equal to $p = \lfloor\rho\rfloor$, and
the function $g(z)$ is a polynomial (not just an arbitrary entire function).

\begin{example}[Hadamard factorisation of $\sin(\pi z)$]\label{ex:hadamard-sin}
$\sin(\pi z)$ has order $\rho = 1$, so $p = 1$. Its zeros are
$a_n = n$ ($n \in \Z \setminus \{0\}$) and it has a simple zero at $0$.
By Hadamard:
\[
  \sin(\pi z) = e^{Q(z)}\, z \prod_{n=1}^{\infty}E_1(z/n)\, E_1(-z/n)
  = e^{Q(z)}\, z\prod_{n=1}^{\infty}\Bigl(1-\frac{z^2}{n^2}\Bigr),
\]
where $\deg Q \leq 1$. From the product formula (Theorem~\ref{thm:sin-product}),
we know the answer is $\pi z \prod(1-z^2/n^2)$, so $e^{Q(z)} = \pi$,
i.e., $Q(z) = \log\pi$ (a constant). This is consistent with $\deg Q \leq 1$.
\end{example}

\begin{example}[Hadamard factorisation of $1/\Gamma(z)$]\label{ex:hadamard-gamma}
\index{gamma function!factorisation}
The reciprocal gamma function $1/\Gamma(z)$ is entire of order $1$, with
simple zeros at $z = 0, -1, -2, \ldots$. The Hadamard (= Weierstrass) product gives:
\[
  \frac{1}{\Gamma(z)} = z\, e^{\gamma z}\prod_{n=1}^{\infty}
  \Bigl(1 + \frac{z}{n}\Bigr)e^{-z/n},
\]
where $\gamma = \lim_{n\to\infty}(1+1/2+\cdots+1/n - \log n)$ is the
Euler--Mascheroni constant\index{Euler--Mascheroni constant}.
\end{example}

\section{Jensen's Formula}\label{sec:jensen}

Jensen's formula connects the distribution of zeros of a holomorphic
function to its growth, providing a bridge between the material of this
chapter and the preceding one.

\begin{theorem}[Jensen's Formula]\label{thm:jensen}
\index{Jensen's formula}
Let $f$ be holomorphic on a neighbourhood of $\overline{D(0,R)}$,
with $f(0) \neq 0$. Let $a_1, \ldots, a_N$ be the zeros of $f$ in
$D(0,R)$ (counted with multiplicity). Then
\begin{equation}\label{eq:jensen}
  \log\abs{f(0)} = -\sum_{k=1}^{N}\log\frac{R}{\abs{a_k}}
  + \frac{1}{2\pi}\int_0^{2\pi}\log\abs{f(Re^{i\theta})}\,\dd\theta.
\end{equation}
\end{theorem}

\begin{proof}
First assume $f$ has no zeros on $\abs{z}=R$. Define
\[
  g(z) = f(z)\prod_{k=1}^{N}\frac{R^2 - \conj{a_k}z}{R(z-a_k)}.
\]
Each factor $\frac{R^2-\conj{a_k}z}{R(z-a_k)}$ is a (rotated) Blaschke factor
with modulus $1$ on $\abs{z}=R$ and a simple zero at $z = R^2/\conj{a_k}$
(outside $\overline{D(0,R)}$) and a simple pole at $z = a_k$. Thus $g$
is holomorphic and nonvanishing on $\overline{D(0,R)}$, and
$\abs{g(z)} = \abs{f(z)}$ on $\abs{z} = R$.

Since $g$ is holomorphic and nonzero on $\overline{D(0,R)}$, $\log\abs{g}$
is harmonic there, so the mean value property gives:
\[
  \log\abs{g(0)} = \frac{1}{2\pi}\int_0^{2\pi}\log\abs{g(Re^{i\theta})}\,\dd\theta
  = \frac{1}{2\pi}\int_0^{2\pi}\log\abs{f(Re^{i\theta})}\,\dd\theta.
\]
Now $g(0) = f(0)\prod_{k=1}^{N}\frac{R^2}{R(-a_k)}
= f(0)\prod_{k=1}^{N}\frac{-R}{a_k}$, so
\[
  \abs{g(0)} = \abs{f(0)}\prod_{k=1}^{N}\frac{R}{\abs{a_k}}.
\]
Taking logarithms:
$\log\abs{f(0)} + \sum_k \log(R/\abs{a_k}) = \frac{1}{2\pi}\int_0^{2\pi}\log\abs{f(Re^{i\theta})}\,\dd\theta$.
\end{proof}

\begin{remark}[Counting function]\label{rem:counting-function}
Jensen's formula can be rewritten using the \textbf{counting function}
$n(r) = \#\{k : \abs{a_k} \leq r\}$ as
\[
  \log\abs{f(0)} = -\int_0^R \frac{n(t)}{t}\,\dd t
  + \frac{1}{2\pi}\int_0^{2\pi}\log\abs{f(Re^{i\theta})}\,\dd\theta.
\]
This form makes it clear that faster growth of $f$ on circles forces
more zeros inside those circles — a quantitative version of the principle
that entire functions of large order must have many zeros.
\end{remark}


\section{Exercises}\label{sec:exercises-ch10}

\begin{exercise}[Entire functions omitting values]\label{exer:entire-omit}
Let $f$ be entire. Show that if $f$ omits two distinct values
$a, b \in \C$, then $f$ is constant.

\emph{Hint:} Consider $g(z) = (f(z)-a)/(b-a)$, which omits $0$ and $1$.
Apply Little Picard.
\end{exercise}

\begin{exercise}[Order computation]\label{exer:order-comp}
Compute the order of each of the following entire functions:
\begin{enumerate}[label=(\alph*)]
  \item $f(z) = e^{z^3 + z}$;
  \item $f(z) = \cos(\sqrt{z})$ (interpret as the entire function
    $\sum_{n=0}^{\infty}(-1)^n z^n/(2n)!$);
  \item $f(z) = \sum_{n=0}^{\infty} z^n / (n!)^2$.
\end{enumerate}
\end{exercise}

\begin{exercise}[Weierstrass product for $\cos$]\label{exer:cos-product}
Derive the product formula
\[
  \cos(\pi z) = \prod_{n=0}^{\infty}
  \Bigl(1 - \frac{4z^2}{(2n+1)^2}\Bigr)
\]
from the product formula for $\sin(\pi z)$ using
$\cos(\pi z) = \sin(2\pi z)/(2\sin(\pi z))$.
\end{exercise}

\begin{exercise}[Mittag-Leffler for $1/\sin^2$]\label{exer:mittag-leffler-csc2}
Verify the partial fraction expansion
$\displaystyle \frac{\pi^2}{\sin^2(\pi z)} = \sum_{n=-\infty}^{\infty}\frac{1}{(z-n)^2}$
by showing that the difference of the two sides is an entire bounded
function, hence constant, and evaluate the constant.
\end{exercise}

\begin{exercise}[Jensen's formula application]\label{exer:jensen-app}
Let $f(z) = \sin(\pi z)$. Apply Jensen's formula with $R = n + 1/2$
(where $n$ is a positive integer) to deduce the asymptotic formula
$n(r) \sim 2r$ for the number of zeros of $\sin(\pi z)$ in $\abs{z} \leq r$.
\end{exercise}

\begin{exercise}[Elliptic function of order 2]\label{exer:elliptic-order2}
Let $f$ be an elliptic function of order $2$ with periods $\omega_1, \omega_2$.
\begin{enumerate}[label=(\alph*)]
  \item Show that $f$ takes every value in $\C \cup \{\infty\}$ exactly
    twice (counted with multiplicity) in each fundamental parallelogram.
  \item Show that if $f$ has a single double pole at $z = 0 \pmod{\Lambda}$,
    then $f(z) = A\,\wp(z) + B$ for constants $A, B$.
\end{enumerate}
\end{exercise}

\begin{exercise}[Growth and zeros]\label{exer:growth-zeros}
Let $f$ be an entire function of order $\rho < \infty$, and let
$(a_n)$ be its nonzero zeros. Show that
$\sum_{n=1}^{\infty}\abs{a_n}^{-\alpha}$ converges for every
$\alpha > \rho$.

\emph{Hint:} Use Jensen's formula to relate $n(r)$ to $\log M(r)$,
then integrate by parts.
\end{exercise}

\begin{exercise}[Liouville for meromorphic functions]\label{exer:liouville-mero}
Let $f$ be a meromorphic function on $\C$ satisfying
$\abs{f(z)} \leq C\abs{z}^N$ for $\abs{z}$ large and some constants
$C, N$. Show that $f$ is a rational function of degree at most $N$.

\emph{Hint:} If $f$ has poles $b_1,\ldots,b_M$ in $\C$, consider
$g(z) = f(z)\prod(z-b_j)^{n_j}$, which is entire.
\end{exercise}

\begin{exercise}[Hadamard for $\cos$]\label{exer:hadamard-cos}
Use the Hadamard factorisation theorem to show that
\[
  \cos(z) = \prod_{n=0}^{\infty}\Bigl(1 - \frac{4z^2}{(2n+1)^2\pi^2}\Bigr).
\]
Verify that this is consistent with $\cos(z)$ having order $1$ and zeros
at $z = (n+1/2)\pi$, $n \in \Z$.
\end{exercise}

\begin{exercise}[The Weierstrass $\sigma$-function]\label{exer:sigma-function}
\index{Weierstrass!sigma-function}
The \textbf{Weierstrass $\sigma$-function} is defined by
\[
  \sigma(z) = z\prod_{\substack{\omega\in\Lambda\\\omega\neq 0}}
  \Bigl(1-\frac{z}{\omega}\Bigr)\exp\!\Bigl(\frac{z}{\omega}+\frac{z^2}{2\omega^2}\Bigr).
\]
Show that $\sigma$ is an entire function with simple zeros at the lattice
points and that $\wp(z) = -(\log\sigma(z))''$.
\end{exercise}


%% ═══════════════════════════════════════════════════════════════════
%%  BIBLIOGRAPHY
%% ═══════════════════════════════════════════════════════════════════
\begin{thebibliography}{99}

\bibitem{ahlfors1979}
  L.\,V.\ Ahlfors,
  \emph{Complex Analysis},
  3rd ed., McGraw-Hill, New York, 1979.

\bibitem{conway1978}
  J.\,B.\ Conway,
  \emph{Functions of One Complex Variable},
  2nd ed., Graduate Texts in Mathematics \textbf{11}, Springer-Verlag,
  New York, 1978.

\bibitem{stein2003}
  E.\,M.\ Stein and R.\ Shakarchi,
  \emph{Complex Analysis},
  Princeton Lectures in Analysis~II, Princeton University Press, 2003.

\bibitem{cartan1995}
  H.\ Cartan,
  \emph{Elementary Theory of Analytic Functions of One or Several
  Complex Variables},
  Dover Publications, New York, 1995 (reprint of the 1963 edition).

\bibitem{rudin1987}
  W.\ Rudin,
  \emph{Real and Complex Analysis},
  3rd ed., McGraw-Hill, New York, 1987.

\bibitem{berenstein1991}
  C.\,A.\ Berenstein and R.\ Gay,
  \emph{Complex Variables: An Introduction},
  Graduate Texts in Mathematics \textbf{125}, Springer-Verlag,
  New York, 1991.

\bibitem{remmert1991}
  R.\ Remmert,
  \emph{Theory of Complex Functions},
  Graduate Texts in Mathematics \textbf{122}, Springer-Verlag,
  New York, 1991.

\bibitem{boas1954}
  R.\,P.\ Boas,
  \emph{Entire Functions},
  Academic Press, New York, 1954.

\bibitem{lang1999}
  S.\ Lang,
  \emph{Complex Analysis},
  4th ed., Graduate Texts in Mathematics \textbf{103}, Springer-Verlag,
  New York, 1999.

\bibitem{needham1997}
  T.\ Needham,
  \emph{Visual Complex Analysis},
  Clarendon Press, Oxford, 1997.

\end{thebibliography}

%% ═══════════════════════════════════════════════════════════════════
%%  INDEX
%% ═══════════════════════════════════════════════════════════════════
\cleardoublepage
\printindex

\end{document}
